\documentclass[12pt,oneside]{book}

% ==================== METADATOS ====================
\newcommand{\universidad}{Universidad de Buenos Aires (UBA)}
\newcommand{\facultad}{Facultad de Derecho}
\newcommand{\materia}{Derecho Procesal Civil y Comercial}
\newcommand{\titulo}{Derecho Procesal Civil y Comercial}
\newcommand{\subtitulo}{Prueba, sentencia, impugnación, ejecución y sistemas cautelares}
\newcommand{\autor}{Isaac Edgar Camacho Ocampo}
\newcommand{\anio}{2020--2025}
% TODO: verificar consistencia entre fuentes originales (los archivos fuente indican años 2020, 2024 y 2025; el curso figura como <<Derecho Procesal Civil>> y como <<Derecho Procesal Civil y Comercial>>).

% ==================== PAQUETES ====================
\usepackage[utf8]{inputenc}
\usepackage[T1]{fontenc}
\usepackage[spanish,es-nodecimaldot]{babel}
\usepackage[
    a4paper,
    top=2.5cm,
    bottom=2.5cm,
    left=3cm,
    right=2.5cm,
    headheight=15pt
]{geometry}
\usepackage{amsmath}
\usepackage{amssymb}
\usepackage{mathtools}
\usepackage{graphicx}
\usepackage{float}
\usepackage{caption}
\usepackage{subcaption}
\usepackage{booktabs}
\usepackage{array}
\usepackage{longtable}
\usepackage{tabularx}
\usepackage{enumitem}
\usepackage{xcolor}
\usepackage[most]{tcolorbox}
\usepackage{fancyhdr}
\usepackage{ragged2e}
\usepackage[
    colorlinks=true,
    linkcolor=blue!50!black,
    citecolor=green!40!black,
    urlcolor=blue!60!black,
    pdftitle={Derecho Procesal Civil y Comercial: prueba, sentencia, impugnación, ejecución y sistemas cautelares},
    pdfauthor={Isaac Edgar Camacho Ocampo},
    pdfsubject={Derecho Procesal Civil y Comercial},
    bookmarks=true
]{hyperref}

% ==================== CONFIGURACIÓN GENERAL ====================
\setcounter{tocdepth}{2}
\setlength{\parindent}{0pt}
\setlength{\parskip}{0.5em}
\clubpenalty=10000
\widowpenalty=10000
\setlist[itemize]{topsep=0.3em, itemsep=0.2em, parsep=0.1em}
\setlist[enumerate]{topsep=0.3em, itemsep=0.2em, parsep=0.1em}
\renewcommand{\arraystretch}{1.15}
\captionsetup{font=small, labelfont=bf}

% Tipos de columna usados en las tablas
\newcolumntype{L}[1]{>{\raggedright\arraybackslash}p{#1}}
\newcolumntype{Y}{>{\raggedright\arraybackslash}X}

% ==================== ENCABEZADOS Y PIES ====================
\pagestyle{fancy}
\fancyhf{}
\fancyhead[L]{\small\nouppercase{\leftmark}}
\fancyhead[R]{\small\materia}
\fancyfoot[C]{\thepage}
\renewcommand{\headrulewidth}{0.4pt}
\renewcommand{\footrulewidth}{0pt}
\fancypagestyle{plain}{%
    \fancyhf{}
    \fancyfoot[C]{\thepage}
    \renewcommand{\headrulewidth}{0pt}
}

% ==================== COMANDOS MATEMÁTICOS ====================
\newcommand{\abs}[1]{\left|#1\right|}
\newcommand{\norm}[1]{\left\lVert#1\right\rVert}
\newcommand{\vect}[1]{\mathbf{#1}}
\newcommand{\deriv}[2]{\frac{d#1}{d#2}}
\newcommand{\pderiv}[2]{\frac{\partial#1}{\partial#2}}

% ==================== CAJAS ACADÉMICAS ====================
\newtcolorbox{definition}[1]{enhanced, breakable, title={Definición: #1}, fonttitle=\bfseries\small,
    colback=blue!3, colframe=blue!50!black, boxrule=0.8pt, arc=2mm}
\newtcolorbox{important}{enhanced, breakable, title={Importante}, fonttitle=\bfseries\small,
    colback=yellow!8, colframe=orange!70!black, boxrule=0.8pt, arc=2mm}
\newtcolorbox{example}[1]{enhanced, breakable, title={Ejemplo: #1}, fonttitle=\bfseries\small,
    colback=green!4, colframe=green!40!black, boxrule=0.8pt, arc=2mm}
\newtcolorbox{formula}[1]{enhanced, breakable, title={Regla: #1}, fonttitle=\bfseries\small,
    colback=purple!4, colframe=purple!50!black, boxrule=0.8pt, arc=2mm}
\newtcolorbox{exercise}[1]{enhanced, breakable, title={Ejercicio: #1}, fonttitle=\bfseries\small,
    colback=gray!5, colframe=gray!60!black, boxrule=0.8pt, arc=2mm}
\newtcolorbox{note}{enhanced, breakable, title={Nota}, fonttitle=\bfseries\small,
    colback=white, colframe=gray!50, boxrule=0.6pt, arc=2mm}
% Texto normativo (artículos y normas transcriptos en las fuentes)
\newtcolorbox{norma}[1]{enhanced, breakable, blanker, borderline west={2pt}{0pt}{blue!50!black},
    left=10pt, before skip=8pt, after skip=8pt, fontupper=\small,
    before upper={\textbf{#1}\par\smallskip}}
\newtcolorbox{normativa}[1]{enhanced, breakable, blanker, borderline west={2pt}{0pt}{blue!50!black},
    left=10pt, before skip=8pt, after skip=8pt, fontupper=\small,
    before upper={\textbf{#1}\par\smallskip}}
\newtcolorbox{articulo}[1]{enhanced, breakable, blanker, borderline west={2pt}{0pt}{blue!50!black},
    left=10pt, before skip=8pt, after skip=8pt, fontupper=\small,
    before upper={\textbf{#1}\par\smallskip}}

% ==================== CUADROS CORNELL ====================
% Tabla real (tabularx): columna izquierda estrecha con preguntas, columna derecha amplia con
% respuestas y última fila completa con el resumen.
\newcommand{\cornelltitle}[1]{\par\bigskip\noindent\textbf{\sffamily Cuadro Cornell: #1}\par\nopagebreak\smallskip}
\newcommand{\cqrow}[2]{\textbf{#1} & #2 \\[0.5em]}
\newcommand{\cornellblock}[1]{%
  \noindent\begin{tabularx}{\textwidth}{@{}>{\raggedright\arraybackslash}p{0.27\textwidth}!{\vrule width 0.4pt}>{\raggedright\arraybackslash}X@{}}
  \toprule
  \textbf{Preguntas / palabras clave} & \textbf{Notas / respuestas} \\
  \midrule
  #1
  \bottomrule
  \end{tabularx}\par}
\newcommand{\cornellblocksum}[2]{%
  \noindent\begin{tabularx}{\textwidth}{@{}>{\raggedright\arraybackslash}p{0.27\textwidth}!{\vrule width 0.4pt}>{\raggedright\arraybackslash}X@{}}
  \toprule
  \textbf{Preguntas / palabras clave} & \textbf{Notas / respuestas} \\
  \midrule
  #1
  \midrule
  \multicolumn{2}{@{}p{\textwidth}@{}}{\textbf{Resumen:} #2} \\
  \bottomrule
  \end{tabularx}\par}

\begin{document}

\begin{titlepage}
\thispagestyle{empty}
\begin{center}
\vspace*{1cm}
{\large \textsc{\universidad}}\\[0.2cm]
{\large \textsc{\facultad}}
\vspace{3cm}

{\Huge \textbf{\titulo}}
\vspace{0.8cm}

{\Large \subtitulo}
\vspace{1.2cm}

\rule{0.70\textwidth}{0.5pt}\\[0.4cm]
{\large Material de estudio}\\[0.4cm]
\rule{0.70\textwidth}{0.5pt}
\vfill
{\large \autor}\\[0.3cm]
{\large \anio}
\end{center}
\vfill
\begin{flushleft}
\textit{Este es un modesto aporte a los estudiantes de ``Derecho Procesal Civil y Comercial''; no pretende sustituir el material oficial de la materia.}
\end{flushleft}
\end{titlepage}

\frontmatter
\tableofcontents
\mainmatter

\chapter*{Introducción general}
\addcontentsline{toc}{chapter}{Introducción general}
\markboth{Introducción general}{Introducción general}
Este volumen reúne y ordena el desarrollo del proceso civil y comercial argentino desde la producción de la prueba hasta la ejecución de la sentencia y los sistemas cautelares. La lectura sigue la progresión natural del proceso: primero la \textbf{prueba y la conclusión de la causa} (Parte~I); luego las formas en que el proceso termina y la \textbf{sentencia} (Parte~II); después los mecanismos para atacar las decisiones, es decir, la \textbf{impugnación y los recursos} (Parte~III); a continuación la \textbf{ejecución de la sentencia} y el juicio ejecutivo (Parte~IV); y, finalmente, los \textbf{sistemas cautelares} (Parte~V), que permiten asegurar el resultado del proceso en cualquiera de sus etapas.

Cada capítulo se cierra con uno o más cuadros Cornell, que formulan preguntas de recuperación, sus respuestas y un resumen, para facilitar el estudio activo.

\part{La prueba y la conclusión de la causa}
\markboth{Parte I}{Parte I}
Esta parte recorre el tramo del proceso civil que va desde la producción de los medios de prueba hasta el momento en que el juez se dispone a sentenciar. Los dos primeros capítulos abordan dos medios de prueba fundamentales del proceso civil y comercial argentino: la \textbf{prueba testimonial} (declaración de testigos) y la \textbf{prueba de confesión} (absolución de posiciones). En ambos casos se analiza el régimen legal aplicable conforme al Código Procesal Civil y Comercial de la Nación (CPCCN), los deberes y derechos de los declarantes, las formas válidas de ofrecimiento, producción y valoración de la prueba, y los supuestos de caducidad que operan objetivamente ante la inactividad procesal de las partes. Se incluyen referencias al Código Civil y Comercial de la Nación (CCyC), jurisprudencia de la Corte Suprema de Justicia de la Nación (CSJN) y ejemplos prácticos del ejercicio forense.

El proceso judicial puede compararse con una investigación periodística: primero hay que conseguir los testimonios de quienes presenciaron los hechos (prueba testimonial) y luego obtener reconocimientos de los propios involucrados (prueba de confesión). Cada medio de prueba sigue tres etapas: \textbf{ofrecimiento}, \textbf{producción} y \textbf{valoración}.

Los tres capítulos siguientes se ocupan de la prueba pericial, la prueba indiciaria y el reconocimiento judicial. Luego se estudia cómo concluye la etapa probatoria, lo que conduce a los alegatos y a la etapa conclusional, y se cierra con los sistemas de valoración de la prueba, es decir, con el modo en que el juez aprecia la prueba producida al dictar sentencia.

\chapter[Prueba testimonial]{Prueba testimonial (arts.~426 a 456 CPCCN)}\label{cap:testimonial}

% =====================================================

La prueba testimonial comienza a regularse a partir del artículo 426 del CPCCN.

\section{Concepto de testigo}

\begin{definition}{Testigo}
El \textbf{testigo} es un tercero que no es parte en el proceso, convocado para que preste declaración sobre hechos de su conocimiento. Es convocado durante la etapa probatoria para que relate lo que sabe sobre los hechos controvertidos o ventilados en la causa.
\end{definition}

\section{Objeto del testimonio y valor según el tipo de percepción}

El objeto del testimonio son los \textbf{hechos alegados por las partes} en sus escritos postulatorios. El testigo puede conocer esos hechos a través de distintas vías:
\begin{itemize}
    \item haberlos visto directamente;
    \item haberlos escuchado;
    \item haberlos conocido a través del relato de otra persona.
\end{itemize}

No todos los testimonios tienen la misma eficacia probatoria: no es lo mismo quien vio un hecho directamente que quien lo conoce de oídas. La valoración del juez al dictar sentencia será distinta en cada caso. Sin embargo, todos esos tipos de testimonio son admisibles.

\begin{norma}{Artículo 426 CPCCN --- Procedencia de la prueba testimonial}
\textit{Toda persona mayor de 14 años podrá ser propuesta como testigo y tendrá el deber de comparecer y declarar, salvo las excepciones establecidas por la ley.}
\end{norma}

\section{Deberes del testigo}

Sobre todo testigo recaen tres deberes:

\begin{center}
\renewcommand{\arraystretch}{1.3}
\begin{tabularx}{\textwidth}{
    >{\raggedright\arraybackslash}p{0.17\textwidth}
    >{\raggedright\arraybackslash}X
    >{\raggedright\arraybackslash}X}
\toprule
\textbf{Deber} & \textbf{Alcance} & \textbf{Excepciones} \\
\midrule
\textbf{Comparecer} & Es una carga pública. Toda persona convocada como testigo debe presentarse ante el juzgado. & Cargo que implique declarar por escrito; causa justificada debidamente acreditada; irregularidad en la notificación (menos de 3 días de anticipación). \\
\addlinespace
\textbf{Declarar} & El testigo tiene la obligación de prestar declaración testimonial. & Secreto profesional (militar, científico, artístico); declaración que perjudique el honor del testigo o lo exponga a enjuiciamiento penal. \\
\addlinespace
\textbf{Decir verdad} & Antes de declarar, el testigo presta juramento o promesa de decir verdad. Se le informan las consecuencias penales del falso testimonio. & No hay excepciones al juramento, pero si el testigo miente incurre en el delito de falso testimonio. \\
\bottomrule
\end{tabularx}
\end{center}

\subsection*{Regla de los 70 kilómetros}

Respecto del deber de comparecer, si el domicilio del testigo se encuentra dentro de un radio de \textbf{70 kilómetros} del tribunal, tiene la obligación de concurrir personalmente. Si su domicilio está fuera de ese radio, se lo convoca mediante la \textbf{Ley 22.172}, a través de un oficio dirigido al juzgado de turno en la sede de su domicilio. Ese oficio debe contener los datos personales del testigo y el interrogatorio completo, ya que la parte contraria también tiene derecho a controlar este medio de prueba y a formular sus propias preguntas.

\section{Testigos excluidos: art. 427 CPCCN y vinculación con el CCyC}

El artículo 427 establece que no pueden ser ofrecidos como testigos los consanguíneos o afines en línea directa de las partes, ni el cónyuge aunque estuviere separado legalmente, salvo en los casos de reconocimiento de firmas. Los parientes colaterales (hermanos, tíos, sobrinos) sí pueden ser testigos.

\begin{norma}{Artículo 427 CPCCN --- Testigos excluidos}
\textit{No podrán ser ofrecidos como testigos los consanguíneos o afines en línea directa de las partes, ni el cónyuge aunque estuviere separado legalmente, salvo si se trata de reconocimiento de firmas.}
\end{norma}

\begin{norma}{Artículo 711 del Código Civil y Comercial de la Nación --- Testigos en juicios de familia}
\textit{Los parientes y allegados a las partes pueden ser ofrecidos como testigos en juicios de familia.}

El fundamento es que, en los conflictos familiares, muchas veces quienes mejor conocen la problemática son precisamente los parientes o allegados, y excluirlos generaría una imposibilidad de probar los hechos alegados (por ejemplo, en casos de violencia familiar).
\end{norma}

\section{Los tres momentos de la prueba testimonial}

Es importante distinguir tres momentos en todo medio probatorio:
\begin{itemize}
    \item \textbf{Ofrecimiento}: se realiza en los escritos postulatorios (demanda y contestación de demanda).
    \item \textbf{Producción}: se lleva a cabo durante la etapa probatoria, una vez abierta la causa a prueba.
    \item \textbf{Valoración}: las partes valoran la prueba en los alegatos; el juez la valora al dictar sentencia.
\end{itemize}

\section{Ofrecimiento de la prueba testimonial}

Para ofrecer un testigo en el escrito postulatorio, el código exige informar:
\begin{itemize}
    \item nombre completo del testigo;
    \item ocupación;
    \item domicilio.
\end{itemize}

El \textbf{artículo 333 CPCCN} agrega un elemento adicional: argumentar el \textbf{objeto de la citación}, es decir, explicar sintéticamente por qué se convoca a esa persona y sobre qué hechos declarará.

\begin{example}{Objeto de la citación}
``Se ofrece a Juan Pérez porque estuvo presente en el momento del accidente'', o ``porque es el médico que estuvo a cargo del tratamiento de la parte actora''.
\end{example}

Se pueden ofrecer \textbf{hasta ocho testigos}. Si la complejidad del caso lo amerita, pueden admitirse más, previa decisión judicial sobre la necesidad de la ampliación.

\section{Producción de la prueba testimonial: citación a audiencia}

Una vez que la causa está abierta a prueba y el juez ha fijado las audiencias de testigos (en la etapa probatoria, luego de la audiencia preliminar), se debe notificar a los testigos. Los medios de notificación admitidos son:
\begin{itemize}
    \item cédula de notificación;
    \item telegrama colacionado;
    \item acta notarial;
    \item carta documento.
\end{itemize}

La notificación debe cursarse con al menos \textbf{tres días} de anticipación a la audiencia. La cédula debe informar tanto la \textbf{audiencia principal} (primera) como la \textbf{audiencia supletoria} (segunda), con la advertencia de que si el testigo no concurre a la primera sin causa justificada, será convocado a la segunda con citación de la fuerza pública (es decir, a través de la policía).

\section{Desarrollo de la audiencia testimonial e interrogatorio preliminar}

Al comenzar la audiencia, el testigo se identifica con su documento nacional de identidad y brinda sus datos personales. Luego se le toma juramento o promesa de decir verdad, explicándole las consecuencias penales del falso testimonio. Inmediatamente después, se procede al interrogatorio preliminar.

\begin{norma}{Artículo 441 CPCCN --- Interrogatorio preliminar (generales de la ley)}
\textit{Aunque las partes no lo pidan, los testigos serán siempre preguntados por: 1) Su nombre, edad, estado, profesión y domicilio; 2) Si es pariente por consanguinidad o afinidad de alguna de las partes y en qué grado; 3) Si tiene interés directo o indirecto en el pleito; 4) Si es amigo íntimo o enemigo de alguna de las partes; 5) Si es dependiente, acreedor o deudor de algunos de los litigantes, o si tiene algún otro género de relación con ellos o algún interés en el pleito.}
\end{norma}

Estas preguntas se denominan, en el lenguaje forense, \textbf{``generales de la ley''}. Su propósito es evaluar si el testigo puede estar afectado por alguna de estas circunstancias, lo que podría incidir en la valoración de su testimonio e incluso en su idoneidad (conforme al art.~456 CPCCN).

\section{Forma de las preguntas al testigo: art. 443}

Luego del interrogatorio preliminar, cada parte puede formular su interrogatorio al testigo, que puede presentarse por escrito o realizarse a viva voz. El artículo 443 establece los requisitos que deben cumplir las preguntas:

\begin{norma}{Artículo 443 CPCCN --- Forma de las preguntas}
\textit{Las preguntas no contendrán más de un hecho, serán claras y concretas. No se formularán las que estén concebidas en términos afirmativos, sugieran la respuesta o sean ofensivas o vejatorias. No podrán contener referencias de carácter técnico, salvo si fuesen dirigidas a personas especializadas.}
\end{norma}

En la práctica forense, la pregunta se formula con la siguiente fórmula ritual: \textit{``Para que jure, diga el testigo si sabe y cómo le consta\ldots{} [el hecho sobre el que se pregunta]''}.

La pregunta debe ser abierta y no sugestiva porque el objetivo es conocer genuinamente lo que el testigo sabe y cómo lo sabe, no confirmar una versión ya armada. Por ejemplo, si se pregunta si el testigo sabe qué sucedió tal día en tal lugar, se busca que el testigo relate libremente, no que responda ``sí'' o ``no'' a una afirmación del abogado.

Si una pregunta viola los requisitos del artículo 443, la parte contraria puede oponerse. Esa oposición genera una \textbf{incidencia} dentro de la audiencia: si el juez hace lugar a la oposición, la pregunta queda rechazada y las costas de la incidencia se imponen a quien formuló la pregunta defectuosa.

\section{Testigo técnico y perito}

La mención a las referencias técnicas del artículo 443 permite distinguir dos figuras que no deben confundirse:

\begin{center}
\renewcommand{\arraystretch}{1.3}
\begin{tabularx}{\textwidth}{
    >{\raggedright\arraybackslash}p{0.2\textwidth}
    >{\raggedright\arraybackslash}X
    >{\raggedright\arraybackslash}X}
\toprule
 & \textbf{Testigo técnico} & \textbf{Perito} \\
\midrule
\textbf{Fungibilidad} & No es fungible. Es una persona específica que conoció los hechos. & Es fungible. Se designa por sorteo de oficio por el juez. \\
\addlinespace
\textbf{Conocimiento del hecho} & Conoció los hechos en el pasado, cuando ocurrieron. & Toma conocimiento del caso en el presente, al momento de su designación. \\
\addlinespace
\textbf{Ejemplo} & El médico que trató al paciente (parte actora) durante su enfermedad. & El perito médico designado por sorteo para realizar la pericia médica en el proceso. \\
\bottomrule
\end{tabularx}
\end{center}

Al testigo técnico, que es una persona especializada, sí se le pueden hacer preguntas de carácter técnico. Al perito, en cambio, se le formulan puntos de pericia en los escritos postulatorios.

\section{Negativa a responder: art. 444}

\begin{norma}{Artículo 444 CPCCN --- Negativa a responder}
\textit{El testigo podrá rehusarse a contestar las preguntas si la respuesta lo expusiera a enjuiciamiento penal o comprometiera su honor, o si no pudiera responder sin revelar un secreto profesional, militar, científico, artístico o industrial.}
\end{norma}

\section{Caducidad de la prueba testimonial: arts. 432 y 437}

Debe distinguirse la negligencia de la caducidad en la producción de la prueba. La \textbf{negligencia} se basa en un elemento subjetivo: el desinterés de la parte en producir su prueba a lo largo del período probatorio. La \textbf{caducidad}, en cambio, opera de forma objetiva: el código establece supuestos específicos que, al configurarse, producen automáticamente la pérdida del derecho a esa prueba.

\begin{norma}{Artículo 432 CPCCN --- Caducidad de la prueba testimonial}
\textit{A pedido de parte y sin sustanciación, se tendrá por desistido del testigo a la parte que lo propuso, en los siguientes supuestos: a) Si no hubiera activado la citación del testigo y este no hubiera comparecido por esta razón; b) Si no habiendo comparecido el testigo a la primera audiencia sin invocar causa justificada, no requiriera oportunamente las medidas de compulsión necesarias; c) Si fracasara la segunda audiencia por motivos no imputables a la parte y no solicitara nueva audiencia dentro del quinto día.}
\end{norma}

\begin{norma}{Artículo 437 CPCCN --- Cuarto supuesto de caducidad}
\textit{Si la parte que ofreció al testigo no concurriera a la audiencia por sí o por apoderado y no hubiere dejado interrogatorio, se la tendrá por desistida de aquel sin sustanciación.}
\end{norma}

\begin{center}
\renewcommand{\arraystretch}{1.3}
\begin{tabularx}{\textwidth}{
    >{\centering\arraybackslash}p{0.06\textwidth}
    >{\raggedright\arraybackslash}X
    >{\raggedright\arraybackslash}X
    >{\raggedright\arraybackslash}p{0.17\textwidth}}
\toprule
\textbf{N.º} & \textbf{Supuesto} & \textbf{Descripción} & \textbf{Norma} \\
\midrule
1 & No se notificó al testigo y no compareció. & La parte no libró la cédula ni ningún medio de notificación al testigo. & Art. 432 inc.~a) \\
\addlinespace
2 & El testigo no compareció a la 1.ª audiencia y no se solicitaron medidas de compulsión. & El testigo faltó sin causa y la parte no gestionó la segunda audiencia con fuerza pública. & Art. 432 inc.~b) \\
\addlinespace
3 & Fracasó la 2.ª audiencia y no se pide una nueva en 5 días. & La segunda audiencia fracasó por causas no imputables a la parte, pero esta no solicitó nueva audiencia dentro del quinto día. & Art. 432 inc.~c) \\
\addlinespace
4 & La parte oferente no concurrió y no dejó interrogatorio. & El testigo sí compareció, pero la parte que lo ofreció no se presentó ni dejó el pliego de preguntas. & Art. 437 \\
\bottomrule
\end{tabularx}
\end{center}

El desistimiento del testigo por quien lo ofreció es válido, sin necesidad de consentimiento de la contraparte, siempre que sea previo a la celebración de la audiencia.

\section{Idoneidad del testigo: art. 456}

\begin{norma}{Artículo 456 CPCCN --- Idoneidad de los testigos}
\textit{Dentro del plazo de prueba, las partes podrán alegar y probar acerca de la idoneidad de los testigos. El juez apreciará, según las reglas de la sana crítica y en oportunidad de dictar sentencia definitiva, las circunstancias y motivos que corroboren o disminuyan la fuerza de las declaraciones.}
\end{norma}

Este artículo permite a las partes generar un \textbf{incidente} por separado del proceso principal para demostrar que un testigo no era idóneo para declarar. Se puede ofrecer en ese incidente toda la prueba que se considere necesaria. Es un mecanismo de defensa ante una declaración que, si bien era formalmente admisible, reveló en su contenido elementos que cuestionan la imparcialidad o credibilidad del testigo.

\cornelltitle{prueba testimonial}
\cornellblock{
\cqrow{¿Quién es un testigo?}{Es un tercero ajeno al proceso, convocado en la etapa probatoria para declarar sobre hechos de su conocimiento, ya sea porque los vio, los escuchó o los conoció a través del relato de terceros. No es parte del proceso.
}
\cqrow{¿Cuáles son los deberes del testigo?}{Tres deberes: 1) comparecer (carga pública; se lo puede citar con fuerza pública en caso de incumplimiento); 2) declarar (salvo secreto profesional o riesgo de autoincriminación); 3) decir verdad (juramento previo; el falso testimonio es delito penal).
}
\cqrow{¿Quiénes están excluidos como testigos?}{Conforme al art.~427, los consanguíneos o afines en línea directa y el cónyuge. Excepción: en juicios de familia (art.~711 CCyC) se admiten parientes y allegados, porque son quienes mejor conocen la problemática familiar.
}
\cqrow{¿Cómo se ofrece la prueba testimonial?}{En los escritos postulatorios, indicando nombre completo, ocupación y domicilio del testigo, más el objeto de la citación (art.~333). Se pueden ofrecer hasta 8 testigos; el juez puede admitir más si es necesario.
}
\cqrow{¿Qué son las ``generales de la ley''?}{Son las preguntas del interrogatorio preliminar (art.~441), que se formulan a todo testigo antes del interrogatorio de fondo: datos personales, parentesco con las partes, interés en el pleito, amistad o enemistad, relación de dependencia. Permiten evaluar la idoneidad e imparcialidad del testigo.
}
\cqrow{¿Cómo deben formularse las preguntas a un testigo?}{Deben contener un solo hecho y ser claras y concretas; NO estar formuladas en términos afirmativos, NO sugerir la respuesta, NO ser vejatorias u ofensivas y NO contener referencias técnicas (salvo al testigo técnico). Se formulan así: ``Para que jure, diga el testigo si sabe y cómo le consta\ldots{}''.
}
\cqrow{¿Qué diferencia hay entre testigo técnico y perito?}{El testigo técnico conoció los hechos en el pasado (por ejemplo, el médico tratante) y no es fungible. El perito es fungible, es designado por sorteo de oficio, toma conocimiento en el presente y dictamina sobre puntos de pericia formulados en los escritos postulatorios. Al primero se le pueden hacer preguntas técnicas; al segundo se le formulan puntos de pericia.
}
}
\medskip
\cornellblocksum{
\cqrow{¿Cuáles son los supuestos de caducidad de la prueba testimonial?}{1) No se notificó al testigo (art.~432 a); 2) no se gestionaron medidas de compulsión tras la 1.ª audiencia fallida (art.~432 b); 3) no se pidió nueva audiencia en 5 días tras fracasar la 2.ª (art.~432 c); 4) la parte oferente no concurrió ni dejó interrogatorio (art.~437). Los 3 primeros son objetivos; todos operan a pedido de parte.
% TODO: contenido incompleto o ambiguo en las notas originales (alcance de "objetivos" y "a pedido de parte" respecto del art. 437)
}
}{La prueba testimonial convoca a terceros ajenos al proceso para que declaren sobre hechos de su conocimiento, imponiéndoles los deberes de comparecer, declarar y decir verdad bajo juramento. Se ofrece en los escritos postulatorios, se produce en la etapa probatoria y se valora en los alegatos y en la sentencia. El CPCCN regula minuciosamente el modo de formular las preguntas (art.~443), el interrogatorio preliminar (art.~441) y los supuestos de caducidad (arts.~432 y 437).
}

\chapter[Prueba de confesión]{Prueba de confesión (arts.~404 a 424 CPCCN)}\label{cap:confesion}

% =====================================================

\section{Concepto de confesión: confesión espontánea y provocada}

\begin{definition}{Confesión}
La confesión como prueba consiste en una \textbf{declaración formulada por quien es parte en el proceso}, sobre hechos personales o de su conocimiento, que produce un efecto desfavorable al confesante y favorable para la otra parte.
\end{definition}

Se distingue entre:
\begin{itemize}
    \item \textbf{Confesión espontánea}: ocurre cuando la parte demandada reconoce un hecho en la contestación de demanda. Al reconocer ese hecho, genera plena prueba sobre él, sin necesidad de ningún otro medio probatorio.
    \item \textbf{Confesión provocada (absolución de posiciones)}: es el medio de prueba propiamente dicho. Se produce en audiencia, donde la parte absuelve las posiciones formuladas por la contraria.
\end{itemize}

La confesión en sí no es el medio de prueba: el medio de prueba es la absolución de posiciones, o confesión provocada en juicio.

\section{Momento del ofrecimiento y de la producción}

Al igual que la prueba testimonial, la confesión se ofrece en los escritos postulatorios. Su producción --- la absolución de posiciones --- se lleva a cabo en la audiencia preliminar, conforme al artículo 360 inciso 4 del CPCCN. En la práctica, muchos jueces fijan una audiencia posterior específicamente para este fin, aunque la norma establece la audiencia preliminar como el momento de producción.

\section{Quiénes pueden ser citados a absolver posiciones: art. 405}

\begin{norma}{Artículo 405 CPCCN --- Personas que pueden ser citadas}
\textit{Podrán asimismo ser citados a absolver posiciones, además del actor y el demandado: los representantes de los incapaces por los hechos en que hayan intervenido personalmente en ese carácter; los apoderados por hechos realizados en nombre de sus mandantes, estando vigente el mandato, y por hechos anteriores cuando estuvieran sus representados fuera del lugar en que se sigue el juicio, siempre que el apoderado tuviese facultades para ello y la parte contraria lo consienta; y los representantes legales de las personas jurídicas, sociedades o entidades colectivas que tuviesen facultad para obligarlas.}
\end{norma}

\section{Elección del absolvente en personas jurídicas: art. 406}

Para las empresas y personas jurídicas, el artículo 406 permite oponerse, dentro del quinto día de notificada la audiencia, a que absuelva el representante elegido por el ponente, siempre que:
\begin{enumerate}
    \item se alegue que ese representante no intervino personalmente ni tuvo conocimiento directo de los hechos;
    \item se indique en el mismo escrito el nombre del representante que absolverá las posiciones;
    \item se deje constancia de que esa persona fue notificada de la audiencia y presta su conformidad.
\end{enumerate}

\section{Citación: domicilio constituido y domicilio real}

La notificación de la audiencia de absolución de posiciones debe hacerse por cédula. La notificación varía según la forma de actuación de la parte en el proceso:

\begin{center}
\renewcommand{\arraystretch}{1.3}
\begin{tabularx}{\textwidth}{
    >{\raggedright\arraybackslash}X
    >{\raggedright\arraybackslash}X
    >{\raggedright\arraybackslash}X}
\toprule
\textbf{Modalidad} & \textbf{Tipo de actuación} & \textbf{Domicilio donde se notifica} \\
\midrule
Parte con letrado patrocinante (actúa por derecho propio) & Firma sus presentaciones junto con su abogado patrocinante. & Domicilio constituido (generalmente el estudio del abogado). \\
\addlinespace
Parte representada por apoderado & El apoderado actúa en nombre de la parte. & Domicilio real de la parte (no el domicilio procesal del apoderado). \\
\bottomrule
\end{tabularx}
\end{center}

La cédula debe entregarse con al menos \textbf{tres días} de anticipación. En caso de urgencia debidamente justificada, el juez puede reducir ese plazo mediante resolución, pero la anticipación no puede ser inferior a un día.

\section{Reserva del pliego de posiciones: art. 410}

\begin{norma}{Artículo 410 CPCCN --- Reserva del pliego de posiciones}
\textit{El pliego deberá ser entregado en secretaría hasta media hora antes de la hora fijada para la audiencia, en sobre cerrado, al que se le pondrá cargo. Si la parte que pidió las posiciones no comparece sin justa causa a la audiencia, no hubiese dejado pliego, y comparece el citado, perderá el derecho a exigirlas.}
\end{norma}

El sobre cerrado se presenta en la mesa de entradas del juzgado hasta media hora antes del inicio de la audiencia. El fundamento de este límite temporal es evitar que la parte ponente prepare dos pliegos distintos (uno para el caso de que el absolvente comparezca y otro para el caso de incomparecencia) y elija el que más le convenga según las circunstancias.

\section{Forma de redactar las posiciones: art. 411}

\begin{norma}{Artículo 411 CPCCN --- Forma de las posiciones}
\textit{Las posiciones serán claras y concretas, no contendrán más de un hecho, serán redactadas en forma afirmativa y deberán versar sobre puntos controvertidos que se refieran a la actuación personal del absolvente. El juez podrá modificar de oficio, sin recurso alguno, el orden y los términos de las posiciones propuestas por las partes sin alterar su sentido, y podrá asimismo eliminar las que fueran manifiestamente inútiles o superfluas.}
\end{norma}

La diferencia fundamental con el interrogatorio de testigos es que, mientras las preguntas a los testigos deben ser abiertas y no sugestivas, las posiciones de la confesión deben formularse en términos afirmativos, de manera que el absolvente pueda responder simplemente ``sí'' o ``no'' (aunque luego puede agregar las aclaraciones que estime pertinentes).

La fórmula es: \textit{``Para que jure como que es cierto que usted se encontraba en tal lugar tal día''}. Eso es la posición.

Un aspecto crucial es que el ponente, al formular la posición, \textbf{está reconociendo ese hecho}. Por lo tanto:
\begin{itemize}
    \item si el absolvente contesta ``sí, es cierto'': hay plena prueba sobre ese hecho (ambas partes lo reconocen);
    \item si el absolvente contesta ``no es cierto'': ese hecho deberá probarse por otro medio.
\end{itemize}

\section{Posición impertinente: art. 414}

\begin{norma}{Artículo 414 CPCCN --- Posición impertinente}
\textit{La parte podrá negarse a contestar la posición si esta fuera impertinente. El juez valorará esa oposición al dictar sentencia y, si considera que la posición era pertinente, podrá tener por confeso al absolvente en ese punto.}
\end{norma}

A diferencia de lo que ocurre en la prueba testimonial (donde la oposición a una pregunta genera una incidencia con fundamentos), en la confesión la oposición a una posición no requiere fundamentación: el abogado del absolvente simplemente manifiesta que se opone en los términos del artículo 414 y eso queda asentado en el acta. Sin embargo, si el juez al dictar sentencia considera que la posición era pertinente, tendrá por confeso al absolvente, lo que puede tener consecuencias muy significativas para el resultado del proceso.

\section{Interrogatorio de las partes: art. 415}

\begin{norma}{Artículo 415 CPCCN --- Interrogatorio de oficio}
\textit{El juez podrá interrogar de oficio a las partes en cualquier estado del proceso, y estas podrán hacerse recíprocamente las preguntas y observaciones que juzgaren convenientes en la audiencia que corresponda, siempre que el juez no las declare superfluas o improcedentes.}
\end{norma}

La posibilidad de que las partes se formulen preguntas recíprocas en cualquier audiencia, o de que el juez interrogue de oficio, es un instrumento muy valioso para esclarecer la verdad de los hechos, pero en la práctica se utiliza con poca frecuencia. Con la reforma procesal y el proceso por audiencias este panorama debería cambiar.

\section{Confesión ficta: art. 417}

\begin{norma}{Artículo 417 CPCCN --- Confesión ficta}
\textit{Si el citado no compareciera a declarar dentro de la media hora de fijada para la audiencia, o si habiendo comparecido rehusara a responder o respondiera de una manera evasiva, el juez al sentenciar lo tendrá por confeso sobre los hechos personales, teniendo en cuenta las circunstancias de la causa y las demás pruebas producidas. En caso de incomparecencia del absolvente, aunque no se hubiera extendido acta, se aplicará lo establecido en este párrafo, si el ponente hubiera presentado oportunamente el pliego de posiciones y el absolvente estuviere debidamente notificado.}
\end{norma}

Para que se configure la confesión ficta deben reunirse los siguientes requisitos:
\begin{itemize}
    \item el absolvente debe estar debidamente notificado de la audiencia (con todos los requisitos de forma ya explicados);
    \item el ponente debe haber presentado el pliego en sobre cerrado hasta media hora antes;
    \item el absolvente no comparece dentro de la media hora de la hora fijada, o se niega a contestar, o responde de manera evasiva.
\end{itemize}

Configurados esos supuestos, el juez al dictar sentencia puede tener por confeso al absolvente. La norma aclara que el juez tomará en cuenta ``las circunstancias de la causa y las demás pruebas producidas'': no es una consecuencia automática absoluta, sino que debe apreciarse en el contexto del proceso.

\section{Valor probatorio de la confesión ficta}

Existen distintas interpretaciones doctrinarias y jurisprudenciales sobre el valor de la confesión ficta:

\begin{center}
\renewcommand{\arraystretch}{1.3}
\begin{tabularx}{\textwidth}{
    >{\raggedright\arraybackslash}p{0.25\textwidth}
    >{\raggedright\arraybackslash}X}
\toprule
\textbf{Posición doctrinaria} & \textbf{Contenido} \\
\midrule
\textbf{Primer sector} & La confesión ficta importa plena prueba siempre que no haya otro elemento de juicio que la contraríe. \\
\addlinespace
\textbf{Segundo sector} & La confesión ficta es plena prueba siempre que haya otros medios de prueba que colaboren y la corroboren. \\
\addlinespace
\textbf{Tercer sector} & No tiene valor absoluto en ningún supuesto: siempre debe conjugarse la confesión ficta con el resto de las constancias y medios de prueba del expediente. \\
\bottomrule
\end{tabularx}
\end{center}

\begin{norma}{Fallo de la CSJN --- 6 de febrero de 2001 (LL 2001-B, pág. 959)}
La CSJN habilitó excepcionalmente el recurso extraordinario en una cuestión de prueba --- materia normalmente ajena a la instancia extraordinaria por ser de derecho común --- y revocó la sentencia de cámara que había desestimado el valor de la confesión ficta. El caso involucraba a una empresa de transportes que había sido tenida por confesa en audiencia de absolución de posiciones, y existían además otros elementos de juicio que corroboraban esa conclusión.
\end{norma}

Desde esta perspectiva, no habría razón para que una ficción legal (la configuración de la confesión ficta) tenga más valor que la realidad que surge del resto de las constancias de la causa. Desde un enfoque sistémico del proceso, el juez al dictar sentencia debe valorar la totalidad de las constancias de la causa, sin que la confesión ficta sea automáticamente superior a la realidad probatoria del expediente.

\section{Confesión expresa: art. 423}

\begin{norma}{Artículo 423 CPCCN --- Efectos de la confesión judicial expresa}
\textit{La confesión judicial expresa constituirá plena prueba, salvo cuando: 1) dicho medio de prueba estuviere excluido por la ley respecto de los hechos que constituyen el objeto del juicio, o sobre derechos que el confesante no puede renunciar o transigir válidamente; 2) recayere sobre hechos cuya investigación prohíbe la ley; 3) se opusiere a las constancias de instrumentos fehacientes de fecha anterior agregados al expediente.}
\end{norma}

La confesión judicial expresa constituye plena prueba, salvo en esos tres supuestos. Es decir, cuando hay confesión expresa no es necesario recurrir a la valoración del resto de la prueba (salvo las excepciones). Esta es la diferencia fundamental con la confesión ficta: la ficta es una creación legal que opera ante la inactividad del absolvente; la expresa es una declaración voluntaria y directa del confesante sobre un hecho.

\section{Confesión extrajudicial: art. 424}

La confesión extrajudicial es aquella que se produce fuera del proceso judicial propiamente dicho. Se mencionan dos ejemplos:
\begin{itemize}
    \item \textbf{Declaración ante autoridad policial}: tiene valor y puede hacerse valer en el proceso.
    \item \textbf{Manifestaciones ante el perito designado judicialmente}: si la parte, durante la realización de una pericia (por ejemplo, una pericia médica), le hace saber al perito algún hecho relevante, esas manifestaciones también pueden tener valor confesional.
\end{itemize}

Esto fue objeto de debate, pero se considera que no afecta el derecho de defensa porque ocurre en el marco del proceso y todas las partes tienen posibilidad de controlarlo. En todo caso, el juez lo valorará en la sentencia junto con el resto de los medios de prueba.

\begin{norma}{Artículo 424 CPCCN --- Alcance de la confesión}
\textit{En caso de duda, la confesión deberá interpretarse en favor de quien la hace. La confesión es indivisible, salvo cuando el confesante invoca hechos impeditivos, modificativos o extintivos, o absolutamente separables e independientes; o cuando las circunstancias calificativas expuestas por quien confiesa fueren contrarias a una presunción legal o inverosímiles, y las modalidades del caso hicieran procedente la divisibilidad.}
\end{norma}

\section{Ausencia del absolvente por enfermedad y regla de los 300 kilómetros}

Si el absolvente no puede concurrir a la audiencia por enfermedad, el artículo 418 permite que el juez o un funcionario del juzgado investido de autoridad se traslade al domicilio o lugar donde se encuentre el absolvente para tomar la declaración. La enfermedad debe justificarse con la debida anticipación y por escrito (art.~419).

En cuanto a la distancia, existe una regla análoga a la de los 70 kilómetros para testigos: si el absolvente domicilia dentro de los \textbf{300 kilómetros} del juzgado, debe concurrir ante el juez de la causa. Si su domicilio está a más de 300 kilómetros, se aplica la Ley 22.172 y se libra un oficio al juzgado competente en la zona del domicilio del absolvente, adjuntando el pliego de posiciones.

\section{Desarrollo de la audiencia de absolución de posiciones}

Al momento de celebrarse la audiencia, se procede a la apertura del sobre que contiene el pliego de posiciones. Si ambas partes ofrecieron el medio de prueba, se absolverán primero las posiciones de una de ellas y luego las de la otra. La audiencia puede ser grabada o puede levantarse un acta escrita. Las posiciones (que ya estaban en el sobre cerrado) se leen al absolvente, quien las contesta por sí o por no, pudiendo agregar las aclaraciones que estime pertinentes.

\cornelltitle{prueba de confesión}
\cornellblocksum{
\cqrow{¿En qué consiste la confesión como prueba?}{En una declaración de quien es parte sobre hechos personales o de su conocimiento, con efecto desfavorable al confesante y favorable a la contraria. Puede ser espontánea (en la contestación de demanda) o provocada (absolución de posiciones en audiencia).
}
\cqrow{¿Cómo se redactan las posiciones?}{En forma afirmativa (``Para que jure como que es cierto que usted se encontraba en\ldots{}''), con un solo hecho, sobre la actuación personal del absolvente y de manera que pueda contestarse por sí o por no. El ponente, al formularlas, reconoce el hecho.
}
\cqrow{¿Qué es la confesión ficta?}{Es la ficción legal del art.~417: si el absolvente, debidamente notificado, no comparece, se niega a contestar o responde evasivamente, el juez al dictar sentencia puede tenerlo por confeso en los hechos de las posiciones. No es automática: el juez evalúa las demás constancias de la causa.
}
\cqrow{¿Qué valor probatorio tiene la confesión ficta frente a la confesión expresa?}{La confesión judicial EXPRESA constituye plena prueba (art.~423), salvo tres excepciones legales. La confesión FICTA tiene valor discutido: para un sector doctrinal es plena prueba si no hay elementos en contra; para otro, requiere corroboración de otros medios; para un tercero, siempre debe cotejarse con el resto de la prueba. La CSJN (fallo del 6/2/2001, LL 2001-B, p.~959) revocó una sentencia que había desestimado la confesión ficta cuando había otros elementos corroborantes.
}
\cqrow{¿Cuándo opera la caducidad en la prueba de confesión?}{Cuando el ponente no comparece a la audiencia, no deja el pliego y el absolvente sí comparece: el ponente pierde el derecho a exigir la absolución de posiciones (art.~410 in fine).
}
}{La prueba de confesión provoca una declaración de las propias partes sobre hechos personales: puede ser espontánea (reconocimiento en la contestación de demanda, que genera plena prueba de inmediato) o provocada a través de la absolución de posiciones, cuya forma afirmativa y específica la distingue de la testimonial. La confesión ficta (art.~417) opera ante la incomparecencia o evasiva del absolvente notificado y, si bien genera una presunción legal desfavorable, su valoración debe conjugarse con el resto de las constancias de la causa, tal como lo señaló la Corte Suprema en el fallo del 6 de febrero de 2001. En conjunto, ambos medios de prueba ilustran la tensión permanente del proceso civil entre el formalismo procedimental y la búsqueda de la verdad objetiva de los hechos.
}

\chapter[La prueba pericial]{La prueba pericial: concepto, sujetos y procedimiento}\label{cap:pericial}

Este capítulo y los dos siguientes abordan tres medios de prueba del proceso civil argentino: la \textbf{prueba pericial}, la \textbf{prueba indiciaria} (con las presunciones legales y las presunciones \textit{iuris hominis}) y el \textbf{reconocimiento judicial}.

Respecto de la prueba pericial se estudian su concepto, los sujetos involucrados (el perito designado de oficio y el consultor técnico de parte), el procedimiento completo desde el ofrecimiento hasta la valoración del dictamen, las obligaciones y derechos del perito, las causales de recusación y remoción, el anticipo de gastos y las consecuencias del incumplimiento. Luego se analiza la prueba indiciaria, es decir, cómo el juez puede formar su convicción a partir de indicios acreditados en el expediente cuando no existe prueba directa. Finalmente se estudia el reconocimiento judicial como la prueba más directa: su concepto, su procedimiento y su vinculación con las facultades del artículo 475 del CPCCN.

\textbf{Normas principales:} arts.\ 459, 163 inc.\ 5, 386, 476, 477 y 479 del CPCCN; art.\ 243 del Código Penal.

\begin{example}{Cómo se conectan los tres medios de prueba}
En un juicio por accidente de tránsito, el juez no sabe si el auto tenía los frenos en mal estado, lo que requiere un perito mecánico. Tampoco hay testigos directos del momento de la colisión, por lo que se necesitan indicios: marcas de frenado, velocidad, comportamiento previo del conductor. Finalmente, el juez visita el lugar del accidente en persona (reconocimiento judicial).

Estos tres instrumentos forman el sistema que permite al juez conocer hechos técnicos o no presenciados y dictar sentencia con convicción fundada.
\end{example}

\section{Finalidad de la prueba en el proceso}

Las pruebas tienen por finalidad \textbf{acreditar} la ocurrencia de un hecho de interés para la litis: un hecho que es \textbf{controvertido}, que debe ser \textbf{conducente} y que fue fijado por el juez en la audiencia preliminar. Quien tiene la carga de acreditar ese hecho debe hacerlo a través de los medios de prueba disponibles.

Se sabe que un hecho existió porque deja un registro, una huella en las cosas o en las personas. Para reproducirlo en el proceso es necesario consultar la fuente y hacerlo a través de los medios correspondientes.

Las pruebas estudiadas hasta este punto (documental, informativa, testimonial y confesional) son medios que tienden a reconstruir hechos para formar la convicción del juez. El juez no puede ver todo ni sabe cómo sucedieron las cosas, por lo que debe fijar los hechos a través de la prueba que las partes ofrezcan y produzcan en la etapa probatoria.

\section{Cuándo es necesaria la prueba pericial}

Cuando para conocer los hechos controvertidos no bastan los conocimientos del común de la gente (ni los del juez, incluso el juez laboral) y hace falta un conocimiento especial en alguna ciencia, arte, industria o actividad técnica especializada, las partes deben valerse de la \textbf{prueba pericial}. Para ello se propone una prueba pericial, a fin de que se realice una peritación.

\subsection{Características de la peritación}

\begin{definition}{Peritación}
Actividad que se desarrolla por encargo judicial, realizada por una persona ajena a las partes y al juez, para hacer conocer al juez determinadas cuestiones técnicas o científicas que escapan a su conocimiento común.
\end{definition}

La peritación no se realiza por iniciativa propia: se ofrece como prueba en los escritos constitutivos del proceso (demanda, contestación, reconvención, contestación de excepciones previas) y es el juez quien la ordena.

\section{Quiénes son los peritos}

Si las profesiones están reglamentadas, son peritos quienes tienen \textbf{título habilitante} para ejercer la profesión. Si la profesión no está reglamentada, se debe buscar a personas que conozcan el tema discutido y tengan amplia experiencia en él, que pueden ser nombradas a falta de profesión reglamentada y de título habilitante.

\subsection{Dónde se ubica a los peritos}

Una vez ofrecida la prueba pericial y admitida por el juez, este designa y sortea un perito dentro de la lista de peritos que llevan las cámaras de cada fuero. Quienes quieren actuar como peritos (por ejemplo ingenieros, arquitectos, contadores, agrimensores, médicos de distintas especialidades) se anotan en las listas de peritos judiciales y, en el momento que corresponda, deben ser sorteados.

\begin{note}
La expresión «desinsaculado» proviene de la práctica antigua en que el sorteo del perito se hacía extrayendo un papel o una bolilla de un saco (\textit{insacular}). Por eso se sigue diciendo que «ha sido desinsaculado el perito».
\end{note}

Generalmente interviene un \textbf{perito único designado de oficio} por el juez, cuando existen razones necesarias para que intervengan personas con conocimientos especiales respecto de determinados hechos.

\begin{definition}{Perito}
Tercero ajeno al proceso, con conocimientos especiales, idóneo y, en su caso, con título habilitante, que esclarece determinadas cuestiones técnicas, científicas, artísticas o industriales.
\end{definition}

El perito es un tercero auxiliar del juez. Tiene el deber de ser \textbf{imparcial} y no puede tener otro interés que esclarecer los hechos, de modo que el juez pueda entenderlos y aplicar la norma que rige. La materia escapa al conocimiento del juez, aunque este deba fijar los hechos.

\section{El consultor técnico (perito de parte)}

La misma complejidad que requiere la intervención del perito para ayudar al juez la enfrenta el abogado que prepara la demanda. Así como el abogado es el asistente técnico en el aspecto jurídico (patrocinio jurídico), la parte puede contar con un asistente en materias científicas o técnicas.

\begin{definition}{Consultor técnico (perito de parte)}
Asistente de la parte en materias científicas o técnicas. Tanto el actor como el demandado tienen derecho a designarlo, y puede designarse también en los escritos constitutivos del proceso.
\end{definition}

\subsection{Diferencias entre perito designado de oficio y consultor técnico}

\begin{table}[H]
\centering
\small
\begin{tabularx}{\textwidth}{
    >{\raggedright\arraybackslash}p{0.19\textwidth}
    >{\raggedright\arraybackslash}X
    >{\raggedright\arraybackslash}X
}
\toprule
\textbf{Criterio} & \textbf{Perito (de oficio)} & \textbf{Consultor técnico} \\
\midrule
Designación & Por el juez, de oficio, por sorteo & Por la parte (actor o demandado) \\
Imparcialidad & Debe ser absolutamente imparcial y objetivo & Puede y debe ser parcial; actúa en favor de su cliente \\
% TODO: contenido incompleto o ambiguo en las notas originales (la tabla dice "con matices"; más adelante se afirma que la aceptación del cargo es facultativa)
Obligatoriedad & Obligado a aceptar el cargo (con matices) & Facultativo; la parte puede removerlo en cualquier momento \\
Dictamen / presentación & Obligado a presentar dictamen en tiempo y forma & No está obligado a presentar nada \\
Recusación & Puede ser recusado con causa (art.\ 17 CPCCN) & No puede ser recusado (es de parte, se asume parcialidad) \\
Honorarios & A cargo de la condena en costas & A cargo de la condena en costas \\
Base del dictamen & Expediente + estudios propios + apoyaturas técnicas & Expediente + asesoría al abogado y a la parte \\
\bottomrule
\end{tabularx}
\end{table}

Una diferencia clave es la imparcialidad. El consultor técnico es un perito de parte: puede ser parcial y, de hecho, debería serlo, porque es contratado por su cliente para que lo asesore y fundamente científico-técnicamente la cuestión de modo que lo favorezca. El perito, en cambio, debe ser absolutamente imparcial y objetivo.

El perito puede basarse en estudios que realiza en instituciones especializadas por fuera del propio expediente judicial. Puede valerse incluso, si la cuestión es compleja, de la ayuda de otros profesionales, pero no puede remitirse en su dictamen únicamente a esas otras cuestiones: el dictamen debe ser hecho por el perito, sin perjuicio de que se valga de otras apoyaturas técnicas o de los exámenes y experimentos necesarios.

\section{Actividad de la contraparte al ofrecer la prueba pericial}

Quien ofrece la prueba pericial debe indicar claramente la \textbf{especialidad} del perito y los \textbf{puntos de pericia}, y puede designar en ese mismo acto al consultor técnico.

Cuando el actor ofrece la prueba pericial, se corre traslado al demandado junto con el traslado de la demanda para que ejerza las facultades que le permite el Código Procesal. Conforme al artículo 468 (mencionado también, en las fuentes originales, como art.~478), la contraparte puede:

% TODO: contenido incompleto o ambiguo en las notas originales (numeración del artículo: 468 o 478)
\begin{enumerate}[label=\alph*)]
    \item Impugnar la procedencia de la prueba pericial por no corresponder conforme al artículo 459, es decir, por no requerirse conocimientos especiales técnicos, científicos, etc.\ para acreditar los hechos que se pretenden a través de los puntos de pericia indicados.
    \item Manifestar desinterés en la pericia, es decir, sostener que se va a abstener de participar en ella.
    \item Impugnar la especialidad del perito propuesto.
    \item Impugnar los puntos de pericia por sostener que están errados, equivocados o que no son procedentes.
    \item Proponer otros puntos de pericia distintos.
    \item Designar su propio consultor técnico.
\end{enumerate}

\begin{formula}{Art. 459 CPCCN --- Procedencia de la prueba pericial}
Procederá la prueba pericial cuando la apreciación de los hechos controvertidos requiere conocimientos especiales en alguna ciencia, arte, industria o actividad técnica especializada.
\end{formula}

\subsection{Costas, traslados y simetría}

% TODO: contenido incompleto o ambiguo en las notas originales (el texto fuente dice "no se hace mérito de la policía"; se interpreta como "pericia")
Si la contraria impugna la procedencia (o manifiesta desinterés) y el juez hace lugar a la pericia, pero en la sentencia esta no resulta un elemento determinante (no se hace mérito de la pericia, no fue necesaria), quien impugnó o manifestó desinterés no queda obligado al pago de las costas de la pericia.

Si la contraparte impugna los puntos de pericia y ofrece otros distintos, se debe correr traslado de vuelta a quien ofreció la prueba. Este traslado va \textbf{por nota}, no por cédula.

Todo lo anterior se aplica en forma simétrica: si el demandado ofrece prueba pericial, se corre traslado al actor para que pueda ejercer la misma actividad. La diferencia es que el traslado al demandado (de la pericial ofrecida por el actor) se corre \textbf{por cédula}, mientras que el traslado al actor (de la pericial ofrecida por el demandado) se corre \textbf{por nota}.

\section{La audiencia preliminar y la designación del perito}

Una vez presentadas todas las impugnaciones y traslados, se llega a la audiencia preliminar, donde el juez debe resolver:

\begin{itemize}
    \item Todo planteo relativo a la especialidad del perito.
    \item Las impugnaciones respecto de los puntos de pericia.
    \item La especialidad del perito, si las partes no se pusieron de acuerdo.
\end{itemize}

En la audiencia preliminar, si las partes se ponen de acuerdo, pueden designar ellas mismas al perito. Sin perjuicio de ello, dado que la calidad del perito se orienta a formar la convicción del juez, este podría cuestionar al perito designado si no es el especialista adecuado o si carece de título habilitante.

El juez tiene también la facultad de agregar otros puntos de pericia o modificar los propuestos, de modo que sean adecuados para acreditar los hechos discutidos.

Si las partes no se ponen de acuerdo, el juez determina la especialidad del perito, fija los puntos de pericia, puede agregar otros si lo considera necesario y puede modificar los existentes. El perito debe ser designado de oficio, sorteado de entre los integrantes de la lista de las cámaras del fuero para la especialidad que resulte.

\section{Recusación del perito}

A partir de la designación del perito en la audiencia preliminar corre un plazo de \textbf{cinco días} para que las partes lo recusen si lo estiman necesario.

Las causales de recusación son las mismas que corresponden a la recusación con causa del juez, contenidas en el artículo 17 del Código Procesal.

\begin{formula}{Art. 17 CPCCN --- Causales de recusación con causa (aplicables al perito)}
Como ejemplos se mencionan: que el perito tenga manifiesto interés en el pleito, ser amigo íntimo o enemigo de alguna de las partes, haber sido denunciado por alguna de las partes, entre otras causales equivalentes a las de recusación judicial.
\end{formula}

Además existen causales específicas de la cuestión pericial:

\begin{itemize}
    \item \textbf{Falta de título habilitante:} el perito carece del título y la profesión está reglamentada.
    \item \textbf{Incompetencia del perito:} tiene un título de profesión en general pero se necesita un especialista en determinada materia.
\end{itemize}

\begin{example}{Incompetencia del perito}
Si se designó un médico clínico y hace falta un médico traumatólogo o un médico psiquiatra; o si se designó un ingeniero cuando en realidad corresponde un agrimensor, como en una pericia sobre inmuebles rurales en la que hay que determinar su superficie.
\end{example}

Si el perito recusado acepta las causales, el juez lo remueve directamente y designa otro perito de oficio. Si el perito rechaza las causales, se forma un \textbf{incidente} que resuelve el juez, y la resolución es irrecurrible en cuanto a si corresponde o no la recusación.

\begin{important}
La notificación de la designación del perito va por nota, por lo que puede sorprender a la parte que deba recusar dentro del plazo por alguna de estas causales. Es necesario estar atento.
\end{important}

\section{Aceptación del cargo, obligaciones y sanciones del perito}

Para intervenir válidamente, el perito no solo debe ser designado por el juez sino que debe presentarse y \textbf{aceptar efectivamente el cargo} ante el Oficial Primero.

La aceptación del cargo es facultativa: no es una carga pública, por lo que el perito puede no aceptarlo. Sin embargo, si reiteradamente y sin justificación un perito no acepta el cargo, puede ser removido de la lista de la cámara que corresponde.

Una vez aceptado el cargo, el perito tiene el deber de:

\begin{itemize}
    \item comparecer todas las veces que sea citado al proceso;
    \item presentar el dictamen en tiempo y forma;
    \item concurrir a las audiencias y dar las explicaciones que correspondan.
\end{itemize}

\begin{formula}{Art. 243 del Código Penal --- Sanción penal al perito}
Establece la sanción para el perito que, legalmente citado, no compareciere o rehusare dar las explicaciones que le sean requeridas: pena de prisión de 15 días a un mes e inhabilitación especial de un mes a un año.
\end{formula}

Además de las sanciones penales, el perito puede ser responsabilizado por los daños y perjuicios que ocasione su incumplimiento, lo que se reclamará por vía incidental.

\section{Anticipo para gastos}

Dentro de los \textbf{tres días} de haber aceptado el cargo, el perito puede solicitar un anticipo para gastos cuando la pericia requiere determinados materiales, insumos o gastos de traslado importantes.

\begin{important}
El anticipo para gastos no debe ser un solapado aumento de honorarios. El perito tiene derecho a recibirlo, pero también el deber de rendir cuenta de los gastos realizados y del destino de los fondos en el dictamen correspondiente, de modo que las partes puedan controlar que los gastos eran tales y fueron destinados a la pericia.
\end{important}

La resolución que fija el anticipo para gastos solamente es susceptible del recurso de \textbf{reposición}; no se puede apelar a la cámara. Una vez notificada y firme, el anticipo debe depositarse dentro de los \textbf{cinco días} de ser notificado por cédula. Si la parte obligada no lo deposita oportunamente, se la tiene por desistida de la prueba pericial y pierde la prueba.

\section{Remoción del perito}

Son causales de remoción del perito:

\begin{itemize}
    \item la renuncia del perito durante el proceso;
    \item que, después de haber aceptado el cargo, rehúse presentarse a dar el dictamen;
    \item que no presente el dictamen oportunamente (luego de haber pedido ampliación del plazo y, en su caso, si no está justificada la demora).
\end{itemize}

En caso de remoción, el perito debe, si corresponde, pagar los daños y es pasible de las sanciones civiles y penales que puedan corresponder.

El plazo para presentar el dictamen, salvo disposición en contrario, es de \textbf{15 días}. Corresponde a la parte que ofreció la prueba pericial realizar las diligencias necesarias para que la prueba no caduque, siendo fundamental notificar al perito de su designación para que se presente en el expediente y acepte el cargo.

\cornelltitle{la prueba pericial, sujetos y procedimiento inicial}
\cornellblocksum{
\cqrow{¿Cuándo procede la prueba pericial?}{Cuando la apreciación de los hechos controvertidos requiere conocimientos especiales en alguna ciencia, arte, industria o actividad técnica especializada que exceden la cultura media del juez (art.\ 459 CPCCN).
}
\cqrow{¿Qué es la peritación?}{Es una actividad desarrollada por encargo judicial, realizada por un tercero ajeno a las partes y al juez, con conocimientos especiales, para hacer conocer al juez cuestiones técnicas o científicas que escapan a su conocimiento común.
}
\cqrow{¿Cómo se designa el perito?}{Por sorteo (desinsaculación) de entre los integrantes de la lista de peritos de la cámara del fuero correspondiente. La especialidad y los puntos de pericia se fijan en la audiencia preliminar.
}
\cqrow{¿Cuál es la diferencia clave entre perito y consultor técnico?}{El perito es designado de oficio, debe ser imparcial y está obligado a presentar el dictamen. El consultor técnico es de parte, puede y debe ser parcial, y su intervención es facultativa; puede ser removido por la parte que lo designó en cualquier momento.
}
\cqrow{¿Cuáles son las causales de recusación del perito?}{Las mismas del art.\ 17 CPCCN para la recusación del juez (interés en el pleito, amistad íntima, enemistad, etc.), más las específicas: falta de título habilitante y falta de la especialidad requerida para los puntos de pericia. El plazo es de cinco días desde la designación.
}
\cqrow{¿Qué pasa si el perito no acepta el cargo reiteradamente?}{Puede ser removido de la lista de la cámara correspondiente, para desalentar la conducta de elegir casos en función de los honorarios esperados.
}
\cqrow{¿Qué consecuencias tiene no pagar el anticipo para gastos?}{Si la parte obligada no deposita el anticipo dentro de los 5 días de notificada por cédula, se la tiene por desistida de la prueba pericial y pierde la prueba (la resolución solo admite recurso de reposición).
}
}{La prueba pericial se ofrece con indicación de especialidad y puntos de pericia, se corre traslado a la contraparte y, en la audiencia preliminar, el juez fija la especialidad y los puntos. El perito, tercero imparcial designado de oficio por sorteo, puede ser recusado en cinco días, debe aceptar el cargo, puede pedir un anticipo para gastos (solo reposición; depósito en cinco días bajo pena de pérdida de la prueba) y puede ser removido por renuncia o incumplimiento. El consultor técnico, en cambio, es de parte y puede ser parcial.
}

\chapter[Práctica de la pericia]{Práctica de la pericia: pasos, dictamen y valoración}\label{cap:practica-pericia}

Establecidos los sujetos de la prueba pericial y su procedimiento inicial, corresponde analizar cómo se realiza la pericia, qué ocurre con el dictamen y de qué modo lo valora el juez.

%==================================================================

\section{Presencia de las partes en la realización de la pericia}

Si las partes quieren presenciar la realización de la prueba pericial, deben hacérselo saber al perito en el expediente. En tal caso, el perito tiene la obligación de informar, con antelación suficiente, el expediente, el lugar, el día y la hora en que llevará a cabo la pericia, de modo que las partes, sobre todo los consultores técnicos, puedan concurrir y hacer las observaciones que consideren pertinentes.

\begin{important}
Si las partes solicitan presenciar la pericia y el perito no informa el día en que realizará su tarea, la pericia es \textbf{nula}, porque se viola el principio de bilateralidad de la audiencia, que tiene raigambre constitucional y no puede ser dejado de lado por ningún sistema que cree el legislador para llevar adelante el proceso.
\end{important}

\section{Los pasos para realizar la prueba pericial}

\subsection{Paso 1: determinación del tema pericial}

El perito debe tener muy en claro qué cuestionario se le pide que conteste y qué cuestiones debe resolver, lo que surge de los puntos de pericia. Puede requerir explicaciones al juez y a las partes respecto de esos puntos e, incluso, puede advertir que exceden su especialidad, en cuyo caso deberá indicar cuál sería el perito de la especialidad correcta.

\subsection{Paso 2: observación del campo de trabajo}

Como experto, el perito debe determinar cuál es su campo de trabajo, qué elementos necesita para llevar a cabo la pericia y dónde se pueden encontrar. Tiene facultad de concurrir a los lugares pertinentes y de recurrir a entidades especializadas para satisfacer las dudas y requerimientos que puedan surgir. Debe recolectar los elementos que luego tendrá que acompañar en cuanto resulten necesarios para el dictamen.

\subsection{Paso 3: elección de los medios}

El perito puede valerse de exámenes de personas, de cosas y de estudios especializados.

\begin{important}
La elección de los medios (la selección de determinados medios importantes para el caso) puede ocasionar después una impugnación o, en su caso, la nulidad de la pericia, si esto pone en riesgo la entidad científica o técnica de la pericia.
\end{important}

\begin{example}{Pericia médica}
El médico legista designado para llevar a cabo la pericia necesita valerse de tomografías, resonancias y estudios complejos (de biología molecular, etc.) que escapan incluso al saber del experto. Se está dando así un paso casi fuera de la prueba pericial e invadiendo la esfera de las pruebas científicas.
\end{example}

El autor Falcón propone establecer la \textbf{prueba científica} como medio probatorio autónomo, porque ciertas cuestiones que encuadrarían dentro de ella requieren estudios especializados, equipos especiales, conjuntos de personas preparadas e instituciones de la materia. Sin embargo, el perito puede integrar su pericia con estos exámenes requeridos a otras instituciones, lo que la convierte en una prueba compleja que excede el marco pericial y traspasa quizás la frontera entre la prueba pericial y la científica.

Esta frontera es de dudosa delimitación, porque el conocimiento avanza paulatinamente: tanto el que tiene el perito respecto de cuestiones que antes no podía alcanzar ni interpretar aun siendo experto, como el de la propia ciencia, que avanza al mismo tiempo en cuestiones más novedosas, complejas y dificultosas que requieren interpretación.

\begin{formula}{Art. 476 CPCCN --- Opinión de instituciones}
Permite que se soliciten opiniones a universidades, instituciones o entidades públicas o privadas respecto de cuestiones técnicas o científicas que resulten complejas. Esto se engloba en lo que puede considerarse prueba científica, que para el autor Falcón debería ser considerada un medio autónomo.
\end{formula}

\subsection{Paso 4: agregación y preservación de los elementos}

El perito no solo debe llevar adelante las operaciones que correspondan sino también preservar los medios de los que se valió, para acompañarlos al expediente.

\begin{example}{Prueba informática}
En una pericia sobre cuestiones electrónicas, el perito debe investigar si en determinada computadora existen ciertos mensajes u otros archivos digitales. Debe preservar esa prueba adecuadamente y rescatarla para asegurar su integridad, trasladarla después y acompañarla como fundamento de sus conclusiones.
% TODO: contenido incompleto o ambiguo en las notas originales (el texto fuente dice "desde el peligro [riesgo] que tiene que ir a investigar")
\end{example}

\subsection{Paso 5: realización de las operaciones técnicas y emisión del dictamen}

El perito debe realizar todas las operaciones técnicas que correspondan para evacuar todos los puntos de pericia y sacar las conclusiones correctas en función de lo que se le pide resolver.

Finalmente debe emitir su \textbf{dictamen}, que deberá ser técnico y científico y valerse de las pruebas necesarias. El dictamen tiene que ser acompañado de las pruebas que sean necesarias.

\section{La prueba documental que acompaña el perito}

Existe un punto delicado: el acompañamiento por el perito de prueba documental que \textquotedblleft descubre\textquotedblright{} puede importar, en muchos casos, una agregación extemporánea de esa prueba al proceso.

La prueba documental debe acompañarse junto con la demanda, la contestación, la reconvención y la contestación de excepciones, salvo que se trate de documentos nuevos o recién conocidos. Los documentos posteriores a esas oportunidades o recién conocidos pueden acompañarse en el proceso hasta el llamamiento de autos para sentencia en primera instancia, sin perjuicio de hacerlo en segunda instancia.

% TODO: contenido incompleto o ambiguo en las notas originales (el pasaje sobre la agregación en segunda instancia -"debe apelarse en la sentencia definitiva y fundarse en forma libre"- no resulta claro)

\begin{important}
El perito no puede acompañar prueba documental supliendo la negligencia de la parte que la tenía a su disposición. Esto rige cuando el perito, al realizar la pericia, concurre a bancos, entidades públicas, etc., y solicita documentación que le será necesaria y que luego acompañará como fundamento de su dictamen al presentarlo en el expediente.
\end{important}

\section{Notificación del dictamen y actividad de las partes}

El perito debe presentar el dictamen \textbf{por escrito, con copias para las partes}. Las partes son notificadas por cédula del dictamen pericial y, a partir de ahí, tienen un plazo de \textbf{cinco días} para realizar distintas actividades.

Dentro de los cinco días de notificadas por cédula, las partes pueden:

\begin{enumerate}[label=\alph*)]
    \item Impugnar la pericia por contener falsedades.
    \item Plantear la nulidad de la prueba pericial (no expresamente contemplada en el código, pero disponible por vía incidental).
    \item Pedir una ampliación de la pericia, dado que el perito no contestó todos los puntos solicitados.
    \item Pedir una nueva pericia por ser inconsistente la anterior y no poder ser corregida por ampliación.
    \item Pedir explicaciones del perito.
\end{enumerate}

La nulidad de la pericia puede darse, por ejemplo, cuando se ha ejercido violencia sobre el perito, cuando el perito perdió sus facultades mentales de modo que la pericia no puede considerarse válida, o cuando tiene groseros errores de fundamentación que ya no admiten corrección a través de un pedido de explicaciones. Si prospera la nulidad, se debe realizar una nueva pericia por otros expertos, no por el mismo perito que ya participó.

\begin{important}
Si no se impugna la pericia por inconsistencias o falsedades dentro de los cinco días de notificado el dictamen por cédula, la impugnación puede hacerse hasta el momento de los alegatos. Si tampoco se hizo en esa oportunidad y el juez, que es quien valora la prueba pericial, la utiliza como elemento de convicción para dictar sentencia, es posible cuestionarla ante la cámara, al expresar agravios en el recurso de apelación contra la sentencia definitiva, por las inconsistencias o falsedades de que adolezca, aunque no se haya hecho al correrse el traslado de la pericia.
\end{important}

\section{El artículo 475 CPCCN: medidas complementarias}

El artículo 475 del Código Procesal permite que las partes ofrezcan, o que el juez ordene de oficio, la ejecución de las siguientes medidas:

\begin{formula}{Art. 475 CPCCN --- Medidas complementarias de la prueba pericial}
\begin{itemize}
    \item \textbf{Inc.\ 1:} planos, relevamientos, tomas fotográficas y filmaciones cinematográficas de objetos, documentos o lugares, con el empleo de medios o instrumentos técnicos.
    \item \textbf{Inc.\ 2:} exámenes científicos necesarios para el mejor esclarecimiento de los hechos controvertidos.
    \item \textbf{Inc.\ 3:} reconstrucción de hechos para comprobar si se han producido o pudieron producirse de una determinada manera.
\end{itemize}
\end{formula}

Las medidas del inciso 1 (planos, fotografías, películas) son en realidad \textbf{prueba documental}, dado que la concepción de prueba documental consiste en la representación del pensamiento humano en algún tipo de objeto. El reconocimiento de hechos tampoco es prueba pericial en sí misma: es otro medio de prueba. Estas facultades del artículo 475 se vinculan con el reconocimiento judicial del artículo 479.

\section{Valoración del dictamen pericial (art.\ 477 CPCCN)}

El dictamen pericial \textbf{no es vinculante} para el juez salvo que la ley disponga lo contrario.

\begin{formula}{Art. 477 CPCCN --- Valoración del dictamen pericial}
El juez valorará el dictamen pericial de acuerdo con las reglas de la sana crítica. Podrá apartarse de las conclusiones del perito cuando existan otras pruebas o razones que justifiquen tal apartamiento.
\end{formula}

El código procesal exige al juez que valore el dictamen en función de las reglas de la sana crítica: la prueba pericial no escapa a la valoración que, para todos los medios de prueba, debe hacerse conforme a esas reglas.

\begin{example}{Juicio de demencia}
La ley dispone lo contrario cuando, en el juicio de demencia, establece que el presunto insano debe ser examinado por tres médicos psiquiatras y que, si lo encuentran en sano juicio, el juez no podría declararlo demente. Esta pericia, por la ley y por su resultado, es absolutamente vinculante para el juez.

No sucede lo mismo a la inversa: aun cuando los médicos psiquiatras consideren demente a la persona, ello no vincula al juez. Fundado en otras razones que surjan del expediente y en lo que valore por la regla de la sana crítica, puede apartarse del dictamen y no declarar demente a esa persona.
\end{example}

Para valorar el dictamen, el juez debe ponderar:

\begin{itemize}
    \item cuán completo sea el dictamen pericial;
    \item cuán competentes e idóneos sean los expertos que lo elaboraron;
    \item cuán coherente y lógico sea el dictamen en sí mismo;
    \item cuán fundado esté en cuestiones teóricas, científicas y en estudios;
    \item cuán completo sea en cuanto a que se acompañaron todos los elementos utilizados al elaborar la pericia.
\end{itemize}

\section{El perito árbitro: cuando el dictamen sí es vinculante}

En el proceso arbitral (título tercero del código procesal), lo que deciden los peritos árbitros en el ámbito de su competencia respecto de cuestiones de hecho es \textbf{vinculante} para el juez.

El artículo 516 del código procesal permite que, en la etapa de ejecución de sentencia, cuando existen cuestiones complicadas o de difícil dilucidación (por ejemplo, liquidaciones de sociedades o de sociedad conyugal), se tome como puerta de salida del proceso, como método alternativo de resolución de conflictos, dejar esas cuestiones a lo que resuelvan estos peritos árbitros, que utilizan la misma regla que los amigables componedores.

El perito árbitro vincula al juez en cuanto a la fijación del hecho a través del laudo emitido en el ámbito de la cuestión que se le dejó para resolver, pero no lo limita en cuanto a la aplicación del derecho en ese proceso.

\cornelltitle{práctica de la pericia, dictamen y valoración}
\cornellblocksum{
\cqrow{¿Qué es el principio de bilateralidad en la pericia?}{Si las partes solicitaron presenciar la pericia, el perito tiene la obligación de informarles con antelación suficiente el lugar, día y hora. Si no lo hace, la pericia es nula, porque se viola el principio de bilateralidad de audiencia, de raigambre constitucional.
}
\cqrow{¿Cuáles son los pasos del perito para realizar la pericia?}{1) Determinación del tema pericial (interpretar los puntos de pericia). 2) Observación del campo de trabajo y recolección de elementos. 3) Elección de los medios (con cautela, pues puede dar lugar a impugnación o nulidad). 4) Agregación y preservación de los elementos. 5) Realización de las operaciones técnicas y emisión del dictamen.
}
\cqrow{¿Qué puede hacer la parte en los 5 días posteriores al dictamen?}{Impugnar por falsedades o inconsistencias, plantear la nulidad, pedir ampliación, pedir nueva pericia o pedir explicaciones al perito. Vencido ese plazo precluyen algunas facultades (la impugnación puede hacerse hasta los alegatos).
}
\cqrow{¿Qué medidas complementarias prevé el art.\ 475 CPCCN?}{Planos, relevamientos, fotografías y filmaciones (inc.\ 1, que constituyen en realidad prueba documental); exámenes científicos (inc.\ 2); y reconstrucción de hechos (inc.\ 3).
}
\cqrow{¿Es vinculante el dictamen pericial para el juez?}{No, salvo que la ley disponga lo contrario (por ejemplo, en el juicio de demencia, si los tres médicos psiquiatras dictaminan que la persona está en sano juicio, el juez no puede declararla demente). El juez valora el dictamen según las reglas de la sana crítica y puede apartarse fundadamente.
}
\cqrow{¿Qué pondera el juez al valorar el dictamen?}{Su completitud, la competencia e idoneidad de los expertos, su coherencia lógica, su fundamento teórico, científico y en estudios, y que se hayan acompañado todos los elementos utilizados.
}
\cqrow{¿Cuándo sí es vinculante el dictamen?}{Cuando interviene un perito árbitro (art.\ 516 CPCCN, en etapa de ejecución de sentencia): su fijación de hechos es vinculante para el juez, aunque este conserva la aplicación del derecho.
}
}{La pericia se realiza en cinco pasos: determinación del tema, observación del campo de trabajo, elección de los medios, preservación de los elementos y emisión del dictamen, con aviso previo a las partes que pidieron presenciarla (bajo pena de nulidad por violación de la bilateralidad). Presentado el dictamen y notificado por cédula, las partes tienen cinco días para impugnar, plantear nulidad, pedir ampliación, nueva pericia o explicaciones. El juez valora el dictamen según la sana crítica (art.\ 477) y puede apartarse fundadamente; solo excepcionalmente la ley lo torna vinculante (juicio de demencia), y el perito árbitro sí vincula al juez en cuestiones de hecho.
}

\chapter[Prueba indiciaria y reconocimiento judicial]{Prueba indiciaria y reconocimiento judicial}\label{cap:indiciaria}

%==================================================================

\section{La prueba indiciaria: concepto y necesidad}

La prueba pericial es un medio complejo no solo por la materia que requiere la intervención del perito, sino por la práctica de la pericia, que tiene distintos momentos y puede requerir la participación de terceras personas o de determinadas instituciones.

Se plantean entonces las siguientes cuestiones: ¿en qué condiciones puede el juez apartarse de lo dispuesto por estos expertos? ¿Qué eficacia probatoria tienen estos dictámenes y los razonamientos que hacen los expertos para ilustrar el conocimiento?

La prueba indiciaria es necesaria precisamente cuando falta la demostración acabada, a través de los otros medios probatorios (documental, informativa, pericial, reconocimiento judicial), de la certeza respecto de la existencia del hecho de importancia para la litis. Si los hechos están acreditados categóricamente por los otros medios establecidos en el código, no es necesario recurrir a estas pruebas llamadas indiciarias.

\section{Presunciones: concepto y tipos}

Al hablar de presunciones existe confusión con los términos que utilizan la ley o la jurisprudencia, que hablan de indicios o los asocian a la conclusión que se extrae de ellos.

Sobre la naturaleza de la prueba indiciaria hay dos posturas. Para algunos autores no es un medio de prueba sino que se limita a ser un razonamiento, una forma de razonamiento jurídico. Otros sostienen que, aunque considerada de forma indirecta y a partir de la utilización de principios lógicos, atiende a la demostración de un hecho controvertido de interés para el proceso y, de esta forma, es un medio de prueba.

\begin{table}[H]
\centering
\small
\begin{tabularx}{\textwidth}{
    >{\raggedright\arraybackslash}X
    >{\raggedright\arraybackslash}X
    >{\raggedright\arraybackslash}X
}
\toprule
\textbf{Tipo de presunción} & \textbf{Quién la establece} & \textbf{Admite prueba en contrario} \\
\midrule
Legal -- \textit{iuris et de iure} (de pleno derecho) & El legislador (la ley) & No: es absoluta \\
Legal -- \textit{iuris tantum} (relativa) & El legislador (la ley) & Sí: puede ser desvirtuada \\
Judicial -- \textit{hominis} & El juez & Depende de la valoración del juez \\
\bottomrule
\end{tabularx}
\end{table}

% TODO: contenido incompleto o ambiguo en las notas originales (el texto fuente dice "como dice la solar, 'tiene por cierto'")
En la presunción legal, el legislador le asigna peso a la prueba: acreditados determinados extremos (los hechos que tienen que actuar), el juez \textit{tiene por cierto} el hecho y extrae su consecuencia. En las presunciones \textit{hominis}, en cambio, es el juez quien asigna peso a los indicios, a diferencia de la presunción legal, en la que el peso lo asigna el legislador.

\section[Presunciones judiciales (hominis)]{Las presunciones judiciales (hominis): requisitos del art.\ 163 inc.\ 5}

Las presunciones judiciales (\textit{hominis}) no están establecidas en el código procesal como un medio particular de prueba en el capítulo correspondiente, sino en el artículo 163, referido a la sentencia, particularmente en su inciso 5.

\begin{formula}{Art. 163 inc.\ 5 CPCCN --- Presunciones judiciales}
Las presunciones no establecidas por ley constituirán prueba cuando se funden en hechos reales y probados y cuando, por su número, precisión, gravedad y concordancia, produzcan convicción según las reglas de la sana crítica, de acuerdo con la naturaleza del juicio.
\end{formula}

Los requisitos de los indicios son los siguientes:

\begin{itemize}
    \item \textbf{Hechos reales y probados:} los indicios deben estar acreditados en el expediente. El indicio no puede ser una conjetura: tiene que ser un hecho concreto, real y probado, que por sí solo no basta para tener por acreditado el otro hecho presumido, pero que sumado a otros indicios sí puede lograrlo.
    \item \textbf{Numerosos:} debe existir pluralidad de indicios.
    \item \textbf{Precisos:} deben tener una interpretación en un solo sentido, sin admitir otra.
    \item \textbf{Graves:} tienen peso por sí mismos para determinar, con cierto grado de realidad, la dirección hacia determinado hecho.
    \item \textbf{Concordantes:} son convergentes, en el sentido de que todos apuntan hacia el mismo lado y forman una cadena sin lagunas y sin lugar a dudas respecto de lo que se puede concluir a partir del conjunto.
\end{itemize}

\section{La conducta procesal de las partes como indicio}

El artículo 163 dispone también que la conducta de las partes puede asumir la entidad de indicio. Por ejemplo: una conducta omisiva, la conducta de la parte que destruye documentación, la parte que no contesta la demanda o la que no actúa con buena fe a lo largo del proceso. Esa conducta, sumada a otros elementos de prueba, puede ser un elemento de convicción que complemente las conclusiones del proceso.

\begin{important}
El juez no podría valerse únicamente de la conducta de las partes para condenar a una de ellas o hacer lugar a la pretensión de la contraria. Pero, reunidos determinados elementos probatorios y verificada una conducta desleal, omisiva o inequitativa de la contraria, esa conducta puede ser un indicio del cual el juez se valga para fallar en contra de quien la asume y perjudica a la contraparte.
\end{important}

Asimismo, al dictar sentencia el juez puede considerar la \textbf{conducta desplegada por las partes durante el proceso}. Si una parte está en condiciones de aportar determinado medio de prueba y omite hacerlo sin explicar por qué, lo más razonable es pensar que ese medio de prueba le era desfavorable. No puede premiarse una conducta que obstaculiza la finalidad del proceso, que es averiguar cómo sucedieron los hechos para determinar si corresponde condenar o absolver al demandado.

Las reglas de la sana crítica, a las que remite el artículo 163 inciso 5 del CPCCN para valorar las presunciones, se desarrollan en el capítulo~\ref{cap:valoracion}.

\section{El reconocimiento judicial (art.\ 479 CPCCN)}

\begin{definition}{Reconocimiento judicial}
Prueba más directa y más fiable, en la que el juez se apersona y percibe directamente los hechos, sin que intermedie para su representación ninguna cosa ni persona. No interviene un testigo que declare lo que percibió por sus sentidos ni un dictamen pericial que saque conclusiones: es el propio juez quien toma conocimiento directo en el lugar de los hechos.
\end{definition}

\begin{formula}{Art. 479 CPCCN --- Reconocimiento judicial}
El juez podrá ordenar, de oficio o a petición de parte, el reconocimiento judicial de lugares o de cosas. Se interpreta que también comprende personas, como en los procesos de restricción de la capacidad, donde el juez tiene entrevista con el presunto incapaz. Si el conocimiento excede el propio del juez, puede valerse de un perito para que interprete y señale las distintas circunstancias en juego.
\end{formula}

Es una prueba \textbf{privativa} del juez. Esto no quiere decir que no pueda ser ofrecida por una parte (puede ser ofrecida por una parte o bien ordenada de oficio), sino que el juez es quien puede decidir la utilidad que le significa esa prueba, en función de lo que podrá recabar y del conocimiento que necesita para tener certeza de cómo sucedieron los hechos de interés para el proceso.

En el acto del reconocimiento judicial, el juez puede hacer comparecer a quienes considere pertinentes, puede pedir la concurrencia de testigos y puede valerse también de las medidas previstas en el artículo 475 (planos, croquis, fotografías, filmaciones, reconstrucción de hechos) para aprehender adecuadamente los hechos que deberá fijar después en el proceso judicial.

\section{Forma de la diligencia del reconocimiento judicial}

El juez debe comparecer personalmente, con las personas del tribunal que considere pertinentes. Se debe levantar un \textbf{acta} dejando constancia de las distintas circunstancias percibidas.

\begin{important}
La narración que haga el juez de sus percepciones en el acta debe ser lo más objetiva posible, porque la etapa probatoria no es el momento de valorar la prueba y podría incurrir en prejuzgamiento. Debe dejar constancia objetiva, además, de las aclaraciones y observaciones de las partes, de los consultores técnicos, de los peritos, de las declaraciones eventuales de los testigos y de todas las dificultades que se presenten.
\end{important}
% TODO: contenido incompleto o ambiguo en las notas originales (el texto fuente alude a "autores técnicos [consultores técnicos]" y a "añadir presentadas por el juez" en la enumeración del contenido del acta)

\cornelltitle{prueba indiciaria y reconocimiento judicial}
\cornellblocksum{
\cqrow{¿Qué son las presunciones hominis?}{Son presunciones judiciales: el juez infiere un hecho no acreditado directamente a partir de indicios que son hechos reales, probados en el expediente, numerosos, precisos, graves y concordantes (art.\ 163 inc.\ 5 CPCCN).
}
\cqrow{¿Cuál es la diferencia entre presunción \textit{iuris et de iure} e \textit{iuris tantum}?}{La \textit{iuris et de iure} (de pleno derecho) no admite prueba en contrario: la conclusión es inexorable. La \textit{iuris tantum} (relativa) admite prueba en contrario: la presunción cede si se prueba lo contrario.
}
\cqrow{¿Qué son las reglas de la sana crítica?}{Un razonamiento libre de error y de vicio, sistematizado por Falcón en ocho reglas, con el que el juez determina cuál de las posiciones de las partes es la correcta y puede demostrar, de manera controlable por las partes, cómo formó su convicción (art.\ 386 CPCCN).
}
\cqrow{¿Qué es el reconocimiento judicial?}{La prueba más directa y fiable: el juez percibe personalmente y sin intermediarios los hechos en el lugar donde ocurrieron (lugares, cosas y personas). Se labra un acta objetiva. El juez puede valerse de peritos para interpretar lo observado y de las facultades del art.\ 475 CPCCN.
}
\cqrow{¿Por qué el acta del reconocimiento judicial debe ser objetiva?}{Porque en la etapa probatoria el juez no debe valorar todavía la prueba, ya que incurriría en prejuzgamiento. El acta debe dejar constancia objetiva de lo percibido y de las aclaraciones de las partes, consultores técnicos, peritos y testigos.
}
}{Cuando no hay prueba directa, el juez puede inferir un hecho controvertido a partir de indicios reales, probados, numerosos, precisos, graves y concordantes (presunciones \textit{hominis}, art.\ 163 inc.\ 5), a diferencia de las presunciones legales (\textit{iuris et de iure}, que no admiten prueba en contrario, e \textit{iuris tantum}, que sí la admiten). La convicción se forma según las reglas de la sana crítica (art.\ 386). El reconocimiento judicial (art.\ 479) es la prueba más directa y fiable: el juez percibe lugares, cosas y personas sin intermediarios, puede valerse de peritos y de las medidas del art.\ 475. Todo el andamiaje probatorio sirve a un único fin: formar la convicción del juez para que pueda dictar una sentencia fundada y justa.
}

\chapter{Los alegatos}\label{cap:alegatos}

Concluida la producción de la prueba, la etapa final del proceso civil ordinario atraviesa tres momentos: los alegatos, la etapa conclusional (``autos para sentencia'') y la valoración de la prueba por el juez al dictar sentencia. Este capítulo se ocupa del primero; los dos siguientes, de los restantes.

Los \textbf{alegatos} son el acto procesal escrito (o, en el futuro, oral) mediante el cual las partes analizan la prueba producida y procuran formar la convicción del juez antes del dictado de la sentencia. La \textbf{etapa conclusional} es el momento en que el juez cierra el debate, con sus distintas vías de acceso, sus efectos y la regulación de la morosidad judicial. Los \textbf{sistemas de valoración de la prueba} explican cómo aprecia el juez la prueba al dictar sentencia, con especial atención al sistema de la sana crítica (art.~386 CPCCN) y sus reglas según Falcón, ilustrado con ejemplos de sentencias reales.

\begin{note}
\textbf{Analogía para comprender el recorrido.} Puede imaginarse un partido de fútbol: la etapa probatoria es el partido en sí (los jugadores ---las partes--- muestran sus jugadas ---las pruebas---). Los alegatos son el análisis posterior al partido: cada equipo explica al árbitro ---el juez--- por qué ganó. Luego el árbitro se retira (autos para sentencia) y delibera utilizando las reglas del reglamento ---la sana crítica--- para dictar el resultado final: la sentencia.
\end{note}

\section{Concepto de alegato}

\begin{definition}{Alegato}
Acto procesal en virtud del cual se concluye con la etapa probatoria y mediante el cual las partes se expresan respecto del mérito de la prueba que se produjo en el expediente.
\end{definition}

En este acto procesal, cada parte analiza las pruebas producidas para arribar a una conclusión y trata de formar la convicción del juez:

\begin{itemize}
    \item Si es el \textbf{actor}: para demostrar por qué es viable hacer lugar a su pretensión.
    \item Si es el \textbf{demandado}: para demostrar por qué debe rechazarse la demanda y la pretensión invocada por el actor.
\end{itemize}

Esto muestra que todas las etapas del proceso están vinculadas entre sí, más allá de que rija un sistema de \textbf{preclusión}, en virtud del cual se avanza por etapas y no pueden realizarse los actos procesales que no se ejercieron en la etapa para la cual estaban destinados.

\section{Naturaleza jurídica: acto escrito y facultativo}

El alegato presenta dos características fundamentales:

\begin{enumerate}[label=\Alph*)]
    \item \textbf{Es escrito.} El Código Procesal Civil y Comercial de la Nación (CPCCN) es fundamentalmente un código escrito. La mayoría de los actos son escritos, con algunas implementaciones de oralidad en las audiencias (preliminar, confesional, testimonial) y eventualmente en audiencias conciliatorias (art.~36 CPCCN). Por eso, el alegato es un acto procesal escrito.
    \item \textbf{Es facultativo.} No es una carga ni un deber, y no trae aparejada sanción ni consecuencia alguna para la parte que no lo presente.
\end{enumerate}

Debe recordarse la distinción entre ambas categorías: el \textbf{deber} es aquel que trae aparejada una sanción, mientras que la \textbf{carga} ---imperativo del propio interés--- es aquella cuya inobservancia, si bien no trae aparejada una sanción, puede traer aparejada una consecuencia (por ejemplo, quien no contesta la demanda puede ser declarado rebelde). En el caso de los alegatos no existe consecuencia alguna.

\section{Utilidad del alegato}

Aunque su presentación sea facultativa, el alegato tiene una doble utilidad.

\subsection{Formar la convicción del juez}

La finalidad inmediata es demostrarle al juez que lo invocado en los escritos introductorios quedó acreditado de determinado modo.

\subsection{Base para la expresión de agravios en apelación}

Cuando se interpone un recurso de apelación contra la sentencia de primera instancia, debe fundarse mediante un escrito llamado \textbf{expresión de agravios}, que consiste en una crítica concreta y razonada de la sentencia impugnada.

Si la parte ya realizó un alegato en el que efectuó su propia valoración de la prueba a la luz de los hechos y arribó a su conclusión, ese alegato puede servirle de base para atacar la sentencia en la que considera que se incurrió en un error, ya sea porque se fijaron mal los hechos o porque se valoró mal la prueba.

\section{Procesos en los que el alegato está excluido}

El legislador excluyó la facultad de alegar en los siguientes procesos, por razones de lógica vinculadas a su extensión o características:

\begin{table}[H]
\centering
\begin{tabularx}{\textwidth}{>{\raggedright\arraybackslash}p{4.2cm} >{\raggedright\arraybackslash}X}
\toprule
\textbf{Proceso} & \textbf{Razón de la exclusión} \\
\midrule
Proceso sumarísimo & Es mucho menos extenso que el ordinario; no tiene sentido alegar. \\
Juicio ejecutivo & No hay valoración de prueba general; la única prueba se vincula a las excepciones. \\
Proceso sucesorio & Es un proceso voluntario; no corresponde el alegato. \\
Proceso voluntario en general & Por sus particularidades, el legislador lo excluyó expresamente. \\
\bottomrule
\end{tabularx}
\end{table}

\section{Regulación legal: artículo 482 CPCCN}

El alegato se encuentra regulado en el artículo 482 del CPCCN, denominado ``Agregación de pruebas -- Alegatos''. Producida la prueba, el expediente se pone a disposición de las partes para alegar. El artículo regula la agregación de los cuadernos de prueba al expediente principal y establece los plazos para retirar el expediente en préstamo y para presentar los alegatos.

Actualmente el sistema de cuadernos de prueba prácticamente no existe debido a la digitalización del expediente: toda la prueba se produce dentro del expediente digital.

\subsection{Procedimiento para llegar a la etapa de alegatos}

Una vez producida toda la prueba (o decretada la negligencia o caducidad de alguna), cualquiera de las partes puede peticionar la clausura de la etapa probatoria y que se pongan los autos para alegar. El juez dicta entonces una resolución cuyo texto típico es:

\begin{quote}
``Buenos Aires, 10 de junio de 2020. Siendo exacto lo expuesto y no existiendo medidas pendientes de producción, decreto clausurada la etapa probatoria y pongo los autos en Secretaría para alegar en los términos del artículo 482 del Código Procesal. El orden establecido para retirar en préstamo el expediente será el siguiente: primero la actora, luego el demandado, luego la citada en garantía. Notifíquese.''
\end{quote}

\subsection{Plazos: individual y común}

El artículo 482 establece dos plazos distintos:

\begin{table}[H]
\centering
\begin{tabularx}{\textwidth}{>{\raggedright\arraybackslash}p{2.3cm} >{\raggedright\arraybackslash}p{2.6cm} >{\raggedright\arraybackslash}p{2.6cm} >{\raggedright\arraybackslash}X}
\toprule
\textbf{Tipo} & \textbf{Finalidad} & \textbf{Duración} & \textbf{Particularidades} \\
\midrule
Plazo individual & Retirar en préstamo el expediente & 6 días por parte & Comienza a correr desde que la resolución queda firme y notificada. El orden lo fija el juez: actor $\rightarrow$ demandado $\rightarrow$ citada en garantía. \\
\addlinespace
Plazo común & Presentar el escrito del alegato & 6 días $\times$ n.º de partes & Vence el mismo día para todos. Se computa multiplicando 6 días por el número de partes (si actúan con la misma representación, cuentan como una). \\
\bottomrule
\end{tabularx}
\end{table}

\begin{formula}{Plazo común para alegar}
\[
\text{Plazo común} = 6\ \text{días} \times \text{n.º de partes}
\]
\end{formula}

\begin{example}{Cómputo del plazo común}
Actor y demandado: 2 partes $\times$ 6 días $=$ 12 días de plazo común para alegar. Si hay 3 partes: $3 \times 6 = 18$ días. Si el demandado y la citada en garantía tienen el mismo abogado, cuentan como una sola parte.
\end{example}

\subsection{Consecuencia por no devolver el expediente en término}

Si la parte que retiró el expediente en préstamo no lo devuelve dentro de los 6 días, sin necesidad de intimación previa, se tiene por no presentado el alegato.

\begin{important}
Para la devolución del expediente no rige el plazo de gracia de las dos primeras horas, porque no se trata de presentar un escrito sino de devolver el expediente. Por ello es fundamental que el expediente sea devuelto dentro de los seis días, para no quedar sujeto a la interpretación del tribunal.
\end{important}

\subsection{Cómputo del plazo: firmeza de la resolución}

Los plazos comienzan a correr luego de que la resolución de clausura quede firme (notificada y no impugnada). El plazo varía según quién la haya firmado:

\begin{itemize}
    \item Si la firma el \textbf{juez}: la resolución queda firme a los \textbf{5 días} de notificada.
    \item Si la firma el \textbf{secretario o prosecretario}: el mecanismo de impugnación es el del art.~38 ter CPCCN, y el plazo es de \textbf{3 días}.
\end{itemize}

Recién una vez firme la resolución corren simultáneamente el plazo individual para retirar el expediente y el plazo común para alegar.

\section{Contenido estructural del alegato}

No existe en la ley un artículo que regule expresamente el contenido del alegato (a diferencia de la demanda ---art.~330--- o la contestación ---art.~356---). No obstante, se recomienda el libro \textit{Cómo se hace un alegato}, de Rojas y Falcón (Ed. Abeledo Perrot), que enseña de manera teórico-práctica la estructura correcta.

La estructura recomendada es:

\begin{enumerate}
    \item \textbf{Suma o sumario} (``Presenta alegato''), encabezado y datos del expediente.
    \item \textbf{Descripción y relato de los hechos:} cuáles invocó el actor en la demanda, cuál versión dio el demandado al contestar y qué hechos quedaron controvertidos (por ejemplo, cuál de los rodados violó la luz del semáforo).
    \item \textbf{Análisis estático de las pruebas:} se analiza una por una la prueba producida, marcando lo favorable y atacando la prueba de la contraria (por ejemplo, cuestionando la idoneidad de testigos que están dentro de las generales de la ley por ser amigos o empleados de la otra parte).
    \item \textbf{Análisis dinámico:} síntesis de qué hechos quedaron acreditados con qué prueba, para formar convicción en el juez.
    \item \textbf{Conclusión y firma.}
\end{enumerate}

En cuanto al contenido de derecho, existe una corriente \textbf{restrictiva} (solo hechos y prueba) y una corriente \textbf{amplia} (puede incluir argumentaciones jurídicas), dado el principio \textit{iura novit curia}, en virtud del cual el juez conoce el derecho y puede apartarse del invocado por las partes.

\section{El alegato y la implementación de la oralidad}

Existe un anteproyecto de nuevo Código Procesal en el Congreso que regula el proceso ordinario como un \textbf{proceso por audiencias}, con mayor oralidad:

\begin{enumerate}
    \item \textbf{Primera etapa escrita:} demanda y contestación.
    \item \textbf{Audiencia preliminar:} el juez ordena la producción de la prueba escrita (pericial, informes) y fija la audiencia de vista de causa.
    \item \textbf{Audiencia de vista de causa:} se produce toda la prueba no escrita (testigos, partes, peritos que explican su dictamen). Una vez terminada la prueba oral, cada parte tiene \textbf{20 minutos} para alegar oralmente.
\end{enumerate}

El alegato oral requiere otra preparación de los operadores jurídicos: técnicas de litigación oral y un estudio mucho más detenido del expediente antes de la audiencia, ya que habrá que tener analizado todo lo ocurrido hasta ese momento y, además, hablar en el alegato de la prueba producida en esa misma audiencia.

La finalidad del alegato es la misma en ambos sistemas: comunicarse con el juez, demostrar cómo la prueba acredita los hechos y persuadirlo para que, al dictar sentencia, haga viable la pretensión. Lo que cambia es la forma (de escrita pasa a oral), lo que permite una mayor inmediación con la jurisdicción.

\cornelltitle{los alegatos}
\cornellblock{
\cqrow{¿Qué es el alegato y cuál es su función principal?}{Es el acto procesal (escrito u oral) mediante el cual las partes analizan la prueba producida en la etapa probatoria y procuran formar la convicción del juez. Tiene dos utilidades centrales: (1) demostrar al juez que los hechos invocados quedaron acreditados, para obtener una sentencia favorable; y (2) sentar las bases para la expresión de agravios en el recurso de apelación, permitiendo atacar los errores en los que pueda incurrir el juez al valorar la prueba.
}
\cqrow{¿Es obligatorio presentar el alegato?}{No. Es un acto facultativo: no constituye ni una carga ni un deber. Quien no lo presenta no sufre consecuencia alguna (no hay rebeldía ni sanción). Su omisión solo perjudica en la medida en que el juez carezca de los argumentos de esa parte al momento de valorar la prueba.
}
\cqrow{¿En qué procesos no se puede alegar?}{El alegato está excluido en: (a) proceso sumarísimo (por su menor extensión); (b) juicio ejecutivo (no hay prueba general que valorar); (c) proceso sucesorio; (d) procesos voluntarios. El legislador los excluye por razones lógicas vinculadas a las particularidades de cada uno.
}
\cqrow{Art.~482 CPCCN: ¿cuáles son los dos plazos y cómo se computan?}{
}
\cqrow{¿Cuándo comienzan a correr los plazos del art.~482?}{Luego de que la resolución de clausura quede firme: si la firmó el juez, a los 5 días de la notificación; si la firmó el secretario o prosecretario, a los 3 días (art.~38 ter CPCCN). Recién una vez firme la resolución, corren simultáneamente el plazo individual para retirar el expediente y el plazo común para alegar.
}
\cqrow{¿Cuál es la estructura recomendada de un alegato?}{1.º Suma o encabezado. 2.º Relato de los hechos invocados y de cuáles quedaron controvertidos. 3.º Análisis estático de cada prueba producida (propia y de la contraria), con posibilidad de atacar la idoneidad de testigos. 4.º Análisis dinámico: qué hechos quedaron acreditados con qué prueba. 5.º Conclusión y firma.
}
}
\medskip
\cornellblocksum{
\cqrow{¿Qué implica la oralidad en los alegatos según el anteproyecto?}{En el proceso por audiencias propuesto, luego de la audiencia de vista de causa (donde se produce toda la prueba oral), cada parte tiene 20 minutos para alegar verbalmente. Esto requiere otra preparación: dominio de técnicas de litigación oral y análisis detenido del expediente antes de la audiencia, ya que deberá mencionarse también la prueba producida ese mismo día.
}
}{el alegato ---acto procesal escrito y facultativo regulado en el art.~482 CPCCN--- cierra la etapa probatoria y le permite a cada parte presentar al juez su análisis de la prueba producida, con la doble utilidad de persuadirlo en el caso concreto y de sentar las bases para una eventual expresión de agravios si la sentencia resulta desfavorable. Sus plazos son individuales para retirar el expediente en préstamo (6 días por parte, en el orden fijado por el juez) y comunes para presentar el escrito (6 días por el número de partes), y comienzan a correr luego de que quede firme la resolución de clausura de la etapa probatoria.
}

\chapter[Etapa conclusional]{La etapa conclusional: autos para sentencia}\label{cap:conclusional}

%==============================================================

\section{Concepto}

Así como el proceso se inicia con la promoción de la demanda, en algún momento debe finalizar. En líneas generales, finaliza con el dictado de la sentencia en la \textbf{etapa conclusional} del proceso.

Este es el momento de máxima expresión de la jurisdicción: el juez crea la \textbf{norma individual} para el caso, que no es otra cosa que la sentencia, en función de los hechos que tuvo por ciertos y de la aplicación de las normas abstractas al caso concreto.

No obstante, puede suceder que no se llegue a esta etapa si el proceso termina antes por \textbf{modos anormales}: acuerdo conciliatorio, transacción, desistimiento del proceso por el actor, allanamiento del demandado o caducidad de instancia (art.~310 CPCCN). Estos modos anormales se tratan por separado.

\section{Vías de acceso a la etapa conclusional}

No siempre se llega a la etapa conclusional por el mismo camino. Hay distintas vías.

\subsection{Vía normal: luego de la etapa probatoria y los alegatos}

Producida la prueba, clausurada la etapa probatoria y vencido el plazo para alegar (hayan presentado o no los alegatos las partes), se entra directamente en la etapa conclusional.

\subsection{Cuestión de puro derecho (art.~359, primer párrafo, CPCCN)}

El primer párrafo del art.~359 CPCCN establece que, contestado el traslado de la demanda o la reconvención, vencidos los plazos para hacerlo y resueltas las excepciones, si la cuestión pudiera ser resuelta como de puro derecho así se decidirá, y firme que se encuentre la providencia se llamará autos para sentencia.

Esto ocurre cuando no hay hechos controvertidos, sino que el debate se refiere únicamente a cuáles son las normas aplicables al caso. En ese supuesto, finalizada la etapa introductoria y antes de la audiencia preliminar, el juez declara que la cuestión es de puro derecho y llama autos para sentencia.

\subsection{Supuestos en la audiencia preliminar (art.~360 inc.~6 y arts.~361--362 CPCCN)}

En la audiencia preliminar también puede resolverse la cuestión como de puro derecho en los siguientes casos:

\begin{itemize}
    \item Las partes reconocen recíprocamente, en el acto de la audiencia, los hechos que antes eran controvertidos según los escritos de demanda y contestación, con lo que ya no hay hechos controvertidos que requieran prueba.
    \item Una de las partes se opuso a la apertura a prueba y el juez hace lugar a esa oposición.
    \item Las partes desisten de todas las pruebas ofrecidas (por ejemplo, porque entienden que con la prueba documental acompañada es suficiente).
    \item Ya se realizó prueba anticipada (art.~326 CPCCN) y el juez entiende que con esa prueba es suficiente.
\end{itemize}

\section{El auto para sentencia}

\begin{definition}{Auto para sentencia}
Resolución que da inicio a la etapa conclusional. Puede dictarla el juez de oficio o a petición de parte, mediante un escrito de dos líneas solicitando que el expediente pase a dictado de sentencia.
\end{definition}

Su texto típico es:

\begin{quote}
``Buenos Aires, desde junio de 2020. Autos para sentencia.''
\end{quote}
% TODO: contenido incompleto o ambiguo en las notas originales (la fórmula "desde junio de 2020" figura así en el apunte).

\subsection{Efectos del auto para sentencia}

\begin{enumerate}
    \item \textbf{Cierre del debate:} las partes ya no pueden realizar ningún acto procesal, presentar escritos ni introducir nuevos documentos. Toda la actividad queda clausurada.
    \item \textbf{Límite para documentos nuevos:} se vincula con el art.~335 CPCCN (documentos de fecha posterior). Solo podían acompañarse hasta que el expediente tuviera auto para sentencia. Una segunda oportunidad se presenta ante la Cámara (art.~260 CPCCN), cuando el expediente se eleva para resolver el recurso de apelación.
    \item \textbf{Improcedencia de la caducidad de instancia:} una vez dictado el auto para sentencia, ya no corre el plazo de caducidad de instancia (art.~310 CPCCN), porque las partes no pueden realizar ninguna actividad procesal.
\end{enumerate}

El art.~313 inc.~4 CPCCN establece que es improcedente la caducidad de instancia si se hubiese llamado autos para sentencia, salvo si se dispusiese prueba de oficio, cuando la producción dependiera de la actividad de las partes.

\subsection{Excepción: medidas para mejor proveer (art.~36 inc.~4 CPCCN)}

\begin{definition}{Medidas para mejor proveer}
Diligencias necesarias para esclarecer la verdad de los hechos controvertidos, que el juez puede ordenar de oficio dejando sin efecto el auto para sentencia, siempre que no afecten el derecho de defensa de las partes y no suplan la negligencia u omisiones de las partes.
\end{definition}

\begin{example}{Límite de las medidas para mejor proveer}
En un juicio de mala praxis en el que la parte actora no ofreció una prueba pericial médica, el juez no puede, como medida para mejor proveer, ordenar la realización de un dictamen pericial médico. Sí puede, en cambio, en un accidente de tránsito en el que un dictamen pericial médico resulta oscuro en algunas cuestiones, dejar sin efecto el auto para sentencia y fijar una audiencia convocando a las partes y al perito para que brinden aclaraciones.
\end{example}

\begin{example}{Oficio a la municipalidad}
En un accidente de tránsito donde se discute el funcionamiento de un semáforo, el juez puede librar un oficio a la municipalidad correspondiente para que informe si el semáforo de determinada intersección el día del accidente funcionaba o estaba fuera de servicio.
\end{example}

\begin{important}
Si el juez ordena la medida pero deja la diligencia a cargo de las partes (por ejemplo, que la parte impulse el diligenciamiento del oficio), vuelve a computarse el plazo de caducidad de instancia mientras no se conteste el oficio y la parte no peticione que el expediente regrese a sentencia.
\end{important}

\subsection{Excepción: cambio de juez}

Si el juez que intervenía se jubila o toma licencia y quien dictará sentencia es otro magistrado, se deja sin efecto el auto para sentencia para notificar a las partes quién es el nuevo juez y permitir que ejerzan, si corresponde, la recusación con causa. Una vez vencido el plazo sin cuestionamientos, el expediente vuelve a estar en estado de sentencia.

\section{Plazos para dictar sentencia y notificación}

Una vez dictado el auto para sentencia, comienza a correr para el juez el plazo para dictar sentencia:

\begin{itemize}
    \item \textbf{Primera instancia} (proceso ordinario): 40 días.
    \item \textbf{Cámara} (tribunal colegiado): 60 días desde que la sala dicta su propio auto para sentencia (por ser tres jueces que deben votar).
\end{itemize}

La sentencia debe ser notificada \textbf{de oficio} por el juzgado o tribunal a todas las partes del proceso (actor, demandado, terceros, citadas en garantía) y a los profesionales a quienes se hayan regulado honorarios. Se notifica la parte dispositiva (la decisión final, también llamada ``fallo'').

\textbf{Excepción: demandado rebelde.} Aunque durante el proceso las notificaciones al rebelde se hacen por ministerio de ley, la sentencia sí debe notificársele. Esta cédula se dirige al domicilio real denunciado y es en formato papel, no electrónico. Su confección suele quedar a cargo de la parte actora.

\section{Morosidad judicial (arts.~167 y 168 CPCCN)}

El código prevé consecuencias para el juez que no dicta sentencia en los plazos establecidos. Se distinguen dos supuestos:

\begin{table}[H]
\centering
\begin{tabularx}{\textwidth}{>{\raggedright\arraybackslash}p{5cm} >{\raggedright\arraybackslash}X}
\toprule
\textbf{Supuesto de morosidad} & \textbf{Consecuencia} \\
\midrule
Demora en providencias simples o interlocutorias & Pesará negativamente en la calificación del juez (mal desempeño de funciones). No hay multa. \\
\addlinespace
Demora en el dictado de la sentencia definitiva & El juez puede ser sancionado con una multa de hasta el 15\% de su remuneración básica, y/o el expediente puede ser remitido a otro juez para que dicte sentencia. \\
\bottomrule
\end{tabularx}
\end{table}

Si el juez advierte que no va a poder dictar sentencia en término, debe presentar una nota a su superior (la Cámara, si es juez de primera instancia; la Corte, si es integrante de una sala de Cámara) con \textbf{diez días de anticipación}, explicando los motivos. Si esos motivos son atendibles, se concederá una prórroga o se evitará la multa y, de corresponder, se determinará si es más idóneo que la dicte otro juez.

\cornelltitle{la etapa conclusional}
\cornellblock{
\cqrow{¿Qué es la etapa conclusional y a través de qué vías se accede a ella?}{Es la etapa en que el juez dicta la sentencia (norma individual). Se accede por: (1) la vía normal: luego de la etapa probatoria y vencido el plazo para alegar; (2) la cuestión de puro derecho (art.~359, párr.~1.º): cuando no hay hechos controvertidos; (3) los supuestos del art.~360 inc.~6 y de los arts.~361 y 362: reconocimiento recíproco en la audiencia, oposición a la apertura a prueba, desistimiento de prueba, prueba anticipada suficiente.
}
\cqrow{¿Qué es el auto para sentencia y qué efectos produce?}{Es la resolución que da inicio a la etapa conclusional. Produce: (a) el cierre del debate (no más escritos ni documentos nuevos); (b) la improcedencia de la caducidad de instancia (art.~313 inc.~4); (c) el inicio del plazo para que el juez dicte sentencia (40 días en primera instancia; 60 días en Cámara). El juez puede dejarlo sin efecto para dictar medidas para mejor proveer (art.~36 inc.~4) o ante el cambio de magistrado.
}
\cqrow{¿Qué son las medidas para mejor proveer?}{Diligencias ordenadas de oficio por el juez (art.~36 inc.~4 CPCCN) para esclarecer la verdad de los hechos controvertidos, siempre que no afecten el derecho de defensa ni suplan la negligencia de las partes. Ejemplos: convocar al perito para aclarar un dictamen oscuro; librar oficio a la municipalidad para informar si un semáforo funcionaba. Si la diligencia queda a cargo de las partes, vuelve a correr la caducidad de instancia.
}
\cqrow{¿Cómo se notifica la sentencia?}{De oficio, por el juzgado o tribunal, a todas las partes (actor, demandado, terceros, citadas en garantía) y a los profesionales a quienes se hayan regulado honorarios; se notifica la parte dispositiva. Al demandado rebelde se le notifica por cédula en papel dirigida al domicilio real denunciado, cuya confección suele quedar a cargo de la parte actora.
}
}
\medskip
\cornellblocksum{
\cqrow{¿Qué consecuencias tiene la morosidad judicial en el dictado de sentencia?}{Si el juez no dicta sentencia en el plazo establecido (40 días en primera instancia, 60 en Cámara), puede ser sancionado con una multa de hasta el 15\% de su remuneración básica, y/o el expediente puede ser remitido a otro juez. Si advierte que no podrá cumplir el plazo, debe notificar a su superior con 10 días de anticipación; si los motivos son atendibles, puede obtenerse una prórroga.
}
}{una vez vencido el plazo para alegar, el proceso entra en la etapa conclusional, que se materializa mediante el auto para sentencia: resolución que cierra el debate, impide nuevas presentaciones de las partes, hace improcedente la caducidad de instancia y activa el plazo del juez para dictar sentencia (40 días en primera instancia, 60 en Cámara). A esta etapa se puede arribar también directamente desde la etapa introductoria si la cuestión es de puro derecho ---sin hechos controvertidos--- o por distintos supuestos que pueden darse en o en torno a la audiencia preliminar.
}

\chapter[Valoración de la prueba]{Sistemas de valoración de la prueba}\label{cap:valoracion}

%==============================================================

\section{Concepto de valoración de la prueba}

\begin{definition}{Valoración de la prueba}
Valorar la prueba significa apreciarla; en sentido figurado, ponerle un precio. Llegado el momento de la sentencia, el juez valora la prueba producida en la etapa probatoria.
\end{definition}

En este punto cobra especial importancia la \textbf{inmediación}.

\begin{definition}{Inmediación}
Contacto directo del juez con la producción de la prueba.
\end{definition}

No es lo mismo que el testigo haya declarado ante un funcionario del juzgado que transcribió su declaración en un acta, a que el testigo haya declarado en presencia del juez, con una declaración filmada que el juez puede reproducir al dictar sentencia. La inmediación permite una sentencia más eficaz, porque el juez aprecia de mejor modo la prueba.

\section{Sistemas de apreciación de la prueba (según Falcón)}

\subsection{Sistema de apreciación prohibida o excluida}

La \textbf{apreciación prohibida} refiere a los casos en que el legislador estableció normativamente la prohibición de ciertos medios de prueba. Por ejemplo, el art.~427 CPCCN establece quiénes no pueden ser testigos (parientes ascendientes, descendientes). Si declararan, el juez no puede apreciar esa declaración.

La \textbf{apreciación excluida} refiere a la facultad del juez (art.~386, párr.~2.º, CPCCN) de no estar obligado a ponderar una por una todas las pruebas: puede excluir aquellas que no considere trascendentes y detenerse en las que resulten más eficientes para formar su convicción.

\subsection{Sistema de la prueba tasada o prueba legal}

En este sistema, el legislador establece de antemano qué valor debe asignársele a cada prueba. Tiene un origen histórico, cuando se consideraba que tal vez el juez no estaba lo suficientemente formado para apreciar la prueba.

\begin{example}{Regla tasada}
Una regla tasada podría ser: ``un hecho solo puede probarse con dos o más testigos; un solo testigo constituye semiplena prueba''.
\end{example}

Es un sistema que puede dar más certezas, pero también mayor injusticia: muchas veces los testigos no se consiguen, y tasar las pruebas puede resultar injusto porque estas pueden ser más o menos valiosas en función de los hechos del caso concreto. Aplicado de modo único y aislado, puede no ser el más eficaz; prácticamente ya no se utiliza de modo único.

\subsection{Sistema de apreciación no tasada o libre convicción}

Dentro de este sistema existe un espectro:

\begin{itemize}
    \item En un extremo, la \textbf{apreciación arbitraria}, donde el juez valora la prueba sin seguir ningún tipo de regla.
    \item En el otro extremo, la \textbf{libre convicción}, donde el juez valora la prueba siguiendo las reglas de la lógica, la experiencia y la intuición (entendida como sentido común).
\end{itemize}

\subsection{Sistema de apreciación conjunta}

No es un sistema puro. En él se aplican, según la prueba de que se trate, criterios de los distintos sistemas anteriores: a veces la prueba tasada, a veces la prohibida o excluida, a veces la libre convicción.

\section{La sana crítica: el sistema del Código (art.~386 CPCCN)}

\begin{formula}{Art. 386 CPCCN --- Regla de valoración de la prueba}
Salvo disposición legal en contrario, los jueces formarán su convicción respecto de la prueba de conformidad con las reglas de la sana crítica. No tendrán el deber de expresar en la sentencia la valoración de todas las pruebas producidas, sino únicamente de las que fueren esenciales y decisivas para el fallo de la causa.
\end{formula}


\begin{definition}{Sana crítica}
Sistema de apreciación de la prueba que es científico y que se basa en las reglas de la experiencia y la lógica. También se lo denomina conjunto de reglas que permiten esclarecer la verdad al momento de dictar sentencia, siempre librado de todo vicio y error. Es el sistema adoptado por el legislador en el CPCCN.
\end{definition}

La sana crítica es, además, el razonamiento que sigue el juez, libre de error y de vicio, para determinar cuál de las dos posiciones asumidas por las partes en el proceso es la correcta. Está sistematizada de modo que el juez pueda demostrar cómo formó su convencimiento de manera concreta, lógica y razonada, y que las partes puedan controlar la construcción de la sana crítica y, si necesitan impugnarla, tengan los fundamentos para hacerlo.

\begin{note}
La expresión «salvo disposición legal en contrario» del artículo 386 remite a la prueba legal tasada: si el legislador le asignó valor a la prueba, el juez no puede apartarse de ella y debe comenzar por esas pruebas legales tasadas. Fuera de ese supuesto, el juez forma su convicción a través de las reglas de la sana crítica.
\end{note}

Si el juez comete un error al valorar la prueba o al fijar los hechos, la sentencia podrá ser objeto de recurso de apelación, para que la Cámara subsane ese error.

\section{Las reglas de la sana crítica según Falcón}

Falcón enumera las siguientes reglas que el juez debe seguir al momento de dictar sentencia:

\begin{longtable}{>{\raggedright\arraybackslash}p{1.6cm} >{\raggedright\arraybackslash}p{3.8cm} >{\raggedright\arraybackslash}p{7.6cm}}
\toprule
\textbf{Regla} & \textbf{Denominación} & \textbf{Desarrollo} \\
\midrule
\endfirsthead
\toprule
\textbf{Regla} & \textbf{Denominación} & \textbf{Desarrollo} \\
\midrule
\endhead

\textbf{1} & \textbf{Solo se prueban los hechos afirmados} & El juez, al dictar sentencia, no puede tener en cuenta otros hechos que no sean los afirmados por las partes. Los hechos no afirmados no existen. Está vinculado con el principio de congruencia. \\
\addlinespace
\textbf{2} & \textbf{Los hechos a probar deben ser controvertidos} & Solo los hechos afirmados y además controvertidos determinan el ámbito probatorio. El juez valorará la prueba producida para acreditar esos hechos. \\
\addlinespace
\textbf{3} & \textbf{Primero debe aplicarse la prueba tasada (si existe)} & Si hay alguna norma que le da cierto valor a determinada prueba, el juez debe aplicarla primero. Ejemplo: art.~579 CCyCN; en un juicio de filiación, la negativa del supuesto progenitor a realizarse la prueba genética se considera indicio grave en su contra. \\
\addlinespace
\textbf{4} & \textbf{Análisis estático de cada medio probatorio} & El juez analiza individualmente todos los medios de prueba producidos para ver cuál es más fiable. El código no establece un orden de valor, pero la doctrina entiende que en principio el orden es: 1.º documental, 2.º informativa, 3.º confesional, 4.º pericial, 5.º testimonial. \\
\addlinespace
\textbf{5} & \textbf{Análisis dinámico y en conjunto de la prueba} & Una vez analizada cada prueba estáticamente, el juez hace un análisis conjunto a la luz de los hechos del caso. Este análisis puede llevar a dar mayor valor a un medio probatorio distinto del orden estático. Ejemplo: en un accidente de tránsito, la declaración de un testigo presencial puede tener más valor que una exposición policial que no aclara cómo ocurrió el hecho. \\
\addlinespace
\textbf{6} & \textbf{Presunciones judiciales y conducta de las partes} & Para obtener la amalgama que estructura el relato de manera coordinada, el juez puede utilizar presunciones judiciales y considerar la conducta desarrollada por las partes a lo largo del proceso (art.~163 inc.~5 CPCCN). \\
\addlinespace
\textbf{7} & \textbf{En caso de duda: carga probatoria} & En caso de duda al momento de valorar la prueba, el juez debe tener por no acreditado el hecho y aplicar las reglas de la carga de la prueba (art.~377 CPCCN): quién tenía la carga de probar ese hecho y si produjo o no la prueba. \\
\addlinespace
\textbf{8} & \textbf{Todo debe volcarse por escrito en la sentencia} & Toda la actividad valorativa debe expresarse por escrito en la sentencia, explicando los fundamentos, argumentos y pasos seguidos para fijar ciertos hechos y para obtener la certeza en el pronunciamiento. La convicción no solo tiene que ser del magistrado, sino también de terceros; de allí la importancia de la construcción de la sentencia. \\
\bottomrule
\end{longtable}

\section{Ejemplos de sentencias reales (accidentes de tránsito)}

Se presentan dos sentencias reales de expedientes de accidentes de tránsito para ejemplificar las reglas de la sana crítica.

\subsection{Ejemplo 1: sentencia con valoración de testigos presenciales}

La sentencia se estructura en tres partes.

\textbf{A) Resultandos (descripción lineal y objetiva).} El juez relata cronológicamente todo lo que consta en el expediente:

\begin{quote}
``A fojas 1/40 se presenta la actora con el patrocinio letrado de tales abogados, promueve demanda contra el demandado Juárez y, si está en garantía, a determinada compañía de seguros. La pretensión tiene por objeto que se condene al demandado a pagar la suma de \$100.000 o lo que más o menos resulte de la prueba del expediente, más intereses y costas. Relata que el día 1.º de diciembre de 2020, cuando era tal hora, iba en su vehículo y [relata la ocurrencia del accidente]. A fojas 75/80 se presenta el demandado, contesta demanda y niega todos los hechos invocados. A fojas 100/120 se presenta la aseguradora, reconoce que existía una póliza de seguros vigente, pero se adhiere a la contestación del demandado. A fojas 98 se celebra la audiencia preliminar. A fojas 157 se clausura el período probatorio. A fojas 160/161 obra el alegato de la parte actora; las restantes partes no presentaron alegatos. A fojas 163 se llama autos para sentencia.''
\end{quote}

\textbf{B) Considerandos (análisis y valoración de la prueba).} El juez aplica las reglas de la sana crítica. En primer lugar, determina qué debe probarse y a quién correspondía la carga:

\begin{quote}
``Efectuada la síntesis de lo actuado, lo que corresponde es determinar si se probó la ocurrencia del hecho, cuya carga recaía sobre la parte actora (artículo 377 del Código Procesal).''
\end{quote}

Luego analiza la prueba testimonial:

\begin{quote}
``A tales fines, la actora ofreció y produjo las testimoniales de los señores Fernández y Palombo, quienes declararon en mi presencia en el marco de un proyecto de procesos orales firmados [indica ver el sobre reservado en el expediente a fojas 159]. Ambos testigos dieron versión concreta de lo acontecido, señalaron dónde se encontraban, qué fue lo que observaron. Dijeron que estaban fumando en el balcón que da a la calle Laprida cuando se produjo el accidente, que vieron cómo fue la colisión [transcribe las partes trascendentes de los testimonios].''
\end{quote}
% TODO: contenido incompleto o ambiguo en las notas originales (las expresiones "si está en garantía" y "proyecto de procesos orales firmados" figuran así en el apunte).

Y valora la credibilidad de los testigos:

\begin{quote}
``No detecté en aquel momento ningún elemento que me hiciera dudar de su veracidad en el momento en que declararon ante él. Todo lo contrario, los recuerdo como testigos confiables que relataron lo que vieron. Ahora, al repasar sus testimonios gracias a las posibilidades que ofrece el registro fílmico, tengo la misma impresión. A ello corresponde añadir que nadie cuestionó el contenido de sus declaraciones ni si no fueron impugnadas por la otra parte por considerar que no eran idóneos o que tenían contradicciones o que habían faltado a la verdad.''
\end{quote}

\textbf{C) Parte dispositiva o fallo.} El juez concluye:

\begin{quote}
``Por ende, tendré por acreditada la materialidad del hecho.''
\end{quote}

Tenido por acreditado el hecho, el juez determina quién es el responsable y pasa a analizar los rubros indemnizatorios reclamados a la luz de la prueba producida, para fijar la procedencia de cada uno y los montos.

\subsection{Ejemplo 2: sentencia con impugnación de testigo (concubina del demandado)}

En otra sentencia, el juez explicita que aplica la sana crítica y explica por qué excluye un testimonio:

\begin{quote}
``Las pruebas han de apreciarse de conformidad a las reglas de la sana crítica (artículo 386 del Código Procesal). Cabe aclarar que los jueces no estamos obligados a ponderar una por una [las pruebas] ---lo que establece el artículo 386---.''
\end{quote}

El juez menciona luego la declaración de una testigo ofrecida por el demandado:

\begin{quote}
``A fojas 210 obra la declaración testimonial de la señora González, quien brindó la siguiente versión de los hechos [la transcribe]. Sin embargo, su declaración fue impugnada por la parte actora a fojas [tanto], ya que la testigo dijo ser pareja/concubina del demandado, afirmando también que al momento de prestar declaración estaba esperando un hijo del demandado.''
\end{quote}

El CPCCN excluye expresamente como testigo al cónyuge, no a la concubina. Sin embargo, en este caso, durante la audiencia se le hicieron preguntas para acreditar sus datos personales y la testigo reconoció ser concubina del demandado y estar esperando un hijo suyo. Por eso la actora impugnó esa declaración.

El juez resuelve:

\begin{quote}
``Lo importante de un testigo es que la declaración sea idónea para crear su convicción respecto de lo que sucedió en la realidad de los hechos, y que no rige aquí la regla de que el testigo único es un testigo nulo. Pero dice: el hecho de que la testigo sea concubina del demandado obliga a examinarla con el mayor rigor, por lo que, ante la ausencia de otros elementos probatorios que precisen las declaraciones que ella dijo y que se habían atacado a través de la impugnación, no puede tenerse en cuenta su declaración.''
\end{quote}

El juez realiza así un \textbf{análisis dinámico}: se aparta de esa declaración no solo por la relación de la testigo con el demandado, sino también porque no existen otros elementos probatorios que corroboren lo que ella dijo. Por eso no resulta suficiente para formar convicción sobre lo ocurrido.

\cornelltitle{valoración de la prueba}
\cornellblock{
\cqrow{¿Cuáles son los 4 sistemas de valoración de la prueba?}{Según Falcón: (1) prohibida o excluida (el legislador prohíbe ciertos medios o el juez puede excluir los no trascendentes ---art.~386, párr.~2.º---); (2) tasada o legal (la ley predetermina el valor de cada prueba; hoy casi no se usa en forma pura); (3) libre convicción (sin reglas ---arbitraria--- o con reglas de lógica y experiencia); (4) conjunta (mezcla de los anteriores según la prueba de que se trate).
}
\cqrow{¿Qué es la sana crítica y cuál es su base legal?}{Es el sistema adoptado por el CPCCN (art.~386). Es un sistema científico de apreciación de la prueba basado en las reglas de la lógica, la experiencia y el sentido común (intuición). Permite al juez esclarecer la verdad al dictar sentencia, libre de todo vicio y error. Los jueces no están obligados a valorar todas las pruebas producidas, solo las esenciales y decisivas.
}
\cqrow{¿Cuáles son las 8 reglas de la sana crítica según Falcón?}{(1) Solo se prueban los hechos afirmados. (2) Esos hechos deben ser además controvertidos. (3) Primero se aplica la prueba tasada, si existe. (4) Análisis estático de cada medio probatorio (orden de fiabilidad doctrinario: documental, informativa, confesional, pericial, testimonial). (5) Análisis dinámico y conjunto (puede variar el orden según el caso). (6) Uso de presunciones judiciales y de la conducta de las partes (art.~163 inc.~5). (7) En caso de duda, aplicar la carga probatoria (art.~377). (8) Todo debe volcarse por escrito en la sentencia, con fundamentos.
}
\cqrow{¿Por qué el análisis dinámico puede alterar el orden de fiabilidad de la prueba?}{Porque el valor de una prueba depende de los hechos del caso concreto. Ejemplo: aunque doctrinariamente la documental está por encima de la testimonial, en un accidente de tránsito donde lo controvertido es cómo ocurrió el hecho, la declaración de un testigo presencial que vio la colisión tendrá más valor que una simple exposición policial que no aclara los detalles de la ocurrencia del accidente.
}
}
\medskip
\cornellblocksum{
\cqrow{¿Cómo se aplicaron las reglas de la sana crítica en las sentencias reales analizadas?}{En la primera sentencia, el juez acreditó el hecho con la testimonial de testigos presenciales, que declararon en su presencia, cuyo contenido nadie cuestionó. En la segunda, descartó la declaración de la concubina del demandado por la falta de elementos que la corroboraran, aplicando el rigor que impone su relación con una de las partes.
}
}{al momento de dictar sentencia, el juez valora la prueba conforme al sistema de la sana crítica (art.~386 CPCCN): un sistema científico basado en la lógica, la experiencia y el sentido común que, según Falcón, se articula en ocho reglas, que van desde tomar en cuenta solo los hechos afirmados y controvertidos, pasando por el análisis estático y dinámico de la prueba, hasta la exigencia de volcar por escrito y con fundamentos toda la actividad valorativa en la sentencia. Dos sentencias reales de accidentes de tránsito ilustran cómo estas reglas se aplican: el hecho se acredita con la testimonial de testigos presenciales (dando mayor valor dinámico a esa prueba sobre la exposición policial) y se descarta la declaración de la concubina del demandado por falta de corroboración.
}

\part{La terminación del proceso y la sentencia}
\markboth{Parte II}{Parte II}
Esta parte estudia las distintas formas en que concluye un proceso civil y las resoluciones con que el tribunal se pronuncia. Primero se examinan los \textbf{modos anormales de terminación} ---es decir, aquellos que no son la sentencia definitiva--- y la caducidad de la instancia. Luego se analiza la clasificación de las \textbf{resoluciones judiciales}, el contenido de la sentencia definitiva, la actuación del juez con posterioridad a la sentencia, la \textbf{cosa juzgada} y la relación entre la sentencia civil y la sentencia penal.


\begin{note}
Para comprender cómo se relacionan estos temas puede imaginarse el proceso judicial como un partido de fútbol: las resoluciones judiciales son las decisiones del árbitro, desde una simple advertencia hasta la tarjeta roja o el gol válido. Cada decisión tiene su nombre, su nivel de complejidad y el recurso (la herramienta) que el litigante puede usar para impugnarla. Cuando el partido termina y el resultado es definitivo, sin posibilidad de revancha, se está ante la cosa juzgada. Y si además hubo una investigación disciplinaria paralela (proceso penal), ambas decisiones deben coordinar sus conclusiones para evitar un <<escándalo lógico>>.
\end{note}

\chapter[Modos anormales de terminación]{Modos anormales de terminación del proceso}\label{cap:modos}

\section{Encuadre: modo normal y modos anormales de terminación}

Este capítulo y el siguiente abordan los \textbf{modos anormales de terminación del proceso civil}, regulados en los artículos 304 a 318 del Código Procesal Civil y Comercial de la Nación (CPCCN). A diferencia del modo normal de terminación ---la sentencia definitiva---, estos institutos permiten que el proceso concluya de manera anticipada o por circunstancias distintas al dictado de una sentencia.

Las etapas del proceso ---la etapa introductoria, la etapa probatoria y la presentación de los alegatos--- conducen a la sentencia, que constituye el \textbf{modo normal} de terminar el proceso (sin perjuicio de los recursos de apelación): se trata de la \textbf{sentencia definitiva}. Los modos anormales, en cambio, definen otras formas en que un proceso puede finalizar.

\begin{note}
\textbf{Aclaración terminológica.} La expresión <<modo anormal>> es la denominación utilizada por el CPCCN; se trata de un término arbitrario adoptado por el Código. Algunos autores prefieren hablar de <<modos extraordinarios>> de terminación del proceso, por oposición al modo <<ordinario>> o <<normal>> (la sentencia). La diferencia es meramente semántica.
\end{note}

El CPCCN prevé \textbf{cinco} modos anormales de terminación del proceso, regulados a partir del artículo 304, que se estudian en el siguiente orden:

\begin{enumerate}
    \item El \textbf{desistimiento} ---del proceso o del derecho---, acto unilateral o bilateral del actor o reconviniente.
    \item El \textbf{allanamiento} ---sometimiento total o parcial del demandado a las pretensiones del actor---.
    \item La \textbf{transacción} ---contrato bilateral con concesiones recíprocas que extingue obligaciones litigiosas---.
    \item La \textbf{conciliación} ---acuerdo intraprocesal celebrado ante el juez---.
    \item La \textbf{caducidad de la instancia} ---extinción del proceso por inactividad de la parte que tiene la carga de impulsarlo, operada en los plazos previstos en el artículo 310 del CPCCN---.
\end{enumerate}

Cada instituto se analiza en sus requisitos, efectos procesales y diferencias con figuras afines, con especial énfasis en la jurisprudencia y la práctica profesional.

\begin{example}{Cómo se relacionan los modos anormales}
Una persona inicia una demanda por cobro de honorarios. Lo normal sería que el proceso avance hasta la sentencia. Pero el actor puede \textbf{desistir} antes de que se notifique al demandado; el deudor puede \textbf{allanarse} y pagar; ambos pueden llegar a una \textbf{transacción}, o el juez puede convocarlos a \textbf{conciliación}. Y si nadie actúa durante 6 meses, opera la \textbf{caducidad de la instancia}. Todos estos son modos en que el proceso puede terminar sin sentencia.
\end{example}

% ----------------------------------------------------------
\section{El desistimiento (art. 304 y ss., CPCCN)}

Desistir, en términos procesales, significa \textbf{renunciar} o \textbf{abandonar}. % TODO: contenido incompleto o ambiguo en las notas originales (las notas originales utilizan también el verbo <<aplicar>> como sinónimo de desistir)
El desistimiento es el primer modo anormal de terminación del proceso previsto en el CPCCN y su nota característica es la renuncia o el abandono de una posición procesal. Puede ser un acto \textbf{unilateral} o \textbf{bilateral}, según la etapa procesal en que se produzca.

\subsection{¿Quién puede desistir?}

El desistimiento corresponde a quien \textbf{articula una pretensión} y no a quien plantea una defensa. Pueden desistir:

\begin{itemize}
    \item el \textbf{actor}, en la demanda principal;
    \item el \textbf{reconviniente}, en caso de existir reconvención.
\end{itemize}

El demandado que se limita a defenderse no puede desistir, ya que no articula una pretensión propia: el desistimiento es un acto de quien plantea la pretensión y no de quien se defiende.

\subsection{Tipos de desistimiento: del proceso y del derecho}

Como modo anormal de terminar el proceso deben distinguirse dos supuestos: el \textbf{desistimiento del proceso} y el \textbf{desistimiento del derecho}. Este último se encuentra en algunos libros como \textbf{desistimiento de la acción}. Ambas modalidades tienen efectos radicalmente distintos.

\begin{table}[H]
\centering
\caption{Desistimiento del proceso y desistimiento del derecho}
\small
\begin{tabularx}{\textwidth}{
    >{\raggedright\arraybackslash}p{0.22\textwidth}
    >{\raggedright\arraybackslash}X
    >{\raggedright\arraybackslash}X
}
\toprule
\textbf{Aspecto} & \textbf{Desistimiento del proceso} & \textbf{Desistimiento del derecho / acción} \\
\midrule
¿Quién puede hacerlo? & Solo el actor (o reconviniente) & Solo el actor (o reconviniente) \\
Conformidad de la contraria & Necesaria si ya se notificó el traslado de demanda & No se requiere (es unilateral siempre) \\
Efecto sobre el derecho & No lo extingue; puede replantearse en otro proceso & Lo extingue definitivamente \\
Efecto de cosa juzgada & No genera cosa juzgada material sobre el derecho & Opera con efecto de cosa juzgada material \\
Posibilidad de nuevo proceso & Sí, siempre que no esté prescripto el derecho & No; el derecho fue renunciado \\
\bottomrule
\end{tabularx}
\end{table}

\subsubsection{Desistimiento del derecho}

Si el actor desiste del \textbf{derecho}, es decir, de la pretensión articulada, \textbf{no se requiere la conformidad de la otra parte}. El efecto que provoca es el de la \textbf{cosa juzgada}: el actor no podrá reclamar ese derecho en otro proceso, porque lo desistió. El desistimiento opera con calidad de cosa juzgada y esa parte no podrá hacer valer ese derecho en un proceso distinto.

Se trata de un acto procesal unilateral. Su fundamento radica en que el actor, al renunciar a su propio derecho sustancial, no perjudica al demandado, quien se ve beneficiado, ya que no podrá ser perseguido nuevamente por esa pretensión. El efecto es análogo al de la cosa juzgada material.

\subsubsection{Desistimiento del proceso}

El desistimiento del \textbf{proceso} puede realizarse de forma unilateral y sin requerir la conformidad de la otra parte \textbf{únicamente antes de notificar la demanda} al demandado. Si ya se notificó el traslado de la demanda, se requiere la \textbf{conformidad del demandado}.

La razón es que este desistimiento no genera el mismo efecto respecto de la cosa juzgada: el derecho no se desistió, sino solamente se renunció a continuar con esa instancia. Por lo tanto, el derecho puede hacerse valer en otro proceso y el actor podría iniciar otro proceso y formular el mismo reclamo. Por esta circunstancia, una vez notificado el traslado de la demanda, el actor solo puede desistir si cuenta con la conformidad del demandado.

\subsection{Legitimación para desistir}

Para desistir se necesita no solo la legitimación sustancial en términos de capacidad, sino también la \textbf{legitimación procesal}. Si quien desiste es un abogado con un poder que no es suficiente para desistir de la acción o del derecho, el desistimiento no es válido: el representante debe estar \textbf{expresamente facultado} para desistir, ya sea del proceso o del derecho.

\begin{important}
\textbf{Requisito de legitimación doble.} Para desistir válidamente se requiere: (1) \textbf{legitimación sustancial} ---capacidad de derecho sobre la pretensión--- y (2) \textbf{legitimación procesal} ---el representante o apoderado debe tener poder especial y expreso que lo autorice a desistir---. El poder general de representación no es suficiente.
\end{important}

\subsection{Oportunidad para desistir}

El desistimiento puede llevarse a cabo \textbf{en cualquier instancia del proceso antes del dictado de la sentencia}. Una vez dictada la sentencia, el proceso terminó, por lo que no se podría desistir del proceso. Además, si la sentencia reconoció un derecho, hizo lugar a la pretensión del actor o la rechazó, tampoco se podría desistir. Una vez dictada la sentencia definitiva, el proceso ha concluido por su modo normal y el desistimiento ya no tiene objeto.

% ----------------------------------------------------------
\section{El allanamiento (art. 307, CPCCN)}

El allanamiento es una de las actitudes posibles del demandado al contestar la demanda y se encuentra contemplado en el artículo 307 del CPCCN.

\begin{articulo}{Art. 307 --- CPCCN}
Regula el allanamiento como actitud del demandado frente a las pretensiones del actor. Prevé los efectos según si el allanamiento es total o parcial, y si viene acompañado o no del cumplimiento de la pretensión.
\end{articulo}

\begin{definition}{Allanamiento}
El allanamiento es un acto jurídico procesal que consiste en el \textbf{sometimiento total o parcial del demandado a las pretensiones del actor}.
\end{definition}

Por oposición al desistimiento, el allanamiento corresponde al \textbf{demandado}. Puede allanarse a las pretensiones del actor el demandado, o bien el \textbf{reconvenido}, que se allana a las pretensiones del reconviniente.

\subsection{Naturaleza: ¿reconocimiento de los hechos o del derecho?}

\begin{important}
Allanarse \textbf{no constituye el reconocimiento ni de la verdad de los hechos ni del derecho}. Es un acto voluntario del demandado que, a través de la actitud que asume, permite el progreso de la pretensión articulada por el actor.
\end{important}

Esta es una distinción esencial: el allanamiento no implica que el demandado reconozca que los hechos alegados por el actor sean verdaderos ni que el derecho invocado sea correcto. Simplemente expresa su voluntad de no oponerse al progreso de la pretensión.

\subsection{Características}

El allanamiento es un \textbf{acto unilateral}: no requiere la conformidad de la otra parte. Si el demandado quiere allanarse, no necesita la conformidad del actor.

\subsection{Modalidades: expreso y tácito}

Se trata de un acto unilateral que puede ser \textbf{expreso} o \textbf{tácito}.

\begin{itemize}
    \item \textbf{Expreso:} debe ser contundente y no generar ninguna duda; la parte demandada manifiesta expresamente que se allana, es decir, que se somete total o parcialmente a las pretensiones de la otra parte.
    \item \textbf{Tácito:} siempre genera mayores dudas. Por ello, el allanamiento tácito exige una actitud del demandado que no genere ninguna duda, que sea \textbf{incontrovertible}, de la que surja que se está sometiendo a las pretensiones del actor aunque no lo manifieste expresamente.
\end{itemize}

\begin{example}{Allanamiento tácito claro}
El demandado, sin oponerse a la pretensión del actor, se presenta y paga, satisfaciendo así la totalidad de la pretensión. Aunque no manifestó expresamente <<me allano>>, el efecto es incontrovertiblemente ese.
\end{example}

\begin{example}{Allanamiento tácito dudoso (demanda de desalojo)}
En una demanda de desalojo, la parte demandada se presenta y deposita en el tribunal las llaves del inmueble. ¿Es un allanamiento tácito? Puede haber otras cuestiones en discusión simultáneamente, por lo que la actitud puede no ser incontrovertible. Se deberá analizar el caso concreto.
\end{example}

\subsection{Allanamiento total y parcial: efectos}

% TODO: contenido incompleto o ambiguo en las notas originales (se plantea qué sucede si el allanamiento es parcial, pero la explicación no se completa; el cuadro siguiente distingue únicamente según se cumpla o no la pretensión)
Los efectos del allanamiento dependen de si, además de allanarse, el demandado cumplió o no con la pretensión.

\begin{table}[H]
\centering
\caption{Efectos del allanamiento según el cumplimiento de la pretensión}
\small
\begin{tabularx}{\textwidth}{
    >{\raggedright\arraybackslash}p{0.20\textwidth}
    >{\raggedright\arraybackslash}X
    >{\raggedright\arraybackslash}X
}
\toprule
\textbf{Supuesto} & \textbf{Allanamiento + cumplimiento de la pretensión} & \textbf{Allanamiento sin cumplimiento} \\
\midrule
¿Qué hace el juez? & Dicta una resolución interlocutoria y manda al archivo las actuaciones (art. 307) & Hace lugar a la pretensión y dicta sentencia definitiva para habilitar la ejecución \\
Tipo de resolución & Resolución interlocutoria & Sentencia definitiva \\
Etapa siguiente & Archivo del expediente & Ejecución de sentencia \\
\bottomrule
\end{tabularx}
\end{table}

El artículo 307 se refiere a la resolución interlocutoria en el caso de que se cumpla la pretensión y a la orden de enviar al archivo las actuaciones.

\subsection{La cosa juzgada en el allanamiento}

La resolución que admite el allanamiento \textbf{hace cosa juzgada}. Pero la cosa juzgada opera con respecto a la \textbf{pretensión} y no con respecto a los \textbf{hechos}, porque el allanamiento no importa el reconocimiento de la verdad de los hechos invocados por el actor ni el reconocimiento del derecho invocado.

En consecuencia, no se podrán invocar esos hechos como verdaderos en otro proceso, como si hubiesen sido reconocidos en una sentencia, porque se trata de otra cosa. Esa es una característica propia de este modo anormal de terminar el proceso.

% ----------------------------------------------------------
\section{La transacción (art. 308, CPCCN; arts. 1641 y ss., CCyCN)}

La transacción es un tercer modo anormal de terminar el proceso. En el CPCCN se regula en el artículo 308.

\begin{articulo}{Art. 308 --- CPCCN}
Las partes podrán hacer valer la transacción del derecho en litigio con la presentación del convenio o suscripción de acta ante el juez. Éste se limitará a examinar la concurrencia de los requisitos exigidos por la ley para la validez de la transacción y la homologará o no. En este último caso, continuarán los procedimientos del juicio.
\end{articulo}

\begin{articulo}{Art. 1641 --- Código Civil y Comercial de la Nación (CCyCN)}
La transacción es un contrato por el cual las partes, para evitar el litigio o ponerle fin, haciéndose concesiones recíprocas, extinguen obligaciones dudosas o litigiosas.
\end{articulo}

\subsection{Concepto y naturaleza}

La transacción proviene del ámbito del derecho civil. Es un \textbf{acto jurídico bilateral} por el cual las partes, haciéndose \textbf{concesiones recíprocas}, extinguen obligaciones litigiosas o dudosas. Define la posibilidad de ponerle fin al litigio, que es el marco que el CPCCN incorpora en el artículo 308 como modo anormal de terminar el proceso. Constituye un medio de extinción de las obligaciones entre las partes que se hacen esas concesiones recíprocas.

Es bilateral (requiere el acuerdo de ambas partes) y puede ser judicial o extrajudicial. En el ámbito procesal, la transacción debe presentarse ante el juez para su homologación.

\subsection{Requisitos formales}

La transacción debe hacerse \textbf{por escrito} y presentarse ante el juez en cuyo proceso están en juego los derechos litigiosos, para que éste la homologue o no. También se puede suscribir un acta ante el juez en el cual tramita la causa.

\begin{articulo}{Art. 1643 --- CCyCN (Forma)}
La transacción debe hacerse por escrito. Si recae sobre derechos litigiosos, sólo es eficaz a partir de la presentación del instrumento firmado por los interesados ante el juez de la causa.
\end{articulo}

\subsection{Materias excluidas de la transacción}

\begin{articulo}{Art. 1644 --- CCyCN (Prohibiciones)}
No puede transigirse sobre derechos en los que está comprometido el orden público, ni sobre derechos irrenunciables. Tampoco pueden ser objeto de transacción los derechos sobre las relaciones de familia o el estado de las personas, excepto que se trate de derechos patrimoniales derivados de aquellos o de otros derechos sobre los cuales expresamente este Código admite pactar.
\end{articulo}

Las cuestiones que afectan el \textbf{orden público} no pueden ser objeto de transacción. Respecto de las \textbf{relaciones de familia}, no es posible, por ejemplo, que las partes acuerden considerarse divorciadas y pacten de ese modo su propio divorcio: esto no funciona así.

Lo que se admite dentro de las relaciones de familia como posible objeto de transacción son las cuestiones derivadas que tengan \textbf{contenido patrimonial}, por ejemplo:

\begin{itemize}
    \item las cuestiones sobre alimentos de los hijos menores;
    \item las cuestiones relativas a la división del régimen de comunidad de bienes, en las que las partes pueden acordar cómo dividir o repartir los bienes conyugales.
\end{itemize}

Tampoco pueden transigirse los \textbf{derechos irrenunciables}, como por ejemplo los derechos laborales o de los trabajadores, que no se pueden pactar en esos términos ni ser objeto de transacción.

% ----------------------------------------------------------
\section{La conciliación (art. 309, CPCCN)}

La conciliación se estudió previamente como sistema de solución alternativa de disputas y también constituye un modo anormal de terminar el proceso. Se encuentra prevista en el artículo 309.

\begin{articulo}{Art. 309 --- CPCCN}
Los acuerdos conciliatorios celebrados por las partes ante el juez y homologados por éste tendrán autoridad de cosa juzgada.
\end{articulo}

\begin{articulo}{Art. 36, inc. 2° --- CPCCN}
En cualquier momento del proceso el juez puede convocar a las partes para alcanzar una conciliación.
\end{articulo}

\begin{articulo}{Art. 360, inc. 1° --- CPCCN}
En la audiencia preliminar, el juez intentará conciliar a las partes y acercarlas a los fines de alcanzar un acuerdo.
\end{articulo}

\subsection{Concepto}

Conciliar significa componer, ponerse de acuerdo. % TODO: contenido incompleto o ambiguo en las notas originales (las notas distinguen la conciliación <<extraprocesal o judicial>> y la <<intraprocesal>> sin precisar cada categoría)
Se distingue la conciliación en la fase extraprocesal o judicial y la intraprocesal, para poner en el centro de la escena la figura del juez como quien convoca la conciliación dentro del proceso.

En los dos supuestos se trata de un \textbf{avenimiento amigable} de las partes. En la conciliación intraprocesal, las diferencias de las partes se arreglan ante la presencia del magistrado o a instancia suya, porque es el propio magistrado quien las convoca.

\subsection{Efecto de cosa juzgada}

Si se alcanza la conciliación y se arriba a un acuerdo, el proceso termina por un modo anormal: no por vía de sentencia, sino por la realización de ese acuerdo. El acuerdo \textbf{extingue las pretensiones antagónicas de las partes} y produce el efecto de \textbf{cosa juzgada} una vez que está \textbf{homologado por el juez}.

\subsection{Conciliación y transacción: diferencias esenciales}

La conciliación y la transacción tienen muchos puntos en común, porque en ambas se alcanza un acuerdo que el juez podrá homologar, pero presentan distinciones importantes.

\begin{table}[H]
\centering
\caption{Diferencias entre transacción y conciliación}
\small
\begin{tabularx}{\textwidth}{
    >{\raggedright\arraybackslash}p{0.19\textwidth}
    >{\raggedright\arraybackslash}X
    >{\raggedright\arraybackslash}X
}
\toprule
\textbf{Aspecto} & \textbf{Transacción} & \textbf{Conciliación} \\
\midrule
Iniciativa & Surge de las partes (voluntad propia) & A instancia o convocatoria del juez \\
Ámbito & Judicial o extrajudicial; derechos litigiosos o dudosos & Siempre intraprocesal (dentro del proceso) \\
Naturaleza jurídica & Contrato (arts. 1641 y ss. CCyCN) & Acuerdo procesal ante el magistrado \\
Estado del derecho & Puede recaer sobre derechos dudosos o litigiosos & Solo sobre derechos ya litigiosos (proceso iniciado) \\
Homologación & Puede o no ser homologado por el juez & Se celebra ante el juez y produce cosa juzgada al homologarse \\
\bottomrule
\end{tabularx}
\end{table}

La diferencia vinculada con la \textbf{voluntad} es la siguiente: en la transacción, la voluntad de las partes se manifiesta en el sentido de querer llevar a cabo esa avenencia ante el juez, de modo que la necesidad de transigir surge de las partes. En cambio, la conciliación surge a raíz de una \textbf{propuesta del juez}: se realiza luego de una convocatoria efectuada por el juez en la audiencia preliminar, en virtud de las facultades previstas en el artículo 36, a lo largo de todo el proceso.

% ----------------------------------------------------------

\cornelltitle{modos anormales de terminación: desistimiento, allanamiento, transacción y conciliación}
\cornellblock{
\cqrow{¿Qué es un modo anormal de terminación del proceso?}{Todo aquel que no es la sentencia definitiva. El CPCCN regula cinco supuestos (arts. 304--318): desistimiento, allanamiento, transacción, conciliación y caducidad de la instancia. Algunos autores los denominan también <<modos extraordinarios>>.
}
\cqrow{¿Quién puede desistir y de qué?}{Solo quien articula una pretensión: el actor o el reconviniente. Puede desistir del proceso (renuncia a continuar esa instancia; no extingue el derecho) o del derecho/acción (renuncia al derecho sustancial; genera cosa juzgada y no puede replantearse). El desistimiento del proceso requiere la conformidad del demandado si ya fue notificado el traslado de la demanda.
}
\cqrow{¿Cuál es la diferencia entre desistir del proceso y desistir del derecho?}{El desistimiento del derecho es siempre unilateral y opera con efecto de cosa juzgada material: el derecho queda extinguido y no puede reclamarse en otro proceso. El desistimiento del proceso no genera cosa juzgada material sobre el derecho: el actor puede iniciar otro proceso, siempre que el derecho no esté prescripto.
}
\cqrow{¿Hasta cuándo puede desistirse?}{En cualquier instancia del proceso antes del dictado de la sentencia. Dictada la sentencia, el proceso terminó por su modo normal y el desistimiento ya no tiene objeto.
}
\cqrow{¿Qué legitimación se requiere para desistir?}{Doble legitimación: sustancial (capacidad sobre el derecho en cuestión) y procesal (el representante debe tener poder especial y expreso para desistir; el poder general no alcanza).
}
\cqrow{¿Qué es el allanamiento y a quién corresponde?}{Es el acto jurídico procesal del demandado (o reconvenido) por el cual se somete total o parcialmente a las pretensiones del actor. Es unilateral (no requiere la conformidad del actor), puede ser expreso o tácito, y no implica reconocimiento de la verdad de los hechos ni del derecho.
}
\cqrow{¿Cómo opera el allanamiento con cumplimiento y sin cumplimiento de la pretensión?}{Con cumplimiento: el juez dicta resolución interlocutoria y manda al archivo las actuaciones. Sin cumplimiento: hace lugar a la pretensión y dicta sentencia definitiva para habilitar la ejecución de sentencia.
}
}
\medskip
\cornellblocksum{
\cqrow{¿Qué cosa juzgada produce el allanamiento?}{La resolución que lo admite hace cosa juzgada respecto de la pretensión, pero no respecto de los hechos (que no fueron reconocidos como verdaderos). No pueden invocarse en otro proceso como hechos reconocidos en una sentencia.
}
\cqrow{¿Qué es la transacción y qué la distingue de la conciliación?}{La transacción es un contrato bilateral en que las partes se hacen concesiones recíprocas para extinguir obligaciones litigiosas o dudosas (art. 1641 CCyCN; art. 308 CPCCN). Surge de la voluntad de las partes y puede ser judicial o extrajudicial. La conciliación, en cambio, surge de la convocatoria del juez y es siempre intraprocesal (art. 309 CPCCN).
}
\cqrow{¿Qué materias no pueden ser objeto de transacción?}{Derechos en que esté comprometido el orden público, derechos irrenunciables (como los laborales) y derechos sobre relaciones de familia o estado de las personas. Excepción: derechos patrimoniales derivados (alimentos, división de bienes conyugales).
}
}{Los modos anormales de terminación del proceso, regulados en los arts. 304 a 318 del CPCCN, se oponen al modo normal, que es la sentencia definitiva. El desistimiento es un acto unilateral del actor o reconviniente con efectos distintos: el desistimiento del derecho genera cosa juzgada material, mientras que el del proceso solo extingue la instancia sin afectar el derecho. El allanamiento es el acto unilateral del demandado por el que se somete total o parcialmente a las pretensiones del actor; no importa reconocimiento de hechos ni de derecho y produce cosa juzgada limitada a la pretensión. La transacción (contrato bilateral con concesiones recíprocas, arts. 1641 y ss. CCyCN y art. 308 CPCCN) y la conciliación (acuerdo intraprocesal ante el juez, art. 309 CPCCN) son institutos afines que difieren en la iniciativa (partes frente a juez) y en el ámbito (judicial o extrajudicial frente a siempre intraprocesal).
}

\chapter{La caducidad de la instancia}\label{cap:caducidad}

\section{Marco normativo y concepto}

La caducidad de la instancia está prevista en el CPCCN en los artículos 310 a 318.

\begin{important}
\textbf{Advertencia sobre derecho comparado provincial.} La regulación de la caducidad de la instancia varía entre provincias: no todos los códigos la tratan de igual modo, y el modo en que opera en el proceso judicial depende de cómo esté prevista en el Código Procesal de que se trate. El análisis de este capítulo se circunscribe al sistema del CPCCN, aplicable en el fuero federal y en la Ciudad Autónoma de Buenos Aires.
\end{important}

La caducidad de la instancia es el modo anormal de terminación del proceso que ocupa la mayor cantidad de artículos del capítulo y, en la práctica profesional, en la casuística y en la interpretación de los jueces en los casos concretos, es el instituto que ha tenido más diversas interpretaciones en la jurisprudencia.

\begin{definition}{Caducidad de la instancia}
Extinción del proceso por la inactividad de la parte que tiene la carga de impulsarlo, operada en los plazos previstos en el artículo 310 del CPCCN.
\end{definition}

\subsection{Nota esencial: la inactividad procesal}

La caducidad de la instancia \textbf{no cancela la pretensión}. Entre los modos anormales de terminación del proceso, esto sucede solamente con el desistimiento del proceso y con la caducidad; % TODO: contenido incompleto o ambiguo en las notas originales (las notas afirman que esto <<pasa solamente>> con el desistimiento del proceso y que en los demás modos <<no sucede>>, sin precisar la comparación con la caducidad)
en los demás no ocurre. Por lo tanto, la misma pretensión puede ser promovida en otro proceso judicial, siempre y cuando el derecho \textbf{no esté alcanzado por la prescripción de la acción}. Son términos que muchas veces se confunden.

\subsection{Ámbito de aplicación: procesos dispositivos}

La caducidad de la instancia es aplicable a los \textbf{procesos dispositivos}, porque en ellos el impulso del proceso está a cargo de las partes. La inactividad de la parte que tiene la carga de impulsarlo genera que, una vez vencido el plazo establecido por el Código, el tribunal pueda declarar, de oficio o a pedido de la parte contraria, el cese del curso de esa instancia y, por ende, la extinción del proceso.

Para que se configure la caducidad deben verificarse, entonces:

\begin{itemize}
    \item un proceso dispositivo, en el que las partes tienen la carga de impulsar el proceso y no lo hacen;
    \item el vencimiento de los plazos previstos en el Código;
    \item inactividad procesal durante todos esos plazos, es decir, que la parte no haya impulsado el proceso.
\end{itemize}

\subsection{Fundamento del instituto}

El Código prevé la caducidad de la instancia como modo anormal de terminar el proceso por dos razones:

\begin{enumerate}
    \item El \textbf{interés público}: que los procesos judiciales no permanezcan paralizados indefinidamente en el tiempo, ocupando casilleros en los tribunales. Si la parte no activa el avance del proceso, se puede presumir su \textbf{abandono}.
    \item La relación procesal no solamente vincula a las partes ---actor y demandado--- sino que también incluye al juez, y el órgano jurisdiccional no puede quedar al arbitrio de los tiempos que las partes quieran manejar en sus procesos judiciales. Por eso es necesario poner un límite y un marco a esta circunstancia.
\end{enumerate}

\subsection{Distinción con la prescripción de la acción}

La prescripción de la acción se funda en una presunción de abandono del derecho: durante el tiempo en que la acción podría haberse hecho efectiva a través de una demanda y un reclamo judicial no se lo hizo, lo que impide hacerlo con posterioridad al vencimiento del plazo de prescripción. De modo análogo, la inactividad dentro del proceso judicial de la parte que tiene la carga de impulsarlo importa la presunción de abandono de esa instancia judicial.

\begin{table}[H]
\centering
\caption{Prescripción de la acción y caducidad de la instancia}
\small
\begin{tabularx}{\textwidth}{
    >{\raggedright\arraybackslash}p{0.21\textwidth}
    >{\raggedright\arraybackslash}X
    >{\raggedright\arraybackslash}X
}
\toprule
\textbf{Aspecto} & \textbf{Prescripción de la acción} & \textbf{Caducidad de la instancia} \\
\midrule
Naturaleza & Instituto de derecho sustancial & Instituto de derecho procesal \\
¿Qué se extingue? & El derecho / la acción & La instancia (el proceso en curso) \\
¿Desde cuándo corre? & Desde que nació el derecho exigible & Desde el último acto impulsor del proceso \\
¿Cancela la pretensión? & Sí (si prescribió, no se puede reclamar) & No (puede iniciarse nuevo proceso si el derecho no prescribió) \\
¿Con qué tipo de días? & Los que disponga cada norma sustancial & Días corridos (incluye inhábiles, salvo ferias judiciales) \\
\bottomrule
\end{tabularx}
\end{table}

Son dos institutos diversos: la prescripción de la acción y la caducidad de la instancia. La caducidad de la instancia \textbf{se abre con la promoción de la demanda}, sin importar que se haya notificado o no el traslado de la demanda.

% ----------------------------------------------------------
\section{Plazos de caducidad (art. 310, CPCCN)}

\begin{articulo}{Art. 310 --- CPCCN (Plazos)}
Se producirá la caducidad de la instancia cuando no se instare su curso dentro de los siguientes plazos:
\begin{enumerate}[label=\arabic*°]
    \item Seis meses, en primera instancia.
    \item Tres meses, en segunda o ulterior instancia y en cualquiera de las instancias en el juicio sumarísimo, en el juicio ejecutivo, en las ejecuciones especiales y en los incidentes.
    \item En el que se opere la prescripción de la acción, si fuese menor a los indicados precedentemente.
    \item Un mes, en el incidente de caducidad de instancia.
\end{enumerate}
\end{articulo}

\begin{table}[H]
\centering
\caption{Plazos de caducidad según el art. 310 del CPCCN}
\small
\begin{tabularx}{\textwidth}{
    >{\raggedright\arraybackslash}p{0.12\textwidth}
    >{\raggedright\arraybackslash}X
}
\toprule
\textbf{Plazo} & \textbf{Supuesto} \\
\midrule
6 meses & Primera instancia \\
3 meses & Segunda o ulterior instancia; juicio sumarísimo; juicio ejecutivo; ejecuciones especiales; incidentes \\
Prescripción & El plazo en que se opere la prescripción de la acción, si fuese menor a los indicados \\
1 mes & Incidente de caducidad de instancia \\
\bottomrule
\end{tabularx}
\end{table}

\begin{example}{Plazo de un mes del incidente de caducidad (inc. 4°)}
Cuando la parte contraria acusa la caducidad de la instancia, fundamenta en el caso concreto por qué se encuentra configurada: indica cuál fue el último acto impulsor y que desde allí se encuentra vencido el plazo de seis meses en primera instancia sin que exista ningún acto procesal que tienda a impulsar el proceso. De esa presentación el juez da traslado a la otra parte, que la contesta, y recién después se resuelve, haciendo lugar o rechazando la caducidad.

El inciso 4° dispone que ese incidente ---el traslado, la contestación y la resolución--- también tiene un plazo de caducidad: \textbf{un mes}. Por ejemplo: una vez acusada la caducidad, el juez ordenó el traslado a la otra parte. Si quien acusó la caducidad no le notifica a la otra parte ---no le envía la cédula electrónica a los fines de notificar ese traslado---, entonces, si pasa un mes, quien debía contestar el traslado pero jamás fue notificado podría acusar la caducidad de ese incidente de caducidad de instancia.
\end{example}

% ----------------------------------------------------------
\section{Cómputo del plazo (art. 311, CPCCN)}

\begin{articulo}{Art. 311 --- CPCCN (Cómputo)}
Los plazos señalados en el artículo precedente se computarán desde la fecha de la última petición de las partes o resolución o actuación del juez, secretario u oficial primero, que tenga por efecto impulsar el procedimiento; correrán durante los días inhábiles salvo los que correspondan a las ferias judiciales.
\end{articulo}

El \textbf{último acto impulsor} es el último acto tomado como válido que significa un impulso procesal. No cualquier acto lo es: por ejemplo, la presentación para autorizar a una persona que trabaja en un estudio jurídico a visualizar el expediente no es un acto impulsor que pueda tomarse como válido en términos de impulso procesal. Debe tratarse de un acto que tenga concretamente por finalidad impulsar el proceso hacia las etapas subsiguientes, es decir, hacerlo avanzar y seguir procesando.

Se toma ese último acto válido que signifique impulso procesal y, a partir del \textbf{día siguiente}, se contabiliza el plazo. Los plazos del artículo 310 corren durante los días inhábiles, salvo los que correspondan a las ferias judiciales, según dispone el artículo 311.

\begin{important}
\textbf{Días inhábiles y ferias.} Los plazos de caducidad de instancia corren en \textbf{días corridos} (incluyendo feriados y días inhábiles), a diferencia de los plazos procesales ordinarios, que se cuentan en días hábiles. La única excepción: no se computan los días correspondientes a la feria de verano (enero) ni a la feria de invierno (julio, 15 días).
\end{important}

% ----------------------------------------------------------
\section{Litisconsorcio (art. 312, CPCCN)}

\begin{articulo}{Art. 312 --- CPCCN (Litisconsorcio)}
El impulso del procedimiento por uno de los litisconsortes beneficiará a los demás.
\end{articulo}

El impulso del procedimiento por uno de los litisconsortes se considera que beneficia a todos los demás, porque la \textbf{instancia es indivisible}: no puede haber caducidad de instancia para un litisconsorte y no para el otro. Por lo tanto, si uno impulsa, en términos del sistema dispositivo, el impulso de la instancia beneficia a todos los demás.

% ----------------------------------------------------------
\section{Supuestos de improcedencia (art. 313, CPCCN)}

\begin{articulo}{Art. 313 --- CPCCN (Improcedencia)}
No se producirá la caducidad:
\begin{enumerate}[label=\arabic*°]
    \item En los procedimientos de ejecución de sentencia, salvo si se tratare de incidentes que no guardaren relación estricta con la ejecución procesal forzada propiamente dicha.
    \item En los procesos sucesorios y, en general, en los voluntarios, también salvo los incidentes que en ellos se susciten.
    \item Cuando los procesos estuvieren pendientes de alguna resolución y la demora en dictarla fuere imputable al tribunal.
    \item Si se hubiere llamado autos para sentencia, salvo si se dispusiere prueba de oficio; cuando su producción dependiere de la actividad de las partes, la carga de impulsar el procedimiento existirá desde el momento en que éstas toman conocimiento de las medidas ordenadas.
\end{enumerate}
\end{articulo}

\begin{itemize}
    \item \textbf{Ejecución de sentencia:} no corresponde acusar la caducidad de instancia, salvo si se tratara de incidentes que no guarden relación estricta con la ejecución procesal forzada propiamente dicha.
    \item \textbf{Procesos sucesorios y voluntarios:} tampoco procede la caducidad en procesos sucesorios ni, en general, en cualquier proceso voluntario.
    \item \textbf{Demora imputable al tribunal:} cuando los procesos estuvieran pendientes de alguna resolución y la demora en dictarla fuera imputable al tribunal, la inactividad debe poder imputarse a la parte y no a una demora que corresponda al tribunal.
\end{itemize}

\begin{example}{Demora imputable al tribunal}
Ya se encuentra dictado el auto para sentencia y éste fue firmado, pero la sentencia no se dictó. No es algo que pueda imputarse a la parte, sino que corresponde al propio tribunal.
\end{example}

% ----------------------------------------------------------
\section{Legitimación y modo de operarse (arts. 314 a 316, CPCCN)}

% TODO: contenido incompleto o ambiguo en las notas originales (el encabezado de las notas menciona los arts. 314 a 316, pero solo se transcriben los arts. 315 y 316)
\begin{articulo}{Art. 315 --- CPCCN (Legitimación para pedir la declaración de caducidad)}
Sin perjuicio de lo dispuesto en el artículo siguiente, la declaración de caducidad podrá ser pedida en primera instancia por el demandado; en el incidente, por el contrario a quien lo hubiese promovido; en el recurso, por la parte recurrida. La petición deberá formularse antes de consentir el solicitante cualquier actuación del tribunal o de la parte posterior al vencimiento del plazo legal, y se sustanciará únicamente con un traslado de la parte contraria.
\end{articulo}

\begin{articulo}{Art. 316 --- CPCCN (Modo de operarse)}
La caducidad será declarada de oficio, sin otro trámite que la comprobación de los plazos señalados en el artículo 310, pero antes de que cualquiera de las partes impulse el procedimiento.
\end{articulo}

La caducidad puede pedirse en primera instancia por el \textbf{demandado}; en el incidente, por el \textbf{contrario a quien lo hubiese promovido}; y en el recurso, por la \textbf{parte recurrida}.

La petición de caducidad debe formularse \textbf{antes de consentir cualquier acto} del tribunal o de la otra parte que pueda considerarse impulsor. Si el plazo de seis meses venció pero, luego de ese vencimiento en primera instancia, se produjo un acto que fue consentido, ya no puede acusarse la caducidad por los seis meses transcurridos anteriormente. Es decir, la caducidad puede acusarse siempre que no se consienta ningún acto impulsor posterior al vencimiento, aunque ya se haya operado el vencimiento del plazo.

Asimismo, el pedido de caducidad en la segunda instancia importa el \textbf{desistimiento del recurso interpuesto por el peticionario}, en el caso de que aquel prospere.

\begin{example}{Caducidad de la segunda instancia con recursos de ambas partes}
Las dos partes interpusieron recursos de apelación contra la sentencia de primera instancia, y la parte demandada acusa la caducidad de la segunda instancia respecto del recurso que interpuso la parte actora. Si esa caducidad se considera configurada y se declara la caducidad de la segunda instancia, se tiene por desistido el recurso de apelación interpuesto por el demandado. Es decir, la Cámara no va a tratar ninguno de los dos recursos.
\end{example}

% ----------------------------------------------------------
\section{Efectos de la caducidad (art. 318, CPCCN)}

\begin{articulo}{Art. 318 --- CPCCN (Efectos)}
La caducidad operada en primera o única instancia no extingue la acción, la que podrá ejercitarse en un nuevo juicio, ni perjudica las pruebas producidas, las que podrán hacerse valer en aquel. La caducidad operada en instancias ulteriores acuerda fuerza de cosa juzgada a la resolución recurrida. La caducidad de la instancia principal comprende la reconvención y los incidentes; pero la de éstos no afecta la instancia principal.
\end{articulo}

\subsection{Caducidad en primera o única instancia}

La caducidad operada en primera o única instancia \textbf{no extingue la acción}: ésta puede ejercitarse en un nuevo juicio, y las pruebas producidas en el proceso que caducó no se perjudican, de modo que pueden hacerse valer en el juicio que se inicie después. Esto no es posible cuando el derecho ya está alcanzado por la \textbf{prescripción}, es decir, cuando la acción se encuentra prescripta.

\subsection{Caducidad en instancias ulteriores}

La caducidad operada en instancias ulteriores \textbf{acuerda fuerza de cosa juzgada a la resolución recurrida}. Si lo que caducó es la segunda instancia, la sentencia definitiva dictada por el juez de primera instancia queda con fuerza de sentencia definitiva pasada en autoridad de cosa juzgada: aunque se hayan interpuesto los recursos de apelación correspondientes, al declararse operada la caducidad en la segunda instancia queda firme la sentencia definitiva de primera instancia.

\subsection{Instancia principal, reconvención e incidentes}

La caducidad de la instancia principal \textbf{comprende la reconvención y los incidentes}: si caduca la instancia principal, se afecta todo lo tramitado (reconvenciones, incidentes, etcétera). En cambio, la caducidad de un incidente generado con motivo del proceso principal \textbf{no} genera que se configure la caducidad para el proceso principal.

% ----------------------------------------------------------

\cornelltitle{caducidad de la instancia}
\cornellblock{
\cqrow{¿Qué es la caducidad de la instancia y cuál es su fundamento?}{Es la extinción de la instancia por inactividad procesal de la parte que tiene la carga de impulsar el proceso, transcurridos los plazos del art. 310 CPCCN. Fundamentos: el interés público en no paralizar indefinidamente los procesos y la tutela del órgano jurisdiccional. Se presume el abandono de la instancia.
}
\cqrow{¿En qué procesos opera y qué la diferencia de la prescripción de la acción?}{Opera en los procesos dispositivos, en los que las partes tienen la carga de impulsar el proceso. La prescripción es un instituto de derecho sustancial que extingue el derecho o la acción; la caducidad es un instituto de derecho procesal que extingue la instancia y no cancela la pretensión.
}
\cqrow{¿Cuáles son los plazos de caducidad (art. 310 CPCCN)?}{6 meses en primera instancia. 3 meses en segunda o ulterior instancia, juicio sumarísimo, ejecutivo, ejecuciones especiales e incidentes. El plazo en que se opere la prescripción, si es menor. 1 mes para el incidente de caducidad de instancia.
}
\cqrow{¿Cómo se computan los plazos de caducidad?}{Desde el último acto impulsor del procedimiento (no cualquier acto, sino uno que tenga por finalidad hacer avanzar el proceso), a partir del día siguiente. Corren en días corridos (incluyendo inhábiles y feriados). Solo se excluyen los días de ferias judiciales (enero y julio). A diferencia de otros plazos procesales, que se computan en días hábiles.
}
\cqrow{¿Qué efecto tiene el impulso de un litisconsorte (art. 312)?}{El impulso del procedimiento por uno de los litisconsortes beneficia a los demás, porque la instancia es indivisible.
}
\cqrow{¿Cuándo no procede la caducidad (art. 313)?}{En ejecución de sentencia (salvo incidentes ajenos a ella), en procesos sucesorios y voluntarios (salvo incidentes), cuando la demora sea imputable al tribunal, y cuando se haya llamado autos para sentencia (salvo que se ordene prueba de oficio a cargo de las partes).
}
}
\medskip
\cornellblocksum{
\cqrow{¿Cuándo debe pedirse la caducidad y qué ocurre si se consiente un acto posterior (arts. 315 y 316)?}{Debe pedirse antes de consentir cualquier actuación del tribunal o de la otra parte posterior al vencimiento del plazo. Puede declararse de oficio, pero antes de que cualquiera de las partes impulse el procedimiento. Si se consiente un acto impulsor posterior al vencimiento, ya no puede acusarse la caducidad por el plazo anterior.
}
\cqrow{¿La caducidad extingue la acción?}{No. La caducidad en primera instancia no extingue la acción: puede iniciarse un nuevo proceso y se pueden hacer valer las pruebas producidas. Excepción: si el derecho ya está prescripto. La caducidad en instancias ulteriores sí acuerda fuerza de cosa juzgada a la resolución recurrida (la sentencia de primera instancia queda firme).
}
\cqrow{¿Qué relación hay entre la caducidad de la instancia principal y la de los incidentes?}{La caducidad de la instancia principal comprende la reconvención y los incidentes. Pero la caducidad de un incidente no afecta a la instancia principal.
}
}{La caducidad de la instancia (arts. 310--318 CPCCN) opera por inactividad de la parte que tiene la carga de impulsar el proceso dispositivo, en los plazos de 6 meses (primera instancia) o 3 meses (segunda instancia, sumarísimo, ejecutivo e incidentes), computados en días corridos con exclusión de las ferias judiciales. No extingue la acción en primera instancia (puede reiniciarse el proceso si el derecho no está prescripto), pero en instancias ulteriores acuerda fuerza de cosa juzgada a la resolución recurrida. Es el instituto con mayor casuística jurisprudencial, por la diversidad de situaciones en torno a qué constituye o no un acto impulsor del procedimiento.
}

\chapter{Resoluciones judiciales}\label{cap:resoluciones}

\section{Concepto y clasificación de las resoluciones judiciales}

El tema de las resoluciones judiciales es de suma importancia, aunque muchas veces no es atendido debidamente en la práctica por los tribunales, dado que no siempre se sigue a rajatabla lo que dice el código.

\begin{definition}{Resoluciones judiciales}
Cuando el Código Procesal habla de resoluciones judiciales hace referencia a las diversas formas en las que se pronuncia el tribunal. El tribunal puede realizar desde un acto sumamente complejo, como la \textbf{sentencia definitiva}, hasta un acto procesal sumamente sencillo, como las \textbf{providencias simples}.
\end{definition}

El código agrupa las resoluciones judiciales de acuerdo con su complejidad y con el objetivo que tiene cada una. La clasificación es la siguiente:

\begin{table}[H]
\centering
\small
\renewcommand{\arraystretch}{1.25}
\begin{tabularx}{\textwidth}{
    >{\raggedright\arraybackslash}p{2.6cm}
    >{\raggedright\arraybackslash}p{1.8cm}
    >{\raggedright\arraybackslash}p{2.3cm}
    >{\raggedright\arraybackslash}p{3.0cm}
    >{\raggedright\arraybackslash}X
}
\toprule
\textbf{Tipo de resolución} &
\textbf{Complejidad} &
\textbf{Plazo} &
\textbf{Firma} &
\textbf{Característica central} \\
\midrule
Providencias simples & Mínima & 3 días & Juez, secretario, pro-secretario o jefe de despacho & Mero trámite del proceso \\
Resoluciones interlocutorias & Media-alta & 10 días (1\textordfeminine{} inst.); 15 días (2\textordfeminine{} inst.) & Juez (exclusivamente) & Sustanciación / motivación de la decisión \\
Sentencias homologatorias & Media & 10 días & Juez (exclusivamente) & Modos anormales de terminación del proceso \\
Sentencia definitiva de 1\textordfeminine{} instancia & Máxima & Plazos del art.~34 CPCC & Juez (exclusivamente) & Pone fin normal al proceso \\
Sentencia definitiva de 2\textordfeminine{} instancia & Máxima & 15 días & Tres vocales & Acuerdo de cámara \\
\bottomrule
\end{tabularx}
\caption{Clasificación de las resoluciones judiciales.}
\end{table}
% TODO: contenido incompleto o ambiguo en las notas originales (el plazo de la sentencia definitiva de primera instancia se remite al art. 34 CPCC, que en otro pasaje de los apuntes se vincula con las providencias simples)

\begin{important}
Saber ante qué tipo de resolución se está es fundamental para el litigante, porque cada resolución judicial tiene determinados remedios o recursos que pueden interponerse contra ella en un plazo específico.
\end{important}

\begin{example}{La analogía del cañón y el mosquito}
No se puede matar un mosquito con un cañón, ni con una pistola ni con una ametralladora, porque lo más probable es que el mosquito se vaya volando. Con el recurso ocurre lo mismo: no se pueden atacar ciertos actos del proceso con un cañón cuando muchas veces lo que se necesita es una simple aclaratoria o una revocatoria. Y la revocatoria o reposición no puede interponerse contra cualquier tipo de resolución judicial. Para saber qué recursos están disponibles para atacar el acto, primero hay que saber qué tipo de acto se tiene delante.
\end{example}

Todas las resoluciones judiciales son actos procesales y, como tales, tienen características definidas: deben ser redactadas en idioma español, en tinta negra, fechadas y firmadas por la autoridad competente del tribunal.

\section{Providencias simples}

El artículo 34 del CPCC establece los plazos para el dictado de las providencias simples. El plazo del tribunal es de \textbf{tres días} para expedirse en aquellas cuestiones que hacen al mero trámite de la causa y que no requieren ningún tipo de complejidad. A partir del artículo 161 del CPCC se establecen las distintas resoluciones judiciales que existen.

Las \textbf{providencias simples} son aquellas que hacen al normal desarrollo y avance del proceso, y que no agregan nada relevante: solo hacen que el proceso siga avanzando con su trámite. Son ejemplos:

\begin{itemize}
    \item Agregar algún documento o prueba que ha sido ordenada.
    \item Tener por contestada la demanda.
    \item Correr una vista o un traslado.
    \item Ordenar una notificación.
    \item Autorizar a sacar fotocopias o a consultar la causa.
\end{itemize}

El Código autoriza al secretario del tribunal, al pro-secretario e incluso al jefe de despacho a firmar determinadas providencias simples, para disminuir la carga de trabajo del juez. En general, pueden firmar las vinculadas a la agregación o devolución de escritos, la devolución de escritos presentados fuera de término, la agregación de pruebas y documentos, y las vistas que deben darse al ministerio público.

Sin embargo, determinadas providencias simples deben ser suscriptas necesariamente por el juez: por ejemplo, el reconocimiento de la condición de parte, tener por contestada la demanda y aquellas que tienen una relevancia mayor en el proceso.

\begin{example}{El daño de una providencia simple mal resuelta}
Si, en lugar de proveerse <<téngase por contestada la demanda>>, el juez resuelve que la contestación fue presentada fuera de término (por considerarla extemporánea), se provoca un daño irreparable al litigante. Por eso es tan importante saber que se está ante una providencia simple: el remedio disponible es la \textbf{revocatoria}, que a veces es el único que puede solucionar ese problema.
\end{example}
% TODO: contenido incompleto o ambiguo en las notas originales (redacción confusa del ejemplo sobre rechazo / extemporaneidad de la contestación de demanda)

\section{Resoluciones interlocutorias}

Quizás lo más difícil para los estudiantes, e incluso para muchos abogados, es diferenciar una providencia simple de una resolución interlocutoria. La sentencia definitiva es clara, porque es la que pone fin de modo normal al proceso. Lo que a veces no resulta tan claro es cómo distinguir una providencia simple de una interlocutoria cuando esta última se ha simplificado demasiado por usos y costumbres.

El elemento diferenciador clave es la \textbf{sustanciación}.

\begin{important}
Cuando el Código Procesal indica que lo que caracteriza a las resoluciones interlocutorias es la <<sustanciación>>, no se refiere al traslado (uso forense coloquial), sino a la \textbf{sustancia}, al contenido de la resolución: la motivación y las razones que el juez tuvo que desarrollar para tomar su decisión.
\end{important}

Las resoluciones interlocutorias, por su propio nombre, \textbf{deciden algo}: pueden hacer lugar o rechazar una excepción, convocar o rechazar la convocatoria de un tercero, u ordenar o rechazar una medida cautelar. Para llegar a esa decisión, el juez debe motivar y dar razones.

La palabra <<interlocutorio>> proviene del latín \textit{inter}, que significa <<en medio del debate>>: son decisiones que se toman en medio de la discusión principal.

Los plazos para dictar estas resoluciones son:

\begin{itemize}
    \item Primera instancia: 10 días.
    \item Segunda instancia: 15 días.
\end{itemize}

\begin{example}{La medida cautelar confirma que <<sustanciación>> no es <<traslado>>}
Cuando una parte pide una medida cautelar no hay traslado previo. Sin embargo, la decisión que toma el juez no es una providencia simple: es una sentencia interlocutoria. Esto demuestra que el código no se refiere a la sustanciación como traslado, sino como sustancia, materia o motivación del acto. De otra manera, habría que concluir que una medida cautelar es una providencia simple, lo cual es incorrecto.
\end{example}

Hay resoluciones interlocutorias que pueden llegar a poner fin al proceso (como la sentencia que hace lugar a una excepción de prescripción o a una caducidad de la instancia), aunque lo hacen de un modo anormal, no del modo normal, que solo se alcanza mediante la sentencia definitiva.

\section{Sentencias homologatorias}

Las \textbf{sentencias homologatorias} son resoluciones vinculadas a determinados mecanismos que constituyen \textbf{modos anormales de terminación del proceso}: un allanamiento, un acuerdo conciliatorio o un desistimiento. El código prevé que se dicte una resolución que no llega a ser una sentencia definitiva, pero que tampoco es una providencia simple.

Tienen menos complejidad en su sustancia, pero no dejan de tenerla, porque el código le exige al juez que evalúe si el acto de la parte es correcto y si reúne las condiciones necesarias para lograr el desistimiento, el allanamiento o el acuerdo conciliatorio.

\begin{important}
Los acuerdos transaccionales homologados tienen efectos de cosa juzgada.
\end{important}

\section{Sentencia definitiva de primera instancia}

La sentencia definitiva es el acto más complejo que el código prevé dentro de las resoluciones judiciales. El artículo 163 del CPCC contiene nueve incisos que describen las distintas condiciones que debe reunir.

\subsection{Estructura formal}

Tanto la resolución interlocutoria como la sentencia definitiva suelen comenzar con la fecha y la palabra <<Autos>> o <<Autos y Vistos>>, seguida del resumen de los planteos de las partes, el <<Y considerando>> (donde el tribunal evalúa los planteos) y, finalmente, la parte resolutiva (el fallo propiamente dicho).

\subsection{Contenido del artículo 163 del CPCC}

\begin{formula}{Art. 163 CPCC --- Contenido de la sentencia definitiva de primera instancia}
El artículo 163 establece nueve incisos que describen las condiciones de la sentencia definitiva. Entre los más relevantes se destacan:
\begin{itemize}
    \item (Inc. 5\textordmasculine) Mencionar quiénes son las partes, sus nombres y apellidos, y la relación sucinta de las cuestiones planteadas.
    \item Los fundamentos y considerandos (motivación de la decisión).
    \item (Inc. 5\textordmasculine, 2.\textordmasculine{} párrafo) Posibilidad de fundar la decisión en \textbf{presunciones} no establecidas por la ley, que constituirán prueba cuando se funden en hechos reales y probados y cuando, por su número, precisión, gravedad y concordancia, produzcan convicción sobre la naturaleza del juicio, de conformidad con la regla de la sana crítica.
    \item (Inc. 6\textordmasculine) La decisión expresa, positiva y precisa, de conformidad con las pretensiones deducidas en el juicio (\textbf{principio de congruencia}).
    \item (Inc. 6\textordmasculine, \textit{in fine}) Mérito de los hechos constitutivos, modificativos o extintivos producidos durante la sustanciación del juicio y debidamente probados (\textbf{hechos sobrevinientes}).
    \item (Inc. 7\textordmasculine) El plazo que tiene la parte para cumplir la condena.
    \item (Inc. 8\textordmasculine) Establecer las costas y honorarios.
\end{itemize}
\end{formula}
% TODO: contenido incompleto o ambiguo en las notas originales (se mencionan nueve incisos pero solo se detallan algunos; el inc. 5.º aparece dos veces)

\subsection{El deber de motivación}

Es de gran importancia el deber de motivar las decisiones judiciales, tanto las resoluciones interlocutorias como las sentencias definitivas. No se trata solo de explicar por qué se llega a una resolución determinada.

\begin{important}
La motivación de la decisión sirve para que las partes y, de modo muy indirecto, el resto de la población controlen los actos de gobierno. Las sentencias y las decisiones judiciales son actos de gobierno: el Poder Judicial es uno de los poderes del Estado.
\end{important}

Además, la motivación es la base del \textbf{control recursivo}: para interponer una apelación contra una sentencia, el código exige una crítica concreta y razonada de la decisión que se quiere modificar. Esa crítica solo puede construirse sobre los considerandos, que son la parte donde el juez brinda las explicaciones de por qué llegó al resultado.

\subsection{Presunciones e indiciaria}

\begin{definition}{Presunción (art. 163 CPCC)}
La presunción no es un medio de prueba, sino una forma de dirimir la cuestión cuando el hecho principal no está probado pero se cuenta con una serie de hechos que, por su importancia, gravedad y concordancia, hacen que deba presumirse que aquel hecho no probado de forma directa ha sucedido. Es un mecanismo de razonamiento que la ley brinda al juez para poder dictar sentencia.
\end{definition}

El juez civil no puede mantener silencio y no dictar sentencia porque un hecho no esté probado: puede recurrir a los razonamientos de la prueba presuncional o indiciaria.

Sobre la conducta desplegada por las partes durante el proceso como indicio, véase el capítulo~\ref{cap:indiciaria}.

\subsection{La congruencia y los hechos sobrevinientes}

El inciso 6\textordmasculine{} del artículo 163 establece la decisión expresa, positiva y precisa de conformidad con las pretensiones deducidas en el juicio. Esto no es más ni menos que el \textbf{principio dispositivo} y la \textbf{congruencia}.

\begin{important}
El juez, al dictar la sentencia, debe hacerse cargo de las cuestiones que han sido introducidas. Pero esto también le marca un límite: no puede dar más de lo que se le pidió, ni algo distinto de lo que se le pidió. Hacerlo implicaría violar la congruencia y el derecho de defensa de la parte demandada.
\end{important}

\begin{example}{La incapacidad no reclamada}
En un proceso donde no se reclama el resarcimiento de determinada incapacidad, el juez no podría concederla aunque le parezca razonable. Si lo hiciera, violaría la congruencia, al otorgar más allá de lo solicitado por los litigantes, y además violaría el derecho de defensa de la parte demandada, que no pudo negar que tal incapacidad existió.
\end{example}

Respecto de los \textbf{hechos sobrevinientes}: en juicios que duran dos, tres o cuatro años pueden suceder hechos que modifiquen la relación jurídica objeto del proceso, la extingan, o que hagan que la realidad exterior exija que la decisión del juez se adecue a ese cambio.

\begin{example}{El cuadro de Picasso destruido}
Si en un proceso se reclama la restitución de un cuadro pintado por Picasso y durante el juicio ese cuadro deja de existir (porque se destruyó o desapareció), el juez no deberá condenar a su restitución en la sentencia, orden que no podría cumplirse. Deberá reemplazarla por una suma de dinero que compense el valor del bien que debía serle restituido a esa parte.
\end{example}

\subsection{Las costas}

El inciso 8\textordmasculine{} del artículo 163 establece que el juez debe, tanto en la sentencia definitiva como en las interlocutorias, establecer las costas y los honorarios.

\begin{definition}{Costas}
Las costas comprenden los gastos en que ha incurrido la parte para llevar adelante su reclamo: no solo los honorarios de su abogado, sino todo otro gasto en el que haya debido incurrir (sellados, obtención de documentos, etc.).
\end{definition}

La regla que establece el código para la condena en costas es la del \textbf{vencido}: la parte vencida en el proceso es la que debe hacerse cargo del pago de las costas. Si la demanda prospera, en general el vencido es el demandado; si la demanda es rechazada, el vencido es el actor.

Como toda regla, tiene excepciones. Una posibilidad es la \textbf{distribución en el orden causado}: cada parte paga sus propios gastos y los comunes, por mitades.

\begin{important}
Hay que ser muy cuidadoso con los argumentos para eximirse de costas. Si se dijera <<se creyó con derecho a litigar>> como fundamento, ese argumento es endeble: si la parte no se hubiera creído con derecho a litigar, ya estaría incurriendo en la sanción por temeridad y malicia. Esta posibilidad es bastante excepcional, especialmente en los casos en que se discuten derechos patrimoniales. En cambio, en cuestiones de familia, como los alimentos, es muy usual la distribución en el orden causado, para no afectar el monto del alimento que se reclama.
\end{important}

\cornelltitle{resoluciones judiciales}
\cornellblock{
\cqrow{¿Qué son las resoluciones judiciales y cuál es su clasificación?}{Son las diversas formas en que se pronuncia el tribunal. Se clasifican en: providencias simples (mero trámite, 3 días; pueden firmarlas el secretario, el pro-secretario o el jefe de despacho); resoluciones interlocutorias (deciden cuestiones con motivación, 10 días en primera instancia, solo el juez); sentencias homologatorias (modos anormales de terminación: allanamiento, desistimiento, conciliación); y sentencias definitivas de primera y segunda instancia (el acto más complejo; art.~163 CPCC).
}
\cqrow{¿Cómo se diferencia una providencia simple de una resolución interlocutoria?}{Por la sustanciación, entendida como <<sustancia>> (contenido y motivación) y no como <<traslado>>. La providencia simple solo impulsa el trámite; la interlocutoria requiere que el juez funde su decisión porque está resolviendo algo (admisión o rechazo de una excepción, medida cautelar, etc.). La medida cautelar, que no tiene traslado previo y es interlocutoria, confirma que la sustanciación no equivale a traslado.
}
\cqrow{¿Por qué es importante saber ante qué tipo de resolución se está?}{Porque cada resolución judicial tiene determinados remedios o recursos disponibles, en plazos específicos. No se puede <<matar un mosquito con un cañón>>: a veces se necesita una simple aclaratoria o una revocatoria, y la revocatoria no procede contra cualquier resolución.
}
\cqrow{¿Qué establece el art.~163 CPCC sobre la sentencia definitiva?}{Establece nueve incisos. Entre ellos: mención de las partes y relación sucinta de las cuestiones planteadas; fundamentos; posibilidad de usar presunciones (prueba indiciaria) basadas en hechos reales y probados que, por su número, precisión, gravedad y concordancia, produzcan convicción según la sana crítica; congruencia (decisión expresa, positiva y precisa conforme a las pretensiones deducidas); hechos sobrevinientes debidamente probados; plazo para cumplir la condena; y costas y honorarios.
}
}
\medskip
\cornellblocksum{
\cqrow{¿Para qué sirve el deber de motivación?}{Para que las partes y, de modo muy indirecto, el resto de la población controlen los actos de gobierno, ya que el Poder Judicial es uno de los poderes del Estado. Además, es la base del control recursivo: la crítica concreta y razonada que exige la apelación solo puede construirse sobre los considerandos.
}
\cqrow{¿Qué es una presunción según el art.~163 CPCC?}{No es un medio de prueba, sino una forma de dirimir la cuestión cuando el hecho principal no está probado pero hay hechos que, por su importancia, gravedad y concordancia, hacen presumir que ocurrió. Es un mecanismo de razonamiento que permite al juez dictar sentencia; el juez civil no puede dejar de sentenciar porque un hecho no esté probado.
}
\cqrow{¿Qué es el principio de congruencia?}{El juez debe hacerse cargo de las cuestiones introducidas en el proceso, pero no puede dar más de lo pedido ni algo distinto. Superar ese límite viola la congruencia y el derecho de defensa del demandado, que no pudo contradecir lo que no se le reclamaba.
}
\cqrow{¿Qué son los hechos sobrevinientes?}{Son los hechos que suceden durante el transcurso del proceso y que modifican o extinguen la relación jurídica objeto del litigio, o que exigen que la decisión se adecue a la realidad exterior. El juez debe ajustar la sentencia a esa nueva realidad (por ejemplo, si el bien reclamado desaparece, se reemplaza su restitución por una suma de dinero equivalente).
}
\cqrow{¿Cuál es la regla en materia de costas y cuáles son sus excepciones?}{Las costas comprenden todos los gastos en que incurrió la parte para sostener su reclamo. La regla es la del vencido. Una excepción es la distribución en el orden causado (cada parte paga sus gastos y los comunes por mitades), muy usual en cuestiones de familia como los alimentos; eximirse de costas alegando <<se creyó con derecho a litigar>> es un argumento endeble.
}
}{Las resoluciones judiciales (providencias simples, interlocutorias, homologatorias y sentencias definitivas) se distinguen por su complejidad, plazo, firmante y sustancia, y de esa clasificación depende el recurso procedente. Las interlocutorias y las sentencias exigen motivación, que permite el control de las partes y de la población. La sentencia definitiva debe ajustarse al art.~163 CPCC: puede apoyarse en presunciones, debe ser congruente, atender a los hechos sobrevinientes y resolver sobre costas y honorarios, con la regla del vencido como pauta general.
}

\chapter[Segunda instancia y actuación posterior]{Segunda instancia y actuación posterior a la sentencia}\label{cap:segunda}

%--------------------------------------------------------------

\section{Sentencia definitiva de segunda instancia}

En la cámara (segunda instancia) hay diferencias en el modo en que deben elaborarse las resoluciones. Los recursos que den lugar a resoluciones interlocutorias, firmadas por los tres vocales de cada sala, deben estar redactados de modo impersonal (en persona gramatical impersonal), aunque con estructura similar a la de las resoluciones interlocutorias de primera instancia.

La \textbf{sentencia definitiva de segunda instancia} es un acto complejo, elaborado por un \textbf{acuerdo} de tres jueces: hay un \textbf{vocal preopinante} y dos vocales que pueden adherir al voto del primero o marcar sus diferencias (disidencias). El fallo puede salir dos a uno, y no siempre el vocal preopinante tendrá el voto que determine el sentido de la decisión.

\begin{table}[H]
\centering
\small
\renewcommand{\arraystretch}{1.25}
\begin{tabularx}{\textwidth}{
    >{\raggedright\arraybackslash}p{2.4cm}
    >{\raggedright\arraybackslash}X
    >{\raggedright\arraybackslash}X
}
\toprule
\textbf{Aspecto} &
\textbf{Resoluciones interlocutorias en 2\textordfeminine{} instancia} &
\textbf{Sentencia definitiva en 2\textordfeminine{} instancia} \\
\midrule
Redacción & Impersonal & Primera persona del singular (cada vocal emite su voto) \\
Firmantes & Tres vocales de la sala & Tres vocales (acuerdo) \\
Estructura & Similar a las interlocutorias de 1\textordfeminine{} instancia & Similar a la sentencia definitiva de 1\textordfeminine{} instancia, pero en votos individuales \\
El <<acuerdo>> & No aplica como tal & Reunión teórica de los tres jueces para discutir y decidir \\
\bottomrule
\end{tabularx}
\caption{Resoluciones de segunda instancia.}
\end{table}

El acuerdo es, en teoría, una reunión de los tres jueces para discutir las decisiones e intercambiar opiniones. En la práctica, sin embargo, no se lleva a cabo en la mayoría de los casos: lo que se hace en las cámaras es trabajar con la circulación escrita de los expedientes.

\begin{important}
No puede dictarse una sentencia de segunda instancia si uno de los jueces no está presente. Pueden verse dos firmas siempre que el tercer juez esté en uso de licencia concedida o esté vacante la vocalía, pero las cámaras no tienen la misma facultad que la Corte federal de dictar decisiones cuando se alcanza la mayoría aunque los otros jueces no las suscriban.
\end{important}

Una vez dictada la sentencia, sea de primera o de segunda instancia, el código establece que esto concluye la competencia del juez o del tribunal para expedirse sobre el caso. El juez ya no puede modificar el fallo, estén o no notificadas las partes.

\section{Actuación del juez con posterioridad a la sentencia}

El artículo 166 del CPCC establece cuál puede ser la actuación posterior del juez una vez dictada la sentencia. Se prevé la posibilidad de enmendar errores puramente numéricos o errores evidentes (\textbf{errores materiales}), pero en ningún caso el tribunal o el juez puede modificar lo sustancial de las decisiones ya tomadas.

\begin{example}{Fallo de la Corte federal sobre el silencio en costas}
Durante mucho tiempo se sostuvo que el silencio sobre costas en una sentencia (es decir, que el tribunal omitiera pronunciarse sobre ellas) implicaba que habían sido fijadas en el orden causado. Esto fue aclarado en un fallo de la Corte federal, según el cual al silencio no puede otorgársele ningún sentido. Lo que debe hacer el tribunal es expedirse sobre un tema sobre el que aún no se había expedido: la condena en costas.
\end{example}

\section{El recurso de aclaratoria}

El segundo párrafo del artículo 166 habilita, a pedido de parte, el denominado \textbf{recurso de aclaratoria}: dentro de los tres días y sin sustanciación pueden corregirse errores materiales y aclararse conceptos oscuros sin alterar lo sustancial de la decisión. El plazo para solicitarla es de tres días en primera instancia y de cinco días en segunda instancia. Su régimen completo se desarrolla en el capítulo~\ref{cap:remedios}.

\cornelltitle{segunda instancia y actuación posterior}
\cornellblocksum{
\cqrow{¿Cómo funciona la sentencia definitiva de segunda instancia?}{Es elaborada por un acuerdo de tres jueces: un vocal preopinante y dos vocales que adhieren o disienten. Puede salir dos a uno. La redacción es en primera persona del singular, pues cada vocal emite su voto individual. Las resoluciones interlocutorias de cámara, en cambio, se redactan de modo impersonal.
}
\cqrow{¿Puede dictarse sentencia de segunda instancia sin la presencia de los tres jueces?}{No, salvo que el tercer juez esté en uso de licencia concedida o la vocalía esté vacante. Las cámaras no tienen la facultad de la Corte federal de dictar decisiones al alcanzarse la mayoría aunque los otros jueces no las suscriban.
}
\cqrow{¿Qué efecto tiene el dictado de la sentencia sobre la competencia del juez?}{Concluye la competencia del juez o del tribunal para expedirse sobre el caso; ya no puede modificar el fallo, estén o no notificadas las partes.
}
\cqrow{¿Qué puede hacer el juez una vez dictada la sentencia (art.~166 CPCC)?}{Solo puede enmendar errores puramente numéricos o errores materiales evidentes, y aclarar conceptos oscuros, sin alterar lo sustancial de lo decidido.
}
\cqrow{¿Qué es el recurso de aclaratoria y en qué plazos procede?}{Es el pedido de la parte para que se corrijan errores materiales y se aclaren conceptos oscuros, dentro de los tres días y sin sustanciación según el art.~166. El plazo para solicitarla es de tres días en primera instancia y cinco días en segunda instancia. No puede modificar lo sustancial de la decisión.
}
\cqrow{¿Qué sentido tiene el silencio sobre costas, según el fallo de la Corte federal?}{Al silencio no puede otorgársele ningún sentido (no implica orden causado). El tribunal debe expedirse sobre un tema sobre el que aún no se había expedido: la condena en costas.
}
}{La sentencia definitiva de segunda instancia es el resultado de un acuerdo de tres jueces (vocal preopinante, adhesiones y disidencias), redactada en primera persona del singular. Una vez dictada la sentencia, concluye la competencia del juez o tribunal: el art.~166 CPCC solo permite corregir errores materiales y aclarar conceptos oscuros, sin alterar lo sustancial, mediante la aclaratoria (3 días en primera instancia, 5 en segunda).
}

\chapter{Cosa juzgada}\label{cap:cosajuzgada}

%--------------------------------------------------------------

\section{Concepto e inmutabilidad}

El concepto de cosa juzgada ya fue visto al estudiar las excepciones procesales en el artículo 347; aquí se profundizan algunos aspectos.

\begin{definition}{Cosa juzgada}
Implica que, contra una decisión de un tribunal de primera o de segunda instancia, no hay más recursos disponibles para interponer, o bien el plazo para apelarla ha transcurrido sin que las partes hayan deducido ninguno. En esas condiciones, la decisión adquiere la \textbf{autoridad de cosa juzgada}.
\end{definition}

\begin{important}
La cosa juzgada significa que dentro del patrimonio de la persona ingresa un activo: el derecho a hacer cumplir lo que emana de esa decisión judicial. La realidad previa deja de existir, y lo único que queda es lo que emana de la sentencia.
\end{important}

El abogado debe controlar muy celosamente lo que surge de la decisión, especialmente lo que el tribunal escribe o asienta después de la palabra <<fallo>>, porque eso es lo que en el futuro podrá exigirse que se cumpla con autoridad de cosa juzgada.

Cuando una sentencia de segunda instancia modifica parcialmente la decisión de primera instancia, la cosa juzgada emana de la \textbf{sumatoria de los dos fallos}: las cuestiones que quedaron firmes por no haber sido apeladas en la sentencia de primera instancia se suman a las que emanan de la sentencia firme de segunda instancia.

\section{Cosa juzgada formal y cosa juzgada material}

No todos los procesos tienen la misma característica en cuanto a la calidad de la cosa juzgada. La doctrina diferencia la \textbf{cosa juzgada formal} de la \textbf{cosa juzgada material}.

\begin{table}[H]
\centering
\small
\renewcommand{\arraystretch}{1.25}
\begin{tabularx}{\textwidth}{
    >{\raggedright\arraybackslash}p{2.6cm}
    >{\raggedright\arraybackslash}p{3.0cm}
    >{\raggedright\arraybackslash}X
    >{\raggedright\arraybackslash}X
}
\toprule
\textbf{Tipo} &
\textbf{Proceso típico} &
\textbf{Qué cuestiones abarca} &
\textbf{Posibilidad de juicio posterior} \\
\midrule
Cosa juzgada material & Proceso ordinario (conocimiento pleno) & Todas las cuestiones debatidas extensamente & No (el fondo fue juzgado definitivamente) \\
Cosa juzgada formal & Juicio ejecutivo / caducidad de instancia & Solo las cuestiones discutibles en ese proceso & Sí, en juicio ordinario posterior (sobre lo no juzgado) \\
\bottomrule
\end{tabularx}
\caption{Cosa juzgada material y formal.}
\end{table}

\begin{example}{El juicio ejecutivo}
En un juicio ejecutivo hay cosa juzgada material si se discutió el pago, porque el pago es una excepción que podía oponerse en ese proceso. En cambio, no habrá cosa juzgada material en relación, por ejemplo, con algún vicio de la voluntad que se hubiera tenido al suscribir el título ejecutivo, porque esa cuestión no podía discutirse en el juicio ejecutivo. Respecto de ese aspecto hay cosa juzgada formal, y las cuestiones cuya discusión no puede darse en el juicio ejecutivo podrán discutirse en el juicio ordinario posterior.
\end{example}

\section{Límites subjetivos de la cosa juzgada}

El inciso 6\textordmasculine{} del artículo 347 del CPCC establece la \textbf{excepción de cosa juzgada}. Lo que ataca la excepción es el derecho de la persona que reclama: ese derecho ya fue discutido y decidido en una sentencia previa que está firme, contra la cual ya no hay más recursos para interponer.
% TODO: contenido incompleto o ambiguo en las notas originales (el apunte dice «dividido en una sentencia previa»; se transcribe como «decidido» por el sentido del pasaje)

Los \textbf{límites subjetivos} indican que la cosa juzgada abarca a las personas que pudieron discutir en el marco del proceso previo: actor, demandado, terceros intervinientes y terceros legitimados sobre los cuales se ha discutido una relación sustancial. No puede afectar a sujetos que no han participado de ese proceso, porque no han tenido la posibilidad de defenderse, ofrecer prueba y contradecir.

Esto se vincula con el \textbf{litisconsorcio}:

\begin{important}
En tanto se trate de un litisconsorcio facultativo, nada impide que el juicio llevado adelante produzca una cosa juzgada material que afecte solo a las partes que han participado, sin relevancia respecto de otros sujetos que podrían haber sido litisconsortes facultativos y que no fueron citados.
\end{important}

\section{Límites objetivos de la cosa juzgada}

Los límites objetivos son más difusos. Nadie duda de que la cuestión debatida en el primer proceso es la que ha sido juzgada, pero la lectura del inciso 6\textordmasculine{} del artículo 347 muestra que el alcance es mucho más expansivo de lo que una primera aproximación indicaría.

\begin{formula}{Art. 347, inc. 6\textordmasculine, CPCC --- Límites objetivos}
Para que sea procedente la excepción de cosa juzgada, el examen integral de las dos contiendas debe demostrar que se trata del \textbf{mismo asunto} sometido a decisión judicial. Pero, además, hay cosa juzgada cuando, por existir \textbf{continencia, conexidad, accesoriedad o subsidiariedad}, la sentencia firme ya haya resuelto lo que constituye la materia o la pretensión deducida en el nuevo juicio que se promueve.
\end{formula}

Esto significa que los límites objetivos de la cosa juzgada no solo comprenden lo debatido en el primer proceso, sino también las cuestiones de continencia, conexidad, accesoriedad y subsidiariedad que podrían haberse planteado en él y no se plantearon.

\begin{example}{El accidente de autos y el lucro cesante}
Si en un juicio por un accidente de autos se reclama solo la reparación del automóvil y se resulta ganador, cabe preguntarse si podría iniciarse un segundo proceso reclamando el lucro cesante (porque el auto era utilizado como taxi o remís). A partir del artículo 347, inc.~6\textordmasculine, se podría concluir que existió lo que se denomina \textbf{cosa juzgada implícita}: el lucro cesante era un rubro que podía pedirse en el primer proceso y no se hizo, por lo que precluyó también la posibilidad de hacerlo con posterioridad, al ser una cuestión accesoria o conexa con la discutida en el primer litigio.
\end{example}

Es importante conocer esto porque la cosa juzgada \textbf{puede ser declarada de oficio por el juez}. Puede suceder que un juez no la advierta y que el demandado tampoco oponga la excepción. En el proceso le va mejor al que más sabe y tiene más herramientas para plantear la discusión de estos supuestos.

\cornelltitle{cosa juzgada}
\cornellblock{
\cqrow{¿Qué es la cosa juzgada?}{Es la autoridad que adquiere una decisión judicial cuando contra ella no hay más recursos disponibles o el plazo para interponerlos ha vencido sin que las partes los dedujeran. Produce dos efectos: (1) ingresa al patrimonio del vencedor el derecho a exigir el cumplimiento; (2) la realidad previa deja de existir y solo queda lo que emana de la sentencia (novación procesal). Los acuerdos transaccionales homologados también tienen efectos de cosa juzgada.
}
\cqrow{¿De dónde emana la cosa juzgada cuando la sentencia de segunda instancia modifica parcialmente la de primera?}{De la sumatoria de ambos fallos: las cuestiones que quedaron firmes por no haber sido apeladas en primera instancia se suman a las que emanan de la sentencia firme de segunda instancia. Por eso el abogado debe controlar lo asentado después de la palabra <<fallo>>.
}
\cqrow{¿Cuál es la diferencia entre cosa juzgada formal y material?}{La cosa juzgada material corresponde al proceso ordinario (conocimiento pleno): abarca todas las cuestiones debatidas y no permite un juicio posterior sobre el fondo. La cosa juzgada formal corresponde, por ejemplo, al juicio ejecutivo o a la caducidad de instancia: abarca solo las cuestiones discutibles en ese proceso y permite un juicio ordinario posterior sobre lo no juzgado (por ejemplo, vicios de la voluntad en el juicio ejecutivo).
}
\cqrow{¿Cuáles son los límites subjetivos de la cosa juzgada?}{Solo alcanza a quienes fueron partes del proceso (actor, demandado, terceros intervinientes, terceros legitimados). No puede afectar a quienes no tuvieron la posibilidad de defenderse, ofrecer prueba y contradecir. Se vincula con el litisconsorcio facultativo: la cosa juzgada material afecta solo a las partes que participaron.
}
\cqrow{¿Cuáles son los límites objetivos de la cosa juzgada (art.~347, inc.~6\textordmasculine)?}{No solo las cuestiones debatidas, sino también las conexas, accesorias, subsidiarias o por continencia que pudieron plantearse y no se plantearon. Puede haber cosa juzgada <<implícita>> respecto de rubros accesorios no reclamados (por ejemplo, el lucro cesante tras un juicio por la reparación del automóvil).
}
}
\medskip
\cornellblocksum{
\cqrow{¿Por qué importa que la cosa juzgada pueda declararse de oficio?}{Porque puede ocurrir que el juez no la advierta y que el demandado no oponga la excepción; en el proceso le va mejor a quien más sabe y tiene más herramientas para plantear estos supuestos.
}
}{Una decisión adquiere autoridad de cosa juzgada cuando ya no hay recursos o ha vencido el plazo para interponerlos; el derecho a su cumplimiento ingresa al patrimonio del vencedor. Puede ser material (proceso ordinario, no permite un juicio posterior sobre el fondo) o formal (limitada a lo discutible en ese proceso, como en el juicio ejecutivo). Sus límites subjetivos alcanzan solo a quienes fueron partes; los objetivos, definidos en el art.~347, inc.~6\textordmasculine, se extienden a cuestiones conexas, accesorias, subsidiarias o por continencia que pudieron plantearse.
}

\chapter[Sentencia civil y sentencia penal]{Relación entre sentencia civil y sentencia penal}\label{cap:civilpenal}

%--------------------------------------------------------------

\section{Prelación de la acción penal sobre la civil}

Vinculado con la cosa juzgada está lo que usualmente se denomina la incidencia de la acción penal sobre la acción civil, regulada en el Código Civil y Comercial (CCyC) a partir del artículo 1774.

\begin{formula}{Art. 1774 y ss. CCyC --- Prelación de la acción penal}
Si la acción penal precede a la acción civil o es intentada durante su curso, el dictado de la sentencia definitiva debe suspenderse en el proceso civil hasta que haya conclusión del proceso penal.

Excepciones a la suspensión:
\begin{itemize}
    \item Si median causas de extinción de la acción penal.
    \item Si la dilación del procedimiento penal provoca en los hechos una frustración efectiva del derecho a ser indemnizado en el juicio civil.
    \item Si la acción civil por reparación del daño está fundada en un factor objetivo de responsabilidad (gran cambio del CCyC: en estos casos nada impide que la sentencia civil se dicte aun cuando no se haya dictado sentencia penal).
\end{itemize}
\end{formula}

La razón de ser de esta coordinación es evitar lo que en general se encuentra en los libros como el <<escándalo jurídico>>, que en definitiva es un escándalo lógico.

\begin{important}
El <<escándalo jurídico>> (o lógico) se produce cuando una sentencia en un fuero dice que el autor del hecho fue una persona y otra sentencia dice que esa persona no fue la autora. La coordinación de los dos ordenamientos busca evitar precisamente esta colisión de decisiones.
\end{important}

\section{Influencia de la sentencia penal sobre la civil}

El Código Civil y Comercial dispone que la inexistencia del hecho, la autoría y la ausencia de dolo penal de responsabilidad determinadas en la sentencia penal no pueden desconocerse en la sentencia civil.

\begin{important}
Si la sentencia penal dice que el sindicado como responsable penal no participó, que no es autor del hecho o que el hecho no ocurrió, esto no podrá ser desconocido en la sentencia civil, que deberá fundarse en consecuencia.
\end{important}

Es posible que una absolución penal por ausencia de delito (porque el accionar no está tipificado en el Código Penal) dé, sin embargo, lugar a la reparación civil. En esos casos no habrá ninguna influencia de las consecuencias penales sobre las civiles.

\section{Sentencia penal posterior a la civil}

Puede ocurrir que, en la acción civil, se demande con base en un factor de atribución objetivo, se permita que la acción civil continúe y se dicte una sentencia civil a la que con posterioridad se agregue una sentencia en sede penal.

\begin{formula}{Art. 1780 CCyC --- Sentencia penal posterior a la sentencia civil}
La sentencia penal posterior a la sentencia civil no produce ningún efecto sobre ella, excepto en el caso de revisión. La revisión procede cuando así se solicite y cuando la sentencia civil asigna alcances de cosa juzgada a cuestiones resueltas por la sentencia penal y esta es revisada respecto de esas cuestiones.
\end{formula}

\begin{example}{Absolución penal posterior por inexistencia del hecho}
Si se dictó una sentencia civil condenando al demandado sobre la base de un factor de atribución objetivo (art.~1774, inc.~c, CCyC) y con posterioridad, en la acción penal, el demandado es absuelto por inexistencia del hecho en que se funda esa condena civil, o por no ser su autor, procede la revisión de la condena civil. Aunque la sentencia penal llegó con posterioridad, afirma que el demandado condenado en sede civil no fue autor del hecho, de modo que de ninguna manera podría quedar condenado.
\end{example}

\cornelltitle{sentencia civil y sentencia penal}
\cornellblocksum{
\cqrow{¿En qué consiste la prelación de la acción penal sobre la civil?}{Según el art.~1774 CCyC, si la acción penal precede a la civil o se intenta durante su curso, la sentencia civil debe suspenderse hasta la conclusión del proceso penal, salvo: (a) extinción de la acción penal; (b) frustración efectiva del derecho a ser indemnizado por la dilación del proceso penal; (c) factor objetivo de responsabilidad (gran innovación del CCyC: en ese caso la sentencia civil puede dictarse sin esperar la penal).
}
\cqrow{¿Qué es el <<escándalo jurídico>> y cómo se evita?}{Es un escándalo lógico: ocurre cuando una sentencia dice que el autor del hecho fue una persona y otra sentencia dice que esa persona no fue la autora. La coordinación de los ordenamientos civil y penal busca evitar esa colisión de decisiones.
}
\cqrow{¿Qué efectos produce la sentencia penal anterior sobre la civil?}{La sentencia penal vincula a la civil: si determina la inexistencia del hecho, la falta de autoría o la ausencia de dolo, esas conclusiones no pueden desconocerse en la sentencia civil. Una absolución penal por ausencia de delito (accionar no tipificado) puede, en cambio, dar lugar a la reparación civil sin influencia de las consecuencias penales.
}
\cqrow{¿Qué efectos produce la sentencia penal posterior a la civil (art.~1780 CCyC)?}{En principio, ninguno sobre la sentencia civil ya dictada, salvo en caso de revisión: procede cuando se solicita y la sentencia civil asignó alcances de cosa juzgada a cuestiones resueltas por la sentencia penal que luego son revisadas (por ejemplo, absolución penal por inexistencia del hecho en que se fundó la condena civil).
}
}{La acción penal tiene prelación sobre la civil: la sentencia civil se suspende hasta la conclusión del proceso penal, salvo extinción de la acción penal, frustración del derecho a ser indemnizado por dilación, o factor objetivo de responsabilidad. La sentencia penal anterior no puede desconocerse en lo relativo a la inexistencia del hecho, la autoría y la ausencia de dolo; la posterior no afecta a la civil, salvo revisión. Todo ello busca evitar el escándalo lógico de fallos contradictorios.
}

\section{Síntesis de la sentencia, la cosa juzgada y su coordinación con el proceso penal}
Las resoluciones judiciales constituyen el lenguaje formal a través del cual el tribunal se comunica con las partes y hace avanzar el proceso. Su clasificación (providencias simples, resoluciones interlocutorias, sentencias homologatorias y sentencias definitivas) no es meramente académica: determina qué recursos pueden interponerse, en qué plazo y ante quién. El deber de motivación no solo garantiza el control de las partes sino también el control democrático de los actos del Poder Judicial.

La sentencia definitiva, regulada en el art.~163 CPCC, debe ser expresa, positiva, precisa y congruente con las pretensiones deducidas, pudiendo el juez valerse de presunciones indiciarias y de la conducta procesal de las partes para forjar su convicción.

Una vez firme, la sentencia adquiere autoridad de cosa juzgada, que puede ser material (definitiva) o formal (limitada a lo discutible en ese proceso), y cuyos límites, objetivos y subjetivos, son más amplios de lo que una primera lectura sugiere, pues alcanzan cuestiones conexas, accesorias y subsidiarias (art.~347, inc.~6\textordmasculine).

La coexistencia del proceso civil y el penal exige una coordinación normativa que, con el Código Civil y Comercial, admite la continuación de la acción civil frente a factores objetivos de responsabilidad sin necesidad de esperar la conclusión del proceso penal, buscando en todos los casos evitar el <<escándalo lógico>> de fallos contradictorios.

\part{Impugnación y recursos}
\markboth{Parte III}{Parte III}
La etapa de impugnación sigue cronológicamente a la etapa de conclusión del proceso, a las resoluciones procesales y a las sentencias. Esta parte desarrolla la \textbf{teoría general de la impugnación}, es decir, el sistema mediante el cual el ordenamiento procesal civil argentino permite atacar las decisiones judiciales erróneas o los actos procesales viciados. Comienza con los fundamentos y las clasificaciones de los actos de impugnación y de los recursos, y con los requisitos que son comunes a todos los recursos. Continúa con la garantía de la doble instancia y los límites del tribunal de alzada, con los remedios que se resuelven ante el mismo juez (aclaratoria y reposición o revocatoria) y con el incidente de nulidad. Finalmente se estudia en detalle el recurso de apelación: su concepto, su concesión, su trámite, sus efectos y las reglas que lo rigen.

\chapter[Teoría general de la impugnación]{Teoría general de la impugnación y clasificación de los recursos}\label{cap:impugnacion}

\section{Ubicación del tema y relación entre impugnación y recurso}

La etapa de impugnación sigue cronológicamente a la etapa de conclusión del proceso, las resoluciones procesales y las sentencias.

Los actos de impugnación y los recursos guardan una relación de \textbf{género y especie}: los \textbf{recursos} son un tipo de acto de impugnación. Existen otros tipos de actos de impugnación, que se desarrollan al analizar las clasificaciones posibles.

La voz \textbf{impugnaci\'on} es un t\'ermino mucho m\'as amplio que el de \textbf{recurso}. Seg\'un Hitter, la impugnaci\'on corresponde m\'as a la Teor\'ia General del Derecho, mientras que el t\'ermino recurso es propio del derecho procesal. La relaci\'on entre ambos es de \textbf{g\'enero a especie}: la impugnaci\'on es el g\'enero y el recurso es una especie de ese g\'enero.

Impugnar quiere decir \textbf{atacar}. Pueden existir actos de impugnaci\'on fuera del proceso judicial y, dentro del proceso, existen distintos tipos de impugnaciones que no siempre son recursos.

\begin{example}{Impugnaci\'on de dictamen pericial}
En la etapa probatoria, si se presenta un dictamen pericial m\'edico y se corre traslado a las partes, \'estas pueden atacarlo mediante un acto procesal escrito de impugnaci\'on, dando los fundamentos de por qu\'e se lo ataca. Ese acto es claramente una impugnaci\'on, pero no es un recurso.
\end{example}

\begin{example}{Impugnaci\'on de liquidaci\'on en ejecuci\'on}
En la etapa de ejecuci\'on del proceso, si se impugna una liquidaci\'on del capital de condena m\'as intereses por entender que los intereses est\'an mal calculados, ese acto tambi\'en es una impugnaci\'on, pero no un recurso.
\end{example}

\begin{important}
S\'olo se est\'a ante un \textbf{recurso} cuando lo que se impugna es una \textbf{resoluci\'on judicial}. Los recursos son mecanismos de impugnaci\'on \textbf{exclusivamente de resoluciones judiciales}: providencias simples, resoluciones interlocutorias y sentencias definitivas (reguladas a partir del art.~160 del C\'odigo Procesal).
\end{important}

\section{Etimología de la palabra «impugnación»}

La palabra \textbf{impugnación} proviene del latín \textit{pugnus}, que significa «puño». De \textit{pugnus} derivaron varias palabras que aluden a la lucha de puños: el pugilato (enfrentamiento a puñetazos), el púgil (el boxeador), la pelea. La noción que subyace a la impugnación es, por lo tanto, la de \textbf{combate, lucha o ataque}.

Aplicado al proceso judicial, cuando se habla de una teoría de la impugnación o de actos de impugnación se hace referencia a todo lo que tiene que ver con atacar o combatir un procedimiento viciado o una decisión equivocada desde la perspectiva de quien impugna. La impugnación es, entonces, todo mecanismo procesal que ataca o combate cualquier acto procesal o procedimiento que lesione el interés de cualquiera de las partes o de un tercero.

\section{Fundamento de la impugnación: el error humano}

El ordenamiento procesal contempla mecanismos de impugnación porque \textbf{existe el error}. Los jueces pueden equivocarse y cometer errores. Del mismo modo, las partes, en los actos procesales que impulsan, pueden cometer irregularidades y realizarlos sin cumplir adecuadamente con las formas. Estos errores de procedimiento o de interpretación de la ley son inevitables, porque están dentro de la naturaleza humana. Los actos de impugnación constituyen la alternativa que busca mitigar los defectos y deficiencias que pueda haber en un proceso o en una resolución judicial.

\textbf{Calamandrei} sosten\'ia que el fundamento de los recursos es, ni m\'as ni menos, el \textbf{error humano}: los jueces son humanos y pueden equivocarse; si se equivocan, debe existir un sistema para subsanar ese error y evitar que se consolide una injusticia.

\begin{note}
\textbf{San Agustín y Cicerón.} La frase «errar es humano» se atribuye habitualmente a San Agustín, aunque tiene antecedentes en Cicerón. Lo grave (lo «diabólico», según San Agustín) es persistir en el error. Precisamente porque el hombre reconoce que se equivoca y puede dar marcha atrás tienen sentido los mecanismos recursivos y toda la teoría de la impugnación.
% TODO: contenido incompleto o ambiguo en las notas originales (la transcripción original menciona «teoría de la evolución», corregida en el propio apunte entre corchetes como «impugnación»)
\end{note}

\begin{note}
\textbf{Francesco Carnelutti.} El procesalista italiano Francesco Carnelutti decía que la amenaza del error pende como \textbf{espada de Damocles} sobre el proceso. Según la leyenda, Damocles tenía sobre su cabeza una espada afilada que pendía de una crin de caballo y que en cualquier momento podía ceder. El error es así el fundamento clásico de toda la teoría de la impugnación.
\end{note}

\textbf{Couture} sosten\'ia que el t\'ermino recurso se vincula con \textit{``volver al camino de partida''}: al recurrir, el tribunal superior vuelve a recorrer un camino ya transitado en la instancia anterior, en la que fue dictada la resoluci\'on impugnada, para revisarlo y, en funci\'on de esa revisi\'on, confirmar la resoluci\'on o modificarla total o parcialmente.

\section{Definición de actos procesales de impugnación}

\begin{definition}{Actos procesales de impugnación (Palacio)}
Los actos procesales de impugnación son aquellos que están dirigidos \textbf{directa e inmediatamente} a provocar la \textbf{modificación o sustitución, total o parcial, de una resolución judicial}, en el mismo proceso en el que ella fue dictada.
\end{definition}

\begin{example}{El sistema de correcciones como un partido mal arbitrado}
Imaginemos un partido de fútbol mal arbitrado. El árbitro comete un error (el juez se equivoca). Los jugadores pueden pedir explicaciones (aclaratoria), pedir que el árbitro reconsidere su decisión en el momento (revocatoria), apelar ante el tribunal deportivo superior (apelación) o, si hubo una grave irregularidad de procedimiento, impugnar la validez del partido entero (incidente de nulidad). Todo ese sistema de correcciones forma la teoría de la impugnación.
\end{example}

\cornelltitle{concepto y fundamento de la impugnación}
\cornellblocksum{
\cqrow{¿Qué relación existe entre los actos de impugnación y los recursos?}{Los actos de impugnación son el género y los recursos son una especie dentro de ese género. Existen otros actos de impugnación que no son recursos.
}
\cqrow{¿Qué significa la palabra «impugnación» según su etimología?}{Proviene del latín \textit{pugnus} (puño) y alude a la lucha, al combate y al ataque. En el proceso, se trata de todo mecanismo que ataca o combate un acto procesal o procedimiento que lesione el interés de las partes o de un tercero.
}
\cqrow{¿Por qué existen los mecanismos de impugnación?}{Porque existe el error: los jueces y las partes pueden equivocarse, y esos errores de procedimiento o de interpretación de la ley son inevitables. Carnelutti decía que la amenaza del error pende como espada de Damocles en el proceso. La impugnación busca mitigar esas deficiencias.
}
\cqrow{¿Cómo se definen los actos procesales de impugnación?}{Son los dirigidos directa e inmediatamente a provocar la modificación o sustitución, total o parcial, de una resolución judicial, en el mismo proceso en el que ella fue dictada (Palacio).
}
}{La impugnación es el sistema que permite atacar actos o resoluciones viciados o equivocados. Su fundamento es la falibilidad humana. El género es el acto de impugnación y los recursos son una de sus especies.
}

\section{Clasificaciones de los actos de impugnación y de los recursos}

Con carácter metodológico general, cabe advertir que se trata de \textbf{clasificaciones posibles}: no hay clasificaciones verdaderas o falsas, buenas o malas. Toda clasificación es una manera de comprender, abarcar y categorizar distintos conceptos. En tanto ayude a comprenderlos mejor, la clasificación es buena; cuando dificulta la comprensión de lo que se quiere categorizar, deja de serlo.


\subsection{Primera clasificación: incidente de nulidad, recursos y otros medios}

La primera clasificación distingue dos grandes grupos, a los que se agregan otros medios menos frecuentes.

\begin{table}[H]
\centering
\small
\begin{tabularx}{\textwidth}{L{0.18\textwidth} >{\raggedright\arraybackslash}X >{\raggedright\arraybackslash}X}
\toprule
 & \textbf{Incidente de nulidad} & \textbf{Recursos} \\
\midrule
\textbf{Normativa y finalidad} &
Arts. 169--174 CPCCN. Corrige errores o vicios en la tramitación del proceso (vicios \textit{in procedendo}). &
Apuntan a la corrección de las decisiones impugnadas. \\
\midrule
\textbf{Ejemplo / tipos} &
Ejemplo típico: notificación defectuosa del traslado de la demanda en un domicilio incorrecto. &
\begin{itemize}[leftmargin=*,nosep]
\item Aclaratoria
\item Revocatoria o reposición
\item Apelación
\item Queja por apelación denegada
\item Recurso extraordinario federal
\item Apelación ordinaria ante la Corte Suprema
\item Inaplicabilidad de ley
\end{itemize} \\
\bottomrule
\end{tabularx}
\end{table}

En el proceso ordinario habitual, los recursos más frecuentes son la \textbf{aclaratoria}, la \textbf{revocatoria} y la \textbf{apelación}. La \textbf{queja} surge como consecuencia de la denegación de la apelación.

\subsection{Otros actos de impugnación (menos frecuentes)}

Si bien no son tan frecuentes, los siguientes mecanismos también integran el sistema de impugnación.

\subsubsection{Replanteo de prueba ante la alzada}

Cuando en el juicio ordinario el juez de primera instancia deniega ciertos medios de prueba, el litigante puede impugnar esa denegatoria en la etapa de revisión de la sentencia. Si la impugnación es adecuada, puede replantear la prueba y producirla ante el tribunal de alzada.

\subsubsection{Acción autónoma de revisión de la cosa juzgada (cosa juzgada írrita)}

Una vez que la sentencia pasa en autoridad de cosa juzgada, la jurisprudencia ha abierto, de modo \textbf{excepcionalísimo} y muy restrictivo, la posibilidad de revisar lo decidido por sentencia firme cuando existen circunstancias excepcionales.

\begin{example}{Juicios de filiación y prueba de ADN}
En otro tiempo no existía la prueba de ADN, por lo que las sentencias dictadas en juicios de filiación no podían valerse de ese medio de prueba. Con el advenimiento de la prueba biológica de altísima precisión que es el ADN, que permite despejar toda duda sobre la cuestión, se abrió la posibilidad de revisar esas sentencias firmes, aun cuando hubieran pasado en autoridad de cosa juzgada. También se aplica cuando, por maniobras fraudulentas, se ocultaron pruebas esenciales durante el proceso.
\end{example}

\subsubsection{Nulidad del laudo arbitral}

Los arts. 767--771 CPCCN prevén las hipótesis de nulidad del laudo arbitral o laudo de amigables componedores como mecanismo de impugnación.
% TODO: contenido incompleto o ambiguo en las notas originales (el docente menciona el art. 771 y, al mismo tiempo, los arts. 767-771)

\subsubsection{Juicio ordinario posterior al juicio ejecutivo}

Permite revisar, en un proceso de conocimiento pleno, una sentencia dictada en juicio ejecutivo, en el que el acotado marco procedimental no permitió ejercer en plenitud las defensas disponibles.

\subsection{Segunda clasificación: remedios y recursos}

La doctrina distingue también entre \textbf{remedios} y \textbf{recursos} en sentido estricto.

\begin{table}[H]
\centering
\small
\begin{tabularx}{\textwidth}{L{0.18\textwidth} >{\raggedright\arraybackslash}X >{\raggedright\arraybackslash}X}
\toprule
 & \textbf{Remedios} & \textbf{Recursos (\textit{stricto sensu})} \\
\midrule
\textbf{Orientación} &
Reparar errores procesales (vicios \textit{in procedendo}). &
Apuntan a que un tribunal jerárquicamente superior reexamine el acto impugnado. \\
\midrule
\textbf{Quién resuelve} &
\textbf{El mismo juez que dictó el acto.} &
Un tribunal jerárquicamente superior. \\
\midrule
\textbf{Vicios que se atacan} &
Vicios \textit{in procedendo}. &
Vicios \textit{in iudicando}: errores en la decisión judicial (interpretación de la ley o de los hechos). \\
\midrule
\textbf{Incluyen} &
Incidente de nulidad, acción autónoma de revisión de cosa juzgada, juicio ordinario posterior, aclaratoria* y revocatoria*. &
Apelación, queja, recurso extraordinario federal, apelación ordinaria ante la Corte, inaplicabilidad de ley. \\
\bottomrule
\end{tabularx}
\end{table}

(*) La aclaratoria y la revocatoria presentan aspectos que podrían involucrarlas dentro de la categoría de remedios, aunque el Código las trate y denomine como «recursos». Técnicamente, en doctrina, sería más adecuado clasificarlas dentro de los remedios.

\subsection{Clasificaci\'on general de los recursos}

Es importante conocer el cuadro completo de los recursos.

\begin{table}[H]
\centering
\begin{tabularx}{\textwidth}{
    >{\raggedright\arraybackslash}p{0.27\textwidth}
    >{\raggedright\arraybackslash}p{0.22\textwidth}
    >{\raggedright\arraybackslash}X
}
\toprule
\textbf{Tipo} & \textbf{Categor\'ia} & \textbf{Ejemplos} \\
\midrule
\textbf{Horizontales} & Remedios & Aclaratoria, revocatoria \\
\addlinespace
\textbf{Verticales ordinarios} & Apelaci\'on & Recurso de apelaci\'on (proceso ordinario) \\
\addlinespace
\textbf{Verticales extraordinarios (comunes)} & Regulados por ley & Inaplicabilidad de ley (fallos plenarios), Recurso Extraordinario Federal (Ley 48) \\
\addlinespace
\textbf{Verticales extraordinarios (excepcionales)} & Creaci\'on pretoriana & Sentencia arbitraria, gravedad institucional, \textit{per saltum} (ya regulado) \\
\addlinespace
\textbf{Especiales} & Queja & Cuando se deniega el recurso de apelaci\'on o el extraordinario \\
\addlinespace
\textbf{Parajudiciales} & Fuera del proceso & Recursos administrativos, recursos arbitrales \\
\bottomrule
\end{tabularx}
\caption{Clasificaci\'on general de los recursos.}
\end{table}

\subsubsection{Recursos horizontales: aclaratoria y revocatoria}

Son horizontales porque se interponen, tramitan y resuelven \textbf{ante el mismo juez} que dict\'o la resoluci\'on que se impugna: se interpone, se funda y es ese mismo juez quien la resuelve. Se dice que, m\'as que recursos, son \textbf{remedios}, porque no existe un nuevo tribunal que vuelva a recorrer el camino de una instancia anterior.

Una caracter\'istica fundamental que los diferencia de la apelaci\'on es que, en los horizontales, el recurso se \textbf{interpone y se funda en el mismo acto}. Esto no sucede en el recurso de apelaci\'on, donde existen dos momentos diferenciados.

\subsubsection{Recursos verticales ordinarios: la apelaci\'on}

La apelaci\'on es un recurso vertical porque se interpone ante el juez de primera instancia (el que dict\'o la resoluci\'on) para que sea la \textbf{c\'amara}, que se encuentra por encima, la que resuelva. Tiene \textbf{dos momentos diferenciados}: la interposici\'on y la fundamentaci\'on, que se desarrollan m\'as adelante.

\cornelltitle{clasificación de los actos de impugnación}
\cornellblocksum{
\cqrow{¿Qué grupos de actos de impugnación distingue la primera clasificación?}{El incidente de nulidad (arts. 169--174 CPCCN, que corrige vicios en la tramitación del proceso) y los recursos (aclaratoria, revocatoria o reposición, apelación, queja por apelación denegada, recurso extraordinario federal, apelación ordinaria ante la Corte Suprema e inaplicabilidad de ley). Se agregan otros medios menos frecuentes.
}
\cqrow{¿Cuáles son los recursos más frecuentes en el proceso ordinario?}{La aclaratoria, la revocatoria y la apelación. La queja surge como consecuencia de la denegación de la apelación.
}
\cqrow{¿Qué otros actos de impugnación menos frecuentes existen?}{El replanteo de prueba ante la alzada, la acción autónoma de revisión de la cosa juzgada, la nulidad del laudo arbitral (arts. 767--771 CPCCN) y el juicio ordinario posterior al juicio ejecutivo.
}
\cqrow{¿Cuándo procede, excepcionalmente, la revisión de la cosa juzgada?}{De modo excepcionalísimo y muy restrictivo, ante circunstancias excepcionales: por ejemplo, la posibilidad de contar con la prueba de ADN en juicios de filiación o el ocultamiento de pruebas esenciales mediante maniobras fraudulentas.
}
\cqrow{¿Cuál es la diferencia entre remedios y recursos \textit{stricto sensu}?}{Los remedios reparan errores procesales (vicios \textit{in procedendo}) y los resuelve el mismo juez que dictó el acto. Los recursos en sentido estricto apuntan a que un tribunal superior reexamine el acto impugnado y atacan los vicios \textit{in iudicando} (errores en la interpretación de la ley o de los hechos).
}
\cqrow{¿Dónde ubica la doctrina a la aclaratoria y a la revocatoria?}{Aunque el Código las denomine «recursos», técnicamente sería más adecuado clasificarlas dentro de los remedios.
}
}{Los actos de impugnación pueden clasificarse en incidente de nulidad, recursos y otros medios menos frecuentes. Otra clasificación distingue remedios (resueltos por el mismo juez, orientados a vicios \textit{in procedendo}) y recursos en sentido estricto (sometidos a un tribunal superior, dirigidos a vicios \textit{in iudicando}).
}

\cornelltitle{fundamentos y clasificación de los recursos}
\cornellblocksum{
\cqrow{¿Qu\'e diferencia hay entre impugnaci\'on y recurso?}{La impugnaci\'on es el g\'enero (atacar un acto, dentro o fuera del proceso); el recurso es una especie de impugnaci\'on y s\'olo procede contra resoluciones judiciales (providencias simples, interlocutorias y sentencias definitivas). Una impugnaci\'on de dictamen pericial o de liquidaci\'on es una impugnaci\'on pero no un recurso.
}
\cqrow{¿Cu\'al es el fundamento de los recursos?}{El error humano (Calamandrei): los jueces pueden equivocarse y debe existir un sistema que subsane el error y evite consolidar una injusticia. Couture vincula el t\'ermino recurso con ``volver al camino de partida'', es decir, que el tribunal superior revise el camino ya transitado para confirmar o modificar la resoluci\'on.
}
\cqrow{¿Por qu\'e la aclaratoria y la revocatoria son recursos horizontales?}{Porque se interponen, tramitan y resuelven ante el mismo juez que dict\'o la resoluci\'on. Por eso se las considera m\'as bien remedios. Se interponen y fundan en el mismo acto.
}
\cqrow{¿Por qu\'e la apelaci\'on es un recurso vertical ordinario?}{Porque se interpone ante el juez de primera instancia, pero lo resuelve la c\'amara, que est\'a por encima. Tiene dos momentos diferenciados: interposici\'on y fundamentaci\'on.
}
}{la impugnaci\'on es el g\'enero y el recurso la especie, limitada a resoluciones judiciales. El fundamento de los recursos es el error humano. Los recursos horizontales se resuelven ante el mismo juez y se interponen y fundan en el mismo acto; la apelaci\'on es un recurso vertical ordinario, que resuelve la c\'amara y tiene dos momentos.
}

\chapter[Requisitos comunes y recurso indiferente]{Requisitos comunes de los recursos y recurso indiferente}\label{cap:requisitos}

\section{Legitimación}

Las \textbf{partes} son las primeras legitimadas para hacer un planteo recursivo: el actor y el demandado. También lo están:
\begin{itemize}
\item los \textbf{terceros}, ya sea por intervención voluntaria u obligada;
\item los \textbf{ministerios públicos} que intervengan en el caso: el fiscal, en cuestiones en las que está comprometido el orden público, y el defensor de menores, en los casos en que hay niños, niñas o adolescentes en el proceso;
\item las \textbf{personas ajenas al proceso} que se ven afectadas en un interés legítimo ante una decisión o actuación procesal; entre ellas, los auxiliares de la justicia que intervienen en el proceso, por ejemplo los peritos cuando se regulan sus honorarios.
\end{itemize}

\begin{example}{Terceros con legitimación recursiva}
Reparticiones públicas o empresas privadas que, en el proceso, incumplen mandatos judiciales (como trabar embargo sobre haberes o no contestar un pedido de informes) y a las que se les impone una multa. Esos terceros, ajenos al proceso, pueden contar con legitimación para plantear un recurso, justamente porque lo decidido en el proceso los afecta de modo directo.
\end{example}

\section{El agravio, gravamen o perjuicio}

Este requisito es fundamental. Los términos \textbf{agravio}, \textbf{gravamen} y \textbf{perjuicio} son sinónimos utilizados indistintamente en doctrina y jurisprudencia.

\begin{definition}{Agravio (gravamen o perjuicio)}
Es la \textbf{insatisfacción total o parcial} de cualquiera de las pretensiones planteadas en virtud de lo decidido; o bien, para la parte demandada, el \textbf{rechazo de las defensas} opuestas. Se trata de la derrota que la decisión judicial provoca frente a los postulados sostenidos en las presentaciones de quien recurre, sea que su postulado apunte a satisfacer una pretensión o a oponer una defensa.
\end{definition}

Sin agravio no hay legitimación para recurrir. El rechazo de la pretensión para el actor, o el rechazo de las defensas para el demandado, constituyen el agravio que habilita el recurso.

\section{El plazo perentorio}

En todos los recursos existe un plazo para plantearlo, que siempre es \textbf{perentorio}, es decir, \textbf{fatal}. El derecho a recurrir debe ejercerse dentro del plazo establecido por el código. Si se deja pasar ese plazo sin plantear el recurso, se pierde la posibilidad futura de formular ese planteo.

\section{Admisibilidad y fundabilidad}

Todos los recursos presentan dos facetas.

\subsection{Admisibilidad}

\begin{definition}{Admisibilidad de un recurso}
Alude a toda la cuestión \textbf{formal} requerida para el planteo impugnativo: el plazo en que debe presentarse; los recaudos formales (si debe ser por escrito o puede ser verbal); qué debe contener el escrito y si debe ser fundado o si basta con el planteo recursivo; y todo lo que tiene que ver con el aspecto formal.
\end{definition}

La admisibilidad también indica si un recurso procede contra determinados actos y contra cuáles no. Es lo que debe analizarse \textbf{antes} de pasar al examen de la sustancia o de lo medular del planteo recursivo.

\subsection{Fundabilidad}

\begin{definition}{Fundabilidad de un recurso}
Apunta al \textbf{fondo} de la cuestión, al aspecto sustancial de la impugnación: la motivación que lleva al sujeto a formular su planteo, la crítica concreta que hace a la decisión del juez o el vicio que señala en alguna actuación procesal. Contempla los agravios, gravámenes o perjuicios que provoca el acto impugnado, en qué afecta a quien recurre y de qué manera se siente perjudicado por lo decidido.
\end{definition}

\begin{important}
\textbf{Orden de análisis.} El análisis de fundabilidad viene \textbf{después} del de admisibilidad. Primero se analiza lo formal: si el recurso procede formalmente contra determinado acto, si se planteó a tiempo y si se planteó de la forma adecuada según el código. Luego se analiza si quien recurre tiene razón en la cuestión que critica y en los motivos que le causan el perjuicio o agravio por el cual se alza contra la decisión o actuación judicial.
\end{important}

\cornelltitle{requisitos comunes de los recursos}
\cornellblocksum{
\cqrow{¿Quiénes tienen legitimación para recurrir?}{Las partes (actor y demandado), los terceros intervinientes (voluntarios u obligados), los ministerios públicos (fiscal, defensor de menores) y las personas ajenas al proceso cuyo interés legítimo resulta afectado (por ejemplo, peritos cuyos honorarios son regulados, o empresas sancionadas por incumplir mandatos judiciales).
}
\cqrow{¿Qué es el agravio o gravamen?}{Es la insatisfacción total o parcial del postulado esgrimido en el proceso (pretensión o defensa). También se denomina gravamen o perjuicio. El rechazo de la pretensión para el actor, o de las defensas para el demandado, constituye el agravio que habilita el recurso. Sin agravio no hay legitimación para recurrir.
}
\cqrow{¿Qué características tiene el plazo para recurrir?}{Es perentorio o fatal: si se deja pasar sin plantear el recurso, se pierde la posibilidad futura de formularlo.
}
\cqrow{¿Cuál es la diferencia entre admisibilidad y fundabilidad de un recurso?}{La admisibilidad se refiere al aspecto formal: plazo, forma de presentación, contenido del escrito, tipo de resolución recurrible. Se analiza primero. La fundabilidad se refiere al fondo: si los agravios son válidos y si la crítica a la decisión judicial es correcta. Se analiza después de superado el juicio de admisibilidad.
}
}{Todo recurso requiere tres presupuestos comunes: legitimación, agravio y plazo perentorio. Presenta dos dimensiones: la admisibilidad formal, que se analiza primero, y la fundabilidad sustancial, que se analiza después.
}

\section{El recurso indiferente}

El recurso indiferente \textbf{no se encuentra tipificado} en el Código Procesal ni en el ordenamiento argentino; no está legislado ni tiene recepción normativa. Algunos autores se refieren a este instituto tomándolo del derecho alemán, en el que la estructura recursiva es bien distinta de la argentina.

Su característica central es que el tribunal \textbf{deja de lado los recaudos de admisibilidad} que tienen todos los planteos impugnativos (plazo, forma, características, tipo de resolución atacada, etc.) en pos de garantizar ciertos derechos fundamentales que están siendo conculcados por alguna decisión judicial. El tribunal de segunda instancia deja de lado los recaudos formales y pasa a conocer y resolver el fondo del asunto, lo sustancial, sin cortapisas ni trabazones de índole formal.

\begin{definition}{Recurso indiferente}
Es aquel que, sin ser el que la ley prescribe expresamente para el caso, o que siéndolo se han omitido elementos formales, produce no obstante los mismos efectos respecto de la posibilidad (procedibilidad) de la vía recursiva que el recurso correctamente articulado.
\end{definition}

\subsection{Casos YPF y Tejida (CSJN)}

\begin{example}{Caso YPF (CSJN), derecho previsional}
El caso que dio andamiaje a este instituto en el ordenamiento argentino fue fallado por la Corte Suprema y es conocido como caso YPF. En el ámbito provincial se había desestimado el planteo de una persona en verdadero estado de vulnerabilidad, que tenía el beneficio de una pensión derivada de su situación de discapacidad. Por defectos en la manera en que realizó los planteos y dedujo los recursos en sede provincial, se había rechazado la posibilidad de revisar y acoger la pretensión. La Corte, dejando de lado esas cuestiones formales, admitió la viabilidad del planteo.
\end{example}

\begin{example}{Caso Tejida (CSJN, marzo de 2018)}
La Corte dijo que, dada la índole peculiar de ciertas pretensiones, compete a los jueces la búsqueda de soluciones que se avengan a éstas, para lo cual deben encauzar los trámites por vías expeditivas, a fin de evitar que la incorrecta utilización de las formas pueda conducir a la frustración de derechos tutelados constitucionalmente, cuya suspensión a las resultas de nuevos trámites es inadmisible.
\end{example}

\cornelltitle{recurso indiferente}
\cornellblocksum{
\cqrow{¿Qué es el recurso indiferente y está regulado en el CPCCN?}{Es una categoría doctrinal que no está tipificada en el Código ni legislada en el ordenamiento argentino. Algunos autores la toman del derecho alemán. Habilita a los tribunales a dejar de lado los requisitos formales de admisibilidad para garantizar derechos fundamentales gravemente afectados, y a resolver el fondo del asunto.
}
\cqrow{¿Qué efectos produce el recurso indiferente?}{Aunque no sea el recurso que la ley prescribe expresamente (o se hayan omitido elementos formales), produce los mismos efectos respecto de la procedibilidad de la vía recursiva que el recurso correctamente articulado.
}
\cqrow{¿Qué casos le sirven de sustento?}{El caso YPF (CSJN), en el que la Corte, dejando de lado cuestiones formales, admitió el planteo de una persona vulnerable con una pensión derivada de su discapacidad; y el caso Tejida (CSJN, marzo de 2018), sobre la búsqueda de soluciones por vías expeditivas para evitar la frustración de derechos constitucionalmente tutelados.
}
}{El recurso indiferente es una construcción jurisprudencial excepcional, no tipificada, para los casos en que las formas frustrarían derechos fundamentales: el tribunal deja de lado los recaudos de admisibilidad y resuelve el fondo.
}

\chapter[Doble instancia y límites del tribunal de alzada]{La garantía de la doble instancia y los límites del tribunal de alzada}\label{cap:doble}

\section{La garantía de la doble instancia}

Cabe preguntarse si existe en el ordenamiento argentino una garantía procesal, amparada por la Constitución Nacional, de que en todo proceso existirá la doble instancia; es decir, la posibilidad de que, por las vías recursivas, toda decisión pueda ser revisada por un órgano diferente de aquel que dictó la resolución.

\subsection{Perspectiva histórica}

La cuestión no es nueva en la historia del proceso. El tema de la doble instancia, incluso la posibilidad de revisiones ulteriores más allá de dos instancias, se remonta a muy vieja data. En el sistema nacional, desde tiempos de la colonia, era reconocida la posibilidad de recurrir a la audiencia, que era la instancia superior a las decisiones de los jueces de primera instancia. Mucho antes, desde la Edad Media y aun desde la época del Imperio Romano, se observa la posibilidad de recurrir las decisiones del pretor ante órganos creados especialmente por el emperador o, en épocas monárquicas, por los reyes, que delegaban en funcionarios de jerarquía la posibilidad de revisar las decisiones de los jueces inferiores.

\subsection{La Constitución Nacional de 1853: ausencia de garantía expresa}

En 1853, la Constitución Nacional no introdujo la doble instancia como garantía del debido proceso. Otros aspectos muy claros están presentes en el art. 18 desde aquella época, pero la doble instancia no está introducida como garantía constitucional: ni siquiera en el aspecto más garantista del ordenamiento, que es el de la persecución de los delitos, ni mucho menos en cuestiones no penales (civiles).

\begin{note}
\textbf{Art. 18 de la Constitución Nacional.} Establece que ningún habitante de la Nación puede ser penado sin juicio previo fundado en ley anterior al hecho del proceso, y que es inviolable la defensa en juicio de la persona y de los derechos. Este artículo, que establece las garantías del debido proceso, no menciona expresamente la doble instancia.
% TODO: contenido incompleto o ambiguo en las notas originales (la cita del art. 18 figura abreviada y con la expresión «hecho del proceso»)
\end{note}

\subsection{El Pacto de San José de Costa Rica y la reforma de 1994}

Recién con la \textbf{Convención Americana de Derechos Humanos} (Pacto de San José de Costa Rica) aparece en el ordenamiento una norma que prevé específicamente el derecho de recurrir del fallo ante juez o tribunal superior. Esto ocurre a raíz de la incorporación del tratado a la Constitución Nacional en 1994, en virtud del \textbf{art. 75, inciso 22}.

\begin{formula}{Art. 8º, inc. 2º, ap. h) CADH}
Establece el derecho de recurrir del fallo ante juez o tribunal superior. El inciso segundo se refiere a las personas inculpadas de un delito, por lo que la garantía está claramente enmarcada en el ámbito penal.
\end{formula}

El art. 75 inc. 22 de la Constitución Nacional otorga jerarquía constitucional a los tratados internacionales de derechos humanos, entre ellos la Convención Americana de Derechos Humanos.

\subsection{El caso Giroldi (CSJN)}

La Corte Suprema ratificó la garantía en materia penal en el precedente \textbf{Giroldi}. Consideró que la incorporación de la Convención Americana de Derechos Humanos a la Constitución Nacional conformaba una garantía que integraba el ordenamiento constitucional y que, por ende, debía respetarse con rango constitucional la doble instancia en materia penal. En materia penal se habla, más que de doble instancia, de \textbf{doble conforme}, para que la decisión condenatoria adquiera validez desde el plano de esta garantía.

\subsection{Materia no penal: ausencia de garantía constitucional}

\begin{important}
La Corte Suprema ha sostenido sistemáticamente que la garantía de la doble instancia \textbf{no tiene jerarquía constitucional en materia no penal}. En materia no penal no hay garantía constitucional de doble instancia, salvo cuando las leyes específicamente lo establezcan.
\end{important}

El \textbf{art. 242 CPCCN} establece, como recaudo de admisibilidad del recurso de apelación, que el monto discutido en el proceso supere cierta suma de dinero (aproximadamente trescientos mil pesos, según las fuentes originales). Toda discusión por una cuestión de menor cuantía no admite apelación, salvo que se trate de cuestiones arancelarias u honorarios.
% TODO: verificar el monto vigente; las fuentes originales indican una cifra aproximada.

Frente al interrogante de si no se priva de la doble instancia a quienes discuten montos menores a 300.000 pesos, la Corte ha dicho que las limitaciones vinculadas con la cuantía establecidas por el legislador son absolutamente razonables y no violan ninguna pauta propia del debido proceso, precisamente porque en materia no penal no existe garantía constitucional de doble instancia.

\section{Límites del conocimiento del tribunal de alzada}

Los tribunales de segunda instancia (cámaras de apelación, tribunales de alzada) sólo son competentes para intervenir por vía de revisión de las decisiones de un juez de primera instancia; no tienen competencia originaria.

\subsection{Primer límite: los escritos constitutivos del proceso}

Al igual que el juez de primera instancia, el tribunal de alzada está limitado por el \textbf{principio de congruencia}, que lo obliga a mantenerse dentro del marco de lo que las partes han postulado en los escritos constitutivos del proceso (pretensión y defensa).

\begin{formula}{Art. 277 CPCCN}
Ante los jueces de segunda instancia no se pueden introducir nuevas pretensiones o defensas no sometidas al conocimiento del juez de primera instancia. Las pretensiones, postulaciones y defensas deducidas ante el juez de primera instancia son las que pueden reeditarse ante el tribunal, con ampliación de fundamentos, pero no pueden plantearse nuevos tópicos no propuestos en primera instancia.
\end{formula}

\subsubsection{Hechos sobrevinientes (\textit{iura novit curia} / hechos nuevos)}
% TODO: contenido incompleto o ambiguo en las notas originales (el apunte rotula este tema como «iura novit curia / hechos nuevos»; anuncia dos tipos de hechos pero sólo se desarrolla uno)

Existen excepciones a este límite. Se anuncian dos tipos de hechos nuevos, de los cuales las fuentes desarrollan el siguiente: los \textbf{hechos posteriores a la sentencia de primera instancia}. Por haber ocurrido con posterioridad, no pudieron ser planteados ante el juez de primera instancia, y el art. 277 los admite.

\begin{formula}{Art. 163, inc. 6º, segunda parte, CPCCN}
Siempre el juez, sea de primera o de segunda instancia, podrá hacer mérito de los hechos constitutivos, modificativos o extintivos de las pretensiones, producidos durante la sustanciación del juicio y debidamente probados.
\end{formula}

\subsection{Segundo límite: los agravios expresados y sus contestaciones}

Además del límite de los escritos constitutivos, el tribunal de alzada está acotado por los propios recursos planteados por las partes: los recursos establecen el límite dentro del cual puede actuar, y los agravios planteados por las partes marcan ese límite de actuación.

\begin{table}[H]
\centering
\small
\begin{tabularx}{\textwidth}{>{\raggedright\arraybackslash}X >{\raggedright\arraybackslash}X}
\toprule
\textbf{Primera limitación} & \textbf{Segunda limitación} \\
\midrule
Los escritos constitutivos del proceso (pretensiones y defensas planteadas en primera instancia). &
Las expresiones de agravios y sus contestaciones (los recursos planteados por las partes). \\
\bottomrule
\end{tabularx}
\end{table}

\section{El principio de \textit{reformatio in peius}}

Este principio tiene íntima relación con los límites del tribunal de alzada.

\begin{definition}{Reformatio in peius}
La sentencia de segunda instancia \textbf{jamás podrá empeorar la situación del único apelante}.
\end{definition}

\begin{example}{Accidente de tránsito y único apelante}
En un juicio de accidente de tránsito, el juez de primera instancia fija una indemnización de 500.000 pesos en favor del actor. La parte demandada no apela, porque considera la suma razonable (o incluso baja). El actor sí apela, porque considera baja la indemnización. El recurso llega a la cámara.

El tribunal de segunda instancia advierte que, a su criterio, la indemnización fue excesiva y que el actor merecía mucho menos. Sin embargo, como el único apelante es el actor (la parte demandada consintió la sentencia), el tribunal de alzada no puede bajar la indemnización: no existe agravio que permita reducirla. A lo sumo podrá confirmar la sentencia.

\textbf{Conclusión:} cuando se apela y se es el único apelante, nunca se puede perder: o se gana más, o se sigue con lo mismo que dio el juez de primera instancia.
\end{example}

El límite por monto del artículo 242 CPCCN como recaudo de admisibilidad de la apelación se analiza con mayor detalle en el capítulo~\ref{cap:apelacion}.

\cornelltitle{doble instancia y límites del tribunal de alzada}
\cornellblocksum{
\cqrow{¿Existe en Argentina una garantía constitucional de doble instancia en materia civil?}{No. En materia no penal, la Corte Suprema sostiene sistemáticamente que no existe garantía constitucional de doble instancia, salvo que las leyes específicamente la establezcan.
}
\cqrow{¿Cómo se incorporó la garantía de doble instancia en materia penal?}{La Constitución de 1853 no la preveía. Surgió con el Pacto de San José de Costa Rica (art. 8º inc. 2º ap. h), incorporado a la Constitución por el art. 75 inc. 22 en 1994, y fue ratificada por la Corte en el caso Giroldi. En materia penal se habla de doble conforme.
}
\cqrow{¿Qué establece el art. 242 CPCCN y por qué no viola el debido proceso?}{Exige que el monto discutido supere cierta suma para admitir la apelación, salvo cuestiones arancelarias u honorarios. La Corte ha dicho que estas limitaciones de cuantía son razonables, porque en materia no penal no existe garantía constitucional de doble instancia.
}
\cqrow{¿Cuáles son los límites del conocimiento del tribunal de alzada?}{Una doble limitación: los escritos constitutivos del proceso (art. 277 CPCCN: no se pueden introducir nuevas pretensiones o defensas) y las expresiones de agravios y sus contestaciones. Se admiten hechos sobrevinientes (art. 163 inc. 6º, segunda parte).
}
\cqrow{¿Qué es el principio de \textit{reformatio in peius}?}{Prohíbe que el tribunal de segunda instancia empeore la situación del único apelante. Si sólo apeló el actor por considerar baja la indemnización y la demandada consintió la sentencia, el tribunal de alzada no puede bajar la indemnización: a lo sumo la confirma o la aumenta.
}
}{La garantía de la doble instancia existe sólo en materia penal con rango constitucional (CADH, art. 75 inc. 22 CN); en materia civil el legislador puede restringirla. El tribunal de alzada tiene una doble limitación (escritos constitutivos y agravios expresados) y no puede perjudicar al único apelante.
}

\chapter[Aclaratoria y reposición]{Remedios ante el mismo juez: aclaratoria y reposición o revocatoria}\label{cap:remedios}

\section{El recurso de aclaratoria}

\begin{definition}{Aclaratoria (Palacio)}
Es un \textbf{remedio} que se concede a las partes para obtener que el mismo juez o tribunal que dictó una resolución subsane las deficiencias materiales o conceptuales que contenga, o la integre de conformidad con las peticiones oportunamente formuladas.
\end{definition}

El \textbf{art. 166, incisos 1º y 2º, CPCCN} regula las actividades del juez después del dictado de la sentencia. El inciso 2º establece que la parte puede plantear la aclaratoria dentro del tercer día de notificada la resolución. El inciso 1º faculta al juez a corregir de oficio cualquier defecto o deficiencia material con anterioridad a la notificación a las partes. En el mismo sentido, el \textbf{art. 36, inciso 6º, CPCCN} establece, entre los deberes y facultades del juez, la facultad de corregir de oficio cualquier defecto o deficiencia material de una resolución, con anterioridad a la notificación a las partes.

\subsection{Supuestos que habilitan la aclaratoria}

\begin{table}[H]
\centering
\small
\begin{tabularx}{\textwidth}{L{0.28\textwidth} >{\raggedright\arraybackslash}X}
\toprule
\textbf{Supuesto} & \textbf{Descripción / ejemplo} \\
\midrule
\textbf{Errores materiales} &
Confundir el nombre de las partes (por ejemplo, escribir «parte actora» cuando se quería decir «parte demandada»); discordancia entre cifras en guarismos y en letras; errores aritméticos en sumas o cálculos. \\
\midrule
\textbf{Errores de copia o transcripción} &
Párrafos «pegados» de otra resolución, especialmente al armar resoluciones con procesadores de texto. Hoy son menos frecuentes que en la época de las máquinas de escribir. \\
\midrule
\textbf{Contradicción entre considerandos y parte resolutiva} &
Los considerandos fundamentan la concesión de un rubro pero la parte resolutiva lo omite, o viceversa. \\
\midrule
\textbf{Conceptos oscuros} &
Discordancia entre la idea que se quiso expresar y las palabras utilizadas, que torna necesaria una aclaración, siempre que ello no implique una modificación sustancial de lo decidido. \\
\midrule
\textbf{Omisión de tratar pretensiones deducidas} &
El juez no trató en la resolución alguna de las pretensiones debatidas en el proceso. La aclaratoria es la oportunidad para ampliar y completar la resolución. \\
\bottomrule
\end{tabularx}
\end{table}

\subsection{Características procesales de la aclaratoria}

\subsubsection{Plazo}

La parte tiene \textbf{tres (3) días} para plantear la aclaratoria, contados desde la notificación de la sentencia o resolución que se pretende aclarar. En el caso de aclaratoria contra sentencia de segunda instancia, el plazo es de \textbf{cinco (5) días}.

\subsubsection{Forma}

Se plantea por escrito, y la fundamentación se hace en el mismo escrito de interposición. Si se trata de una providencia no notificada por cédula, los tres días comienzan a correr desde el martes o viernes posterior al dictado (notificación por ministerio de la ley).

\subsubsection{Sustanciación}

Según el Código, la aclaratoria es resuelta por el juez \textbf{sin traslado previo} a la otra parte (sin sustanciación).

\subsubsection{Relación con el recurso de apelación}

\begin{important}
La aclaratoria \textbf{no interrumpe el plazo para apelar}. La apelación debe interponerse autónomamente dentro del quinto día, sin perjuicio de la aclaratoria planteada. Lo conveniente es plantear la aclaratoria y, antes del quinto día, plantear la apelación; si la aclaratoria suple los defectos y torna innecesaria la apelación, se desistirá eventualmente de ésta. No puede perderse la apelación por haber planteado una aclaratoria.

No cabe la apelación subsidiaria de ese planteo: es aclaratoria o apelación, pero no las dos cosas juntas. \textbf{El único recurso que admite la apelación en subsidio es la revocatoria.}
\end{important}

\cornelltitle{recurso de aclaratoria}
\cornellblocksum{
\cqrow{¿Para qué sirve el recurso de aclaratoria y cuándo no procede?}{Sirve para que el mismo juez corrija errores materiales (nombres, sumas, cálculos), aclare conceptos oscuros o integre la resolución cuando omitió tratar alguna pretensión. No procede para modificar sustancialmente lo decidido: para ello existen la revocatoria o la apelación.
}
\cqrow{¿Qué supuestos habilitan la aclaratoria?}{Errores materiales, errores de copia o transcripción, contradicción entre considerandos y parte resolutiva, conceptos oscuros y omisión de tratar pretensiones deducidas.
}
\cqrow{¿Cuál es el plazo para plantearla?}{Tres días desde la notificación de la resolución; cinco días si se plantea contra una sentencia de segunda instancia.
}
\cqrow{¿Cómo se plantea y cómo se resuelve?}{Por escrito, con la fundamentación en el mismo escrito. El juez la resuelve sin traslado previo a la otra parte.
}
\cqrow{¿Qué relación tiene con el recurso de apelación?}{No interrumpe el plazo para apelar: la apelación debe interponerse autónomamente dentro del quinto día. No cabe la apelación subsidiaria de la aclaratoria; el único recurso que admite apelación en subsidio es la revocatoria.
}
}{La aclaratoria permite al mismo juez corregir errores formales o integrar omisiones sin alterar la esencia de lo decidido. Se plantea por escrito dentro de los tres días (cinco contra sentencia de segunda instancia), no se sustancia y no interrumpe el plazo para apelar.
}

\section{El recurso de reposición o revocatoria}

\begin{definition}{Revocatoria o reposición}
Es el \textbf{remedio} procesal tendiente a que el mismo juez o tribunal que dictó una resolución subsane los agravios que ella ocasiona, \textbf{por contrario imperio}. Se le pide al juez o tribunal que dictó la resolución que la modifique, la cambie o la revise, porque esa resolución causa perjuicios que justifican que corrija la decisión y la modifique en favor de quien hace el planteo.
\end{definition}

Los \textbf{arts. 238 a 241 CPCCN} regulan el recurso: procedencia (art. 238), forma y trámite (arts. 239 y 240) y efectos (art. 241). Se plantea ante el mismo juez que dictó la decisión y es resuelto por ese mismo juez.

\subsection{Recaudos de admisibilidad}

\subsubsection{Tipo de resolución atacable}

Procede únicamente contra \textbf{providencias simples}. Es inadmisible contra sentencias interlocutorias o sentencias definitivas.

\subsubsection{Jueces ante quienes procede}

Procede contra decisiones de jueces de primera instancia y también contra decisiones del presidente de una cámara de apelación (providencias simples de mero trámite adoptadas durante la sustanciación del juicio en segunda instancia, con anterioridad al dictado de la sentencia de alzada).

Cuando la revocatoria se plantea contra una providencia del presidente de la cámara, resuelve el tribunal en pleno (los tres jueces). Esas decisiones del tribunal en pleno hacen ejecutoria: no pueden ser recurridas.

\subsubsection{Providencias que causen o no gravamen irreparable}

\begin{formula}{Art. 238 CPCCN}
El recurso de reposición o revocatoria procede contra providencias simples, causen o no gravamen irreparable.
\end{formula}

El \textbf{gravamen irreparable} es aquel que no puede ser suplido por lo que se decida en la sentencia: la providencia no puede ser corregida durante el curso del proceso ni con la sentencia, sino que, una vez adoptada la decisión (aunque sea en una providencia simple), marca un cambio de rumbo tal que no hay vuelta atrás.

\begin{example}{Contestación de demanda tenida por no presentada}
La providencia simple por la cual se tiene por no contestado el traslado de la demanda en término, por considerarse extemporánea la contestación, deja al demandado sin posibilidad de defensa si no se corrige por vía de revocatoria. Eso es irreparable desde el plano de la defensa. Por ello, causa gravamen irreparable y es susceptible tanto de revocatoria como de apelación.

En cambio, una providencia que ordena agregar copia de un escrito antes de correr traslado puede no causar gravamen irreparable (la cuestión se soluciona presentando la copia), pero aun así es admisible la revocatoria.
\end{example}

\begin{note}
Mientras la revocatoria procede contra providencias simples, causen o no gravamen irreparable, la apelación exige que la providencia simple cause gravamen irreparable para ser admisible.
\end{note}

\subsection{Plazo}

El plazo para plantear la revocatoria es de \textbf{tres (3) días} desde la notificación de la providencia que se pretende impugnar, ya sea por ministerio de la ley o de manera electrónica (art. 135 CPCCN).

\textbf{Excepción:} si la decisión fue adoptada en una audiencia, la revocatoria debe plantearse \textbf{en el mismo acto de la audiencia, verbalmente}. No puede presentarse días después.

\begin{example}{Audiencia preliminar}
En una audiencia preliminar el juez adopta una serie de decisiones. La parte que no esté de acuerdo con alguna de ellas tiene el derecho de plantear la revocatoria en ese mismo momento. No puede hacerlo después de la audiencia, a los tres días.
\end{example}

\subsection{Forma de interposición y fundamentación}

La revocatoria se interpone por escrito (en la mayoría de los casos) o verbalmente (si la providencia fue adoptada en audiencia). Se funda en el mismo escrito o en el mismo acto de su interposición. Se exponen los fundamentos, los agravios y los motivos que llevan a criticar el acto cuya modificación se pretende. El \textbf{art. 239 CPCCN} regula la forma de interposición, incluida la posibilidad de interposición verbal cuando la decisión se adopta en el marco de una audiencia.

\subsection{Trámite y resolución}

Si la decisión atacada fue el resultado de una \textbf{petición formulada por la otra parte}, cuando se plantea la revocatoria el juez debe oír previamente a esa otra parte: debe correrse traslado del planteo de revocatoria a la parte que había pedido la decisión que ahora se impugna. Si la decisión no se originó en la petición de una de las partes, sino que fue una providencia dictada de oficio por el juez, \textbf{no es necesario} correr traslado a la otra parte.

El \textbf{art. 240 CPCCN} regula el trámite, incluidos los supuestos en que procede el traslado a la otra parte y aquellos en que no. El traslado, cuando corresponde, cuenta con un plazo de tres días.

\subsection{Efectos: la ejecutoria y sus excepciones}

Lo que el juez decida respecto de la revocatoria \textbf{hace ejecutoria}: queda firme lo decidido y resulta ejecutable. No es impugnable por vía de otro recurso.

Excepciones a la ejecutoria:
\begin{enumerate}
\item Si la revocatoria fue planteada con \textbf{apelación subsidiaria}: si el juez rechaza la revocatoria, no hace ejecutoria, porque queda pendiente el recurso de apelación interpuesto en subsidio.
\item Si el juez \textbf{admite} la revocatoria: esa decisión puede ser apelada por la contraparte.
\end{enumerate}

\cornelltitle{recurso de reposición o revocatoria}
\cornellblocksum{
\cqrow{¿Qué es la revocatoria o reposición?}{Es el remedio por el cual se pide al mismo juez o tribunal que dictó una resolución que la modifique por contrario imperio, porque le causa un perjuicio al recurrente. Está regulada en los arts. 238 a 241 CPCCN.
}
\cqrow{¿Contra qué tipo de resoluciones procede?}{Sólo contra providencias simples, causen o no gravamen irreparable (art. 238 CPCCN). No procede contra sentencias interlocutorias ni definitivas. Procede contra decisiones de jueces de primera instancia y del presidente de una cámara de apelación. La apelación, en cambio, exige que la providencia simple cause gravamen irreparable.
}
\cqrow{¿Cuál es el plazo y qué excepción tiene?}{Tres días desde la notificación. Si la decisión se adoptó en una audiencia, debe plantearse verbalmente en el mismo acto.
}
\cqrow{¿Cómo se tramita?}{Se interpone y funda por escrito (o verbalmente en audiencia). Si la decisión atacada se originó en una petición de la otra parte, se corre traslado a ésta; si fue dictada de oficio, no es necesario el traslado.
}
\cqrow{¿Qué efectos tiene lo resuelto y qué excepciones existen?}{Hace ejecutoria: queda firme y no es impugnable por otro recurso. No hace ejecutoria si se planteó con apelación subsidiaria y se rechazó la revocatoria, ni si se admite la revocatoria, pues la contraparte puede apelar.
}
}{La revocatoria permite que el mismo juez revea providencias simples. Se plantea dentro de los tres días (o en la audiencia, si allí se adoptó la decisión), admite apelación en subsidio y, como regla, lo resuelto hace ejecutoria.
}

\chapter{El incidente de nulidad}\label{cap:nulidad}

El incidente de nulidad no es un recurso, pero resulta fundamental en el proceso.

\begin{definition}{Incidente de nulidad}
Es el \textbf{remedio} procesal que apunta a la privación de efectos imputada a los actos del proceso que adolecen de algún vicio en sus elementos esenciales y que, por ende, carecen de aptitud para cumplir el fin a que están destinados.
\end{definition}

\begin{example}{Nulidad del traslado de la demanda}
La notificación del traslado de la demanda está dotada, en virtud del art. 339 CPCCN, de una serie de recaudos que hacen a la garantía del derecho de defensa. Es un acto trascendental. Si este acto, por algún vicio o defecto, no cumple su propósito (si el demandado no se entera de que fue emplazado a comparecer a un proceso), el proceso sigue adelante sin que él pueda ejercer su defensa. El incidente de nulidad es el mecanismo que permite atacar ese acto irregular.
\end{example}

\section{Principio general: conservación de los actos procesales}

El proceso tiende a que los actos procesales cumplidos sean válidos. El principio general es el de \textbf{conservación de los actos procesales}: se reputan válidos, salvo que se reúnan los criterios que habilitan la declaración de nulidad de un acto.

\section{Los cinco principios que rigen las nulidades procesales}

Los cinco principios están regulados en los arts. 169 a 174 CPCCN.

\subsection{Principio de especificidad (art. 169, primer párrafo)}

Para que exista nulidad tiene que haber una \textbf{irregularidad formal}, una contradicción a una norma que, de no haber sido observada, apareje la nulidad. Tiene que haber una pauta formal vulnerada. Si en el proceso no se violó ninguna norma formal, no se daría este principio, necesario para que haya nulidad.

\begin{formula}{Art. 169, primer párrafo, CPCCN}
Ningún acto procesal será declarado nulo si la ley no prevé expresamente esa sanción. Sin embargo, la nulidad procederá cuando el acto carezca de los requisitos indispensables para la obtención de su finalidad.
\end{formula}

\subsection{Principio de finalidad o conservación (art. 169, segundo párrafo)}

El acto viciado debe carecer de los requisitos indispensables para obtener su finalidad. Pueden existir hipótesis en las que, pese a la irregularidad del acto cumplido y atacado, éste haya cumplido la finalidad a la que estaba destinado. En ese caso prevalece la \textbf{validez} del acto.

\begin{formula}{Art. 169, segundo párrafo, CPCCN}
Si el acto, no obstante su irregularidad, ha cumplido la finalidad a la que estaba destinado, no procede la declaración de nulidad.
\end{formula}

\subsection{Principio de subsanación (arts. 170 y 172)}
% TODO: contenido incompleto o ambiguo en las notas originales (el apunte rotula este principio con los arts. 170 y 172, pero sólo transcribe el art. 170; el art. 172 se transcribe como principio de trascendencia)

La nulidad es improcedente si el acto viciado fue \textbf{convalidado} o purgado, o fue conocido por quien la invoca y éste no hizo ningún planteo. Se presume que el acto ha sido convalidado si no se promueve el incidente de nulidad dentro del \textbf{quinto día} de conocido.

\begin{formula}{Art. 170 CPCCN}
La nulidad no podrá ser declarada cuando el acto haya sido consentido, aunque fuere tácitamente, por la parte interesada en la declaración. Se considerará que hay consentimiento tácito cuando no se promoviere incidente de nulidad dentro del quinto día de conocido el acto nulo.
\end{formula}

\begin{example}{Emplazado anoticiado de manera irregular pero con conocimiento real}
Una persona se muda de domicilio pero mantiene contacto con quien vive allí. La cédula de traslado de la demanda llega al domicilio anterior. La persona que recibe la cédula le avisa al demandado, quien toma conocimiento pleno. Si el demandado no hace nada y deja transcurrir el plazo de cinco días, el acto queda saneado y ya no procede el incidente de nulidad. El acto irregular cumplió su finalidad (el emplazado supo que era demandado) y fue convalidado.
\end{example}

\subsection{Principio de trascendencia (art. 172)}

Quien plantea el incidente debe \textbf{invocar y demostrar un perjuicio cierto} del que derive su interés en formular el planteo, dimensionando las defensas que no ha podido oponer como consecuencia del acto viciado. La nulidad debe dejar desnuda ante el juez una situación de indefensión de quien la invoca.

\begin{formula}{Art. 172 CPCCN}
El incidente de nulidad es inadmisible cuando quien lo deduce no exprese el perjuicio sufrido del que derive el interés en obtener la declaración y las defensas que no ha podido oponer.
\end{formula}

\begin{important}
\textbf{«Si hay indefensión, hay nulidad. Si no hay indefensión, no hay nulidad.»} El mecanismo de las nulidades apunta a garantizar la defensa, no la regularidad de los actos o el cumplimiento de un formalismo. No existe la nulidad por la mera forma.
\end{important}

\subsection{Principio de protección (art. 171)}

Nadie puede invocar la nulidad cuando fue partícipe o gestor de la nulidad que invoca. Por ejemplo, el actor no puede plantear la nulidad del traslado de la demanda en el domicilio que él mismo denunció y en el que postuló que se hiciera la notificación.

\begin{formula}{Art. 171 CPCCN}
La parte que ha dado lugar a la nulidad no podrá pedir la invalidez del acto realizado.
\end{formula}

\cornelltitle{incidente de nulidad}
\cornellblocksum{
\cqrow{¿Qué es el incidente de nulidad?}{Es el remedio procesal que apunta a la privación de efectos de los actos del proceso que adolecen de algún vicio en sus elementos esenciales y que, por ello, carecen de aptitud para cumplir su fin. Ataca vicios procedimentales (\textit{in procedendo}), como la notificación irregular del traslado de la demanda.
}
\cqrow{¿Qué principio general rige los actos procesales?}{La conservación de los actos procesales: se reputan válidos salvo que se reúnan los criterios que habilitan la declaración de nulidad.
}
\cqrow{¿Qué establece el principio de especificidad?}{Debe haber una norma formal vulnerada (art. 169, primer párrafo). Ningún acto será declarado nulo si la ley no prevé expresamente esa sanción, salvo que el acto carezca de los requisitos indispensables para su finalidad.
}
\cqrow{¿Qué establece el principio de finalidad?}{Si el acto, pese a su irregularidad, cumplió la finalidad a la que estaba destinado, no hay nulidad (art. 169, segundo párrafo).
}
\cqrow{¿Qué establece el principio de subsanación?}{Si el acto fue convalidado, no procede la nulidad (art. 170). Hay consentimiento tácito si no se promueve el incidente dentro del quinto día de conocido el acto.
}
\cqrow{¿Qué establece el principio de trascendencia?}{Quien plantea el incidente debe invocar y demostrar un perjuicio cierto y las defensas que no pudo oponer (art. 172). Debe existir indefensión real: sin indefensión, no hay nulidad.
}
\cqrow{¿Qué establece el principio de protección?}{Quien generó el vicio no puede invocarlo (art. 171).
}
}{El incidente de nulidad ataca vicios procedimentales y exige los cinco principios: especificidad, finalidad, subsanación, trascendencia y protección. La clave es la trascendencia: sin indefensión real, no hay nulidad.
}

\chapter[Recurso de apelación: concepto, procedencia y forma]{El recurso de apelación: concepto, procedencia y forma}\label{cap:apelacion}

% --------------------------------------------------

\section{Definici\'on del recurso de apelaci\'on}

\begin{definition}{Recurso de apelaci\'on}
El recurso de apelaci\'on es un \textbf{acto procesal de impugnaci\'on} que interponen las partes o terceros ante el mismo juez que dict\'o una resoluci\'on, que debe ser una providencia simple, interlocutoria o definitiva, en tanto y en cuanto cause \textbf{gravamen irreparable}, con el objetivo de que el tribunal que est\'a por encima de ese juzgado la revoque total o parcialmente.
\end{definition}

\section{Legitimados para interponer el recurso}

No son s\'olo las partes en sentido estricto (actor o demandado): cualquier sujeto procesal con \textbf{inter\'es acreditado} puede interponer el recurso. Ejemplos:

\begin{itemize}
    \item El Ministerio P\'ublico, si est\'a representando los intereses de un menor o de un incapaz.
    \item Un perito que considera exiguos sus honorarios regulados.
    \item El abogado interviniente, respecto de la regulaci\'on de sus honorarios.
\end{itemize}

El requisito es que exista un \textbf{gravamen irreparable}: concreto, actual, que no se pueda reparar luego con el dictado de la sentencia definitiva.

\section{Resoluciones apelables: la regla de la apelabilidad}

En el proceso ordinario la regla es la \textbf{apelabilidad}: todas las resoluciones son apelables, salvo que el C\'odigo establezca lo contrario. Esto es diferente a lo que ocurre en el proceso ejecutivo o en el sumar\'isimo, donde la apelabilidad est\'a muy restringida.

\begin{important}
\textbf{Inapelable} no es lo mismo que \textbf{irrecurrible}. Si el C\'odigo dice ``irrecurrible'', la resoluci\'on no resiste ning\'un tipo de recurso (ni siquiera revocatoria). Cuando dice ``inapelable'', se refiere \'unicamente a que no resiste el recurso de apelaci\'on; otros recursos podr\'ian interponerse.
\end{important}

\subsection{Art\'iculo 379 CPCCN: inapelabilidad de resoluciones sobre prueba}

% TODO: contenido incompleto o ambiguo en las notas originales (la redacci\'on del segundo p\'arrafo del art. 379 figura as\'i en el apunte)
\begin{formula}{Art. 379 CPCCN}
\textit{Ser\'an inapelables las resoluciones del juez sobre producci\'on, denegaci\'on y sustanciaci\'on de las pruebas. Si se hubiere negado alguna medida, la parte interesada podr\'a solicitar a la c\'amara que la diligencia cuando el expediente le fuere remitido para que conozca del recurso contra la sentencia definitiva.}
\end{formula}

\begin{example}{Denegaci\'on de prueba testimonial en audiencia preliminar}
En la audiencia preliminar, el juez decide denegar la producci\'on de la prueba testimonial por considerarla innecesaria en funci\'on de los hechos controvertidos. La parte actora o demandada no puede apelar esa resoluci\'on, porque el art.~379 la declara inapelable. Sin embargo, podr\'a replantear la cuesti\'on ante la c\'amara cuando el expediente sea elevado para resolver el recurso contra la sentencia definitiva.
\end{example}

\subsection{Art\'iculo 242 CPCCN: l\'imite por monto}

\begin{formula}{Art. 242 CPCCN (actualizado por Acordada 41/2019, vigente desde el 1.º de enero de 2020)}
\textit{Ser\'an inapelables las sentencias definitivas y las dem\'as resoluciones, cualquiera fuere su naturaleza, que se dicten en los procesos en los que el valor cuestionado no supere los \$\,300.000 (pesos trescientos mil). Esta inapelabilidad no rige en materia de alimentos ni respecto de los honorarios regulados en ese proceso.}
\end{formula}

Muchos alumnos pueden tener c\'odigos procesales desactualizados: el monto se va actualizando mediante acordadas de la C\'amara, siendo la \'ultima la Acordada 41/2019.

Se ha cuestionado la constitucionalidad de este art\'iculo por limitar el acceso a la segunda instancia. La Corte Suprema ha interpretado que la \textbf{doble instancia} como principio constitucional-convencional (Pacto de San Jos\'e de Costa Rica) rige s\'olo en materia penal. En materia civil, el legislador puede establecer este l\'imite sin que resulte inconstitucional ni anticonvencional. La finalidad es que los juicios de menor trascendencia econ\'omica tramiten en una sola instancia para ser resueltos m\'as r\'apidamente.

\section{Revocatoria con apelaci\'on en subsidio}

La \'unica combinaci\'on posible de recursos es la \textbf{revocatoria con apelaci\'on en subsidio}, y s\'olo procede ante una \textbf{providencia simple que cause gravamen irreparable}. En ese supuesto, el recurrente interpone la revocatoria y, para el caso de que el juez de primera instancia la rechace, subsidiariamente est\'a la apelaci\'on para ir a la c\'amara.

\begin{example}{Providencia simple que causa gravamen irreparable}
Si se contest\'o la demanda en t\'ermino pero el juzgado la tuvo por extempor\'anea por haber computado mal el plazo de 15 d\'ias (sin tener en cuenta un feriado), eso causa un gravamen irreparable. En ese supuesto: (a) puede interponerse una revocatoria con apelaci\'on en subsidio, o (b) directamente un recurso de apelaci\'on.
\end{example}

\section{Errores que comprende el recurso de apelaci\'on}

\begin{formula}{Art. 253 CPCCN}
El recurso de apelaci\'on comprende tanto los errores \textit{in iudicando} (error en el juzgamiento: fijaci\'on de los hechos o aplicaci\'on del derecho) como los errores \textit{in procedendo} (error en el procedimiento de construcci\'on de la sentencia). Por ello, en el proceso ordinario no existe un recurso de nulidad separado: el recurso de nulidad est\'a comprendido dentro del recurso de apelaci\'on.
\end{formula}

\section{Plazo y forma de interposici\'on}

\begin{table}[H]
\centering
\begin{tabularx}{\textwidth}{
    >{\raggedright\arraybackslash}p{0.25\textwidth}
    >{\raggedright\arraybackslash}p{0.18\textwidth}
    >{\raggedright\arraybackslash}X
}
\toprule
\textbf{Proceso} & \textbf{Plazo} & \textbf{Observaci\'on} \\
\midrule
Proceso ordinario & \textbf{5 d\'ias} & Desde la notificaci\'on (c\'edula o ministerio de ley) \\
\addlinespace
Proceso sumar\'isimo & \textbf{3 d\'ias} & Plazo especial previsto por el c\'odigo \\
\addlinespace
Verbal (en audiencia) & \textbf{Inmediato} & Se deja constancia en el acta \\
\bottomrule
\end{tabularx}
\caption{Plazo de interposici\'on del recurso de apelaci\'on seg\'un el tipo de proceso.}
\end{table}

\section{Los dos momentos del recurso de apelaci\'on}

A diferencia de los recursos horizontales, el recurso de apelaci\'on tiene \textbf{dos momentos diferenciados}.

\begin{table}[H]
\centering
\begin{tabularx}{\textwidth}{
    >{\raggedright\arraybackslash}X
    >{\raggedright\arraybackslash}X
}
\toprule
\textbf{1.º momento: interposici\'on} & \textbf{2.º momento: fundamentaci\'on} \\
\midrule
Dentro de los 5 d\'ias de notificada la resoluci\'on. S\'olo se declara que se interpone el recurso de apelaci\'on contra la resoluci\'on de fecha tal, por causar gravamen irreparable. \textbf{Nada m\'as.} &
Una vez concedido el recurso, en funci\'on de c\'omo fue concedido. Se hace una cr\'itica concreta y razonada de la resoluci\'on, explicando a la c\'amara qu\'e se modifica total o parcialmente y d\'onde est\'a el error del juez de primera instancia. \newline
\textbf{Sentencia definitiva:} expresi\'on de agravios. \newline
\textbf{Otras resoluciones:} memorial. \\
\bottomrule
\end{tabularx}
\caption{Momentos del recurso de apelaci\'on.}
\end{table}

\begin{important}
Si no se fundamenta en tiempo y forma, se tiene por \textbf{desierto} el recurso: es como si no se hubiera interpuesto y la resoluci\'on queda firme.
\end{important}

\cornelltitle{el recurso de apelación}
\cornellblock{
\cqrow{¿Qu\'e es el recurso de apelaci\'on?}{Es un acto procesal de impugnaci\'on que interponen las partes o terceros ante el mismo juez que dict\'o una resoluci\'on (providencia simple, interlocutoria o definitiva), en tanto cause gravamen irreparable, para que el tribunal superior la revoque total o parcialmente.
}
\cqrow{¿Qui\'enes est\'an legitimados para apelar y qu\'e requisito se exige?}{Cualquier sujeto procesal con inter\'es acreditado: por ejemplo, el Ministerio P\'ublico que representa a un menor o incapaz, un perito respecto de sus honorarios o el abogado respecto de la regulaci\'on de los suyos. Se exige un gravamen irreparable, concreto y actual, que no pueda repararse con la sentencia definitiva.
}
\cqrow{¿Cu\'al es la regla sobre apelabilidad y qu\'e diferencia hay entre inapelable e irrecurrible?}{En el proceso ordinario la regla es la apelabilidad. ``Inapelable'' significa que no procede la apelaci\'on, aunque podr\'ian interponerse otros recursos; ``irrecurrible'' significa que la resoluci\'on no resiste ning\'un recurso.
}
\cqrow{¿Qu\'e establece el art.~379 CPCCN?}{Son inapelables las resoluciones del juez sobre producci\'on, denegaci\'on y sustanciaci\'on de las pruebas. La parte interesada puede replantear la cuesti\'on ante la c\'amara cuando el expediente sea remitido para conocer del recurso contra la sentencia definitiva.
}
\cqrow{¿Qu\'e establece el art.~242 CPCCN y por qu\'e no es inconstitucional?}{Son inapelables las resoluciones en procesos cuyo valor cuestionado no supere \$\,300.000 (Acordada 41/2019), salvo en materia de alimentos y honorarios regulados. La Corte Suprema interpret\'o que la doble instancia como principio constitucional-convencional rige s\'olo en materia penal, por lo que en materia civil el l\'imite es v\'alido.
}
\cqrow{¿Qu\'e errores comprende la apelaci\'on?}{Los errores \textit{in iudicando} (fijaci\'on de hechos o aplicaci\'on del derecho) y los errores \textit{in procedendo} (procedimiento de construcci\'on de la sentencia). El recurso de nulidad est\'a comprendido en la apelaci\'on (art.~253).
}
\cqrow{¿Cu\'al es el plazo para interponer la apelaci\'on?}{5 d\'ias en el proceso ordinario, desde la notificaci\'on; 3 d\'ias en el proceso sumar\'isimo; de forma inmediata en el proceso verbal en audiencia, con constancia en el acta.
}
}
\medskip
\cornellblocksum{
\cqrow{¿Cu\'ales son los dos momentos de la apelaci\'on?}{Interposici\'on (dentro de los 5 d\'ias, s\'olo se declara que se apela) y fundamentaci\'on (cr\'itica concreta y razonada, que es expresi\'on de agravios si se apela la sentencia definitiva y memorial en los dem\'as casos). Si no se funda en tiempo y forma, el recurso se tiene por desierto.
}
}{la apelaci\'on procede ante resoluciones que causan gravamen irreparable y su regla es la apelabilidad en el proceso ordinario, con l\'imites como los arts.~379 y 242. Comprende errores de juzgamiento y de procedimiento, se interpone en 5 d\'ias y se funda en un segundo momento; sin fundamentaci\'on el recurso queda desierto. La \'unica combinaci\'on posible de recursos es revocatoria con apelaci\'on en subsidio ante providencias simples.
}

\chapter[Concesión, trámite y efectos de la apelación]{Concesión, trámite y efectos del recurso de apelación}\label{cap:concesion}

Interpuesto el recurso de apelación, corresponde analizar cómo se concede, qué trámite sigue y qué efectos produce. Este capítulo reúne la forma de concesión del recurso (libre o en relación), su trámite (inmediato o diferido) y sus efectos (suspensivo o no suspensivo).

\section{Las tres consecuencias de la concesi\'on}

Al interponerse el recurso de apelaci\'on y ser concedido, se generan tres consecuencias: una forma, un tr\'amite y un efecto.

\begin{table}[H]
\centering
\begin{tabularx}{\textwidth}{
    >{\raggedright\arraybackslash}X
    >{\raggedright\arraybackslash}X
    >{\raggedright\arraybackslash}X
}
\toprule
\textbf{Forma} & \textbf{Tr\'amite} & \textbf{Efecto} \\
\midrule
\textbf{Libre} (amplia) vs. \textbf{en relaci\'on} (restringida) &
\textbf{Inmediato} (regla general) vs. \textbf{diferido} (excepci\'on legal) &
\textbf{Suspensivo} (regla general) vs. \textbf{no suspensivo} o devolutivo (excepci\'on) \\
\bottomrule
\end{tabularx}
\caption{Consecuencias de la concesi\'on del recurso.}
\end{table}

\section{Forma de concesi\'on: libre vs. en relaci\'on}

\begin{table}[H]
\centering
\begin{tabularx}{\textwidth}{
    >{\raggedright\arraybackslash}p{0.22\textwidth}
    >{\raggedright\arraybackslash}X
    >{\raggedright\arraybackslash}X
}
\toprule
 & \textbf{Concesi\'on libre} & \textbf{Concesi\'on en relaci\'on} \\
\midrule
\textbf{Supuesto} & S\'olo cuando se apela la sentencia definitiva de un proceso ordinario & Todos los dem\'as casos \\
\addlinespace
\textbf{Plazo para fundar} & 10 d\'ias (expresi\'on de agravios) & 5 d\'ias (memorial) \\
\addlinespace
\textbf{D\'onde se funda} & Ante la c\'amara & Ante el juzgado de primera instancia \\
\addlinespace
\textbf{Actividad adicional} & S\'i: art.~260 (5 d\'ias paralelos) & No \\
\addlinespace
\textbf{Traslado de fundamentos} & 10 d\'ias (notificaci\'on por ministerio de ley) & 5 d\'ias (notificaci\'on por ministerio de ley) \\
\addlinespace
\textbf{Efecto} & Siempre suspensivo & Depende de cada caso \\
\addlinespace
\textbf{Tr\'amite} & Siempre inmediato & Depende de cada caso \\
\bottomrule
\end{tabularx}
\caption{Concesi\'on libre y concesi\'on en relaci\'on.}
\end{table}

\subsection{Procedimiento cuando se concede en relaci\'on}

\begin{enumerate}
    \item El recurrente queda notificado de la resoluci\'on (por c\'edula o por ministerio de ley).
    \item Dentro de los 5 d\'ias interpone el recurso de apelaci\'on.
    \item El juez dicta una resoluci\'on concediendo en relaci\'on el recurso (se notifica por ministerio de ley).
    \item Desde que queda notificada esa concesi\'on, el recurrente tiene 5 d\'ias para presentar el memorial (fundamentaci\'on).
    \item Del memorial se corre traslado a la otra parte por 5 d\'ias (notificaci\'on por ministerio de ley).
    \item Vencido ese plazo o contestado el traslado, el expediente se eleva a la c\'amara.
    \item La c\'amara, a trav\'es de la sala correspondiente, resuelve directamente el recurso.
\end{enumerate}

\begin{note}
\textbf{Principio de prevenci\'on:} si a lo largo del proceso el expediente ya fue elevado a la c\'amara para resolver otro recurso, se remite a la misma sala que ya intervino. Si es la primera vez, se sortea la sala.
\end{note}

\subsection{Procedimiento cuando se concede libremente}

Este procedimiento es totalmente distinto al anterior.

\begin{enumerate}
    \item Se dicta la sentencia definitiva en un proceso ordinario.
    \item Se notifica por c\'edula. Dentro de los 5 d\'ias se interpone el recurso.
    \item El juez concede libremente todos los recursos interpuestos y el expediente es elevado directamente a la c\'amara.
    \item La sala dicta una resoluci\'on que dice, literalmente: \textit{``P\'onganse los autos a los fines de los art\'iculos 259 y 260 del C\'odigo Procesal. Notif\'iquese.''}
    \item Las partes son notificadas por c\'edula de esa resoluci\'on. Desde esa notificaci\'on corren en paralelo dos plazos:
    \begin{itemize}
        \item 10 d\'ias (art.~259): para expresar agravios.
        \item 5 d\'ias (art.~260): para realizar la actividad adicional en c\'amara.
    \end{itemize}
    \item De la expresi\'on de agravios se corre traslado a la otra parte por 10 d\'ias (notificaci\'on por ministerio de ley).
    \item El tribunal dicta resoluci\'on de autos para sentencia y vota para resolver el recurso.
\end{enumerate}

\subsection{Art\'iculo 260: actividad adicional en c\'amara}

% TODO: contenido incompleto o ambiguo en las notas originales (el apunte transcribe "elevaci\'on" con errores de tipeo en el art. 260)
\begin{formula}{Art. 260 CPCCN, facultades de las partes en segunda instancia (concesi\'on libre)}
\textit{Dentro de los cinco d\'ias de la notificaci\'on de la providencia que dispone la elevaci\'on, los litigantes podr\'an:}
\begin{enumerate}[label=\arabic*.º]
    \item \textit{Fundar los recursos concedidos con tr\'amite diferido.}
    \item \textit{Replantear la prueba denegada o declarada negligente.}
    \item \textit{Acompa\~nar documentos de fecha posterior o desconocidos al inicio del proceso.}
    \item \textit{Invocar hechos nuevos.}
    \item \textit{Pedir la absoluci\'on de posiciones sobre hechos no debatidos en la confesional de primera instancia.}
    \item \textit{Pedir apertura a prueba para la producci\'on del replanteo o de los hechos nuevos.}
\end{enumerate}
\end{formula}

\begin{example}{Replanteo de prueba (art.~260 en relaci\'on con art.~379)}
En la audiencia preliminar el juez deneg\'o la prueba testimonial por considerarla innecesaria. Luego dict\'o sentencia rechazando la demanda por no tener por acreditada la materialidad del accidente. Si bien en su momento esa resoluci\'on era inapelable (art.~379), ahora, al estar el expediente en c\'amara, el recurrente puede replantear la prueba denegada, fundando c\'omo la declaraci\'on de esos testigos presenciales podr\'ia acreditar la existencia del accidente y cambiar el resultado de la sentencia.
\end{example}

\section{L\'imites de la alzada al resolver}

La c\'amara tiene tres l\'imites para resolver:

\begin{enumerate}
    \item \textbf{Expresi\'on de agravios o memorial:} la c\'amara s\'olo entiende en aquello que haya sido materia de cr\'itica concreta y razonada. Lo que no fue objeto de agravios queda firme.
    \item \textbf{Principio de congruencia:} el tribunal no puede considerar hechos que no hayan sido invocados en la demanda, la contestaci\'on, la reconvenci\'on o la contestaci\'on de la reconvenci\'on.
    \item \textbf{\textit{Reformatio in peius}:} la c\'amara no puede empeorar la situaci\'on del recurrente. S\'olo puede modificar en lo que fue impugnado.
\end{enumerate}

\cornelltitle{concesión del recurso}
\cornellblocksum{
\cqrow{¿Qu\'e tres consecuencias genera la concesi\'on del recurso?}{Una forma (libre o en relaci\'on), un tr\'amite (inmediato o diferido) y un efecto (suspensivo o no suspensivo).
}
\cqrow{¿Cu\'ando se concede libremente?}{S\'olo cuando se apela la sentencia definitiva de un proceso ordinario. En todos los dem\'as casos se concede en relaci\'on.
}
\cqrow{¿C\'omo es el procedimiento cuando se concede en relaci\'on?}{Se interpone en 5 d\'ias; el juez concede en relaci\'on; desde la notificaci\'on de la concesi\'on el recurrente tiene 5 d\'ias para presentar el memorial ante el juzgado de primera instancia; se corre traslado por 5 d\'ias; luego el expediente se eleva a la c\'amara, que resuelve directamente. Rige el principio de prevenci\'on en la elecci\'on de la sala.
}
\cqrow{¿C\'omo es el procedimiento cuando se concede libremente?}{Concedido el recurso, el expediente se eleva a la c\'amara y la sala ordena poner los autos a los fines de los arts.~259 y 260. Desde la notificaci\'on por c\'edula corren en paralelo 10 d\'ias para expresar agravios y 5 d\'ias para la actividad adicional. De la expresi\'on de agravios se corre traslado por 10 d\'ias y el tribunal dicta autos para sentencia y vota.
}
\cqrow{¿Qu\'e puede hacerse en c\'amara seg\'un el art.~260?}{Dentro de los 5 d\'ias: fundar recursos concedidos con tr\'amite diferido, replantear prueba denegada o declarada negligente, acompa\~nar documentos posteriores o desconocidos, invocar hechos nuevos, pedir absoluci\'on de posiciones sobre hechos no debatidos en primera instancia y pedir apertura a prueba para el replanteo o los hechos nuevos.
}
\cqrow{¿Cu\'ales son los l\'imites de la alzada?}{La expresi\'on de agravios o memorial (lo que no fue criticado queda firme), el principio de congruencia (no se consideran hechos no invocados en los escritos postulatorios) y la prohibici\'on de \textit{reformatio in peius} (no se puede empeorar la situaci\'on del recurrente).
}
}{la concesi\'on determina forma, tr\'amite y efecto. La forma libre se reserva a la sentencia definitiva del proceso ordinario (fundamentaci\'on ante la c\'amara en 10 d\'ias, con actividad adicional del art.~260); en los dem\'as casos se concede en relaci\'on (memorial ante primera instancia en 5 d\'ias). La c\'amara resuelve dentro de los l\'imites de los agravios, la congruencia y la prohibici\'on de \textit{reformatio in peius}.
}

\section{Tr\'amite: inmediato vs. diferido}

La regla general es que el recurso se concede con \textbf{tr\'amite inmediato}: una vez interpuesto y concedido, de modo continuo se computa el plazo para fundar, se corre traslado y la c\'amara resuelve.

La excepci\'on es el \textbf{tr\'amite diferido}, cuando la ley lo dispone: el momento de fundamentaci\'on, sustanciaci\'on y resoluci\'on se difiere para cuando el expediente es elevado a la c\'amara por el recurso contra la sentencia definitiva.

\begin{important}
Advertencia terminol\'ogica: el legislador usa la expresi\'on ``efecto diferido'' en el c\'odigo, pero en realidad se trata de un tema de tr\'amite, no de efecto. Cuando el c\'odigo dice ``efecto diferido'' debe leerse ``tr\'amite diferido''.
\end{important}

\subsection{Supuestos de tr\'amite diferido}

\begin{itemize}
    \item \textbf{Art\'iculo 366, hechos nuevos:} la resoluci\'on que rechaza el planteo de hechos nuevos puede apelarse con tr\'amite diferido. La fundamentaci\'on se hace cuando el expediente llega a la c\'amara para resolver el recurso contra la sentencia definitiva. Finalidad: evitar que el expediente sea elevado varias veces antes de la sentencia.
    \item \textbf{Art\'iculo 69, costas y honorarios incidentales:} las apelaciones sobre costas y honorarios regulados a lo largo del proceso se conceden con tr\'amite diferido para el mismo momento. Finalidad: celeridad y econom\'ia procesal.
\end{itemize}

\section{Efectos: suspensivo vs. no suspensivo (devolutivo)}

El efecto hace referencia a qu\'e sucede con la resoluci\'on impugnada: si se suspenden o no sus efectos. La regla general es el \textbf{efecto suspensivo}: al interponerse el recurso, se suspende la ejecutoriedad de la resoluci\'on hasta que la c\'amara resuelva. La raz\'on es que la c\'amara puede modificar o revocar la resoluci\'on, por lo que tiene sentido no ejecutarla mientras est\'a siendo revisada.

\subsection{Origen hist\'orico del t\'ermino ``devolutivo''}

El t\'ermino ``devolutivo'' que usa el c\'odigo tiene origen en el derecho romano: el \'unico con jurisdicci\'on era el emperador, quien la delegaba en los jueces. Cuando una parte impugnaba una decisi\'on, se ``devolv\'ia'' la jurisdicci\'on al emperador. En el sistema actual, todos los jueces (desde el de primera instancia hasta la Corte Suprema) tienen la misma funci\'on jurisdiccional; lo que var\'ia es el \'ambito de su competencia. Por eso el t\'ermino ``devolutivo'' puede llevar a confusi\'on.

\begin{important}
Cuando el c\'odigo dice ``devolutivo'', debe leerse \textbf{``no suspensivo''}.
\end{important}

\subsection{Supuestos de efecto no suspensivo}

\begin{itemize}
    \item \textbf{Alimentos (sentencia que admite):} la sentencia que fija alimentos a cargo del demandado puede ejecutarse aunque est\'e pendiente el recurso de apelaci\'on del demandado ante la c\'amara.
    \item \textbf{Medidas cautelares (art.~198):} la resoluci\'on que concede una medida cautelar (por ejemplo, un embargo) puede apelarse por el afectado, pero el recurso tiene efecto no suspensivo: el embargo se traba y contin\'ua trab\'andose sin importar que el recurso est\'e pendiente ante la c\'amara.
    \item \textbf{Beneficio de litigar sin gastos (art.~81):} la resoluci\'on que concede el beneficio tiene efecto no suspensivo: se aplica de inmediato, aunque haya recurso en tr\'amite.
\end{itemize}

\section{Art\'iculo 250: procedimiento seg\'un el tipo de resoluci\'on}

Cuando el recurso se concede con efecto no suspensivo, el procedimiento var\'ia seg\'un el tipo de resoluci\'on impugnada.

\begin{table}[H]
\centering
\begin{tabularx}{\textwidth}{
    >{\raggedright\arraybackslash}X
    >{\raggedright\arraybackslash}X
}
\toprule
\textbf{La resoluci\'on es una sentencia} & \textbf{La resoluci\'on es una interlocutoria} \\
\midrule
El expediente se eleva a la c\'amara para que pueda revisar todo el proceso y resolver el recurso. Para ejecutar lo resuelto en la sentencia, se forma un \textbf{incidente de ejecuci\'on de sentencia} que queda en primera instancia (con copia certificada de la sentencia). \newline
\textit{Expediente principal $\rightarrow$ c\'amara. Incidente de ejecuci\'on $\rightarrow$ primera instancia.} &
El expediente queda en primera instancia para continuar el proceso (la sentencia definitiva a\'un no fue dictada). Se forma un \textbf{incidente de apelaci\'on} con las copias indicadas por el juez, que se eleva a la c\'amara. El apelante tiene 5 d\'ias para acompa\~nar esas copias. \newline
\textit{Expediente principal $\rightarrow$ primera instancia. Incidente de apelaci\'on $\rightarrow$ c\'amara.} \\
\bottomrule
\end{tabularx}
\caption{Procedimiento del art\'iculo 250 seg\'un el tipo de resoluci\'on (efecto no suspensivo).}
\end{table}

Si el apelante no acompa\~na las copias para formar el incidente dentro de los 5 d\'ias, se declara la \textbf{deserci\'on} del recurso: es como si no se hubiera interpuesto.

\cornelltitle{trámite y efectos de la apelación}
\cornellblocksum{
\cqrow{¿Qu\'e es el tr\'amite diferido?}{Es la excepci\'on legal a la regla del tr\'amite inmediato: la fundamentaci\'on, sustanciaci\'on y resoluci\'on del recurso se difieren para cuando el expediente sea elevado a la c\'amara por el recurso contra la sentencia definitiva.
}
\cqrow{¿Cu\'ales son los supuestos de tr\'amite diferido y su finalidad?}{El art.~366 (rechazo del planteo de hechos nuevos) y el art.~69 (costas y honorarios incidentales). Su finalidad es la celeridad y la econom\'ia procesal, evitando que el expediente se eleve varias veces antes de la sentencia.
}
\cqrow{¿Qu\'e significa ``efecto diferido'' en el c\'odigo?}{Se trata de un tema de tr\'amite y no de efecto: donde el c\'odigo dice ``efecto diferido'' debe leerse ``tr\'amite diferido''.
}
\cqrow{¿Qu\'e es el efecto suspensivo?}{Al interponerse el recurso se suspende la ejecutoriedad de la resoluci\'on impugnada hasta que la c\'amara resuelva. Es la regla general, porque la c\'amara puede modificar o revocar la resoluci\'on.
}
\cqrow{¿Qu\'e es el efecto no suspensivo y c\'omo debe leerse ``devolutivo''?}{Es la excepci\'on: la resoluci\'on puede ejecutarse aunque el recurso est\'e pendiente. El t\'ermino ``devolutivo'' proviene del derecho romano (devoluci\'on de la jurisdicci\'on al emperador) y debe leerse como ``no suspensivo''. Ejemplos: alimentos, medidas cautelares (art.~198) y beneficio de litigar sin gastos (art.~81).
}
\cqrow{¿Qu\'e ocurre, seg\'un el art.~250, con el efecto no suspensivo?}{Si se apela una sentencia, el expediente va a la c\'amara y se forma un incidente de ejecuci\'on que queda en primera instancia. Si se apela una interlocutoria, el expediente queda en primera instancia y se forma un incidente de apelaci\'on que se eleva a la c\'amara, con copias que el apelante debe acompa\~nar en 5 d\'ias; de lo contrario, el recurso se declara desierto por deserci\'on.
}
}{el tr\'amite es inmediato como regla y diferido por excepci\'on legal (arts.~366 y 69); el efecto es suspensivo como regla y no suspensivo (devolutivo) por excepci\'on (alimentos, cautelares, beneficio de litigar sin gastos). Con efecto no suspensivo, el art.~250 organiza el expediente seg\'un se apele una sentencia o una interlocutoria, y la falta de copias en 5 d\'ias produce la deserci\'on.
}

\chapter[Reglas y mecanismos especiales de la apelación]{Reglas, ejercicios y mecanismos especiales de la apelación}\label{cap:reglas}

\section{Las seis reglas del recurso de apelación}

Las reglas estudiadas permiten sintetizar el r\'egimen del recurso de apelaci\'on.

\begin{longtable}{
    >{\raggedright\arraybackslash}p{0.13\textwidth}
    >{\raggedright\arraybackslash}p{0.22\textwidth}
    >{\raggedright\arraybackslash}p{0.50\textwidth}
}
\toprule
\textbf{N.º} & \textbf{Regla} & \textbf{Contenido} \\
\midrule
\endfirsthead
\toprule
\textbf{N.º} & \textbf{Regla} & \textbf{Contenido} \\
\midrule
\endhead
\bottomrule
\endlastfoot

\textbf{Regla 1} & \textbf{Alcance del recurso} & El recurso de apelaci\'on comprende tanto los errores \textit{in iudicando} (fijaci\'on de hechos o aplicaci\'on del derecho) como los errores \textit{in procedendo} (procedimiento de construcci\'on de la sentencia). El recurso de nulidad est\'a incluido dentro del de apelaci\'on (art.~253 CPCCN). \\
\addlinespace
\textbf{Regla 2} & \textbf{Plazo de interposici\'on} & El plazo para interponer el recurso es de 5 d\'ias, salvo que la ley disponga lo contrario (por ejemplo, proceso sumar\'isimo: 3 d\'ias). \\
\addlinespace
\textbf{Regla 3} & \textbf{Forma de concesi\'on} & El recurso se concede siempre en relaci\'on, salvo un \'unico supuesto en que se concede libremente: cuando lo que se apela es la sentencia definitiva dictada en un proceso ordinario. \\
\addlinespace
\textbf{Regla 4} & \textbf{Libre: siempre suspensivo e inmediato} & Cuando el recurso se concede libremente (sentencia definitiva de proceso ordinario), siempre es con efecto suspensivo y tr\'amite inmediato. \\
\addlinespace
\textbf{Regla 5} & \textbf{Tr\'amite} & El recurso se concede siempre con tr\'amite inmediato, salvo que la ley disponga expresamente el tr\'amite diferido (por ejemplo, arts.~366 y 69 CPCCN). \\
\addlinespace
\textbf{Regla 6} & \textbf{Efecto} & El recurso se concede siempre con efecto suspensivo, salvo que la ley disponga lo contrario (efecto no suspensivo o devolutivo; por ejemplo, alimentos, medidas cautelares y beneficio de litigar sin gastos). \\
\end{longtable}

\cornelltitle{las seis reglas de la apelación}
\cornellblocksum{
\cqrow{¿Cu\'ales son las reglas sobre alcance y plazo?}{La apelaci\'on comprende errores \textit{in iudicando} e \textit{in procedendo}, con el recurso de nulidad incluido (art.~253). El plazo es de 5 d\'ias, salvo disposici\'on legal en contrario, como en el proceso sumar\'isimo (3 d\'ias).
}
\cqrow{¿Cu\'ales son las reglas sobre concesi\'on libre?}{El recurso se concede siempre en relaci\'on, salvo cuando se apela la sentencia definitiva de un proceso ordinario, en cuyo caso se concede libremente, siempre con efecto suspensivo y tr\'amite inmediato.
}
\cqrow{¿Cu\'ales son las reglas sobre tr\'amite y efecto?}{El tr\'amite es siempre inmediato salvo disposici\'on expresa de tr\'amite diferido (arts.~366 y 69). El efecto es siempre suspensivo salvo disposici\'on legal en contrario (alimentos, medidas cautelares, beneficio de litigar sin gastos).
}
}{cada regla del recurso de apelaci\'on admite excepciones legales: plazo de 5 d\'ias (salvo 3 en el sumar\'isimo), concesi\'on en relaci\'on (salvo la libre para la sentencia definitiva del proceso ordinario), tr\'amite inmediato (salvo diferido) y efecto suspensivo (salvo no suspensivo).
}

\section{Admisibilidad vs. fundabilidad: \textquestiondown qui\'en analiza qu\'e?}

\begin{exercise}{Admisibilidad y fundabilidad}
\textbf{Pregunta:} \textit{¿Cu\'al va a ser el juez o tribunal que analiza, en el caso del recurso de apelaci\'on, los requisitos de admisibilidad del recurso y los de fundabilidad?}

\medskip
\textbf{Respuesta:}
\begin{itemize}
    \item \textbf{Admisibilidad $\rightarrow$ juez de primera instancia:} analiza si el recurso fue interpuesto en el plazo legal, si la resoluci\'on es susceptible de apelaci\'on y si quien lo interpone es el legitimado. Lo hace al momento de conceder o denegar el recurso.
    \item \textbf{Fundabilidad $\rightarrow$ c\'amara:} analiza si la expresi\'on de agravios o el memorial constituyen una cr\'itica concreta y razonada, y decide si hace lugar o no al recurso (confirma, modifica o revoca la resoluci\'on).
\end{itemize}
\end{exercise}

\section{Caso pr\'actico: contestaci\'on de demanda declarada extempor\'anea}

\begin{exercise}{Contestaci\'on de demanda extempor\'anea}
\textbf{Supuesto:} se act\'ua como abogados del demandado en un proceso ordinario de da\~nos y perjuicios por mala praxis m\'edica. Se tienen 15 d\'ias para contestar demanda y se la contesta el d\'ia 16, dentro del plazo de gracia (las dos primeras horas del d\'ia, antes de las 9:30 a.\,m.). El d\'ia siguiente, el juzgado emite esta resoluci\'on:

\begin{quote}
\textit{``Buenos Aires, 22 de junio de 2020. DESGLOSESE por extempor\'aneo.''} Firma: Mar\'ia Eugenia G\'omez, Secretaria.
\end{quote}

El juzgado no tuvo en cuenta un feriado intermedio y por eso comput\'o mal el plazo. \textbf{¿Qu\'e se hace?}
\end{exercise}

Pueden plantearse las siguientes respuestas:

\begin{itemize}
    \item Interponer una revocatoria (art.~238 y ss.): respuesta incorrecta.
    \item Interponer revocatoria con apelaci\'on en subsidio: respuesta incorrecta.
    \item Interponer directamente un recurso de apelaci\'on: respuesta incorrecta.
\end{itemize}

\begin{important}
Todas son incorrectas porque la resoluci\'on \textbf{no est\'a firmada por el juez, sino por la secretaria del juzgado}. El secretario y el prosecretario act\'uan por delegaci\'on. S\'olo el juez tiene jurisdicci\'on y s\'olo el juez puede revocar por contrario imperio. La resoluci\'on no es susceptible ni de revocatoria ni de apelaci\'on. \textbf{La v\'ia correcta es la revocaci\'on nominada del art\'iculo 38 ter del CPCCN.}
\end{important}

\section{Revocaci\'on nominada: art\'iculo 38 ter CPCCN}

\begin{formula}{Art. 38 ter CPCCN}
\textit{Dentro del plazo de tres d\'ias desde que se notifican de la resoluci\'on dictada por el secretario o prosecretario administrativo, las partes pueden requerir al juez que deje sin efecto lo dispuesto por el secretario o prosecretario administrativo.}
\end{formula}

Diferencias con la revocatoria del art.~238:

\begin{table}[H]
\centering
\begin{tabularx}{\textwidth}{
    >{\raggedright\arraybackslash}p{0.25\textwidth}
    >{\raggedright\arraybackslash}X
    >{\raggedright\arraybackslash}X
}
\toprule
 & \textbf{Revocatoria (art.~238)} & \textbf{Revocaci\'on nominada (art.~38 ter)} \\
\midrule
\textbf{Plazo} & 5 d\'ias & \textbf{3 d\'ias} \\
\addlinespace
\textbf{Fundamentaci\'on} & Simult\'anea a la interposici\'on & Simult\'anea a la interposici\'on \\
\addlinespace
\textbf{Resoluci\'on que ataca} & Providencias simples dictadas por el juez & Resoluciones dictadas por secretario o prosecretario \\
\addlinespace
\textbf{Qui\'en resuelve} & El mismo juez & El juez (revisa lo dispuesto por el auxiliar) \\
\bottomrule
\end{tabularx}
\caption{Revocatoria y revocaci\'on nominada.}
\end{table}

\section{Reposici\'on \textit{in extremis}: creaci\'on pretoriana}

La reposici\'on \textit{in extremis} es un mecanismo impugnatorio \textbf{no regulado en el c\'odigo}, creado por la jurisprudencia. No debe confundirse ni con la revocatoria del art.~238 ni con la revocaci\'on del art.~38 ter.

\begin{important}
Requisitos de la reposici\'on \textit{in extremis}:
\begin{itemize}
    \item Se trata de una resoluci\'on que \textbf{no} es una providencia simple.
    \item No existe otro remedio ni recurso posible.
    \item El error es de \'indole \textbf{sustancial} (de fondo), no meramente formal.
    \item El error es \textbf{notorio}: de mantenerse, se consolidar\'ia una injusticia muy grande para alguno de los litigantes.
    \item Se interpone y se funda en el mismo acto. Plazo: 3 d\'ias.
\end{itemize}
\end{important}

\begin{example}{Sentencia inapelable por monto con error sustancial}
Una sentencia definitiva dictada en un proceso donde la demanda es de \$\,200.000 (menor al m\'inimo del art.~242) es inapelable. Si esa sentencia tiene un error sustancial en la fijaci\'on de los hechos que constituye una injusticia notoria, la parte afectada puede interponer la reposici\'on \textit{in extremis} dentro de los 3 d\'ias, fund\'andola en ese mismo acto.
\end{example}

\subsection{Jurisprudencia sobre la reposici\'on \textit{in extremis}}

\begin{note}
\textbf{C\'amara Civil, Sala E (sumario citado en las notas).} Las resoluciones del tribunal de alzada no son, en principio, susceptibles de recurso de reposici\'on, por no revestir el car\'acter de providencias simples. Este principio reconoce excepciones en circunstancias excepcionales: cuando media un error esencial en la apreciaci\'on de los antecedentes del caso y el mantenimiento de la situaci\'on conducir\'ia a un resultado injusto que es deber de los jueces modificar, para preservar a los litigantes.
\end{note}

La Corte Suprema de Justicia tambi\'en ha dicho que sus decisiones no son susceptibles de recurso de reposici\'on, aunque ese principio puede reconocer excepciones cuando se trata de situaciones serias e inequ\'ivocas que demuestran con nitidez manifiesta el error que se pretende subsanar.

\cornelltitle{ejercicios y mecanismos especiales de la apelación}
\cornellblocksum{
\cqrow{¿Qu\'e analiza el juez de primera instancia y qu\'e la c\'amara?}{El juez de primera instancia analiza la admisibilidad (plazo legal, resoluci\'on susceptible de apelaci\'on y legitimado) al conceder o denegar el recurso. La c\'amara analiza la fundabilidad: si la expresi\'on de agravios o el memorial son una cr\'itica concreta y razonada, y decide si confirma, modifica o revoca.
}
\cqrow{¿Por qu\'e no procede revocatoria ni apelaci\'on contra una resoluci\'on firmada por el secretario?}{Porque el secretario y el prosecretario act\'uan por delegaci\'on: s\'olo el juez tiene jurisdicci\'on y s\'olo \'el puede revocar por contrario imperio.
}
\cqrow{¿Qu\'e es la revocaci\'on nominada del art.~38 ter?}{Es el mecanismo por el cual, dentro de los 3 d\'ias de notificada la resoluci\'on dictada por el secretario o prosecretario, las partes piden al juez que la deje sin efecto. Se interpone y funda en el mismo acto.
}
\cqrow{¿Qu\'e diferencia hay entre el art.~38 ter y la revocatoria del art.~238?}{La revocatoria (art.~238) tiene plazo de 5 d\'ias y ataca providencias simples dictadas por el juez; la revocaci\'on nominada (art.~38 ter) tiene plazo de 3 d\'ias y ataca resoluciones del secretario o prosecretario. En ambas la fundamentaci\'on es simult\'anea a la interposici\'on.
}
\cqrow{¿Cu\'ando procede la reposici\'on \textit{in extremis}?}{Cuando la resoluci\'on no es una providencia simple, no existe otro remedio o recurso posible, el error es sustancial y notorio, y su mantenimiento consagrar\'ia una injusticia grande. Se interpone y funda en el mismo acto, en 3 d\'ias. Es de creaci\'on pretoriana y no est\'a regulada en el c\'odigo.
}
\cqrow{¿Qu\'e han dicho la C\'amara Civil (Sala E) y la Corte Suprema sobre la reposici\'on?}{Que sus resoluciones no son, en principio, susceptibles de reposici\'on, pero ese principio admite excepciones ante un error esencial o ante situaciones serias e inequ\'ivocas que demuestran con nitidez el error que se pretende subsanar.
}
}{la admisibilidad se analiza en primera instancia y la fundabilidad en la c\'amara. Contra resoluciones del secretario o prosecretario procede la revocaci\'on nominada del art.~38 ter (3 d\'ias, fundamentaci\'on simult\'anea). La reposici\'on \textit{in extremis}, de creaci\'on pretoriana, opera cuando no existe otro remedio, el error es sustancial y notorio y su mantenimiento implicar\'ia una injusticia grave. Identificar correctamente el tipo de resoluci\'on y qui\'en la firm\'o es la clave para elegir la v\'ia impugnativa correcta.
}

\section{Síntesis general de la impugnación y los recursos}

\begin{important}
Esta parte desarrolla la \textbf{teoría general de la impugnación}, que constituye el sistema mediante el cual el ordenamiento procesal civil argentino permite atacar las decisiones judiciales erróneas o los actos procesales viciados. El fundamento último es la falibilidad humana: los jueces y las partes pueden equivocarse, y ese error requiere mecanismos de corrección.
\end{important}

\cornelltitle{general de la impugnación y los recursos}
\cornellblock{
\cqrow{¿Qué es un acto de impugnación y cuál es la diferencia con un recurso?}{Los actos de impugnación son el género: todo mecanismo procesal dirigido a atacar un acto o procedimiento que lesione el interés de las partes o de terceros. Los recursos son una especie dentro de ese género. Otros tipos son el incidente de nulidad, la acción autónoma de revisión de cosa juzgada, el juicio ordinario posterior, etc.
}
\cqrow{¿Por qué existen los mecanismos de impugnación?}{Porque el error es inherente a la condición humana. Los jueces y las partes pueden equivocarse. Carnelutti decía que la amenaza del error pende como espada de Damocles en el proceso. La impugnación existe para mitigar esas deficiencias.
}
\cqrow{¿Cuál es la diferencia entre vicio \textit{in procedendo} y vicio \textit{in iudicando}?}{El vicio \textit{in procedendo} es un defecto en el procedimiento (por ejemplo, una notificación irregular) y se ataca por incidentes de nulidad o remedios. El vicio \textit{in iudicando} es un error en la decisión judicial (interpretación de la ley o de los hechos) y se ataca por recursos en sentido estricto, ante un tribunal superior.
}
\cqrow{¿Quiénes tienen legitimación para recurrir y qué es el agravio?}{Las partes, los terceros intervinientes, los ministerios públicos y las personas ajenas al proceso con interés legítimo afectado. El agravio es la insatisfacción total o parcial del postulado esgrimido; sin agravio no hay legitimación para recurrir.
}
\cqrow{¿Existe garantía constitucional de doble instancia en materia civil?}{No. En materia no penal no existe garantía constitucional de doble instancia. En materia penal sí, desde la incorporación del Pacto de San José de Costa Rica (art. 8º inc. 2º ap. h) a la Constitución vía art. 75 inc. 22 CN, ratificada por el caso Giroldi.
}
\cqrow{¿Qué es la \textit{reformatio in peius}?}{Prohíbe que el tribunal de segunda instancia empeore la situación del único apelante.
}
\cqrow{¿Qué es el recurso indiferente?}{Una categoría doctrinal no tipificada en el CPCCN que habilita a los tribunales a dejar de lado los requisitos formales de admisibilidad para garantizar derechos fundamentales gravemente afectados. Se aplica excepcionalmente y tiene sustento en los casos «YPF» y «Tejida» de la CSJN.
}
}
\medskip
\cornellblocksum{
\cqrow{¿Para qué sirve la aclaratoria y contra qué procede la revocatoria?}{La aclaratoria sirve para que el mismo juez corrija errores materiales, aclare conceptos oscuros o integre la resolución; no procede para modificar sustancialmente lo decidido. La revocatoria procede sólo contra providencias simples, causen o no gravamen irreparable (art. 238 CPCCN), y no contra sentencias interlocutorias ni definitivas.
}
\cqrow{¿Cuáles son los cinco principios de las nulidades procesales?}{1) Especificidad (art. 169, primer párrafo): debe haber una norma formal vulnerada. 2) Finalidad (art. 169, segundo párrafo): si el acto cumplió su fin, no hay nulidad. 3) Subsanación (art. 170): si el acto fue convalidado (no se plantea en cinco días), no procede. 4) Trascendencia (art. 172): debe existir indefensión real; sin indefensión, no hay nulidad. 5) Protección (art. 171): quien generó el vicio no puede invocarlo.
}
}{Esta parte desarrolla la teoría general de la impugnación. El género es el acto de impugnación (que abarca el incidente de nulidad, los recursos y otros medios atípicos) y la especie más frecuente son los recursos. La doctrina diferencia remedios (resueltos por el mismo juez, orientados a vicios procedimentales) y recursos en sentido estricto (sometidos a un tribunal superior, dirigidos a vicios en la decisión). Todo recurso requiere legitimación, agravio y plazo perentorio, y presenta dos dimensiones: admisibilidad formal y fundabilidad sustancial. La garantía de la doble instancia existe con rango constitucional sólo en materia penal. El tribunal de alzada tiene una doble limitación (los escritos constitutivos y los agravios expresados) y no puede perjudicar al único apelante. El recurso indiferente es una construcción jurisprudencial excepcional. La aclaratoria corrige errores formales o integra omisiones sin alterar la esencia de lo decidido; la revocatoria permite que el mismo juez revea providencias simples. El incidente de nulidad ataca vicios procedimentales y exige los cinco principios; sin indefensión real, no hay nulidad.
}

\part{Ejecución de la sentencia y juicio ejecutivo}
\markboth{Parte IV}{Parte IV}
Esta parte aborda la ejecución de sentencia en el proceso civil y comercial argentino (Libro Tercero del Código Procesal Civil y Comercial de la Nación). Se estudia su naturaleza jurídica como etapa del proceso de conocimiento, sus diferencias con el juicio ejecutivo, los requisitos para ejecutar una sentencia (que debe estar consentida o ejecutoriada), el procedimiento según el tipo de obligación (dar dinero líquido, ilíquido, hacer, no hacer o dar cosas), las excepciones admisibles conforme al artículo 506 del CPCCN y los títulos que habilitan este proceso especial de ejecución.

\textbf{Hilo conductor.} Así como en un partido de fútbol primero se define el reglamento, luego se juega y finalmente se ejecuta el resultado (el equipo ganador cobra el trofeo), el proceso judicial tiene una etapa de conocimiento donde se discute y un fallo que solo cobra vida útil cuando se ejecuta. El recorrido de esta parte corresponde a ese último tramo: cómo hacer efectiva la sentencia.

Luego del estudio de la ejecución de la sentencia, la parte desarrolla el juicio ejecutivo y el cumplimiento de la sentencia de remate.

\textbf{El juicio ejecutivo.} Se desarrolla el juicio ejecutivo como proceso de conocimiento sumario: concepto, diferencia con el proceso ordinario, títulos ejecutivos (completos e incompletos), preparación de la vía ejecutiva, embargo ejecutivo, excepciones admisibles, sentencia de trance y remate, apelación, costas y ejecuciones especiales.

\textbf{Cumplimiento de la sentencia de remate.} Se examina la etapa coercitiva del proceso: embargo ejecutorio, subasta de bienes muebles e inmuebles, designación del martillero, publicación de edictos, informes registrales, base de la subasta, acto de remate, aprobación, depósito del saldo, levantamiento de gravámenes, distribución del producido, sobreseimiento, nulidades y adquisición del título.

\begin{example}{del cheque no pagado}
Un vecino firmó un cheque para pagar una refacción pero luego no lo pagó. El acreedor no necesita iniciar un largo juicio ordinario: lleva directamente ese cheque al tribunal (juicio ejecutivo), el juez embarga de inmediato y, si el deudor no paga ni opone excepciones válidas, se dicta sentencia de trance y remate. Una vez firme esa sentencia comienza el cumplimiento de la sentencia de remate: el acreedor ejecuta los bienes del deudor ---por ejemplo, su departamento--- a través de una subasta judicial dirigida por un martillero, y el producido de esa subasta se aplica al crédito.
\end{example}

Ambos son los dos momentos de un mismo proceso: primero se discute, con limitaciones, si se debe, y luego se cobra por la fuerza.

\chapter[Naturaleza y encuadre de la ejecución]{Naturaleza jurídica y encuadre de la ejecución de sentencia}\label{cap:ejec-naturaleza}

\section{Ubicación normativa en el Código Procesal}

La ejecución de sentencia se encuentra regulada en el Libro Tercero del CPCCN, que trata de los procesos de ejecución. Dentro de este libro se distinguen tres títulos:

\begin{itemize}
    \item \textbf{Título I:} ejecución de sentencias.
    \item \textbf{Título II:} juicio ejecutivo.
    \item \textbf{Título III:} ejecuciones especiales (hipotecaria, prendaria, fiscal y comercial).
\end{itemize}

\section{La ejecución de sentencia no es un proceso independiente}

Si bien el Código incluye a la ejecución de sentencia como un proceso de ejecución más, y por ello cierta doctrina la considera un proceso independiente del principal en el cual fue dictada la sentencia, la posición que aquí se sostiene es que la ejecución de sentencia es tan solo \textbf{una etapa más dentro del proceso de conocimiento} en el que fue dictada.

El proceso de conocimiento comprende las siguientes etapas:

\begin{enumerate}
    \item \textbf{Etapa introductoria:} demanda, contestación, reconvención y excepciones.
    \item \textbf{Etapa probatoria:} medios de prueba para formar convicción en el juez.
    \item \textbf{Etapa decisoria:} valoración de la prueba y dictado de la sentencia.
    \item \textbf{Etapa recursiva:} mecanismos impugnativos.
\end{enumerate}

Una vez agotada la etapa recursiva se accede a la ejecución. El hecho de que el Código la trate como un proceso de ejecución es lo que permite a la doctrina contraria sostener su independencia, pero para la posición aquí adoptada se trata de una etapa del proceso de conocimiento.

\section{Diferencias entre ejecución de sentencia y juicio ejecutivo}

La \textbf{intimación de pago} no es un requisito esencial en la ejecución de sentencia, mientras que sí lo es en el juicio ejecutivo, que es un proceso de ejecución obviamente independiente. La ejecución de sentencia no requiere intimación de pago porque ello se cubre con la notificación de la sentencia al demandado. En cambio, el juicio ejecutivo exige la intimación de pago porque es la que permite al ejecutado oponer excepciones para defenderse: al ser intimado de pago puede deducir las excepciones del artículo 544, mediante las cuales puede cuestionar el título base de la ejecución.

\begin{formula}{Art. 544 CPCCN}
Excepciones admisibles en el juicio ejecutivo. El ejecutado puede oponer, al ser intimado de pago, estas excepciones para cuestionar el título base de la ejecución: falsedad e inhabilidad de título, entre otras.
\end{formula}

Ambos procesos tienen en común que en ellos se ejecutan títulos, pero con notables diferencias. En el juicio ejecutivo se ejecutan \textbf{títulos ejecutivos}, que si bien documentan derechos, presumiblemente no tienen la certeza que posee el título ejecutorio, que es la sentencia.

\begin{table}[H]
\centering
\small
\begin{tabularx}{\textwidth}{
    >{\raggedright\arraybackslash}p{0.19\textwidth}
    >{\raggedright\arraybackslash}X
    >{\raggedright\arraybackslash}X
}
\toprule
\textbf{Aspecto} & \textbf{Juicio ejecutivo} & \textbf{Ejecución de sentencia} \\
\midrule
Tipo de título &
Título ejecutivo (derecho presunto). &
Título ejecutorio: sentencia (derecho cierto). \\[0.3em]

Intimación de pago &
Esencial. Habilita las excepciones del art.~544. &
No es requisito. La sentencia ya fue notificada. \\[0.3em]

Obligación que instrumenta &
Solo obligación de dar cantidad de dinero líquida o fácilmente liquidable. &
Obligaciones de dar, hacer o no hacer. \\[0.3em]

Embargo &
Optativo (puede trabarse o no). &
Absolutamente esencial para la ejecución. \\[0.3em]

Excepciones admisibles &
Art.~544 (más amplias, pueden abrirse a prueba). &
Art.~506 (sumamente limitadas, solo hechos posteriores a la sentencia). \\[0.3em]

Apelaciones durante el proceso &
Concedidas con trámite diferido hasta apelar la sentencia de remate. &
Concedidas con trámite diferido hasta apelar la sentencia de venta. \\
\bottomrule
\end{tabularx}
\caption{Diferencias entre el juicio ejecutivo y la ejecución de sentencia.}
\end{table}

\section{La sentencia como título ejecutorio: el derecho cierto}

La sentencia también constituye un título, pero un título que contiene un \textbf{derecho cierto}: ese derecho ha sido discutido con la suficiente amplitud en el juicio de conocimiento en el que se dictó la sentencia, y la \textbf{cosa juzgada} impide la discusión de hechos anteriores a ella.

Quien tiene a su favor una sentencia tiene, por ello, un título que contiene un derecho cierto. Por eso se lo llama \textbf{título ejecutorio}: trae aparejada una coacción inmediata para agredir el patrimonio del condenado que no cumple la sentencia. El condenado dispone de algunas defensas, pero sumamente limitadas, que se canalizan mediante la oposición de las excepciones que consagra el artículo 506 del CPCCN.

\begin{formula}{Art. 506 CPCCN}
Excepciones admisibles en la ejecución de sentencia. Solo se consideran legítimas: 1) falsedad de la ejecutoria; 2) prescripción de la ejecutoria; 3) pago; 4) quita, espera o remisión. En todos los casos, los hechos deben ser posteriores al dictado de la sentencia y probarse con constancias del expediente o documentos emanados del ejecutante.
\end{formula}

\cornelltitle{naturaleza y encuadre de la ejecución de sentencia}
\cornellblocksum{
\cqrow{¿Dónde se regula la ejecución de sentencia?}{En el Libro Tercero del CPCCN (procesos de ejecución), Título I. El Título II regula el juicio ejecutivo y el Título III las ejecuciones especiales (hipotecaria, prendaria, fiscal y comercial).
}
\cqrow{¿Por qué la ejecución de sentencia no es un proceso independiente?}{Porque es la culminación lógica del proceso de conocimiento: una etapa más dentro de él. Está íntimamente ligada a la sentencia que se ejecuta, es competente el mismo juez que conoció en la causa y la intimación de pago no es un requisito esencial, ya que se cubre con la notificación de la sentencia. El Código la trata como proceso de ejecución, y ello permite a la doctrina contraria sostener su independencia.
}
\cqrow{¿Qué diferencia al título ejecutorio del título ejecutivo?}{El título ejecutivo (juicio ejecutivo) contiene un derecho \textbf{presunto}, que puede discutirse con mayor amplitud (art.~544). El título ejecutorio (sentencia) contiene un derecho \textbf{cierto}, porque ya fue debatido con amplitud en el proceso de conocimiento. La cosa juzgada impide discutir hechos anteriores al dictado de la sentencia.
}
\cqrow{¿Qué papel cumple el embargo en cada proceso?}{En el juicio ejecutivo el embargo es optativo; en la ejecución de sentencia es absolutamente esencial.
}
\cqrow{¿Qué excepciones se admiten en cada proceso?}{En el juicio ejecutivo, las del art.~544 (más amplias, pueden abrirse a prueba). En la ejecución de sentencia, las del art.~506 (sumamente limitadas, solo hechos posteriores a la sentencia).
}
}{La ejecución de sentencia es una etapa del proceso de conocimiento y no un proceso autónomo. Ejecuta un título ejecutorio (la sentencia) que contiene un derecho cierto amparado por la cosa juzgada, a diferencia del juicio ejecutivo, que ejecuta títulos con derecho presunto, exige intimación de pago y admite las excepciones más amplias del art.~544. En la ejecución de sentencia el embargo es esencial y las excepciones se limitan a las del art.~506.
}

\chapter{Requisitos para ejecutar la sentencia}\label{cap:ejec-requisitos}

\section{La sentencia consentida}

La principal condición que debe cumplir la sentencia para ser ejecutada, conforme al artículo 499, es estar \textbf{consentida o ejecutoriada}.

\begin{definition}{Sentencia consentida}
Una sentencia está consentida cuando ha sido debidamente notificada a las partes y no ha sido impugnada. El consentimiento puede ser \textbf{expreso} (por ejemplo, mediante un escrito en el que se manifiesta que se consiente la sentencia) o \textbf{tácito} (cuando ha vencido el plazo para impugnarla sin que ninguna de las partes lo haya hecho).
\end{definition}

\begin{formula}{Art. 499 CPCCN}
Requisitos para ejecutar una sentencia. La sentencia debe estar consentida o ejecutoriada. Vencido el plazo de cumplimiento fijado en la sentencia ---o al día siguiente de quedar firme si el juez no fijó plazo especial--- y a pedido de parte, el juez procede a la ejecución.
\end{formula}

\section{La sentencia ejecutoriada}

\begin{definition}{Sentencia ejecutoriada}
Una sentencia está ejecutoriada cuando ha sido impugnada mediante los mecanismos impugnativos que admitía y ya no admite nuevos recursos.
\end{definition}
% TODO: contenido incompleto o ambiguo en las notas originales (la frase original sobre los recursos que ya no admite resulta ilegible; se adoptó la formulación del cuadro Cornell del apunte)

Tanto si ha sido consentida como si ha quedado ejecutoriada, la sentencia se encuentra \textbf{firme}.

\section{Plazo vencido, incumplimiento e instancia de parte}

Una vez firme la sentencia, y vencido el plazo para su cumplimiento, el juez procede a la ejecución a pedido de parte. El plazo puede ser fijado expresamente por el juez en la sentencia; si nada se dice al respecto, el cumplimiento es exigible al día siguiente de aquel en que la sentencia quedó firme.

Se requiere \textbf{inescindiblemente} el pedido de la parte vencedora para proceder a ejecutar la sentencia, ya que rige el \textbf{principio dispositivo}.

Además, para poder iniciar la etapa de ejecución la sentencia debe hallarse \textbf{incumplida} por el condenado. Si vencido el plazo la sentencia está prácticamente cumplida y se promueve la ejecución, esta será rechazada porque ha habido cumplimiento. Se entiende que hay cumplimiento aun cuando este se produzca por la aplicación de \textbf{astreintes} (multa progresiva a favor del vencedor que la parte puede reclamar para que el ejecutado cumpla su obligación); en ese caso no se trata de ejecución de sentencia sino de cumplimiento de la sentencia.

\begin{important}
Requisitos para ejecutar una sentencia: (1) sentencia consentida o ejecutoriada; (2) vencimiento del plazo de cumplimiento; (3) incumplimiento del condenado; (4) instancia de parte.
\end{important}

\section{Excepción: la ejecución provisoria}

Como excepción, el Código permite en algunas situaciones ejecutar la sentencia aunque no esté firme. Se trata de una \textbf{ejecución provisoria}, sujeta a lo que en definitiva suceda, generalmente por vías recursivas.

El principio es que todos los recursos de apelación se conceden con \textbf{efecto suspensivo} (apelada la sentencia, se suspende su cumplimiento), salvo que la ley establezca lo contrario. Existe la posibilidad de que la apelación se conceda con \textbf{efecto devolutivo}, y en esos casos la sentencia puede ejecutarse aunque esté apelada.

\subsection{Alimentos}

\begin{formula}{Art. 647 CPCCN}
Efecto del recurso en materia de alimentos. El recurso de apelación contra la sentencia que condena al pago de alimentos se concede con efecto devolutivo. Esto permite al alimentado ejecutar la sentencia de primera instancia aunque esté apelada, formándose incidente de apelación con el testimonio de la sentencia de primera instancia.
\end{formula}

El alimentado beneficiado con una sentencia que condena al alimentante a pagar alimentos puede ejecutarla aunque el alimentante la apele, dado que el recurso se concede con efecto devolutivo. Se protege así un derecho esencial: la necesidad del alimentado de contar con los alimentos para su subsistencia.

\subsection{Proceso sumarísimo}

\begin{formula}{Art. 498 CPCCN}
Proceso sumarísimo. El recurso de apelación contra la sentencia definitiva se concede con efecto devolutivo, salvo que su cumplimiento pudiere ocasionar al condenado un perjuicio irreparable, en cuyo caso el juez podrá concederlo con efecto suspensivo.
\end{formula}

Por lo tanto, la sentencia del juicio sumarísimo se cumple aunque esté apelada, salvo perjuicio irreparable para el condenado.

\subsection{Recurso extraordinario}

Una situación similar se presenta cuando se interpone recurso extraordinario contra una sentencia de cámara que confirma la sentencia de primera instancia. Si las dos sentencias son coincidentes, aunque la de segunda instancia haya sido objeto de recurso extraordinario, se permite la ejecución de la sentencia dando \textbf{fianza} de responder por si la Corte revocara luego la sentencia a través del recurso extraordinario.

\section{Ejecución parcial de la sentencia}

La segunda parte del artículo 499 permite ejecutar parcialmente una sentencia cuando una de sus partes ha sido consentida, aunque se haya interpuesto recurso respecto de otros rubros. Se puede ejecutar parcialmente la sentencia respecto de la parte que quedó firme.

No se trata de una excepción a la ejecutoriedad de la sentencia no consentida: lo que se permite es que, habiéndose cuestionado una parte del fallo, la otra parte, independiente y que puede cuantificarse separadamente, haya quedado firme y pueda llevarse adelante su ejecución.

\begin{example}{Accidente de tránsito}
El juez hizo lugar a los rubros daño emergente, lucro cesante y daño moral, y la parte demandada apela solamente los rubros relativos al daño moral y al lucro cesante. Firme la suma fijada por daño emergente, se puede proceder a ejecutar parcialmente la sentencia por los rubros que quedaron firmes.
\end{example}

Para ello se libra un \textbf{testimonio} en el que consta que determinado rubro, por determinado importe, ha quedado firme, de modo que pueda ser ejecutado en primera instancia, mientras el expediente sube a la cámara para conocer del recurso por los demás rubros cuestionados.

\cornelltitle{requisitos para ejecutar la sentencia}
\cornellblocksum{
\cqrow{¿Cuándo se considera consentida una sentencia?}{Cuando fue debidamente notificada y no fue impugnada en el plazo, o cuando fue consentida expresamente por escrito. El consentimiento puede ser expreso o tácito. % TODO: contenido incompleto o ambiguo en las notas originales (el cuadro Cornell del apunte agrega un tercer supuesto: apelación cuyo recurso se abandonó, por ejemplo por no fundarlo en tiempo; no figura en la transcripción original)
Se agrega, como tercer supuesto, que la sentencia fue apelada pero el recurso se abandonó (por ejemplo, por no fundarlo en tiempo).
}
\cqrow{¿Qué significa que la sentencia esté ejecutoriada?}{Que fue impugnada por todos los recursos que admitía y ya no admite nuevos embates: se agotaron los mecanismos impugnativos disponibles.
}
\cqrow{¿Cuáles son los requisitos para ejecutar una sentencia?}{1) Sentencia consentida o ejecutoriada. 2) Vencido el plazo de cumplimiento (fijado por el juez o el día siguiente a quedar firme). 3) Incumplimiento del condenado. 4) Instancia de parte. Rige el principio dispositivo: sin pedido de parte no hay ejecución.
}
\cqrow{¿Cuándo se puede ejecutar una sentencia que no está firme?}{Excepcionalmente, mediante ejecución provisoria, cuando la apelación se concede con efecto devolutivo: a) alimentos (art.~647 CPCCN); b) proceso sumarísimo (art.~498, salvo perjuicio irreparable); c) recurso extraordinario contra sentencia de cámara confirmatoria de la de primera instancia, dando fianza de responder si la Corte revoca la sentencia.
}
\cqrow{¿Cuál es la diferencia entre ejecución total y ejecución parcial?}{La ejecución parcial procede cuando solo una parte del fallo quedó firme y el resto fue apelado. Se libra testimonio acreditando que ese rubro quedó consentido, se ejecuta en primera instancia y el expediente sube a cámara por los rubros cuestionados. Por ejemplo: en un accidente de tránsito, el daño emergente queda firme mientras se apelan el lucro cesante y el daño moral.
}
}{Para ejecutar una sentencia debe estar consentida o ejecutoriada, haber vencido el plazo de cumplimiento, haber incumplido el condenado y existir pedido de parte. Excepcionalmente puede ejecutarse una sentencia no firme (alimentos, proceso sumarísimo, recurso extraordinario con doble sentencia coincidente y fianza), y puede ejecutarse parcialmente el rubro que quedó firme.
}

\chapter[Competencia y procedimiento de la ejecución]{Competencia y procedimiento general de la ejecución}\label{cap:ejec-procedimiento}

\section{Competencia: el juez que dictó la sentencia}

\begin{formula}{Art. 501 CPCCN}
Competencia para la ejecución. Inc.~1: el juez que dictó la sentencia es competente para entender en su ejecución. Inc.~2: si el objeto de la ejecución lo impone, puede encomendarse a un juez de otra competencia territorial la realización de determinadas diligencias, lo que no importa una excepción a las reglas de competencia. Inc.~3: es competente el juez del proceso principal cuando existan conexiones directas entre causas sucesivas.
\end{formula}

El artículo 501 establece que resulta competente el juez que dictó la sentencia. Ello se vincula con la idea de que el proceso no debe entenderse aislado de la ejecución de sentencia: el proceso principal tiene una etapa más, la de ejecución, y por lógica consecuencia es competente para entender en ella quien entendió en todo el proceso hasta el dictado de la sentencia.

\subsection{Delegación a un juez de otra competencia territorial}

El inciso 2 permite que, si así lo impone el objeto de la ejecución, se encomiende a un juez de otra competencia territorial una determinada diligencia para ejecutar la sentencia. Esto no es una excepción a las reglas de competencia, sino una encomienda que un juez de una jurisdicción hace a otro: una delegación de la competencia para una diligencia determinada.

\begin{example}{Inmueble en extraña jurisdicción}
Debe otorgarse la posesión de un inmueble ubicado en extraña jurisdicción, o llevarse adelante un desalojo en extraña jurisdicción. El juez que dictó la sentencia de desalojo y tiene asiento en Buenos Aires no puede trasladarse a ejecutarla a la ciudad de Corral; por delegación interviene el juez con competencia territorial en relación con el lugar del inmueble respecto del cual debe darse la tenencia a la parte actora.
\end{example}
% TODO: contenido incompleto o ambiguo en las notas originales (el nombre de la ciudad «Corral» figura así en el apunte)

\subsection{Conexión directa entre causas sucesivas}

El inciso 3 establece que es competente el juez del principal si median conexiones directas entre causas sucesivas.

\begin{example}{Desalojo y daños y perjuicios}
Se inició como proceso principal un desalojo y, como consecuencia de él, se inició luego un proceso de daños y perjuicios derivados de esa ocupación. El juez del desalojo puede entender en ambos.
\end{example}

\section{Condena a pagar cantidad líquida}

Debe diferenciarse si la sentencia condena al pago de una cantidad líquida o de una cantidad ilíquida.

\begin{definition}{Condena a cantidad líquida}
La sentencia condena al pago de una cantidad líquida cuando esta está expresada numéricamente o es fácilmente determinable mediante una simple operación aritmética.
\end{definition}

\begin{example}{Capital más interés}
Si la sentencia condena a pagar un capital de 500.000 pesos con más el interés del 40\,\% anual, la sentencia se entiende líquida aunque no esté expresado su monto final, dado que puede calcularse con una simple operación aritmética que incluso puede realizar el juez para determinar el monto.
\end{example}

\section{La traba del embargo}

El embargo es absolutamente esencial: la ejecución de sentencia pretende, ante el incumplimiento del condenado, individualizar bienes del deudor para proceder a su venta y realización efectiva, eventualmente a través de la subasta. Los tipos de embargo (preventivo, ejecutivo y ejecutorio) se distinguen en el capítulo~\ref{cap:embargo}.

El juez tiene todo el poder del Estado, \textbf{imperium}, para hacer cumplir coercitivamente sus decisiones. El obligado que no cumple voluntariamente es embargado; una vez afectado un bien de su patrimonio al cumplimiento de la sentencia mediante un embargo ejecutorio, y previos los trámites que correspondan para ejercer algún tipo de defensa, ese bien es subastado para pagar al acreedor vencedor.

\subsection{Bienes inmuebles}

Para pedir la traba de un embargo sobre un inmueble debe acompañarse un \textbf{informe de dominio} donde conste la titularidad registral en cabeza del deudor. Acreditada la titularidad dominial, el juzgado ordena trabar embargo por la cantidad líquida que consta en la sentencia, más la suma que presupueste el juzgado para responder a intereses y costas.

Si el inmueble está en la misma jurisdicción, el embargo se diligencia por \textbf{oficio} firmado por el secretario al Registro de la Propiedad Inmueble. Si está en extraña jurisdicción (por ejemplo, un inmueble en la ciudad de La Plata con embargo ordenado por un juez de la ciudad de Buenos Aires), se traba mediante el llamado \textbf{testimonio ley 22.172}.

\begin{formula}{Ley 22.172}
Comunicaciones entre jueces de distintas jurisdicciones. Establece las formas en las que se comunican los jueces de distintas jurisdicciones para llevar a cabo determinadas diligencias, en este caso la inscripción del embargo ante el Registro de la Propiedad Inmueble de otra provincia. El testimonio debe ser firmado por el juez y por el secretario, tener sello de agua y otras características establecidas en la ley.
\end{formula}

\subsection{Bienes muebles registrables}

Si se trata de un mueble registrable (por ejemplo, un automóvil, una embarcación o una aeronave), se embarga por \textbf{oficio} firmado por el secretario. Se acompaña un informe de dominio expedido por los registros, como el del automotor, que acredite la titularidad en cabeza del demandado; el juez ordena el embargo y el oficio para que se inscriba ante el registro correspondiente.

Si es en extraña jurisdicción, también se tramita por testimonio de la ley 22.172, firmado por el juez y el secretario, con constancia en el texto del testimonio de los recaudos que exige la ley.

\subsection{Bienes muebles no registrables}

Son muebles no registrables, por ejemplo, las máquinas y determinadas herramientas que no sean de uso indispensable según la profesión del deudor (las que sí lo sean serían inembargables). Se embargan mediante \textbf{mandamiento de embargo} firmado por el secretario. Debe combinarse la fecha con el oficial de justicia, encargado de llevarlo a cabo, quien se constituye en el domicilio del deudor o donde estén los bienes y traba el embargo mediante un acta que labra en el momento.

\section{Efectos del embargo}

El embargo concede \textbf{prioridad como primer embargante} para valerse del producido del bien embargado cuando se efectúe la subasta y se convierta en dinero. Importa una \textbf{indisponibilidad}, una afectación del bien al proceso.

\begin{important}
Que el bien esté embargado no significa que no pueda ser vendido. Si el deudor vende el bien indicando que está embargado, un tercero puede adquirirlo, pero deberá soportar las consecuencias del embargo, porque este persigue a la cosa aunque la transfiera el deudor. La indisponibilidad consiste en que no puede disponerse del bien como si estuviera libre de gravámenes.
\end{important}

Si, por el tipo de bien que deba subastarse, se quiere desapoderar directamente al deudor, debe pedirse otra medida además del embargo: el \textbf{secuestro} del bien, para trasladarlo y quitarlo de la esfera de dominio de su dueño.

\section{Citación de venta}

Una vez embargado un bien y afectado al proceso, corresponde continuar la ejecución citando de venta al deudor. Aunque en la etapa de ejecución el conocimiento del juez es mínimo, no se excluye cierto grado de conocimiento a través de las excepciones que el deudor puede oponer al progreso de la ejecución; de ahí la necesidad de la citación de venta.

La secuencia que establece el Código es: \textbf{primero embargar} (sin bienes embargados no hay sobre qué realizar la sentencia) y \textbf{después citar de venta} por cédula al deudor, para que comparezca a oponer excepciones dentro de \textbf{cinco días} desde que fue notificado. En la práctica a veces ambos actos se hacen conjuntamente: por ejemplo, mediante un mandamiento que traba embargo sobre bienes muebles y deja constancia de que se cita de venta al deudor para que oponga las excepciones del artículo 506.

Si el deudor no opone excepciones, el juez dicta directamente la sentencia de venta o resolución de venta mandando llevar adelante la ejecución, y se pasa a las reglas de cumplimiento de la sentencia de remate del juicio ejecutivo (artículos 559 y siguientes), que se tratan en el capítulo~\ref{cap:remate}.

\section{Las excepciones del artículo 506}

\subsection{Falsedad de la ejecutoria}

La falsedad de la ejecutoria supone que se está ejecutando una sentencia adulterada: porque la firma no pertenece al juez, o porque se enmendó (sobreraspado, adulterado) el monto que originalmente constaba en la sentencia.

No es común, pero podría darse cuando se pide un testimonio, que sale de la esfera del expediente y puede ser objeto de adulteración. No puede imponerse al ejecutado el cumplimiento de una obligación que no se corresponde con la que concretamente se dictaminó en la sentencia: debe haber una correspondencia perfecta entre lo que se ejecuta y lo decidido por sentencia firme en la causa.

\subsection{Prescripción de la ejecutoria}

Se aplica el plazo genérico del artículo 2560 del Código Civil y Comercial de la Nación (CCCN), de \textbf{cinco años}, por no existir uno específico. Antes era de diez años por el artículo 4023 del Código de Vélez; ahora se redujo a cinco.

\begin{formula}{Art. 2560 del CCCN}
Plazo genérico de prescripción: cinco años, excepto que esté previsto uno diferente en la legislación local. Es el plazo aplicable para la prescripción de la ejecutoria cuando no existe plazo específico.
\end{formula}

El vencedor no puede dejar pasar el plazo que quiera para ejecutar la sentencia: a partir del día en que resulta exigible el cumplimiento, tiene cinco años para ejecutarla, ya que si no realiza actos que impulsen el proceso de ejecución opera la prescripción de la ejecutoria.

\begin{important}
Aunque haya transcurrido el plazo de prescripción, el vencedor puede deducir la pretensión de ejecución. Si la otra parte no opone la excepción de prescripción de ejecutoria en el plazo que corresponde (cinco días), aunque estuviera vencido el plazo de prescripción liberatoria, la sentencia podrá ejecutarse eficazmente. Rige el principio dispositivo: la prescripción debe ser opuesta por la parte que se beneficiaría con ella y no puede ser decidida de oficio por el juez.
\end{important}

\subsection{Pago}

Si el deudor canceló la obligación que surge de la sentencia, puede defenderse acreditando el pago mediante documentos que emanen del vencedor.

\begin{important}
Todas las excepciones admisibles en la ejecución de sentencia deben fundarse en \textbf{hechos posteriores a la sentencia}, dado que los anteriores están cubiertos por la eficacia de la cosa juzgada y la preclusión impide volver sobre ellos. Además, hay una limitación en los medios de prueba: solo se permiten \textbf{constancias del expediente} o \textbf{documentos emanados del vencedor}, con exclusión de cualquier otro medio probatorio.
\end{important}

\subsection{Quita, espera y remisión}

El inciso 4 del artículo 506 se refiere a las excepciones de quita, espera o remisión.

\begin{itemize}
    \item \textbf{Quita:} después de dictada la sentencia, el vencedor realiza una quita del monto que se le debe. El ejecutado se excepciona acompañando el documento que da cuenta de la quita, y se reduce la ejecución y, eventualmente, el embargo.
    \item \textbf{Espera:} el vencedor otorga un plazo de cumplimiento al ejecutado y, no obstante, ejecuta la sentencia antes de su vencimiento. El ejecutado opone la excepción de espera; si se hace lugar, el ejecutante debe aguardar a que pase el plazo para poder ejecutar la sentencia.
    \item \textbf{Remisión:} consiste, por decirlo así, en el perdón de la deuda. Técnicamente, remisión es devolver al deudor el título donde consta la obligación, pero esto no es posible aquí porque existe un expediente. Si se perdona la deuda, se acompaña el instrumento correspondiente y el juez hace lugar a la excepción, ordena levantar el embargo.
\end{itemize}
% TODO: contenido incompleto o ambiguo en las notas originales (la explicación original de por qué no es posible devolver el título en la remisión resulta ambigua)

\begin{example}{Quita}
La sentencia es de 600.000 pesos por daños y perjuicios y existe un documento emanado del vencedor, posterior a la sentencia, que le hace una quita de 200.000 pesos. El ejecutado puede excepcionarse acompañando ese documento, y la ejecución y, eventualmente, el embargo se reducen a 400.000 pesos.
\end{example}

\section{Excepciones no enumeradas admitidas por la jurisprudencia}

El artículo 506 dice que solo se consideran legítimas las excepciones que enumera, pero la jurisprudencia entiende que no pueden dejarse de lado algunas otras, como la \textbf{inhabilidad de título}. Si el ejecutante pretende ejecutar a quien no es el demandado condenado, el pretendido ejecutado puede oponer la excepción de falta de legitimación pasiva, que se incluye dentro de la inhabilidad de título.

Lo mismo se ha permitido respecto de otras excepciones vinculadas con presupuestos procesales, como la \textbf{incompetencia}: si se va a ejecutar la sentencia ante un juez incompetente, el deudor puede excepcionarse aunque la excepción no esté señalada expresamente entre las del artículo 506.

\section{Trámite de las excepciones y sentencia de venta}

Opuesta alguna excepción, se da traslado al ejecutante por cinco días para que la conteste. Una vez contestada, el juez debe resolver. % TODO: contenido incompleto o ambiguo en las notas originales (el apunte atribuye «resolver» al ejecutado; por el contexto corresponde al juez)

\begin{itemize}
    \item \textbf{Hace lugar a las excepciones:} ordena levantar el embargo. Si es la excepción de pago o de prescripción de ejecutoria, levanta el embargo; si es la de quita, adecua el embargo en la forma que corresponda; si es la de espera, hace lugar a la excepción pero deja trabado el embargo para que el ejecutante pueda continuar cuando se verifique el plazo correspondiente.
    \item \textbf{Rechaza las excepciones:} dicta la sentencia de venta mandando llevar adelante la ejecución.
\end{itemize}

La sentencia de venta puede ser apelada por el ejecutado. El recurso se concede con \textbf{efecto suspensivo y trámite inmediato}, pero si el ejecutante otorga \textbf{fianza} para responder por daños y perjuicios, puede transformar el efecto de la apelación en devolutivo y ejecutar la sentencia no obstante estar apelada por el ejecutado.

Si el juez ordena levantar el embargo, puede apelar el ejecutante. Esa apelación va en relación, con efecto suspensivo y trámite inmediato. Una vez firme la resolución, se continúa con el trámite de cumplimiento de la sentencia de remate del juicio ejecutivo, que es una etapa común al juicio ejecutivo y a la ejecución de sentencia.

\cornelltitle{competencia y procedimiento general de la ejecución}
\cornellblock{
\cqrow{¿Qué juez es competente para ejecutar la sentencia?}{El juez que dictó la sentencia (art.~501 CPCCN), porque la ejecución es una etapa más del proceso principal. Puede encomendar diligencias determinadas a un juez de otra competencia territorial, lo que no es excepción a las reglas de competencia sino una delegación. Si hay conexión directa entre causas sucesivas, es competente el juez del proceso principal.
}
\cqrow{¿Cuándo una condena es a cantidad líquida?}{Cuando la cantidad está expresada numéricamente o es fácilmente determinable mediante una simple operación aritmética, aunque no esté expresado el monto final.
}
\cqrow{¿Por qué el embargo es esencial en la ejecución de sentencia?}{Porque el procedimiento tiene por finalidad afectar bienes del deudor para su posterior subasta. Sin bienes afectados al proceso no puede continuarse la ejecución. En el juicio ejecutivo el embargo es optativo; en la ejecución de sentencia es absolutamente esencial.
}
\cqrow{¿Cómo se traba el embargo según el tipo de bien?}{Inmueble (misma jurisdicción): oficio del secretario al Registro de la Propiedad, con informe de dominio. Inmueble (otra jurisdicción): testimonio ley 22.172, firmado por juez y secretario, con sello de agua. Mueble registrable (automóvil, embarcación, aeronave): oficio del secretario con informe de dominio del registro correspondiente. Mueble no registrable: mandamiento de embargo firmado por el secretario, diligenciado por oficial de justicia mediante acta.
}
\cqrow{¿Qué efectos tiene el embargo?}{Prioridad como primer embargante sobre el producido de la subasta e indisponibilidad (afectación del bien al proceso). El bien puede ser vendido, pero el adquirente soporta el embargo, que persigue a la cosa. Para desapoderar al deudor se requiere además el secuestro.
}
\cqrow{¿Qué es la citación de venta y cuándo se produce?}{Es el acto que habilita al ejecutado a oponer excepciones. Se practica después del embargo (aunque en la práctica puede hacerse en el mismo mandamiento) y se notifica por cédula. El ejecutado tiene cinco días para oponer las excepciones del art.~506 CPCCN.
}
}
\medskip
\cornellblocksum{
\cqrow{¿Cuáles son las excepciones del art.~506 CPCCN?}{1) Falsedad de la ejecutoria (adulteración de la sentencia). 2) Prescripción de la ejecutoria (cinco años, art.~2560 CCCN; debe ser opuesta por la parte, no de oficio). 3) Pago (documentado con constancias emanadas del vencedor). 4) Quita, espera o remisión (documentos posteriores a la sentencia emanados del ejecutante). Condición común: hechos posteriores a la sentencia; prueba solo por constancias del expediente o documentos del ejecutante.
}
\cqrow{¿Qué excepciones admite la jurisprudencia además de las del art.~506?}{Inhabilidad de título (falta de legitimación pasiva: ejecutar a quien no fue condenado) e incompetencia (ejecutar ante juez incompetente).
}
\cqrow{¿Qué pasa si el ejecutado no opone excepciones?}{El juez dicta directamente la sentencia de venta mandando llevar adelante la ejecución, y se aplican las reglas de cumplimiento de la sentencia de remate del juicio ejecutivo (arts.~559 y ss. CPCCN; capítulo~\ref{cap:remate}).
}
\cqrow{¿Cómo se apela la sentencia de venta?}{Con efecto suspensivo y trámite inmediato; si el ejecutante otorga fianza para responder por daños y perjuicios, puede transformarse en devolutivo y ejecutarse la sentencia pese a la apelación.
}
}{Es competente el juez que dictó la sentencia, con posibilidad de delegar diligencias en jueces de otra jurisdicción. Ante una condena líquida se traba embargo (según el tipo de bien) y luego se cita de venta al deudor, que tiene cinco días para oponer las excepciones del art.~506 (falsedad, prescripción, pago, quita, espera o remisión), fundadas en hechos posteriores a la sentencia y probadas solo con constancias del expediente o documentos del ejecutante; la jurisprudencia admite además inhabilidad de título e incompetencia. Sin excepciones, o rechazadas estas, se dicta la sentencia de venta, apelable con efecto suspensivo salvo fianza.
}

\chapter[Condena ilíquida y liquidación]{Condena a cantidad ilíquida y práctica de la liquidación}\label{cap:ejec-iliquida}

\section{Cuándo corresponde practicar liquidación}

Cuando la sentencia no establece un monto determinado que debe abonar el condenado, se plantea cómo proceder. Ocurre, por ejemplo, con los intereses (la sentencia fija la tasa pero no el monto), con los frutos que deben restituirse o con los daños y perjuicios cuyo derecho se reconoce al actor pero que, por distintas cuestiones procesales, el juez no puede fijar en el acto del fallo.

\begin{formula}{Art. 165 CPCCN}
Condena al pago de frutos, intereses o daños. El juez, al pronunciar sentencia, debe fijar el importe de los frutos, intereses, daños y perjuicios con arreglo a lo dispuesto por el artículo 165. Cuando ello no sea posible, el juez debe fijar las bases para que se practique la liquidación conforme a las pautas que establezca.
\end{formula}
% TODO: contenido incompleto o ambiguo en las notas originales (el texto del recuadro del apunte para el art. 165 se autorreferencia)

El artículo 165 exige en primer término que el juez determine la cantidad en forma líquida por estos rubros, pero reconoce que pueden existir impedimentos; en ese caso exige que al menos fije las pautas para que se establezca posteriormente en una liquidación.

\begin{important}
La liquidación debe estar en consonancia y ser correlativa con lo establecido en la sentencia. No pueden incluirse rubros ni irse más allá o por fuera de lo solicitado en las pretensiones articuladas al iniciar el proceso, pues sería nugatoria la ejecución que se llevara a cabo con base en una liquidación que incluyera conceptos indebidos.
\end{important}

\section{Quién practica la liquidación y en qué plazo}

\begin{formula}{Art. 503 CPCCN (liquidación)}
Liquidación. El acreedor podrá practicar la liquidación dentro del plazo de diez días contados desde que la sentencia sea exigible. Dará traslado a la otra parte por cinco días. Si hay conformidad, o si vencido el plazo el deudor no la impugna, el juez aprobará la liquidación presentada, disponiendo que se la tenga como base para el embargo y la subasta.
\end{formula}
% TODO: contenido incompleto o ambiguo en las notas originales (el apunte cita el art. 503 tanto para la liquidación como para la condena a hacer, sección 5.1)

Es una \textbf{facultad del vencedor} practicar la liquidación, para lo cual dispone de diez días contados desde que sea exigible el cumplimiento, es decir, desde que la sentencia quedó firme. Si el juez fijó un plazo especial, el plazo del vencedor para practicar la liquidación comienza al día siguiente del vencimiento del plazo para cumplir.

De la liquidación se da traslado a la otra parte, por cédula y por cinco días. Si la contraparte consiente expresamente o guarda silencio (consentimiento tácito, por no impugnar), el juez aprueba la liquidación y se continúa como en el proceso con cantidad líquida: se traba embargo sobre el monto de la liquidación y lo que eventualmente se presupueste para responder a intereses y costas.

\section{Impugnación de la liquidación}

\begin{formula}{Art. 504 CPCCN}
Impugnación de la liquidación. Si hubiere impugnación a la liquidación, se formará incidente siguiendo las normas de los artículos 178 y siguientes. Las partes deberán ofrecer la prueba de la que intenten valerse en sus respectivos escritos de impugnación y contestación.
\end{formula}

El vencido puede impugnar la liquidación porque entiende que no corresponden los rubros o porque, aun correspondiendo, están mal o defectuosamente calculados. Se forma entonces el incidente que exige el artículo 504, siguiendo las pautas del artículo 178 y siguientes. Las partes deben ofrecer la prueba de la que intenten valerse para fundar la impugnación o la contestación de la impugnación que haga el ejecutante. Vencidos los plazos y contestada la impugnación, el juez resuelve y, en su caso, aprueba la liquidación que corresponda.

\section{Debate doctrinal sobre la apelación de la liquidación aprobada}

Existen dos posiciones doctrinarias sobre el efecto con que se concede la apelación contra la resolución que aprueba la liquidación:

\begin{table}[H]
\centering
\small
\begin{tabularx}{\textwidth}{
    >{\raggedright\arraybackslash}p{0.17\textwidth}
    >{\raggedright\arraybackslash}X
    >{\raggedright\arraybackslash}X
}
\toprule
 & \textbf{Alsina} & \textbf{Palacio} \\
\midrule
Fundamento &
Por tratarse de un incidente, el recurso se concede en relación. &
Se aplica el art.~509, que dispone que todas las apelaciones en ejecución de sentencia se conceden con trámite diferido. \\[0.3em]

Efecto &
Suspensivo y trámite inmediato. &
Trámite diferido. \\[0.3em]

Consecuencias &
Se suspende el cumplimiento de la resolución y el expediente va a la cámara para su resolución; una vez confirmada, vuelve a primera instancia y permite trabar embargo y continuar la ejecución. &
La liquidación aprobada, aunque apelada por el ejecutado, sirve de base para trabar el embargo y continuar la ejecución; la apelación se funda recién junto con una eventual apelación contra la sentencia de venta. \\
\bottomrule
\end{tabularx}
\caption{Posiciones doctrinarias sobre el recurso contra la liquidación aprobada.}
\end{table}

Esta discusión varía fundamentalmente la solución del caso.

\cornelltitle{condena ilíquida y liquidación}
\cornellblocksum{
\cqrow{¿Qué sucede cuando la sentencia condena a una cantidad ilíquida?}{Corresponde practicar liquidación (art.~503 CPCCN). El vencedor tiene diez días desde que la sentencia es exigible. Se da traslado a la otra parte por cinco días. Si hay conformidad o silencio, el juez aprueba la liquidación y se sigue con el trámite de embargo. Si hay impugnación, se forma incidente (art.~504, reglas del art.~178 y ss.).
}
\cqrow{¿Qué debe hacer el juez según el art.~165 cuando no puede fijar el monto?}{Debe fijar las bases o pautas para que se practique la liquidación posteriormente. La liquidación debe ser correlativa con la sentencia y con las pretensiones articuladas al iniciar el proceso.
}
\cqrow{¿Cuál es el debate sobre la apelación de la liquidación aprobada?}{Alsina: al ser un incidente, se concede en relación, con efecto suspensivo y trámite inmediato. Palacio: rige el art.~509 y se concede con trámite diferido, de modo que la liquidación sirve de base para el embargo y la apelación se funda junto con la apelación de la sentencia de venta.
}
}{Si la condena es ilíquida, el vencedor practica la liquidación en diez días desde que la sentencia es exigible; se corre traslado por cinco días y, si hay conformidad o silencio, se aprueba y se sigue con el embargo; si hay impugnación se tramita un incidente (arts.~503 y 504). Sobre el efecto de la apelación de la liquidación aprobada, Alsina y Palacio sostienen soluciones opuestas.
}

\chapter{Obligaciones de hacer, no hacer y dar cosas}\label{cap:ejec-hacer}

\section{Condena a hacer alguna cosa}

\begin{formula}{Art. 503 CPCCN (condena a hacer)}
Si la condena es a hacer alguna cosa y el obligado no la realiza dentro del plazo correspondiente, el ejecutante puede optar por: a) que la cosa sea realizada por un tercero a costa del ejecutado; o b) que se le paguen los daños y perjuicios provenientes de la inejecución.
\end{formula}
% TODO: contenido incompleto o ambiguo en las notas originales (el apunte cita el art. 503 también para la liquidación, sección 4.2)

Si el deudor cumple haciendo aquello a lo que fue obligado dentro del plazo correspondiente, no hay ejecución. Si no lo hace, el ejecutante puede promover la ejecución de sentencia, optando conforme al artículo 503. Debe advertirse que el deudor no puede ser obligado por la fuerza física a hacer algo contra su voluntad. El vencedor, ejecutante, puede elegir entre solicitar que la prestación sea realizada por un \textbf{tercero} o que se obligue al deudor a resarcir los \textbf{daños y perjuicios} provenientes de la inejecución.

Si opta por que lo realice un tercero, hay que estimar el costo de esa actividad (no una liquidación sino una \textbf{estimación}) y dar traslado al vencido para que se expida sobre si ese valor se corresponde con el objeto de la obligación que ahora realizará un tercero a su costa. La opción es del acreedor, pero la actividad es a costa del ejecutado.

\section{Condena a otorgar escritura pública}

\begin{formula}{Art. 502 CPCCN}
Condena a otorgar escritura pública. Si la sentencia condena al otorgamiento de escritura pública y este no es cumplido por el ejecutado, el juez la suscribirá a su costa, sin perjuicio de que si llegara a devenir de imposible cumplimiento por alguna circunstancia, esa escrituración se transforme en la obligación de resarcir los daños y perjuicios provenientes de esa imposibilidad.
\end{formula}

Si la sentencia condena al otorgamiento de escritura pública y el ejecutado no la cumple, el procedimiento continúa con quien fue designado para suscribirla a su costa. Si el cumplimiento deviniera imposible, la condena escrituraria se transforma en la obligación de resarcir los daños y perjuicios derivados de esa imposibilidad, que se fijan por el mismo procedimiento que establece el artículo 513.

\section{Condena a no hacer}

Cuando se trata de una condena a no hacer y el ejecutado la incumple, es decir, hace algo contrario a lo que estaba obligado, el acreedor puede ejercer la opción entre solicitar que se \textbf{reponga el estado de cosas anterior} o que se traduzca en la \textbf{indemnización de daños y perjuicios}, conforme al artículo 513.

\begin{formula}{Art. 513 CPCCN}
Daños y perjuicios por inejecución. Los daños y perjuicios derivados del incumplimiento se determinarán conforme a las pautas de los artículos 503 y 504 (liquidación con vía incidental) o bien a través de juicio ordinario o sumarísimo, según lo determine el juez de manera irrecurrible.
\end{formula}

Los daños y perjuicios se determinan, entonces, mediante el procedimiento de liquidación, con la derivación que hace el Código a las normas de los incidentes, o por juicio ordinario o sumarísimo, según lo determine el juez, y en forma irrecurrible.

\section{Condena a dar cosas: mandamiento de desapoderamiento}

Si la sentencia no se cumple voluntariamente, pueden aplicarse \textbf{astreintes}, es decir, multas progresivas conminatorias a favor de la otra parte para obligar a cumplir al ejecutado.

Si la obligación de dar alguna cosa no es cumplida voluntariamente, el vencedor puede ejecutar la sentencia mediante un \textbf{mandamiento} que ordena librar el juez para \textbf{desapoderar}, aun contra la voluntad del deudor, al vencido de aquello que fue ordenado por sentencia que se restituya al vencedor. En el mismo momento en que se diligencia el mandamiento para desapoderar de la cosa al ejecutado, se lo cita para que oponga las excepciones del artículo 506.

Si no puede cumplirse la entrega de la cosa por alguna circunstancia (porque ha desaparecido o es imposible su cumplimiento), el vencedor puede solicitar que se compense con una suma equivalente al valor de la cosa. La determinación de cuánto vale la cosa y cuánto debe compensarse al ejecutante se realiza también por las pautas de los artículos 503 y 504, por la vía incidental o por juicio ordinario o sumarísimo, según establezca el juez y de manera irrecurrible.

\section{Liquidaciones especiales}

\begin{formula}{Art. 516 CPCCN}
Liquidaciones especiales. Cuando las liquidaciones impliquen cuentas complicadas de lenta o difícil justificación o que requieran conocimientos especiales, podrán someterse a la decisión de peritos árbitros o amigables componedores. El laudo que emitan será vinculante y obligatorio para las partes. Lo mismo se aplica para liquidaciones de sociedades conyugales o determinación del carácter propio o ganancial de bienes.
\end{formula}

Cuando las liquidaciones implican cuentas complicadas de lenta o difícil justificación, o que requieren conocimientos especiales, pueden dejarse a la decisión de peritos árbitros o, en su caso, de amigables componedores. Tanto el perito árbitro como el amigable componedor emiten un laudo que es vinculante y obligatorio para las partes, y que debe respetar incluso el juez en estas cuestiones técnicas fijadas por los peritos en el marco de la competencia que se les atribuye.

\section{Otros títulos ejecutables por este procedimiento}

\begin{formula}{Art. 500 CPCCN}
Títulos ejecutables por el procedimiento de ejecución de sentencia: 1) las sentencias de los tribunales judiciales y los laudos de tribunales arbitrales (aunque estos carecen de imperium); 2) las transacciones o acuerdos homologados; 3) (inciso incorporado por ley 26.589) los acuerdos celebrados en mediación con la firma certificada del mediador, salvo que incluyan menores o incapaces, en cuyo caso requieren homologación judicial; 4) las multas procesales; 5) los honorarios regulados judicialmente.
\end{formula}
% TODO: contenido incompleto o ambiguo en las notas originales (el apunte menciona en la explicación oral el «inciso cuarto» para los acuerdos de mediación y el «inciso segundo» para las multas procesales, lo que no coincide con la numeración del recuadro)

\begin{itemize}
    \item \textbf{Sentencias judiciales y laudos arbitrales:} los tribunales arbitrales tienen jurisdicción pero carecen de imperium para hacer cumplir coercitivamente sus decisiones.
    \item \textbf{Transacciones o acuerdos homologados:} las transacciones son las que realizan las partes extinguiendo obligaciones litigiosas o dudosas; los acuerdos homologados pueden provenir, por ejemplo, de una conciliación llevada a cabo por el juez en la audiencia preliminar.
    \item \textbf{Acuerdos de mediación con firma certificada del mediador:} incorporados por la reforma de la ley 26.589. Su incumplimiento permite saltar todas las etapas del proceso de conocimiento y acceder directamente a la última, la ejecución de sentencia, para exigir el cumplimiento según las pautas ya vistas. Si el acuerdo incluye menores o incapaces, su representante debe solicitar al juez la homologación, con intervención del Ministerio Pupilar, para proteger esos intereses.
    \item \textbf{Multas procesales:} por ejemplo, las impuestas por temeridad y malicia a favor de una de las partes, que esta puede ejecutar por este proceso.
    \item \textbf{Honorarios regulados judicialmente:} su cobro también se ejecuta por este procedimiento.
\end{itemize}

\begin{note}
Paradójicamente, el acuerdo conciliatorio celebrado ante el mismo juez en la audiencia preliminar requiere homologación judicial, mientras que el acuerdo celebrado ante el mediador, con firma certificada, puede ejecutarse sin homologación (salvo que incluya menores o incapaces).
\end{note}

\section{Reflexión final}

El Código establece un procedimiento de ejecución de sentencia \textbf{abreviado}, limitado en cuanto a la apelación y a las defensas, dada la trascendencia del título ejecutorio, que instrumenta un derecho cierto. Permite transitar una etapa con un mínimo de conocimiento y con una mínima posibilidad de prueba por parte del ejecutado, y hacer efectiva la sentencia de modo de satisfacer el derecho que pretendía el actor al iniciar la demanda, cuya pretensión no se agota en una mera sentencia declarativa de la certeza de su derecho, sino que requiere efectivo cumplimiento.

Se recuerda, además, la afirmación de un autor según la cual, para tener derecho, hay que contar con la norma que lo acuerde, saber exponerlo, un juez que quiera reconocerlo y un deudor que pueda pagarlo.

\cornelltitle{hacer, no hacer, dar cosas y otros títulos}
\cornellblock{
\cqrow{¿Qué ocurre con la condena a hacer algo?}{El ejecutante puede optar (art.~503): a) que un tercero realice la cosa a costa del ejecutado (se estima el costo y se da traslado al vencido); b) el pago de daños y perjuicios por inejecución (se cuantifican por liquidación/incidente o juicio ordinario/sumarísimo, según determine el juez, de manera irrecurrible). El deudor no puede ser obligado por la fuerza física a hacer.
}
\cqrow{¿Qué ocurre con la condena a otorgar escritura pública?}{Si el ejecutado no la cumple, el juez la suscribe a su costa (art.~502). Si deviniera de imposible cumplimiento, se transforma en la obligación de resarcir los daños y perjuicios, fijados por el procedimiento del art.~513.
}
\cqrow{¿Qué pasa si la condena es a no hacer algo y se incumple?}{El acreedor puede optar (art.~513): a) la reposición al estado anterior; b) la indemnización por daños y perjuicios, cuantificados por liquidación/incidente o juicio ordinario/sumarísimo según determine el juez de manera irrecurrible.
}
\cqrow{¿Cómo se ejecuta la condena a dar cosas?}{Pueden aplicarse astreintes. Si no se cumple voluntariamente, mediante mandamiento de desapoderamiento librado por el juez, al tiempo que se cita al ejecutado para oponer las excepciones del art.~506. Si la entrega es imposible, se compensa con una suma equivalente al valor de la cosa, determinada por los arts.~503 y 504 o por juicio ordinario o sumarísimo.
}
\cqrow{¿Cómo se resuelven las liquidaciones especiales (art.~516)?}{Cuando implican cuentas complicadas o conocimientos especiales, pueden someterse a peritos árbitros o amigables componedores, cuyo laudo es vinculante y obligatorio para las partes e incluso para el juez.
}
\cqrow{¿Qué títulos habilitan la ejecución de sentencia además de la sentencia judicial?}{Art.~500 CPCCN: 1) laudos arbitrales (el árbitro carece de imperium); 2) transacciones y acuerdos homologados; 3) acuerdos de mediación con firma certificada del mediador (ley 26.589), salvo menores o incapaces, que requieren homologación; 4) multas procesales; 5) honorarios regulados judicialmente.
}
}
\medskip
\cornellblocksum{
\cqrow{Paradoja del art.~500: acuerdo judicial vs.\ acuerdo de mediación}{El acuerdo conciliatorio celebrado ante el mismo juez en la audiencia preliminar requiere homologación judicial para ser ejecutable. En cambio, el acuerdo de mediación con firma certificada del mediador no la requiere (salvo menores o incapaces).
}
}{Según el tipo de obligación, la ejecución adopta distintos carriles: en la condena a hacer, el ejecutante opta entre un tercero a costa del deudor o los daños y perjuicios; en la escritura pública, el juez suscribe; en la condena a no hacer, entre la reposición al estado anterior o la indemnización; en la condena a dar cosas, el mandamiento de desapoderamiento, con compensación en dinero si la entrega es imposible. Además de la sentencia, se ejecutan por este procedimiento laudos, transacciones o acuerdos homologados, acuerdos de mediación certificados, multas procesales y honorarios regulados.
}

\section{Síntesis de la ejecución de sentencia}
La ejecución de sentencia es la etapa final del proceso de conocimiento ---no un proceso autónomo---, mediante la cual el Estado pone su imperium al servicio del vencedor para hacer efectivo un derecho declarado cierto por sentencia firme.

A diferencia del juicio ejecutivo, que ejecuta títulos con derechos presuntos y admite defensas más amplias (art.~544 CPCCN), la ejecución de sentencia opera sobre un título ejecutorio que ya superó el debate exhaustivo del proceso de conocimiento: la cosa juzgada material impide discutir hechos anteriores. Por eso las excepciones admisibles (art.~506 CPCCN) se reducen a la falsedad de la ejecutoria, la prescripción, el pago y la quita, espera o remisión, todas fundadas en hechos \textbf{posteriores} a la sentencia y probables solo con documentos del expediente o del ejecutante.

El procedimiento requiere que la sentencia esté consentida o ejecutoriada, que haya vencido el plazo de cumplimiento y que exista incumplimiento del condenado, todo a instancia de parte. El embargo es esencial: traba bienes para su posterior subasta. Según el tipo de obligación condenada (dar dinero, hacer, no hacer, dar cosas), el proceso adopta distintos carriles: liquidación (si la condena es ilíquida), tercero a costa del deudor o indemnización (si la condena es a hacer o no hacer) y mandamiento de desapoderamiento (si es a dar cosas).

Otros títulos equiparados a la sentencia ---laudos arbitrales, transacciones homologadas, acuerdos de mediación certificados, multas procesales y honorarios--- también se ejecutan por este procedimiento, lo que subraya su carácter de herramienta central en la tutela jurisdiccional efectiva.

\chapter{El juicio ejecutivo}\label{cap:juicioejecutivo}

\chaptermark{El juicio ejecutivo}

%----------------------------------------------------------
\section{Ubicación en el código procesal}

El juicio ejecutivo se localiza en la parte del código referida al proceso de ejecución. Queda integrado, dentro de esa parte, a continuación del tema de la ejecución de sentencias.

El código agrupa tanto la ejecución de sentencias ---el cumplimiento de la decisión de un juicio ordinario de conocimiento pleno--- como el juicio ejecutivo, a partir del artículo 520. Ambos mecanismos comparten luego la regulación referida a la realización de los bienes del deudor.

\begin{quote}
\textbf{Art. 520 CPCCN.} A partir de este artículo el código regula el juicio ejecutivo.
\end{quote}

\begin{important}
El juicio ejecutivo no debe confundirse con el proceso de ejecución de sentencias, ya que responden a una lógica totalmente distinta.
\end{important}

%----------------------------------------------------------
\section{Concepto y naturaleza jurídica}

\subsection{El proceso de conocimiento sumario}

La ejecución de sentencia tiene como antecedente un conocimiento pleno, es decir, exhaustivo en relación con la causa de la obligación. En el juicio ejecutivo, en cambio, se está ante un \textbf{juicio de conocimiento sumario}, de conocimiento acotado o, si se quiere, superficial.

% TODO: contenido incompleto o ambiguo en las notas originales (el apunte menciona "el de la provincia de González"; se interpreta como referencia al código procesal provincial)
No debe confundirse esta noción con el proceso sumario previsto, por ejemplo, en la provincia de Buenos Aires, o el que antes existía en el código procesal civil y comercial. Cuando el código nacional o el código provincial hablan de juicio sumario, se refieren a un juicio de conocimiento pleno. En cambio, cuando se afirma que el juicio ejecutivo es un proceso de conocimiento sumario, se indica que el conocimiento del juez sobre el conflicto está acotado, disminuido o seccionado: no va más allá del título cartular que se lleva a ejecutar.

\subsection{La opción por la rapidez}

El juicio ejecutivo es un proceso en el cual se opta por la rapidez en detrimento de la exhaustividad en el conocimiento. Su razón de ser es el tiempo: el legislador quiso que, en ciertos casos en que existen determinadas prestaciones, el acreedor tenga a su alcance ---ante el incumplimiento del deudor--- una herramienta rápida para recobrar el crédito o las sumas de dinero adeudadas, sin tener que someterse, como regla, a un juicio de conocimiento pleno.

En este proceso no hay traslado de demanda, excepciones como las del juicio ordinario, apertura a prueba, alegatos ni sentencia en los términos del ordinario. Tampoco rige la regla de que todas las decisiones son apelables: sucede lo contrario. Es un proceso al cual sólo se puede acceder ante la presencia de un título ejecutivo.

%----------------------------------------------------------
\section{Diferencia con la ejecución de sentencias}

\begin{table}[H]
\centering
\renewcommand{\arraystretch}{1.25}
\begin{tabularx}{\textwidth}{Y Y}
\toprule
\textbf{Ejecución de sentencias} & \textbf{Juicio ejecutivo} \\
\midrule
Requiere un juicio ordinario previo con conocimiento pleno. & Sólo requiere un título ejecutivo. \\
La sentencia es el antecedente del trámite ejecutorio. & El título cartular es el fundamento directo. \\
El deudor pudo discutir exhaustivamente la causa. & La causa de la obligación no puede discutirse en el proceso. \\
La cosa juzgada es material para todas las cuestiones. & La cosa juzgada es formal respecto de las cuestiones causales. \\
El embargo resultante es «ejecutorio». & El embargo resultante es «ejecutivo». \\
\bottomrule
\end{tabularx}
\end{table}

\subsection{El desdoblamiento del conocimiento del conflicto}

El juicio ejecutivo desdobla el conocimiento del conflicto y opta por un mecanismo procesal ágil y rápido mediante el cual ---con menos cantidad de trámites--- el acreedor, a partir de la existencia de un título que traiga aparejada fuerza ejecutiva, puede hacerse más rápidamente del crédito, en detrimento de las defensas que puede oponer el deudor. Esas defensas sólo pueden referirse a cuestiones que emerjan del título cartular que se ejecuta.

Para resguardar el derecho de defensa del deudor, se le concede la posibilidad de discutir todas las cuestiones que quedan fuera del juicio ejecutivo en un \textbf{juicio ordinario posterior}, en el cual la discusión ya no está acotada a las cuestiones que surjan del título.

\subsection{La cosa juzgada en el juicio ejecutivo}

En alguna medida, se dice que la cosa juzgada que emana del juicio ejecutivo es una \textbf{cosa juzgada formal}: lo es porque existen cuestiones causales que no pueden debatirse en él. Sin embargo, respecto de las cuestiones que sí pueden debatirse en el juicio ejecutivo ---por ejemplo, el pago--- no se trata de cosa juzgada formal, porque el pago es una excepción que sí puede interponerse en el marco del ejecutivo. Para el pago, entonces, el juicio ejecutivo arroja una \textbf{cosa juzgada material}: no podrá discutirse en el ordinario posterior una cuestión que sí podía plantearse en el ejecutivo.

%----------------------------------------------------------
\section{Requisitos para acceder a la vía ejecutiva}

A partir del artículo 523, el código enumera los títulos que tienen fuerza ejecutiva por ley. Para acceder a la vía ejecutiva se requieren dos condiciones esenciales:

\begin{enumerate}
    \item Contar con un título que traiga aparejada la ejecución (fuerza ejecutiva).
    \item Que exista una obligación exigible de dar cantidades líquidas de dinero o fácilmente liquidables.
\end{enumerate}

\subsection{La exigibilidad}

La obligación debe tener un plazo vencido; es decir, debe ser exigible, de modo que el acreedor pueda pedir el cobro. Debe tratarse de una obligación de dar sumas de dinero y no de una obligación de hacer o de no hacer. El juicio ejecutivo queda, por el momento, limitado a obligaciones de dar sumas de dinero.

\subsection{La liquidez o fácil liquidabilidad}

Una suma es \textbf{líquida} cuando la cantidad de dinero adeudada está expresada en números. El código también admite sumas \textbf{fácilmente liquidables}: por ejemplo, si se trata de cánones mensuales de alquileres, la suma líquida puede obtenerse mediante una liquidación sencilla en el marco del proceso ejecutivo.

\subsection{El caso de la moneda extranjera}

Una situación bastante común se presenta en las hipotecas pactadas en moneda extranjera. En Argentina, los dólares tuvieron distintos tratamientos jurídicos a lo largo del tiempo: en un momento se los consideró una obligación de dar cosas y, en otro, una obligación equivalente a dar moneda corriente.

% TODO: contenido incompleto o ambiguo en las notas originales (el apunte indica que actualmente se la trata "como una obligación de dar" sin precisar de qué clase)
Actualmente se la trata como una obligación de dar, por lo que, conforme el artículo 520, corresponde convertir la moneda extranjera a su equivalente nacional para obtener la suma líquida por la cual se puede ejecutar.

%----------------------------------------------------------
\section{Títulos ejecutivos --- Artículo 523}

El artículo 523 realiza una larga enumeración de supuestos en los que se puede acceder a la vía ejecutiva. La vía ejecutiva es siempre una opción para el acreedor: podría optar directamente por el juicio ordinario (art.~521).

\begin{quote}
\textbf{Art. 523 CPCCN --- Títulos con fuerza ejecutiva:}
\begin{enumerate}[label=\arabic*.º),leftmargin=2.2em]
    \item El instrumento público presentado en forma.
    \item El instrumento privado suscripto por el obligado, reconocido judicialmente, o cuya firma estuviese certificada por escribano con intervención del obligado y registrada la certificación en el protocolo.
    \item La confesión de deuda líquida y exigible prestada ante el juez competente para conocer en la ejecución.
    \item La cuenta aprobada o reconocida como consecuencia del procedimiento establecido en el artículo 525.
    \item La letra de cambio, factura de crédito, cobranza bancaria de factura de crédito, vale o pagaré, el cheque y la constancia de saldo deudor en cuenta corriente bancaria, de acuerdo con las condiciones establecidas por las leyes en la materia.
    \item El crédito por alquileres o arrendamiento de muebles.
    \item Los demás títulos que tuvieran fuerza ejecutiva por ley y no estén sujetos a un procedimiento especial.
\end{enumerate}
\end{quote}

\subsection{El inciso 7.º: apertura a títulos especiales}

El código abre la posibilidad de acceder a la vía ejecutiva a cualquier otro título al que una ley especial le dé fuerza ejecutiva. Es el caso, por ejemplo, del cobro de algunas matrículas en ciertos colegios profesionales, las deudas fiscales y otros créditos que gozan de este beneficio: la entidad emite un certificado de deuda ---por ejemplo, las matrículas de los abogados en los colegios profesionales--- y ese instrumento constituye un título ejecutivo que emana directamente del acreedor.

También tiene fuerza ejecutiva el crédito por expensas ---las llamadas \textbf{ejecuciones de expensas}---, que se conforma, por ejemplo, con el certificado de deuda suscripto por el administrador del consorcio.

%----------------------------------------------------------
\section{Títulos incompletos y preparación de la vía ejecutiva}

Dentro de la enumeración de los títulos ejecutivos se encuentran los denominados \textbf{títulos incompletos}: instrumentos que no tienen por sí mismos fuerza ejecutiva, sino que requieren algún trámite previo para lograr la \textbf{preparación de la vía ejecutiva}.

\begin{example}{el contrato de locación}
Cuando el locador quiere ejecutar los alquileres pactados en un contrato de locación, del título no surge cuál es la suma líquida que pretende ejecutar, porque al firmarse el contrato no se sabía en qué mes el deudor dejaría de pagar el canon locativo. Por eso, de modo previo, se cita al inquilino al tribunal para que acompañe el último recibo de pago del alquiler: si no lo trae, se tienen por adeudados todos los posteriores y, con una cuenta sencilla, se obtiene la suma fácilmente liquidable.
\end{example}

Si el contrato de locación tampoco tiene la firma certificada por escribano, el inquilino deberá reconocer o desconocer la firma que se le atribuye. La sanción por no concurrir es que se lo tendrá por reconocido. Si concurre y niega la firma, y luego se prueba mediante pericial que la firma le pertenece, le corresponde una multa del 30\,\% más de la deuda que se reclama.

%----------------------------------------------------------
\section{El embargo ejecutivo}

Cuando se tiene un título completo con fuerza ejecutiva ---o un título incompleto que luego se preparó y adquirió fuerza ejecutiva--- el paso siguiente es directamente la intimación al deudor para que oponga excepciones y el embargo. El deudor es embargado desde el comienzo del proceso ejecutivo.

El \textbf{embargo ejecutivo} es el que se traba a consecuencia de la fuerza ejecutiva de un título cartular.

Los tres tipos de embargo ---preventivo, ejecutivo y ejecutorio--- se distinguen según el momento procesal en que se traban y se desarrollan en el capítulo~\ref{cap:embargo}.

%----------------------------------------------------------
\section{La intimación de pago --- Artículo 531}

La intimación de pago equivale, en el juicio ejecutivo, al traslado de demanda del juicio ordinario, pero se realiza mediante un \textbf{mandamiento}. El mandamiento, que dice «intimación de pago y citación de remate», se libra al domicilio del deudor.

Se otorgan cinco días hábiles judiciales para interponer las excepciones que el deudor se crea con derecho a oponer. El mandamiento puede estar acompañado del embargo domiciliario ---en cuyo caso el oficial de justicia va al domicilio, intima de pago y embarga los muebles embargables--- o ser simplemente un mandamiento de intimación de pago, y el embargo se tramita por otro carril: si se trata de inmuebles, se libra oficio directamente al Registro de la Propiedad Inmueble; si se trata de cuentas bancarias, a la entidad bancaria donde esté registrada la cuenta.

\begin{quote}
\textbf{Art. 531 CPCCN.} Regula la intimación de pago y la forma del mandamiento.

Son irrenunciables para el deudor (no pueden ser objeto de acuerdos procesales): 1) la intimación de pago; 2) la citación para oponer excepciones; 3) la sentencia.
\end{quote}

%----------------------------------------------------------
\section{Ampliación de la ejecución --- Artículos 540 y 541}

En el curso del juicio ejecutivo ---especialmente en las ejecuciones de expensas--- pueden vencer nuevos períodos de la misma obligación mientras el juicio tramita. El código permite ampliar la ejecución a esos nuevos períodos.

\begin{quote}
\textbf{Art. 540 CPCCN.} Ampliación por vencimientos anteriores a la sentencia: se reitera la intimación y los nuevos períodos pasan a integrar la suma líquida reclamada.

\textbf{Art. 541 CPCCN.} Vencimientos posteriores a la sentencia: sólo se permite que el deudor traiga el recibo de pago; no se abre nueva discusión.
\end{quote}

%----------------------------------------------------------
\section{Excepciones admisibles --- Artículo 544}

Una vez notificado, el ejecutado tiene cinco días para interponer las excepciones previstas y limitadas por el artículo 544, que son enunciadas de modo taxativo. Si no opone excepciones, el juez dicta la sentencia denominada de \textbf{trance y remate}.

\begin{quote}
\textbf{Art. 544 CPCCN --- Excepciones admisibles en el juicio ejecutivo:}
\begin{enumerate}[label=\arabic*.º),leftmargin=2.2em]
    \item Incompetencia.
    \item Falta de personería en el ejecutante o en el ejecutado, por carecer de capacidad civil para estar en juicio.
    \item Litispendencia.
    \item Falsedad o inhabilidad del título con el que se pide la ejecución.
    \item Prescripción.
    \item Pago total o parcial documentado.
    \item Compensación de crédito líquido que resulte de documento que traiga aparejada la ejecución.
    \item Quita, espera, remisión, novación, transacción, conciliación o cosa juzgada.
\end{enumerate}
\end{quote}

\begin{note}
La falta de legitimación ---omitida por el código--- se ha incorporado jurisprudencial y doctrinariamente a través de la excepción de inhabilidad de título.
\end{note}

%----------------------------------------------------------
\section{Falsedad e inhabilidad del título (art.~544, inc.~4.º)}

\subsection{Falsedad material y falsedad ideológica}

Es fundamental distinguir entre la falsedad material y la falsedad ideológica.

\begin{table}[H]
\centering
\renewcommand{\arraystretch}{1.25}
\begin{tabularx}{\textwidth}{Y Y}
\toprule
\textbf{Falsedad material (admisible)} & \textbf{Falsedad ideológica (inadmisible en el ejecutivo)} \\
\midrule
Adulteración física del documento. & Cuestiona la causa de la obligación. \\
Ejemplo: tachaduras, sobreescrituras, fecha alterada, nombre agregado. & Ejemplo: alegar que no se recibió el dinero del préstamo. \\
Se puede probar con pericial caligráfica. & Debe discutirse en el juicio ordinario posterior. \\
Cosa juzgada material para ese punto. & Produce sólo cosa juzgada formal en el ejecutivo. \\
\bottomrule
\end{tabularx}
\end{table}

La falsedad a la que se refiere la excepción es la \textbf{adulteración física del documento}: que haya sido tachado, enmendado o sobrescrito, que se le hayan agregado condiciones después de firmado, que se haya agregado un nombre o que se haya cambiado o sobrescrito la fecha. Toda adulteración o falsedad material del documento es una excepción admisible en el marco del juicio ejecutivo. Lo que no es admisible es la llamada \textbf{falsedad ideológica}.

\subsection{El caso extremo: ordinarización del proceso}

\begin{example}{hipoteca firmada bajo amenazas}
Una persona firma una hipoteca bajo coacción con amenazas ---con un revólver--- y no recibe nada a cambio. Luego los acreedores ---en este caso, delincuentes--- ejecutan la hipoteca, subastan la casa y se quedan con ese dinero.
\end{example}

En casos extremos como este, la jurisprudencia ha admitido la \textbf{ordinarización del proceso ejecutivo}: se permite que el juicio ordinario corra de modo paralelo al ejecutivo. En ese ordinario, el deudor podría pedir como medida cautelar la suspensión del juicio ejecutivo o, al menos, que ambas sentencias se dicten de modo conjunto.

\subsection{La inhabilidad del título}

La excepción de inhabilidad de título permite cuestionar que el instrumento presentado no trae aparejada la ejecución, lo cual obliga al juez a examinar nuevamente el título en la sentencia. También se utiliza para plantear la falta de legitimación activa o pasiva: que el acreedor no es la persona que surge del título, o que el ejecutado no es quien debe.

\begin{important}
Requisito formal (art.~544, inc.~4.º): tanto la excepción de falsedad como la de inhabilidad del título sólo pueden interponerse siempre que el ejecutado niegue la existencia de la deuda.
\end{important}

%----------------------------------------------------------
\section{Apertura a prueba y sentencia de trance y remate}

El juez analiza las excepciones, rechaza conforme el artículo 547 las que no fueron autorizadas por ley y evalúa las restantes para verificar si cumplen los requisitos. Si son de puro derecho, las resolverá en la sentencia de trance y remate. Si corresponde, puede abrirse prueba por un plazo breve.

La sentencia del juicio ejecutivo se denomina \textbf{sentencia de trance y remate} y debe dictarse en un plazo de diez días, porque se equipara a una resolución interlocutoria.

\begin{quote}
\textbf{Art. 547 CPCCN.} El juez rechaza \emph{in limine} las excepciones no previstas legalmente.
\end{quote}

%----------------------------------------------------------
\section{Apelación e inapelabilidad en el juicio ejecutivo}

En el juicio ejecutivo la regla es contraria a la del juicio ordinario: la \textbf{inapelabilidad} es la norma y sólo son apelables las decisiones que el código autoriza expresamente. Cualquier apelación que se deduzca en el trámite ---salvo las autorizadas--- será denegada. La regla de la inapelabilidad es mucho más estricta para el ejecutado en el marco del juicio ejecutivo.

\begin{quote}
\textbf{Art. 554 CPCCN --- Casos en que procede la apelación de la sentencia de trance y remate:}
\begin{enumerate}[label=\alph*),leftmargin=2.2em]
    \item Cuando se desestiman sin más las excepciones por no cumplirse los requisitos del art.~544.
    \item Cuando las excepciones hubiesen tramitado como de puro derecho.
    \item Cuando se hubiese producido prueba respecto de las opuestas.
    % TODO: contenido incompleto o ambiguo en las notas originales (la redacción del inciso d) del apunte es confusa: "cuyo ordinario posterior")
    \item Cuando la sentencia versare sobre puntos ajenos al ámbito natural del proceso o que cause gravamen irreparable cuyo ordinario posterior.
\end{enumerate}
\end{quote}

El recurso aplicable es la \textbf{apelación en relación}, más acotada que la apelación libre.

%----------------------------------------------------------
\section{Costas y juicio ordinario posterior}

Las costas siguen la regla general: las soporta la parte vencida. La salvedad es que, si el ejecutado logra que se rechace una parte de la ejecución ---por ejemplo, acreditando el pago parcial documentado---, no se imponen las costas de modo proporcional al porcentaje rechazado: la regla es que prospera la ejecución y no hay eximición parcial de costas.

\begin{quote}
\textbf{Art. 553 CPCCN --- Condiciones para el juicio ordinario posterior:} hasta que no esté concluido el juicio ejecutivo no se podría promover el ordinario posterior (salvo los casos de ordinarización paralela ya vistos).
\end{quote}

%----------------------------------------------------------
\section{Ejecuciones especiales --- Artículo 595 y siguientes}

A partir del artículo 595 el código regula las ejecuciones especiales: distintos supuestos de subastas de bienes muebles, inmuebles y otras cuestiones. Se prevén reglas específicas para algunos tipos de ejecución que se basan en el trámite del juicio ejecutivo pero tienen condiciones especiales.

\subsection{La ejecución hipotecaria}

Mediante la Ley 24.441 se creó la denominada hipoteca privada: una ejecución que se apoya en la jurisdicción para realizar algunos trámites pero que se lleva adelante fundamentalmente de modo privado.

Esta regulación recibe objeciones, porque omite el dictado de una sentencia: si el deudor no opone una excepción, el juez no tiene que dictar sentencia. Ello viola directamente la Constitución Nacional, que establece que nadie puede ser privado de su propiedad sino en virtud de una sentencia previa. Si el deudor no pone excepciones, aun cuando debería haber sentencia, se ejecutan los bienes del deudor.

% TODO: contenido incompleto o ambiguo en las notas originales (frase ininteligible en el apunte: "Esto es algo que el adorno vio")
Esta deficiencia no fue subsanada hasta la fecha.

\cornelltitle{juicio ejecutivo}
\cornellblock{
\cqrow{¿Qué es el juicio ejecutivo y cómo se ubica en el código?}{Es un proceso de conocimiento sumario (acotado) regulado a partir del art.~520 CPCCN, dentro del proceso de ejecución. Prioriza la rapidez sobre la exhaustividad: el acreedor accede directamente al embargo y al cobro sin necesidad de un juicio ordinario previo.
}
\cqrow{¿Cuáles son los requisitos para acceder a la vía ejecutiva?}{1) Tener un título ejecutivo (completo o preparado). 2) Que la obligación sea de dar sumas de dinero. 3) Que sea líquida o fácilmente liquidable. 4) Que sea exigible (plazo vencido).
}
\cqrow{¿Qué son los títulos completos e incompletos?}{Completos: traen aparejada la ejecución por sí solos (art.~523: pagaré, cheque, escritura pública, etc.). Incompletos: requieren un trámite previo de preparación de la vía ejecutiva (por ejemplo, el contrato de locación, que requiere la citación del inquilino para obtener la suma líquida).
}
\cqrow{¿Cuáles son los tres tipos de embargo?}{Preventivo: como medida cautelar (requiere verosimilitud y peligro en la demora). Ejecutivo: en el juicio ejecutivo, por la fuerza del título cartular, sin sentencia previa. Ejecutorio: resultado de una sentencia firme en ejecución.
}
\cqrow{¿Qué excepciones admite el juicio ejecutivo?}{Las del art.~544, taxativas: incompetencia, falta de personería, litispendencia, falsedad o inhabilidad del título, prescripción, pago documentado, compensación con título ejecutivo, y quita, espera, remisión, novación, transacción, conciliación o cosa juzgada. La falta de legitimación se incorpora por vía de la inhabilidad de título.
}
\cqrow{¿Qué diferencia hay entre falsedad material e ideológica?}{Material (admisible en el ejecutivo): adulteración física del documento (tachaduras, sobreescrituras, fecha alterada). Ideológica (inadmisible): cuestiona la causa de la obligación (por ejemplo, «firmé pero no recibí el dinero»); debe discutirse en el juicio ordinario posterior.
}
\cqrow{¿Qué es la cosa juzgada formal en el juicio ejecutivo?}{Las cuestiones causales no debatibles en el ejecutivo no quedan cubiertas por cosa juzgada material: pueden discutirse en el ordinario posterior. En cambio, las cuestiones que sí se debatieron (por ejemplo, el pago) generan cosa juzgada material y no se pueden replantear.
}
}
\medskip
\cornellblocksum{
\cqrow{¿Cómo es la apelación en el juicio ejecutivo?}{La regla es la inapelabilidad (inversa a la del juicio ordinario). Sólo son apelables las resoluciones que el código autoriza expresamente (art.~554). El recurso aplicable es la apelación en relación, más acotada que la libre.
}
}{En el juicio ejecutivo el legislador sacrifica la exhaustividad del conocimiento en aras de la celeridad: el acreedor que cuenta con un título ejecutivo ---completo o debidamente preparado--- puede agredir el patrimonio del deudor de manera inmediata, sin aguardar un proceso ordinario de pleno conocimiento. Las excepciones que el deudor puede oponer son taxativas (art.~544 CPCCN) y excluyen la discusión de la causa de la obligación, reservada para el juicio ordinario posterior. La cosa juzgada resultante es, en consecuencia, mixta: formal respecto de las cuestiones causales no debatidas y material respecto de aquellas que sí pudieron ventilarse (como el pago).
}

\chapter{Cumplimiento de la sentencia de remate}\label{cap:remate}

\chaptermark{Cumplimiento de la sentencia de remate}

%----------------------------------------------------------
\section{La etapa coercitiva del proceso}

Hasta la sentencia definitiva, toda la discusión ---la demanda, la contestación, la prueba, la etapa recursiva--- tuvo como meta obtener un reconocimiento de derechos en sede jurisdiccional. Corresponde ahora llevar a la realidad lo dispuesto en el proceso judicial: los derechos reconocidos deben tener su correlato en prestaciones o conductas de las personas en la vida de relación, como pagar la deuda reconocida judicialmente o hacer o dejar de hacer alguna actividad.

Los elementos de la jurisdicción pueden enunciarse en tres: la \textbf{notio} (conocer el conflicto), la \textbf{documentación} y la \textbf{coercio}, es decir, la coerción: el Estado impone su decisión por la fuerza cuando las personas no han cumplido voluntariamente la decisión jurisdiccional.

\subsection{Títulos asimilables a la sentencia}

Tanto el cumplimiento de la sentencia de trance y remate como el procedimiento de ejecución de sentencia ---que se utiliza cuando el juicio ordinario está concluido--- sirven también para los títulos asimilables a la sentencia:

\begin{itemize}
    \item Acuerdos de mediación.
    \item Honorarios regulados judicialmente.
    \item Sentencias extranjeras reconocidas por el mecanismo de exequátur.
    \item Otros que las leyes hayan asimilado.
\end{itemize}

\subsection{Inapelabilidad del ejecutado --- Artículo 560}

El ejecutado que llegó a esta etapa tuvo amplia posibilidad de defensa durante todo el proceso previo. Por eso, en el cumplimiento de la sentencia de remate la regla es la \textbf{inapelabilidad para el ejecutado}.

\begin{quote}
\textbf{Art. 560 CPCCN.} Son inapelables para el ejecutado todas las resoluciones que se dicten durante el trámite de cumplimiento de la sentencia de remate.

Excepciones: las resoluciones que no puedan constituir objeto del juicio ordinario posterior; las que, debiendo ser objeto del ordinario posterior, ya fueron debatidas; las que se relacionen con el carácter de parte del ejecutado; y los casos del art.~550 y 4, inc.~4 (puntos ajenos al ámbito natural del proceso que causan gravamen irreparable).
\end{quote}
% TODO: contenido incompleto o ambiguo en las notas originales (la remisión "art. 550 y 4, inc. 4" figura así en el apunte)

%----------------------------------------------------------
\section{El embargo ejecutorio: presupuesto del remate}

No es posible avanzar hacia el remate de una cosa si no ha sido embargada previamente. Se trata del \textbf{embargo ejecutorio}, que es el resultado de la traba que se realiza una vez que existe una condena firme en sede jurisdiccional.

Si lo que se va a percibir son sumas de dinero ya embargadas, no hay que hacer ningún remate: con la firmeza de la sentencia de venta se practica una liquidación y se ordena el giro para que el acreedor se haga de ese dinero. El remate cobra relevancia cuando hay un bien ---mueble o inmueble--- que debe subastarse para transformarlo en dinero.

%----------------------------------------------------------
\section{Bienes subastables}

Se puede subastar todo aquello que esté en el comercio, todo aquello que pueda ser vendido.

También es posible subastar derechos hereditarios. Cuanto mayor sea la determinación del acervo hereditario, con mayor precisión se podrá establecer el valor de ese derecho y, de tal manera, fijar la base de la venta.

%----------------------------------------------------------
\section{El martillero: designación y funciones}

El encargado de llevar adelante la subasta a las órdenes del juez es el \textbf{martillero} o corredor inmobiliario, que se anota en las listas de las cámaras para ser designado. Actualmente, en el fuero civil de la Nación, la designación se realiza directamente por sorteo por computadora.

\subsection{Obligaciones del martillero}

\begin{enumerate}
    \item Cumplir las funciones personalmente: no puede delegarlas.
    \item Responder sólo a las órdenes del juez, no de las partes.
    \item Aceptar el cargo obligatoriamente.
    \item Informar el valor fiscal y de mercado del bien a subastar.
    \item Exhibir el inmueble a los eventuales compradores.
    \item Publicar los edictos.
    \item Depositar en el proceso las sumas percibidas en el acto de subasta.
    \item Puede pedir un anticipo de fondos al ejecutante para cubrir los gastos del trámite.
\end{enumerate}

%----------------------------------------------------------
\section{La liga: distorsión del sistema de subastas}

\begin{example}{el fenómeno de la «liga»}
Una serie de personas se dedica a distorsionar lo que se obtiene en las subastas judiciales.

\textbf{Mecanismo:} impiden que concurran eventuales compradores interesados en un bien particular, de modo de intentar adquirirlo por el menor valor posible ---si es posible, por la base---. Luego realizan la «verdadera subasta» fuera del tribunal: el verdadero comprador paga a la liga la diferencia entre lo pagado en sede judicial y el valor real.

\textbf{Impacto:} esa diferencia ---que puede oscilar entre un 20 y un 30\,\% del valor de un inmueble--- es la que el deudor no obtiene. El acreedor cobra menos y el deudor ---si había un remanente--- también pierde.

\textbf{Ejemplo numérico:} si el inmueble vale \$100.000 y se subasta a \$60.000 cuando podría haberse subastado a \$90.000, esos \$30.000 quedan en poder de la liga.
\end{example}

%----------------------------------------------------------
\section{Subastas electrónicas}

Se señala como una necesidad impostergable la modernización del sistema: resulta fuera de contexto que en la práctica se siga sacando a subasta un inmueble en una oficina, con todos los interesados presentes durante unas dos horas para hacer la puja en forma presencial. Los anuncios podrían hacerse por internet en lugar de por edictos, que son costosos, y las subastas podrían realizarse de modo digital o electrónico.

La Provincia de Buenos Aires tiene una ley de subasta electrónica que se está implementando gradualmente. La subasta no dura dos horas sino quince días: distintos interesados toman conocimiento del inmueble por internet, se registran y hacen ofertas, que se actualizan a medida que aparece un mejor postor.

Las ventajas son múltiples: elimina a la liga, obtiene valores mucho más altos y reales y beneficia tanto al acreedor ---que cobra más--- como al deudor ---que puede obtener un remanente luego de pagada la deuda---.

%----------------------------------------------------------
\section{Trámites previos a la subasta de inmuebles}

\subsection{Informes registrales y de deudas}

Cuando el juez ordena la subasta y designa al martillero, ordena también que las partes obtengan informes de todos los entes ---fundamentalmente las municipalidades y los organismos de recaudación provincial, como ARBA en Buenos Aires--- sobre las deudas que tienen los inmuebles. También se solicitan informes de deudas por servicios: agua (AySA), gas, luz, etc.

\begin{itemize}
    \item \textbf{Informes registrales:} se requieren al Registro de la Propiedad Inmueble sobre embargos, inhibiciones y gravámenes (hipotecas) que pesen sobre el inmueble y sobre el ejecutado.
    \item \textbf{Validez de los informes:} 60 días. No es necesario actualizarlos antes de la subasta, salvo que haya transcurrido mucho tiempo por alguna suspensión del proceso.
    \item \textbf{Inmuebles en propiedad horizontal:} debe requerirse también la deuda por expensas.
\end{itemize}

\subsection{Citación de los acreedores con gravámenes}

Este trámite es muy importante, porque no es lo mismo subastar un inmueble que no tiene gravamen que subastar uno que tiene una hipoteca previa. Cuando el inmueble tiene hipotecas o embargos, el acreedor, antes de cobrar, deberá respetar las preferencias y privilegios de esos gravámenes.

Una vez que el inmueble se subasta y el embargo pasa al producido de la subasta ---por subrogación real---, el juez debe repartir ese dinero de acuerdo con los privilegios y preferencias de los gravámenes y embargos. Por eso se notifica a todos los acreedores antes de la subasta, para que puedan presentarse, defender el precio de venta y controlar la legalidad del proceso.

%----------------------------------------------------------
\section{Los edictos y la publicidad --- Artículo 566}

\begin{quote}
\textbf{Art. 566 CPCCN.} El remate se anunciará por dos días en el Boletín Oficial y en otro diario de amplia circulación del lugar donde se ubica el inmueble.
\end{quote}

El edicto debe contener obligatoriamente:

\begin{itemize}
    \item Expediente, carátula y juzgado donde tramita.
    \item Fecha, hora y lugar de la subasta.
    \item Días y horas en que el martillero exhibirá el inmueble a los interesados.
    \item Deudas del inmueble (servicios, tasas, contribuciones municipales, expensas).
    \item Gravámenes que pesan sobre el bien.
    \item Obligación del adquirente de pagar en el acto la seña y la comisión del martillero.
    \item Modalidades especiales fijadas por el juez para la venta.
\end{itemize}

Se observa que la referencia a los diarios como medio de comunicación más relevante resulta anacrónica, dado que poca gente sigue comprando diarios, si se compara con el tiempo en que la norma fue redactada. Sin embargo, la norma continúa vigente y debe cumplirse.

El martillero puede también solicitar al tribunal autorización para realizar publicidad adicional ---por ejemplo, por internet--- para atraer más compradores.

%----------------------------------------------------------
\section{La base de la subasta --- Artículo 578}

El martillero debe informar al juez el valor fiscal del inmueble y su valor de mercado. Muchas veces los valores fiscales están muy desconectados del verdadero valor de mercado de una propiedad, es decir, del valor que podría llegar a obtenerse en una venta.

\begin{formula}{base de la subasta (art.~578 CPCCN)}
Si no hay acuerdo de partes, la base de venta se fijará en \textbf{dos tercios de la valuación fiscal actualizada}.
\end{formula}

\begin{note}
En general, el juez se aparta de ese valor fiscal para impedir que los bienes sean malvendidos, fijando una base menor al valor de mercado para atraer distintos compradores. Las partes ---y los acreedores con gravámenes--- pueden impugnar la tasación informada por el martillero.
\end{note}

%----------------------------------------------------------
\section{El acto de remate}

Hay distintos resultados posibles: que haya un comprador, que el remate fracase por falta de postores o que el propio acreedor ejecutante quiera adquirir el bien.

\subsection{Cuando hay comprador}

Si hay comprador, el martillero le hace firmar en el acto de la subasta el \textbf{boleto de venta judicial}, donde consta el precio de compra, lo que paga de seña, lo que paga de comisión del martillero y algún otro impuesto que corresponda a la venta. El martillero retiene su comisión y deposita el dinero obtenido en el proceso, informando el nombre e identidad del comprador.

El comprador debe constituir domicilio electrónico (art.~41 CPCCN) para ser notificado de las decisiones judiciales.

Secuencia posterior al acto de remate:

\begin{enumerate}
    \item El martillero informa el resultado al tribunal.
    \item Se da conocimiento del resultado a acreedor y deudor, que tienen cinco días para formular objeciones.
    \item Sin objeciones (o desestimadas), el juez aprueba el remate.
    \item El comprador tiene cinco días para depositar el saldo de precio: el 70\,\% restante, dado que el 30\,\% se paga como seña en el acto.
\end{enumerate}

\subsection{Crítica a la exigencia del dinero en efectivo}

Se señala que exigir el 30\,\% en efectivo el día de la subasta está fuera de contexto: obliga a las personas a concurrir con dinero que quizá no usarían, con todos los riesgos que ello trae. Para los fines de la efectividad de la subasta, resulta bastante complejo.

%----------------------------------------------------------
\section{Aprobación y nulidades}

El plazo para plantear nulidades es de cinco días contados desde que se conoce el resultado de la subasta. Para quien estuvo presente, el cómputo se inicia desde el propio acto, independientemente de cuándo el martillero informe al tribunal.

Las nulidades que pueden plantearse en esta etapa son las \textbf{nulidades propias del acto}: las ocurridas ese día, no las nulidades procesales anteriores. Estas últimas ---las que se fueron produciendo desde que se ordenó la subasta, como la base o el edicto--- debían cuestionarse dentro de los cinco días desde que se conocieron en el expediente. En general, la única nulidad que podrá plantearse es algún incumplimiento en el acto propio de la subasta.

\begin{important}
\textbf{Principio de trascendencia:} quien articula la nulidad debe acreditar el perjuicio. Sin perjuicio concreto, no procede la nulidad: la falta menor no hace caer el acto. La trascendencia de la nulidad debe tener entidad suficiente para hacer caer la venta.
\end{important}

%----------------------------------------------------------
\section{El sobreseimiento del juicio --- Artículo 583}

Aun hecha la subasta, el deudor tiene una posibilidad más de recuperar el bien: el ejecutado puede, antes de que el comprador en subasta deposite el saldo de precio, pedir el \textbf{sobreseimiento}. Una vez que el comprador depositó el saldo de precio, se considera perfeccionada la venta.

Para que proceda el sobreseimiento, el deudor debe depositar:

\begin{enumerate}
    \item El capital adeudado.
    \item Los intereses presupuestados.
    \item Las costas del proceso.
    \item Todos los gastos.
    \item Un importe a favor del comprador integrado por la comisión que pagó al martillero, los sellados del boleto y el equivalente a la seña depositada.
\end{enumerate}

%----------------------------------------------------------
\section{Título del adquirente e inscripción registral}

Así como en una compraventa de inmueble fuera de los tribunales hay una escritura, dentro del proceso se extiende un \textbf{testimonio} que contiene toda la referencia y los antecedentes de la subasta judicial y las distintas resoluciones del juez, para que el adquirente en subasta lo inscriba en el Registro de la Propiedad Inmueble. Ese testimonio será el nuevo título del bien. Puede obtenerse sin necesidad de escribano: es un trámite más del proceso, que extiende el secretario del tribunal.

\subsection{Levantamiento de gravámenes}

El adquirente pide el levantamiento de todas las medidas dispuestas contra el inmueble ---hipotecas, embargos--- al solo efecto de escritura, para que el título ingrese y el inmueble quede liberado. Los acreedores que tenían esos gravámenes deben presentarse en el proceso a hacer valer sus preferencias y privilegios en la distribución del producido.

%----------------------------------------------------------
\section{Distribución del producido: la liquidación}

Una vez depositado el saldo de precio, se practica la liquidación: capital, intereses, costas y gastos. Se pone en conocimiento del deudor ---que tiene un plazo para impugnarla--- y, una vez aprobada, se entrega al ejecutante, cancelando el crédito.

El deudor puede tener un \textbf{remanente}: si el producido de la subasta supera la deuda total (capital, intereses y costas), el excedente le corresponde al deudor.

%----------------------------------------------------------
\section{Subasta fracasada y acreedor comprador}

\subsection{Fracaso por falta de postores}

Si no hay postores, se realizará una nueva subasta, en general intentando bajar la base para atraer nuevos interesados (arts.~592 y 593 CPCCN).

\subsection{El acreedor como comprador}

El acreedor puede concurrir a comprar él mismo el inmueble que se subasta y \textbf{compensar su crédito con el precio}. En general, los acreedores ejecutantes piden al juez que, en caso de resultar adquirentes, los exima de depositar la seña. Es un gran beneficio para ese comprador: no tiene que concurrir con el dinero en la mano, expuesto a la inseguridad.

Esta posibilidad también podría ser solicitada por algún acreedor embargante de otro juicio que deba concurrir a la subasta del inmueble en el proceso en que se va a realizar, para ser eximido del depósito de la seña.

%----------------------------------------------------------
\section{Los plenarios: diferencia entre fueros}

\begin{important}
\textbf{Plenario de la Cámara Comercial:} el adquirente debe pagar las deudas que están expresadas en el edicto.

\textbf{Plenario de la Cámara Civil:} el adquirente no tiene que afrontar todas esas deudas, en tanto alcance el dinero que se obtiene en la subasta.

\textbf{Consecuencia práctica:} según se subaste en un juzgado civil o comercial, la ecuación económica para el comprador cambia radicalmente. Un abogado debe verificar los plenarios vigentes antes de aconsejar a su cliente.
\end{important}

%----------------------------------------------------------
\section{Desocupación del inmueble --- Artículo 589}

\begin{quote}
\textbf{Art. 589 CPCCN.} Como regla, no hay desalojo de los ocupantes del inmueble hasta que no se hubiese pagado el saldo de precio y hecho la tradición de la cosa al adquirente. Se exceptúan los casos en que se haya pactado otra cosa, como en las hipotecas, donde el hipotecado generalmente se obliga a desocupar dos días después del auto de subasta.
\end{quote}

\cornelltitle{cumplimiento de la sentencia de remate}
\cornellblocksum{
\cqrow{¿Qué es el cumplimiento de la sentencia de remate?}{Es la etapa coercitiva del proceso: una vez firme la sentencia de trance y remate (o de venta en el ordinario), el Estado impone su decisión por la fuerza ejecutando los bienes del deudor. Presupuesto indispensable: el embargo ejecutorio.
}
\cqrow{¿Cuáles son las funciones del martillero?}{Auxiliar del tribunal: evalúa el bien, informa valor fiscal y de mercado, lo exhibe a los compradores, publica edictos, conduce el acto de remate, suscribe el boleto de venta judicial, retiene su comisión y deposita el producido. No puede recibir órdenes de las partes.
}
\cqrow{¿Qué es «la liga» y cómo opera?}{Un grupo organizado que impide la concurrencia de compradores interesados al acto de remate para adquirir el bien por el mínimo (la base) y luego realiza una subasta privada fuera del tribunal. La diferencia de precio (20--30\,\% del valor) queda en sus manos, en perjuicio del acreedor y del deudor.
}
\cqrow{¿Cuál es la base de la subasta según el art.~578?}{Si no hay acuerdo de partes: dos tercios de la valuación fiscal actualizada. En la práctica, el juez suele apartarse de ese valor fiscal (muy desconectado del mercado) y fija una base menor al valor de mercado para atraer compradores sin que el bien sea malvendido.
}
\cqrow{¿Cuáles son los pasos tras el acto de remate?}{1) El martillero informa al tribunal. 2) Cinco días para objeciones. 3) Aprobación del remate. 4) Cinco días para depositar el saldo de precio. 5) Posibilidad de sobreseimiento por el deudor antes del depósito. 6) Liquidación, aprobación y pago al acreedor. 7) Testimonio e inscripción del título del adquirente.
}
\cqrow{¿Qué es el sobreseimiento del juicio (art.~583)?}{El ejecutado puede recuperar el bien después del remate pero antes de que el comprador deposite el saldo de precio. Para ello debe pagar: capital, intereses, costas, gastos, comisión del martillero, sellados del boleto y el equivalente a la seña.
}
\cqrow{¿Por qué importan los plenarios sobre deudas en las subastas?}{La Cámara Comercial establece que el adquirente paga sólo las deudas expresadas en el edicto. La Cámara Civil establece que todas las deudas se abonan con el producido de la subasta (el adquirente no las asume si alcanza el dinero). La diferencia cambia radicalmente la ecuación económica para el comprador.
}
}{Una vez firme la sentencia de trance y remate, el Estado despliega su poder coercitivo pleno: se traba el embargo ejecutorio, se designa un martillero, se practican los informes registrales, se publican edictos, se notifica a los acreedores concurrentes y se realiza el acto de subasta. El producido se distribuye según los privilegios y preferencias de los gravámenes. El adquirente en subasta obtiene su título a través de un testimonio extendido por el secretario del tribunal, sin necesidad de escritura pública. El sistema presenta dos grandes deficiencias estructurales: la persistencia del mecanismo presencial, que alimenta el fenómeno de la «liga», y la divergencia entre los plenarios de la Cámara Civil y la Cámara Comercial sobre las deudas asumidas por el adquirente, que el operador jurídico debe consultar en cada caso para brindar un asesoramiento económicamente preciso.
}

\part{Sistemas cautelares}
\markboth{Parte V}{Parte V}
Los sistemas cautelares son herramientas procesales que permiten al juez \textbf{conservar} el estado de las cosas mientras dura el proceso, evitando que la sentencia llegue tarde o resulte inútil. Esta parte comienza con los fundamentos, la clasificación, los presupuestos y los caracteres de las medidas cautelares, continúa con el sistema asegurativo genérico y con el ámbito excepcional de la tutela anticipada, y concluye con el estudio de las medidas cautelares en particular: el embargo, el secuestro, la intervención judicial, la anotación de litis, la prohibición de contratar y las medidas de protección de personas y de familia. Las medidas en particular se estudian siempre en conexión con la parte general (verosimilitud del derecho, peligro en la demora y contracautela).

El proceso civil cuenta con diversas herramientas para garantizar la eficacia de la sentencia, del mismo modo en que un médico dispone de distintos instrumentos para tratar diferentes dolencias: desde el embargo, la más usada y análoga a un seguro sobre un bien, hasta las medidas de familia, orientadas a proteger derechos fundamentales de los niños y los vínculos familiares.

\chapter[Fundamentos de los sistemas cautelares]{Fundamentos de los sistemas cautelares}\label{cap:fundamentos-cautelares}

% TODO: verificar consistencia entre fuentes originales (uno de los archivos fuente numera de manera inconsistente la unidad de sistemas cautelares).

\section{Importancia de los sistemas cautelares}

Los sistemas cautelares constituyen un tema de gran trascendencia para el ejercicio profesional, porque son herramientas que posibilitan a los abogados, a quienes trabajan dentro del Poder Judicial y a los funcionarios cumplir con la \textbf{tutela judicial efectiva}.

Sin los sistemas cautelares sería necesario esperar a la sentencia en todos los procesos judiciales, lo que haría que muchos derechos quedaran vulnerados. En ausencia de cautelares, muchas sentencias solamente servirían para «colgarse en un cuadro».

Se trata de un tema de gran importancia para el justiciable y para todo operador jurídico que deba emprender la gestión de un conflicto, ya que estas herramientas permiten, frente a un peligro inminente o una situación afligente, dar una respuesta desde el sistema de justicia en un tiempo distinto del que insumiría un proceso hasta llegar a su fin, es decir, hasta la sentencia.

Los sistemas cautelares son herramientas procesales que permiten al juez \textbf{conservar} el estado de las cosas mientras dura el proceso, evitando que la sentencia llegue tarde o resulte inútil. Puede pensarse el proceso judicial como una carrera: los sistemas cautelares funcionan como cercos que impiden que alguien modifique el camino antes de que se declare al ganador.

\section{La actuación de la ley en el proceso}

Al estudiar la jurisdicción y sus distintas etapas (momentos jurisdiccionales) se distinguen tres modalidades de actuación de la ley en el proceso.

\begin{table}[H]
\centering
\small
\begin{tabularx}{\textwidth}{>{\raggedright\arraybackslash}p{0.2\textwidth} Y Y}
\toprule
\textbf{Modo de actuación} & \textbf{Cuándo ocurre} & \textbf{En qué etapa} \\
\midrule
Para \textbf{conocer} & Cuando la jurisdicción conoce los hechos & Etapa introductoria (escritos postulatorios), probatoria (producción de prueba) y alegatos \\
Para \textbf{ejecutar} & Cuando el juez ejecuta la sentencia & Luego de la sentencia firme y consentida (ejecutoriada) \\
Para \textbf{conservar / asegurar} & Antes o durante el proceso, ante urgencia & En cualquier momento en que se requiera respuesta inmediata \\
\bottomrule
\end{tabularx}
\end{table}

Esta distinción corresponde a las tres finalidades que, según Calamandrei, tiene la actuación de la ley en el proceso: que el juez conozca, que ejecute y que conserve.

La actuación de la ley con \textbf{carácter conservatorio} es la que interesa para el estudio de los sistemas cautelares. Esta actuación abarca tanto la función conservativa como la innovativa:

\begin{itemize}
    \item \textbf{Conservar para mantener un determinado statu quo:} cuando el riesgo proviene de que se altere la situación existente.
    \item \textbf{Conservar alterando:} cuando el riesgo proviene de que se mantenga una determinada situación o statu quo y esa conservación es la que genera el daño.
\end{itemize}

\section{El poder cautelar}

El \textit{iter} procesal no puede soslayarse: un proceso judicial insume necesariamente tiempo. Si se transcurre todo el \textit{iter} procesal hasta llegar a la sentencia, es probable que situaciones vigentes, o que requieran una respuesta rápida, inmediata y efectiva, no tengan solución. Los sistemas cautelares vienen a dar respuesta a esas situaciones.

\begin{definition}{Poder cautelar}
Respuesta efectiva e inmediata que se da en el marco de un proceso judicial prudente, con el fin de respetar la tutela judicial efectiva con relación a los derechos discutidos en ese proceso.
\end{definition}

Sus objetivos son:

\begin{itemize}
    \item Contrarrestar los efectos negativos del transcurso del tiempo en el proceso judicial.
    \item Respetar la tutela judicial efectiva de los derechos discutidos.
    \item Mantener la igualdad de las partes en el proceso.
\end{itemize}

\begin{important}
El legislador crea los sistemas cautelares para contrarrestar el efecto que el tiempo genera en el proceso ---lo que Calamandrei denominaba el ``ordinario iter procesal''--- y para mantener la \textbf{igualdad de las partes ante la jurisdicción}. Se busca dar respuesta a situaciones urgentes, garantizar una tutela judicial efectiva y evitar que la sentencia resulte estéril por tardía.
\end{important}

\begin{example}{Camacho Acosta (CSJN)}
Fallo de la Corte Suprema de Justicia de la Nación que se trabajará más adelante al estudiar la prohibición de innovar. Un trabajador pierde el antebrazo y necesita una prótesis. La prótesis exige una respuesta inmediata de la jurisdicción: no se puede esperar a la sentencia para decidir si la prótesis le correspondía o no. La jurisdicción debe anticipar esa decisión a fin de resguardar la salud y el derecho de esa persona a contar con la prótesis. Es el ejemplo paradigmático de la necesidad de las medidas cautelares.
\end{example}

Otros supuestos frecuentes en los que se requieren sistemas cautelares son:

\begin{itemize}
    \item La protección de una persona.
    \item La conservación de una prueba para el proceso.
    \item Mantener el giro comercial de una empresa.
    \item La aplicación de una norma inconstitucional a un caso concreto (solicitud de medida cautelar para suspender sus efectos mientras tramita el proceso).
\end{itemize}

\section{Cautela y condena}

\textbf{Cautela no es condena}: son dos términos distintos. El término \textit{cautelar} deriva de «cauto», quien observa con cuidado, y también de «caución», vinculado con «precaver» y con lo «precavido». Todos estos términos remiten a la prevención y a la conservación, es decir, a poder cautelar una determinada situación de hecho o de derecho que se dé en un caso concreto.

\begin{important}
La cautela se mueve en un ámbito \textbf{provisorio}, que nada tiene que ver con la condena. La condena se da únicamente en la sentencia, luego de pasar por todas las etapas del proceso.
\end{important}

\begin{table}[H]
\centering
\begin{tabularx}{\textwidth}{>{\raggedright\arraybackslash}p{0.47\textwidth} Y}
\toprule
\textbf{Cautela} & \textbf{Condena} \\
\midrule
Provisoria & Definitiva \\
Anterior a la sentencia & Es la sentencia definitiva \\
Basada en la verosimilitud del derecho & Basada en certeza jurídica plena \\
Puede modificarse o levantarse & No puede modificarse (cosa juzgada) \\
Tramita \textit{inaudita parte} & Con plena bilateralidad \\
Requiere contracautela & No requiere contracautela \\
Sujeta a caducidad & No caduca \\
Instrumental, accesoria & Principal, autónoma \\
\bottomrule
\end{tabularx}
\end{table}

La sentencia de mérito implica el máximo momento de la jurisdicción: fijar los hechos del proceso, elevarlos a las normas abstractas y crear la norma individual para el caso concreto, luego del conocimiento del juez, de los hechos, de la producción de la prueba y de los alegatos. Ese proceso previo a la sentencia es ajeno a la cautela, que tiene otro nivel de conocimiento, inferior, fundado en la verosimilitud del derecho.

\cornelltitle{fundamentos de los sistemas cautelares}
\cornellblocksum{
\cqrow{¿Por qué son importantes los sistemas cautelares?}{Porque permiten cumplir con la tutela judicial efectiva: sin ellos habría que esperar a la sentencia en todos los procesos y muchos derechos quedarían vulnerados; muchas sentencias solo servirían para «colgarse en un cuadro». Permiten dar una respuesta al justiciable, frente a un peligro inminente o una situación afligente, en un tiempo distinto del que insume el proceso hasta la sentencia.
}
\cqrow{¿Cuáles son los tres modos en que la ley actúa en el proceso?}{1) Para \textbf{conocer} (etapas introductoria, probatoria y de alegatos). 2) Para \textbf{ejecutar} (luego de la sentencia firme y consentida). 3) Para \textbf{conservar o asegurar} (ante urgencia, en cualquier momento en que se requiera respuesta inmediata); este último es el que interesa a los sistemas cautelares.
}
\cqrow{¿Qué significa conservar manteniendo el statu quo y conservar alterando?}{Conservar para mantener un statu quo se da cuando el riesgo proviene de que se altere la situación existente. Conservar alterando se da cuando el riesgo proviene de que se mantenga la situación y esa conservación es la que genera el daño.
}
\cqrow{¿Qué es el poder cautelar y cuál es su finalidad?}{Es la respuesta efectiva e inmediata que se da en el marco de un proceso judicial prudente para respetar la tutela judicial efectiva de los derechos discutidos. Busca contrarrestar los efectos negativos del transcurso del tiempo, respetar esa tutela y mantener la igualdad de las partes.
}
\cqrow{¿Qué ilustra el caso Camacho Acosta (CSJN)?}{La necesidad de las medidas cautelares: un trabajador que perdió el antebrazo necesita una prótesis y la jurisdicción debe anticipar la decisión, sin esperar a la sentencia, para resguardar su salud y su derecho.
}
\cqrow{¿Qué significa «cautelar» y en qué se diferencia de la condena?}{Deriva de «cauto» y de «caución» (precaver), lo que remite a prevención y conservación. La cautela es provisoria, anterior a la sentencia, basada en la verosimilitud, modificable, \textit{inaudita parte}, con contracautela, sujeta a caducidad e instrumental. La condena es definitiva, es la sentencia, se basa en certeza plena, no puede modificarse (cosa juzgada), tiene plena bilateralidad, no requiere contracautela, no caduca y es principal y autónoma.
}
}{la ley actúa en el proceso para conocer, ejecutar y conservar o asegurar. El poder cautelar es la respuesta efectiva e inmediata que contrarresta el efecto del tiempo procesal, respeta la tutela judicial efectiva y mantiene la igualdad de las partes. La cautela, provisoria y fundada en la verosimilitud, se distingue de la condena, definitiva y fundada en la certeza plena que solo se alcanza con la sentencia.
}

\cornelltitle{marco general de los sistemas cautelares}
\cornellblocksum{
\cqrow{¿Cuáles son las tres finalidades de la ley en el proceso según Calamandrei?}{La ley actúa para que el juez conozca (etapa introductoria y probatoria), para que ejecute (etapa de ejecución, para que sus decisiones sean cumplidas) y para que conserve (asegurar, proteger y preservar, que es el campo de los sistemas cautelares).
}
\cqrow{¿Qué significa que el juez ``conserve''?}{Conservar puede ser mantener o alterar una situación de hecho o de derecho, según de dónde provenga el perjuicio: se ordena mantener cuando el perjuicio proviene de alterar la situación, y a la inversa.
}
\cqrow{¿Por qué existen los sistemas cautelares?}{Para contrarrestar el efecto del tiempo en el proceso (el ``ordinario iter procesal'' de Calamandrei), mantener la igualdad de las partes ante la jurisdicción, responder a situaciones urgentes con tutela judicial efectiva y evitar una sentencia estéril por tardía.
}
\cqrow{¿Cómo se clasifican los sistemas cautelares?}{En ámbito normal (asegurativo probatorio, asegurativo garantizador, precautorio tradicional y asegurativo genérico), donde la pretensión cautelar no coincide con la sustancial; y ámbito excepcional (tutela anticipada), donde se superpone total o parcialmente con ella.
}
\cqrow{¿Qué presupuestos deben acreditarse para peticionar una medida cautelar?}{Verosimilitud en el derecho (probabilidad, no certeza), peligro en la demora (perjuicio inminente o irreparable), contracautela (garantía por los daños que pudiera causar la traba incorrecta) y los presupuestos procesales del Art.~195 CPCC.
}
\cqrow{¿Qué diferencia hay entre lo cautelar y la sentencia en cuanto a su carácter?}{Lo cautelar es provisorio y puede modificarse, reducirse, ampliarse o dejarse sin efecto según varíen las circunstancias; la sentencia es definitiva.
}
}{Los sistemas cautelares permiten al juez conservar la situación de hecho o de derecho durante el proceso, contrarrestando el efecto del tiempo y preservando la igualdad de las partes. Se dividen en un ámbito normal y un ámbito excepcional (tutela anticipada). Para obtener cualquier medida se acreditan verosimilitud en el derecho, peligro en la demora, contracautela y los presupuestos procesales del Art.~195 CPCC, y lo resuelto es siempre provisorio, a diferencia de la sentencia.
}

\chapter[Clasificación de los sistemas cautelares]{Clasificación de los sistemas cautelares}\label{cap:clasificacion-cautelar}

\section{El tratamiento como sistema}

Los institutos cautelares se estudian como un \textbf{sistema} porque, al recorrer el Código Procesal Civil y Comercial, se encuentra que la ley actúa con carácter conservativo o asegurativo en distintos lugares y con distintas modalidades. Agrupar esas manifestaciones en subsistemas permite comprenderlas en su conjunto y distinguirlas, a los fines de entender el sistema del código procesal en todos aquellos aspectos en los que la ley actúa de este modo.

\section{Ámbito normal y ámbito excepcional}

Los sistemas cautelares se clasifican en dos grandes ámbitos.

\begin{table}[H]
\centering
\begin{tabularx}{\textwidth}{>{\raggedright\arraybackslash}p{0.25\textwidth} Y}
\toprule
\textbf{Ámbito} & \textbf{Tipo de sistema} \\
\midrule
Normal & Sistemas asegurativos probatorios \\
Normal & Sistemas precautorios tradicionales \\
Normal & Sistemas asegurativos garantizadores \\
Normal & Sistemas asegurativos genéricos \\
Excepcional & Tutela anticipada \\
Excepcional & Sistemas protectores \\
\bottomrule
\end{tabularx}
\end{table}

Según otra formulación presente en las fuentes, el criterio que distingue ambos ámbitos es la relación entre la pretensión cautelar y la pretensión sustancial:

\noindent\begin{tabularx}{\textwidth}{@{}>{\raggedright\arraybackslash}p{2.6cm} Y Y@{}}
\toprule
\textbf{Ámbito} & \textbf{Sistemas} & \textbf{Característica principal} \\
\midrule
Normal (típico) &
Asegurativo probatorio; asegurativo garantizador; precautorio tradicional; asegurativo genérico. &
Están regulados en la ley; la pretensión cautelar \textbf{no} coincide con la pretensión sustancial. \\
\addlinespace
Excepcional &
Tutela anticipada. &
La pretensión cautelar se superpone total o parcialmente con la pretensión sustancial. \\
\bottomrule
\end{tabularx}

% TODO: verificar consistencia entre fuentes originales (una fuente incluye los sistemas protectores en el ámbito excepcional; otra menciona únicamente la tutela anticipada).

\section{Ámbito normal}

\subsection{Sistemas asegurativos probatorios}

Estos sistemas propenden a la \textbf{mayor eficacia del proceso}. Su finalidad es resguardar fuentes de prueba que pueden perderse con el transcurso del tiempo: si no se permite a la ley actuar en el proceso de este modo, hay prueba que se perdería.

\begin{normativa}{Artículo 326 --- CPCCN: prueba anticipada}
Regula la \textbf{prueba anticipada}. Es allí donde se encuentra regulado el sistema asegurativo probatorio. La ley actúa en el proceso con carácter asegurativo o conservatorio, permitiendo adelantar el momento en el cual se produce una prueba que, de no producirse en ese momento, se perdería al aguardar la etapa probatoria ordinaria del proceso.
\end{normativa}

Si en la etapa probatoria ordinaria se pierde una prueba porque no se permitió su adelantamiento, la parte que la tenía pierde la posibilidad de acreditar el fundamento de su pretensión, lo que afecta la igualdad de las partes en el proceso.

\subsection{Sistemas precautorios tradicionales}

Agrupan las medidas cautelares clásicas: el embargo, la anotación de litis, el secuestro, entre otras. Apuntan fundamentalmente a resguardar bienes.

\begin{important}
Su finalidad es resguardar el efectivo cumplimiento de la sentencia. Si no se traba un embargo para hacer efectiva la sentencia, esa sentencia solamente servirá para «colgarse en un cuadro», porque no podrá ejecutarse.
\end{important}

\begin{table}[H]
\centering
\begin{tabularx}{\textwidth}{>{\raggedright\arraybackslash}p{0.3\textwidth} Y}
\toprule
\textbf{Sistema} & \textbf{Finalidad} \\
\midrule
Asegurativo probatorio & Propender a la mayor eficacia del proceso (resguardar prueba) \\
Precautorio tradicional & Resguardar el efectivo cumplimiento de la sentencia (resguardar bienes) \\
\bottomrule
\end{tabularx}
\end{table}

\subsection{Sistemas asegurativos garantizadores}

Estos sistemas garantizan a la parte contraria, es decir, a quien puede ser afectado por la medida cautelar, frente a la posibilidad de que la medida sea revocada o no reconocida en la sentencia. El instrumento central es la \textbf{contracautela}.

Si un juez dicta una medida y luego se acredita que no fue correctamente dictada, el sistema asegurativo garantizador, es decir, la contracautela, resguarda los daños y perjuicios que la medida pueda haber ocasionado al afectado.

\begin{normativa}{Artículo 509 --- CPCCN: fianza para ejecutar}
Exige una fianza para poder seguir avanzando con la ejecución, a fin de resguardar los derechos de la parte contraria en caso de que la sentencia sea revocada.
\end{normativa}

\begin{normativa}{Artículo 555 --- CPCCN: juicio ejecutivo}
También exige fianza para ejecutar, por las mismas razones que el artículo 509: resguardar a quien soportó la resolución contraria en caso de reversión del proceso.
\end{normativa}

\subsection{Sistemas asegurativos genéricos}

Son el cuarto tipo dentro del ámbito normal.

\begin{normativa}{Artículo 228 --- CPCCN: inhibición general de bienes}
Es una medida residual que se otorga cuando no existen bienes para embargar. Pertenece al sistema asegurativo genérico dentro del ámbito normal de los sistemas cautelares.
\end{normativa}

\begin{normativa}{Artículo 230 --- CPCCN: prohibición de innovar}
Es la medida elegida como eje explicativo del sistema cautelar, porque permite el acceso a la tutela anticipada (ámbito excepcional) y porque presenta dos facetas: la medida de no innovar y la medida innovativa.
\end{normativa}

\begin{table}[H]
\centering
\small
\begin{tabularx}{\textwidth}{>{\raggedright\arraybackslash}p{0.2\textwidth} Y Y}
\toprule
\textbf{Faceta} & \textbf{Cuándo se aplica} & \textbf{Qué se pide} \\
\midrule
Medida de \textbf{no innovar} & Cuando el perjuicio proviene de la modificación del statu quo & Que no se innove; que se conserve la situación existente \\
Medida \textbf{innovativa} & Cuando el perjuicio proviene de mantener el statu quo & Que se innove; que se modifique la situación para evitar el daño irreparable \\
\bottomrule
\end{tabularx}
\end{table}

En ambos casos se conserva: la medida innovativa se da cuando el perjuicio va a provenir de mantener el statu quo; entonces se pide que se innove, porque si se mantiene la situación como está se genera el daño irreparable que necesita una respuesta inmediata.

\begin{normativa}{Artículo 232 --- CPCCN: medida cautelar genérica o innominada}
Contempla la posibilidad de peticionar medidas cautelares no previstas específicamente en la ley, cuando las circunstancias del caso lo requieren.
\end{normativa}

También se menciona la cautelar autónoma, proveniente del ámbito del contencioso-administrativo. Estos institutos se desarrollan en el capítulo~\ref{cap:generico}.

\section{Ámbito excepcional}

En el ámbito excepcional se encuentran la tutela anticipada y los sistemas protectores.

\begin{definition}{Tutela anticipada}
Pronunciamiento jurisdiccional de carácter provisorio a través del cual se adelanta, en todo o en parte, aquello que debió ser objeto de la sentencia de mérito, a fin de evitar un daño actual o inminente, o una situación afligente. A diferencia de las cautelares ordinarias, puede recaer sobre el mismo objeto de la sentencia definitiva.
\end{definition}

La tutela anticipada se desarrolla en el capítulo~\ref{cap:tutela}.
% TODO: verificar el desarrollo de los sistemas protectores, que las fuentes mencionan sin desarrollarlos en este punto.

\section{Panorama de los sistemas cautelares}

Los sistemas cautelares se desarrollan del siguiente modo:
\begin{itemize}
    \item \textbf{Parte general:} qué son, su tratamiento como sistema, los distintos tipos creados por el legislador, los presupuestos sustanciales y procesales para peticionar y obtener una medida cautelar, su trámite, las facultades del juez, los derechos del demandado y las características propias de las resoluciones dictadas en materia cautelar (capítulos~\ref{cap:fundamentos-cautelares} a~\ref{cap:caracteres}).
    \item \textbf{Sistemas precautorios tradicionales:} embargo, inventario, anotación de litis, secuestro y demás medidas del sistema precautorio tradicional reguladas en el Código Procesal (capítulos~\ref{cap:embargo} a~\ref{cap:proteccion}).
    \item \textbf{Sistema asegurativo probatorio:} es el que el legislador creó a través de la producción de prueba anticipada, regulada en el art.~326 CPCC.
    \item \textbf{Sistema asegurativo garantizador:} vinculado con la contracautela y la fianza (art.~258 CPCC).
    \item \textbf{Sistema asegurativo genérico, ámbito excepcional y tutela anticipada} (capítulos~\ref{cap:generico} y~\ref{cap:tutela}).
\end{itemize}
% TODO: contenido incompleto o ambiguo en las notas originales (el sistema asegurativo garantizador y el Art. 258 CPCC sólo figuran en el índice de una de las fuentes; no se desarrollan).

\cornelltitle{clasificación de los sistemas cautelares}
\cornellblock{
\cqrow{¿Por qué se estudian los sistemas cautelares como un «sistema»?}{Porque en distintos lugares del Código Procesal se verifican manifestaciones en las que la ley actúa con carácter conservatorio o asegurativo. Agruparlas en subsistemas permite comprenderlas en su conjunto, distinguirlas y entender el sistema del código procesal en esos aspectos.
}
\cqrow{¿Cómo se clasifican los sistemas cautelares según su ámbito?}{Ámbito normal: asegurativos probatorios, precautorios tradicionales, asegurativos garantizadores y asegurativos genéricos. Ámbito excepcional: tutela anticipada y sistemas protectores.
}
\cqrow{¿Qué resguardan los sistemas asegurativos probatorios?}{Propenden a la mayor eficacia del proceso, resguardando fuentes de prueba que pueden perderse con el tiempo. Se encuentran regulados en el art.~326 CPCCN (prueba anticipada), que permite adelantar la producción de una prueba que se perdería al aguardar la etapa probatoria ordinaria.
}
\cqrow{¿Cuál es la finalidad de los sistemas precautorios tradicionales?}{Resguardar el efectivo cumplimiento de la sentencia, fundamentalmente mediante el resguardo de bienes (embargo, anotación de litis, secuestro, entre otras). Sin embargo, la sentencia solo serviría para «colgarse en un cuadro».
}
\cqrow{¿Qué garantizan los sistemas asegurativos garantizadores?}{Garantizan a la parte afectada por la medida frente a la posibilidad de que sea revocada o no reconocida en la sentencia. El instrumento central es la contracautela; los arts.~509 y 555 CPCCN exigen fianza para ejecutar.
}
\cqrow{¿Qué comprenden los sistemas asegurativos genéricos?}{La inhibición general de bienes (art.~228, medida residual cuando no hay bienes para embargar), la prohibición de innovar (art.~230) y la medida cautelar genérica o innominada (art.~232). También se menciona la cautelar autónoma del ámbito contencioso-administrativo.
}
\cqrow{¿Cuáles son las dos facetas de la prohibición de innovar?}{Medida de no innovar: cuando el perjuicio proviene de modificar el statu quo (se pide conservar la situación). Medida innovativa: cuando el perjuicio proviene de mantener el statu quo (se pide modificarlo para evitar el daño irreparable).
}
}
\medskip
\cornellblocksum{
\cqrow{¿Qué es la tutela anticipada y por qué pertenece al ámbito excepcional?}{Es un pronunciamiento jurisdiccional provisorio que adelanta, en todo o en parte, lo que debió ser objeto de la sentencia de mérito, para evitar un daño actual o inminente o una situación afligente. A diferencia de las cautelares ordinarias, puede recaer sobre el mismo objeto de la sentencia definitiva.
}
}{los sistemas cautelares se agrupan en un ámbito normal (probatorios, precautorios tradicionales, garantizadores y genéricos) y un ámbito excepcional (tutela anticipada y sistemas protectores). Cada subsistema responde a una finalidad distinta: resguardar prueba, asegurar el cumplimiento de la sentencia, proteger al afectado por la medida o admitir medidas genéricas como la prohibición de innovar, que presenta una faceta de no innovar y otra innovativa.
}

\chapter[Medidas cautelares y presupuestos]{Medidas cautelares tradicionales y sus presupuestos}\label{cap:presupuestos}

\section{Concepto de medidas cautelares}

\begin{definition}{Medidas cautelares}
Actos procesales del órgano jurisdiccional, adoptados en el curso de un proceso o previos a él, para:
\begin{itemize}
    \item Asegurar bienes o pruebas.
    \item Mantener una determinada situación de hecho.
    \item Brindar seguridad o protección a una persona.
\end{itemize}
En todos los casos implican un \textbf{anticipo provisorio}, que no es definitivo, y que constituye uno de los caracteres esenciales de las medidas cautelares.
\end{definition}

\section{Presupuestos de las medidas cautelares}

Las medidas cautelares tienen presupuestos propios que determinan su procedencia. Estos presupuestos se exigen tanto para las medidas del sistema precautorio tradicional como para las del sistema asegurativo genérico.

\begin{table}[H]
\centering
\begin{tabularx}{\textwidth}{>{\raggedright\arraybackslash}p{0.25\textwidth} Y}
\toprule
\textbf{Tipo de presupuesto} & \textbf{Contenido} \\
\midrule
\textbf{Sustanciales} & Verosimilitud en el derecho y peligro en la demora; como consecuencia, la contracautela \\
\textbf{Procesales} & Enumerados en el art.~195 CPCCN: derecho a asegurar, medida solicitada, disposición legal aplicable y requisitos específicos de la medida \\
\bottomrule
\end{tabularx}
\end{table}

\section{Presupuestos sustanciales}

\subsection{Verosimilitud en el derecho (\textit{fumus boni iuris})}

El viejo adagio latino \textit{fumus boni iuris} se traduce como el «humo de buen derecho». Habla de la \textbf{apariencia}: es un juicio de valor que hace la jurisdicción, según el cual lo que aparenta es que la pretensión va a tener un buen resultado en la sentencia. Frente a esa apariencia, y frente al juicio de valor sobre el derecho que invoca quien peticiona la medida, el juez considera que está dado el presupuesto sustancial para admitir la cautelar. No tiene relación con la certeza a la que debe acceder el juez cuando dicta la sentencia definitiva.

\begin{important}
La verosimilitud del derecho \textbf{no es certeza}. Es un conocimiento inferior y provisional, que se basa en la apariencia de que la pretensión del peticionante está jurídicamente fundada. La certeza plena se exige solo para la sentencia definitiva.
\end{important}

\subsection{Peligro en la demora}

Es el temor fundado de que, con el transcurso del tiempo que insume un proceso judicial, el derecho que se invoca se torne estéril, es decir, que no se pueda dictar una sentencia eficaz. Se configura cuando existe un riesgo fundado de que, de no decretarse la cautelar, la sentencia que en su momento se dicte resulte ineficaz o de cumplimiento imposible.

\subsection{La contracautela}

\begin{definition}{Contracautela}
Garantía que debe otorgar quien pide la medida, para responder por los daños y perjuicios y las costas que pueda generar si la medida resulta luego rechazada en la acción de fondo o revocada por un tribunal superior.
\end{definition}

La contracautela guarda relación con los otros presupuestos y funciona como \textbf{vasos comunicantes}:

\begin{formula}{Contracautela como vasos comunicantes}
A mayor verosimilitud en el derecho y mayor peligro en la demora, menor será la contracautela. A menor verosimilitud en el derecho y menor peligro en la demora, mayor será la contracautela.
\end{formula}

\begin{table}[H]
\centering
\small
\begin{tabularx}{\textwidth}{>{\raggedright\arraybackslash}p{0.17\textwidth} Y Y}
\toprule
\textbf{Tipo} & \textbf{En qué consiste} & \textbf{Ejemplo} \\
\midrule
\textbf{Real} & Garantía sobre bienes concretos & Depósito en dinero, hipoteca sobre un bien, embargo sobre un inmueble \\
\textbf{Personal} & Fianza otorgada por una persona con reconocida capacidad económica & Fianza de un tercero solvente \\
\textbf{Juratoria} & Juramento del peticionante de responder por los daños & Supuestos de máxima verosimilitud en el derecho; es la caución menos exigente \\
\bottomrule
\end{tabularx}
\end{table}

\section{Presupuestos procesales: artículo 195 CPCCN}

\begin{normativa}{Artículo 195 --- CPCCN: oportunidad y presupuesto}
Las providencias cautelares podrán ser solicitadas antes o después de deducida la demanda, a menos que de la ley resultare que ésta debe entablarse previamente. El escrito deberá expresar: el derecho que se pretende asegurar, la medida que se pide, la disposición de la ley en que se funde, y el cumplimiento de los requisitos que correspondan en particular a la medida requerida.
\end{normativa}

Los cuatro presupuestos procesales que enumera el artículo 195 son:

\begin{enumerate}
    \item El derecho que se pretende asegurar.
    \item La medida que se pide (embargo, inhibición, secuestro, etc.).
    \item La disposición legal en que se funda.
    \item El cumplimiento de los requisitos que correspondan en particular a la medida requerida.
\end{enumerate}

\cornelltitle{presupuestos de las medidas cautelares}
\cornellblocksum{
\cqrow{¿Qué son las medidas cautelares?}{Actos procesales del órgano jurisdiccional, adoptados en el curso de un proceso o previos a él, para asegurar bienes o pruebas, mantener una determinada situación de hecho o brindar seguridad o protección a una persona. En todos los casos implican un anticipo provisorio, que no es definitivo.
}
\cqrow{¿Cuáles son los presupuestos de las medidas cautelares?}{Sustanciales: verosimilitud en el derecho y peligro en la demora, con la contracautela como consecuencia. Procesales: los enumerados en el art.~195 CPCCN.
}
\cqrow{¿Qué es el \textit{fumus boni iuris}?}{Es el «humo de buen derecho»: la apariencia de que la pretensión del peticionante tiene fundamento jurídico suficiente. Es un juicio de valor de la jurisdicción que no exige certeza plena, sino verosimilitud, y que habilita el dictado de la cautelar.
}
\cqrow{¿Qué es el peligro en la demora?}{Es el temor fundado de que, con el transcurso del tiempo que insume el proceso, el derecho invocado se torne estéril y la sentencia resulte ineficaz o de cumplimiento imposible.
}
\cqrow{¿Qué es la contracautela y de qué tipos puede ser?}{Es la garantía que debe otorgar quien pide la medida para responder por daños, perjuicios y costas si la medida es rechazada en la acción de fondo o revocada por un tribunal superior. Puede ser real (garantía sobre bienes concretos), personal (fianza de una persona con reconocida capacidad económica) o juratoria (juramento de responder por los daños; la menos exigente).
}
\cqrow{¿Cómo funciona la contracautela como vasos comunicantes?}{A mayor verosimilitud y mayor peligro en la demora, menor contracautela; a menor verosimilitud y menor peligro en la demora, mayor contracautela.
}
\cqrow{¿Qué debe expresar el escrito según el art.~195 CPCCN?}{1) El derecho que se pretende asegurar. 2) La medida que se pide. 3) La disposición legal en que se funda. 4) El cumplimiento de los requisitos particulares de la medida requerida.
}
}{las medidas cautelares son actos procesales que implican un anticipo provisorio. Proceden cuando concurren los presupuestos sustanciales (verosimilitud en el derecho y peligro en la demora), de los que se sigue la contracautela, que opera en forma inversa a la intensidad de aquellos, y los presupuestos procesales del art.~195 CPCCN.
}

\chapter[Caracteres y clasificación de las medidas cautelares]{Caracteres particulares y clasificación de las medidas cautelares}\label{cap:caracteres}

\section{Provisoriedad: artículo 202 CPCCN}

\begin{normativa}{Artículo 202 --- CPCCN: provisoriedad}
Las medidas cautelares son provisorias. Subsistirán mientras duren las circunstancias que las determinaron. En cualquier momento en que éstas cesaren, se podrá requerir su levantamiento.
\end{normativa}

Esta característica diferencia esencialmente a la cautela de la sentencia: cuando cambian las circunstancias que dieron origen al dictado de la medida cautelar, ésta puede modificarse o levantarse. La sentencia, en cambio, pone fin al proceso y no puede alterarse.

Lo que se decide en materia cautelar es \textbf{provisorio}: la medida puede dejarse sin efecto, modificarse, reducirse o ampliarse a lo largo del proceso, según se modifiquen o no las circunstancias de hecho y de derecho que existían al momento en que el juez la ordenó.

\begin{important}
No deben confundirse ambos planos: lo que se resuelve en la \textbf{sentencia} es \textbf{definitivo}; lo que se resuelve en materia \textbf{cautelar} es \textbf{provisorio}.
\end{important}

\section{Modificabilidad: artículo 203 CPCCN}

\begin{normativa}{Artículo 203 --- CPCCN: modificación}
El afectado por una medida cautelar puede pedir su sustitución por otra que le resulte menos gravosa, ofreciendo bienes que resulten igualmente suficientes. También puede peticionarse la ampliación o mejora de la medida.
\end{normativa}

\begin{example}{Sustitución de un embargo}
Si se embargó un inmueble y este embargo genera un daño innecesario al afectado, éste puede ofrecer otro bien a embargo que resulte igualmente suficiente al que fue embargado; la medida puede sustituirse para evitar el daño que genera la medida dictada.
\end{example}

\section{Facultades del juez: artículo 204 CPCCN}

\begin{normativa}{Artículo 204 --- CPCCN: facultades del juez}
El juez, para evitar perjuicios o gravámenes innecesarios al titular de los bienes, podrá disponer una medida precautoria distinta de la solicitada, o limitarla, teniendo en cuenta la importancia del derecho que se intentare proteger.
\end{normativa}

El juez tiene plenas facultades para modificar o limitar la medida cautelar solicitada, siempre resguardando el derecho que se intenta proteger. Esto también distingue a la cautelar de la sentencia definitiva.

\section{Caducidad: artículo 207 CPCCN}

\begin{normativa}{Artículo 207 --- CPCCN: caducidad de las medidas cautelares}
Cuando las medidas cautelares se hayan peticionado antes del inicio del proceso principal, a los 10 días siguientes de su traba (o de su inscripción registral o notificación, según corresponda), debe promoverse la demanda. La inhibición general de bienes y el embargo se extinguen a los 5 años desde la fecha de inscripción en el registro respectivo. Para evitar la extinción, debe pedirse la reinscripción de la medida antes del vencimiento de ese plazo, lo que genera un nuevo período de 5 años.
\end{normativa}

Mecanismo de reinscripción: antes de que venzan los 5 años de la inscripción original, se solicita la reinscripción; el oficio correspondiente se libra y se inscribe en el registro, generando un nuevo plazo de 5 años desde esa nueva inscripción.

\section{Tramitación \textit{inaudita parte} y recursos: artículo 198 CPCCN}

\begin{normativa}{Artículo 198 --- CPCCN: cumplimiento de la medida y recursos}
Las medidas precautorias se decretarán y cumplirán sin audiencia de la otra parte (\textit{inaudita parte}). Ningún incidente planteado por el destinatario de la medida podrá detener su cumplimiento. Si el afectado no hubiese tomado conocimiento de las medidas con motivo de su ejecución, se les notificarán personalmente o por cédula dentro de los 3 días. Quien hubiese obtenido la medida será responsable de los perjuicios que generare la demora en notificarla. La providencia que admitiere o denegare una medida cautelar será recurrible por vía de reposición; también será admisible la apelación subsidiaria o directa. El recurso de apelación, en caso de admitirse la medida, será concedido con efecto devolutivo (no suspensivo).
\end{normativa}

La medida cautelar tramita \textit{inaudita parte}, lo que también significa que no se está ante una hipótesis de condena, sino de cautela. Es obligación de quien obtuvo la cautelar notificarla dentro del plazo de 3 días de trabada, bajo responsabilidad por los perjuicios que genere la demora.

\begin{important}
El recurso de apelación en materia cautelar se concede con efecto \textbf{devolutivo} (no suspensivo): la vigencia de la medida cautelar continúa mientras se sustancia el recurso. Aquí radica la importancia de la contracautela: si el tribunal superior revoca la resolución, la contracautela permite responder por los daños durante el tiempo en que la medida estuvo vigente (art.~199 CPCCN).
\end{important}

\section{Clasificación sobre bienes, entes ideales y personas}

Para hacer más didáctico el estudio, las medidas cautelares se clasifican según el objeto sobre el que recaen. Esta clasificación es meramente didáctica y el detalle de cada medida se desarrolla en los capítulos sobre medidas cautelares en particular.

\begin{table}[H]
\centering
\small
\begin{tabularx}{\textwidth}{>{\raggedright\arraybackslash}p{0.2\textwidth} Y Y}
\toprule
\textbf{Objeto} & \textbf{Grado / tipo} & \textbf{Medida específica} \\
\midrule
\textbf{Bienes} & Individualización del bien & Inventario / anotación de litis \\
\textbf{Bienes} & Indisposición del bien & Embargo / prohibición de contratar \\
\textbf{Bienes} & Desapoderamiento del bien & Secuestro / interventor recaudador \\
\textbf{Entes ideales} (empresas / sociedades) & Interferencia simple & Veedores, interventores informativos \\
\textbf{Entes ideales} & Intervención compleja & Administración judicial de la sociedad \\
\textbf{Personas} & Sistema genérico / residual & Guarda de personas (art.~234 CPCCN) \\
\textbf{Personas} & Protección específica & Protección contra violencia familiar (Ley 24.417) \\
\bottomrule
\end{tabularx}
\end{table}

\begin{normativa}{Artículo 234 --- CPCCN: guarda de personas}
Regula la guarda de personas como medida cautelar residual sobre personas. Se aplica cuando las circunstancias del caso lo requieren.
\end{normativa}

\begin{normativa}{Ley 24.417 --- Protección contra la Violencia Familiar}
Esta ley, de aplicación en el ámbito de Capital Federal, prevé medidas cautelares específicas de protección para personas en situación de violencia doméstica.
\end{normativa}

\cornelltitle{caracteres y clasificación de las medidas cautelares}
\cornellblock{
\cqrow{¿Qué establece el art.~202 CPCCN sobre la provisoriedad?}{Las medidas cautelares son provisorias: subsisten mientras duren las circunstancias que las determinaron y, cuando estas cesan, se puede requerir su levantamiento. Esto las diferencia de la sentencia, que pone fin al proceso y no puede alterarse.
}
\cqrow{¿Qué permite el art.~203 CPCCN (modificabilidad)?}{Que el afectado pida la sustitución de la medida por otra menos gravosa, ofreciendo bienes igualmente suficientes, y que se peticione la ampliación o mejora de la medida.
}
\cqrow{¿Qué facultades otorga el art.~204 CPCCN al juez?}{Para evitar perjuicios o gravámenes innecesarios al titular de los bienes, puede disponer una medida distinta de la solicitada o limitarla, teniendo en cuenta la importancia del derecho que se intenta proteger.
}
\cqrow{¿Qué establece el art.~207 CPCCN sobre la caducidad?}{Si la cautelar se peticionó antes del inicio del proceso principal, la demanda debe promoverse a los 10 días siguientes de su traba (o de su inscripción registral o notificación, según corresponda). La inhibición general de bienes y el embargo se extinguen a los 5 años desde su inscripción registral, salvo que se pida la reinscripción antes del vencimiento, lo que genera un nuevo período de 5 años.
}
\cqrow{¿Qué significa que la cautelar tramita \textit{inaudita parte} (art.~198 CPCCN)?}{Que se decreta y cumple sin audiencia de la otra parte, y ningún incidente del destinatario puede detener su cumplimiento. Si el afectado no tomó conocimiento por la ejecución, se le notifica personalmente o por cédula dentro de los 3 días, y quien obtuvo la medida responde por los perjuicios que cause la demora en notificarla.
}
\cqrow{¿Cómo se recurre la resolución cautelar y qué efecto tiene la apelación?}{La providencia que admite o deniega la medida es recurrible por reposición, con apelación subsidiaria o directa. Si se admite la medida, la apelación se concede con efecto devolutivo (no suspensivo): la medida sigue vigente mientras se sustancia el recurso, lo que explica la importancia de la contracautela.
}
}
\medskip
\cornellblocksum{
\cqrow{¿Cómo se clasifican las medidas cautelares según su objeto?}{Sobre bienes (individualización: inventario / anotación de litis; indisposición: embargo / prohibición de contratar; desapoderamiento: secuestro / interventor recaudador). Sobre entes ideales (interferencia simple: veedores e interventores informativos; intervención compleja: administración judicial de la sociedad). Sobre personas (guarda de personas, art.~234 CPCCN; protección contra violencia familiar, Ley 24.417).
}
}{las medidas cautelares son provisorias, modificables y sujetas a caducidad; tramitan \textit{inaudita parte} y su apelación se concede con efecto devolutivo, de modo que la contracautela cobra importancia. Estos caracteres las distinguen de la sentencia definitiva. Además, se clasifican didácticamente según recaigan sobre bienes, entes ideales o personas, y cada medida se estudia en detalle en los capítulos sobre medidas cautelares en particular.
}

\chapter[Sistema asegurativo genérico]{Sistema asegurativo genérico (ámbito normal)}\label{cap:generico}

\section{Caracterización general}

El sistema asegurativo genérico comprende cuatro medidas:

\begin{enumerate}
    \item \textbf{Inhibición general de bienes} (Art.~228 CPCC).
    \item \textbf{Medida cautelar genérica} (Art.~232 CPCC).
    \item \textbf{Prohibición de innovar}, que incluye una faceta innovativa y otra no innovativa (Art.~230 CPCC).
    \item \textbf{Cautela autónoma}, creación pretoriana luego regulada por la Ley 26.854.
\end{enumerate}

\begin{important}
\textbf{Carácter residual o subsidiario.} Las medidas de este sistema sólo prosperan cuando las medidas del sistema precautorio tradicional (embargo, secuestro, etc.) no resultan aptas para la situación de hecho y de derecho que se pretende proteger en el caso concreto.
\end{important}

\section{Inhibición general de bienes (Art.~228 CPCC)}

\begin{quote}
\small
\textbf{Artículo 228 --- Código Procesal Civil y Comercial de la Nación.} \textit{Inhibición general de bienes. En todos los casos en que habiendo lugar a embargo éste no pudiere hacerse efectivo por no conocerse bienes del deudor, o por no cubrir éstos el importe del crédito reclamado, podrá solicitarse contra aquél la inhibición general de enajenar y gravar bienes, la que se deberá dejar sin efecto siempre que presentare a embargo bienes suficientes o diere caución bastante. El que solicitare la inhibición deberá expresar el nombre, apellido y domicilio del deudor, así como todo otro dato que pueda individualizar al inhibido, sin perjuicio de los demás requisitos que impongan las leyes. La inhibición sólo producirá efecto desde la fecha de su anotación salvo para los casos en que el dominio se hubiere transmitido con anterioridad, de acuerdo con lo dispuesto en la legislación general. No concederá preferencia sobre las anotadas con posterioridad.}
\end{quote}

\subsection{Definición y finalidad}

\begin{definition}{Inhibición general de bienes}
Medida cautelar genérica, regulada en un solo artículo (el 228 CPCC), mediante la cual se pretende que el demandado o afectado por la medida \textbf{no venda ni grave} bienes que formen parte de su patrimonio.
\end{definition}

\subsection{Carácter residual: cuándo procede}

La inhibición general de bienes \textbf{no es una medida autónoma o principal} como el embargo, sino residual. Procede en los mismos casos en que procede el embargo, pero únicamente cuando se da alguno de estos supuestos:

\begin{enumerate}
    \item No existen bienes para embargar (el deudor no tiene bienes identificados).
    \item Existen bienes para embargar, pero no resultan suficientes para responder o garantizar el crédito que se pretende asegurar.
\end{enumerate}

\begin{example}{Embargo insuficiente}
Si se debe asegurar una suma equivalente a un millón de pesos y sólo pudieron embargarse quinientos mil pesos en una cuenta bancaria, sin otros bienes para embargar, queda por garantizar el cumplimiento de los quinientos mil pesos restantes. La inhibición general de bienes puede ir junto con el embargo de manera residual, porque los bienes embargados no resultan suficientes.
\end{example}

\subsection{Diferencias con el embargo}

\medskip
\renewcommand{\arraystretch}{1.3}
\noindent\begin{tabularx}{\textwidth}{@{}>{\raggedright\arraybackslash}p{2.8cm} Y Y@{}}
\toprule
\textbf{Aspecto} & \textbf{Embargo} & \textbf{Inhibición general de bienes} \\
\midrule
Determinación &
\textbf{Determinado:} recae sobre un bien identificado y particular. &
\textbf{Indeterminada:} afecta a todos los bienes presentes y futuros del patrimonio del deudor. \\
\addlinespace
Tipo de bienes &
Todo tipo de bienes: muebles, inmuebles, registrables y no registrables. &
Sólo bienes \textbf{registrables} (porque debe inscribirse para producir efectos). \\
\addlinespace
Preferencia temporal &
El primer embargante tiene derecho de preferencia sobre los posteriores. &
No existe derecho de preferencia; no rige el principio de prioridad temporal. \\
\addlinespace
Habilita ejecución forzada &
Sí: permite avanzar hacia el remate y cobro del crédito. &
No: no habilita la ejecución forzada porque no hay bienes identificados de los que cobrar. \\
\addlinespace
Extinción &
5 años desde la inscripción. &
5 años desde la inscripción (igual que el embargo). \\
\addlinespace
Caducidad si es previa al proceso &
10 días para iniciar el proceso principal. &
10 días para iniciar el proceso principal (igual que el embargo). \\
\bottomrule
\end{tabularx}
\renewcommand{\arraystretch}{1}

\subsection{Alcance temporal y efectos}

La inhibición produce efectos desde su inscripción en adelante. Es posible que se desconozcan los bienes que tiene el deudor o demandado, o que éste no tenga bienes al momento de la medida y luego ingresen bienes a su patrimonio, que quedarán alcanzados por la inhibición.

Tanto el embargo como la inhibición se extinguen a los \textbf{cinco años} de su inscripción. Para mantener la validez de la medida, antes de su vencimiento debe peticionarse al juez la \textbf{reinscripción}, a fin de evitar su extinción.

\subsection{Procedimiento}

\begin{enumerate}
    \item Se peticiona que se dicte la inhibición general de bienes en determinados registros, acreditando los presupuestos cautelares.
    \item El juez ordena y decreta la inhibición, identificando al demandado por nombre completo y número de DNI.
    \item El letrado confecciona un oficio dirigido al registro correspondiente (por ejemplo, el Registro de la Propiedad Inmueble), que se deja en el juzgado para su revisión y firma.
    \item El letrado diligencia el oficio ingresándolo al registro para que se inscriba la inhibición.
\end{enumerate}

Suele pedirse informes en los registros pertinentes para conocer qué bienes tiene el demandado a su nombre y con su DNI: el Registro de la Propiedad Inmueble de la Ciudad de Buenos Aires y de la Provincia de Buenos Aires, y el Registro de la Propiedad Automotor, de alcance nacional.

\subsection{Cese de la inhibición}

La inhibición puede levantarse antes de los cinco años en los siguientes casos:

\begin{enumerate}
    \item Si el deudor ofrece un bien a embargo que sea suficiente para garantizar el crédito.
    \item Si el deudor ofrece una garantía suficiente (en el ámbito de una medida cautelar).
\end{enumerate}

El pedido de levantamiento por sustitución, formulado por el deudor, se traslada al actor o acreedor, y luego el juez resuelve; fundamentalmente porque el bien ofrecido a embargo debe ser suficiente para garantizar el crédito en cuestión.

\section{Medida cautelar genérica o innominada (Art.~232 CPCC)}

\begin{quote}
\small
\textbf{Artículo 232 --- Código Procesal Civil y Comercial de la Nación.} \textit{Fuera de los casos previstos en los artículos precedentes, quien tuviere motivos para temer que durante el tiempo anterior al reconocimiento judicial de su derecho, éste pudiera sufrir un perjuicio inminente o irreparable, podrá solicitar las medidas urgentes que, según las circunstancias, fuesen más aptas para asegurar provisionalmente el cumplimiento de la sentencia.}
\end{quote}

La medida cautelar genérica establece que, fuera de los casos previstos en los artículos precedentes (todas las medidas cautelares ya estudiadas), quien tenga motivos para temer que, antes de la sentencia, su derecho pueda sufrir un perjuicio inminente o irreparable, puede solicitar las medidas urgentes que según las circunstancias resulten más aptas para asegurar provisionalmente el cumplimiento de la sentencia.

Se la llama \textbf{genérica} o \textbf{innominada} porque habilita a la jurisdicción a dictar cualquier medida, ya sea sola o aislada, o en complemento de las tradicionales, que mejor responda a las situaciones de hecho y de derecho que se pretende asegurar y proteger antes de la sentencia.

\begin{important}
La medida cautelar genérica rige para asegurar cualquier tipo de bienes, personas o derechos. Es la que más amplía las posibilidades de peticionar una medida cautelar cuando las reguladas anteriormente no resultan adecuadas y suficientes.
\end{important}

\section{Prohibición de innovar (Art.~230 CPCC)}

\subsection{La prohibición de innovar como medida cautelar matriz}

Se entiende que la prohibición de innovar es un sistema asegurativo \textbf{matriz} de los restantes, la ``medida cautelar madre'' de todas las demás. Muchos autores atribuyen ese carácter al embargo, pero se considera más adecuado afirmar que el sistema matriz es la prohibición de innovar, porque permite cautelar no sólo bienes, sino también personas y derechos, según los hechos del caso concreto.

\begin{quote}
\small
\textbf{Artículo 230 --- Código Procesal Civil y Comercial de la Nación.} \textit{Podrá decretarse la prohibición de innovar en toda clase de juicios siempre que: 1. El derecho resulte verosímil. 2. Existiere peligro de que si se mantuviera o alterara, en su caso, la situación de hecho o de derecho, la modificación pudiera influir en la sentencia o convirtiera su ejecución en ineficaz. 3. La cautela no pudiere obtenerse por medio de otra medida precautoria.}
\end{quote}

\subsection{Las dos facetas de la prohibición de innovar}

El inciso 2 del Art.~230 revela las dos facetas de esta medida. Algunos autores las consideran dos vías distintas (medida innovativa y de no innovar); sin embargo, según cómo las reguló el legislador, ambas son facetas de la prohibición de innovar.

\medskip
\renewcommand{\arraystretch}{1.3}
\noindent\begin{tabularx}{\textwidth}{@{}>{\raggedright\arraybackslash}p{2.6cm} Y Y@{}}
\toprule
\textbf{Faceta} & \textbf{Cuándo procede} & \textbf{Qué ordena el juez} \\
\midrule
Innovativa (medida innovativa) &
El perjuicio proviene de \textbf{mantener} el statu quo existente. &
Ordena \textbf{alterar} la situación de hecho o de derecho. \\
\addlinespace
No innovativa (medida de no innovar) &
El perjuicio proviene de \textbf{alterar} el statu quo existente. &
Ordena \textbf{mantener} la situación de hecho o de derecho. \\
\bottomrule
\end{tabularx}
\renewcommand{\arraystretch}{1}

\subsection{Los dos ámbitos en que opera}

Además de sus dos facetas, la prohibición de innovar puede operar en dos ámbitos distintos:

\begin{enumerate}
    \item \textbf{Ámbito normal:} cuando la pretensión cautelar \textbf{no} se superpone con la pretensión sustancial del proceso principal.
    \item \textbf{Ámbito excepcional:} cuando la pretensión cautelar \textbf{se superpone} total o parcialmente con la pretensión sustancial; esto configura la \textbf{tutela anticipada}.
\end{enumerate}

\section{Cautela autónoma}

\subsection{Ubicación en el sistema y origen}

La cautela autónoma es la cuarta medida del sistema asegurativo genérico. Es propia del derecho administrativo y de la justicia en lo contencioso-administrativo; no está regulada en el Código Procesal Civil y Comercial de la Nación. En un primer momento no tuvo regulación y fue una \textbf{creación pretoriana} jurisprudencial; luego recibió una regulación legal bastante cuestionable.

\subsection{En qué consiste}

\begin{definition}{Cautela autónoma}
Petición que el justiciable hace al juez para que \textbf{suspenda los efectos de un acto administrativo} hasta que se resuelva el recurso administrativo interpuesto contra ese acto y que, de algún modo, habilitó la instancia judicial, es decir, hasta que se agote la vía administrativa.
\end{definition}

\subsection{La presunción de legitimidad (Art.~12 Ley 19.549)}

\begin{quote}
\small
\textbf{Artículo 12 --- Ley 19.549 (Procedimiento Administrativo).} \textit{El acto administrativo goza de presunción de legitimidad; su fuerza ejecutoria faculta a la Administración a ponerlo en práctica por sus propios medios, salvo que la ley o la naturaleza del acto exigieren la intervención judicial, e impide que los recursos que interpongan los administrados suspendan su ejecución y efectos, salvo que una norma expresa establezca lo contrario. Sin embargo, la Administración podrá, de oficio o a pedido de parte y mediante resolución fundada, suspender la ejecución por razones de interés público, o para evitar perjuicios graves al interesado.}
\end{quote}

El acto administrativo goza de presunción de legitimidad. Por ello, aunque el administrado interponga un recurso administrativo para que el acto sea revisado, ese recurso no suspende los efectos del acto, que puede ser ejecutado. Esto se vincula con los efectos suspensivos y no suspensivos de los recursos: el recurso administrativo no tiene efecto suspensivo.

Esta situación puede generar un daño o perjuicio irreparable en el administrado; por ejemplo, si el acto ordenó abonar una suma o multa de importancia trascendente. En ese caso, estando pendiente de resolución el recurso administrativo, el administrado puede acudir a la justicia y peticionar la cautela autónoma, es decir, que el juez disponga una medida que suspenda los efectos del acto administrativo hasta que se resuelva el recurso.

\subsection{Por qué se llama ``autónoma''}

Se llama autónoma porque se trata de un procedimiento administrativo y la medida es autónoma respecto de él: no hay un proceso judicial en curso. Es instrumental y accesoria, porque se vincula con la suspensión de los efectos del acto administrativo.

\subsection{Creación pretoriana y trámite original}

Originalmente se promovían amparos para obtener este tipo de cautelas, hasta que la Corte resolvió que el amparo no era la vía idónea, y entonces se creó la cautela autónoma. El trámite original se encuadró en el Art.~230 CPCC (prohibición de innovar) y los presupuestos requeridos eran:

\begin{enumerate}
    \item Verosimilitud del derecho.
    \item Peligro en la demora.
\end{enumerate}

El carácter de la medida era \textbf{más excepcional} que el de las medidas ordinarias, dado que el acto administrativo goza de presunción de legitimidad.

\subsection{La Ley 26.854: medidas cautelares contra el Estado}

Lo estudiado hasta aquí sobre medidas cautelares no rige cuando el Estado está del otro lado. En ese caso debe tenerse en cuenta la \textbf{Ley 26.854}, que regula cuestiones particulares, distintas de las ya vistas, cuando el justiciable pretende una medida cautelar contra el Estado.

\begin{important}
\textbf{Artículo 13 --- Ley 26.854 (suspensión de efectos de un acto estatal).} Regula la cautela autónoma en el fuero contencioso-administrativo. Para peticionar la suspensión de los efectos de un acto estatal ya no bastan la verosimilitud del derecho y el peligro en la demora. Se exige:
\begin{enumerate}
    \item Acreditar sumariamente perjuicios graves que impliquen imposible reparación posterior.
    \item Que tramite \textbf{bilateralizada}: el juez debe convocar previamente a la entidad estatal a presentar un informe.
    \item Que el administrado haya solicitado previamente en sede administrativa la suspensión del acto; sólo ante la negativa o el silencio de cinco días queda habilitada la vía judicial.
\end{enumerate}
\end{important}

En cuanto al trámite, según el Art.~4 de la ley, cualquier medida cautelar peticionada contra el Estado deja de tramitar \textbf{inaudita parte}: siempre, antes de resolver, el juez debe convocar a la entidad estatal, que presenta un informe pormenorizado de la cuestión peticionada, y luego resuelve.

Esta regulación podría ser atacada por ciertas inconstitucionalidades y, al igual que la regulación de la tutela anticipada de urgencia que se analiza más adelante, muchas veces se legisla en contra del justiciable.

\cornelltitle{sistema asegurativo genérico}
\cornellblock{
\cqrow{¿Qué es el carácter residual o subsidiario del sistema asegurativo genérico?}{Las medidas del sistema (inhibición general de bienes, medida cautelar genérica, prohibición de innovar) sólo proceden cuando las medidas del sistema precautorio tradicional (embargo, secuestro, etc.) no resultan aptas para proteger la situación de hecho y de derecho del caso concreto.
}
\cqrow{¿Cuándo procede la inhibición general de bienes?}{En los mismos casos en que procede el embargo, pero sólo cuando no existen bienes para embargar o cuando los bienes embargables no alcanzan para garantizar el crédito. Busca que el demandado no venda ni grave bienes de su patrimonio.
}
\cqrow{¿En qué se diferencia la inhibición del embargo?}{La inhibición es indeterminada (afecta bienes presentes y futuros), alcanza sólo bienes registrables, no otorga preferencia temporal y no habilita la ejecución forzada. El embargo es determinado, recae sobre cualquier tipo de bien, otorga preferencia al primer embargante y habilita el remate. Ambos se extinguen a los 5 años de su inscripción y, si son previos al proceso, caducan a los 10 días si no se inicia el proceso principal.
}
\cqrow{¿Cómo cesa o se mantiene la inhibición?}{Se levanta antes de los 5 años si el deudor ofrece un bien suficiente a embargo o una garantía suficiente; se corre traslado al acreedor y el juez resuelve. Para mantenerla debe reinscribirse antes de su vencimiento.
}
\cqrow{¿Qué es la medida cautelar genérica o innominada?}{Es la del Art.~232 CPCC: permite solicitar las medidas urgentes que según las circunstancias sean más aptas para asegurar provisionalmente el cumplimiento de la sentencia, fuera de los casos previstos en los artículos anteriores, ya sea de forma aislada o complementando las tradicionales. Sirve para cautelar bienes, personas o derechos.
}
\cqrow{¿Cuáles son las dos facetas de la prohibición de innovar?}{Innovativa: el perjuicio proviene de mantener el statu quo, por lo que el juez ordena alterarlo. No innovativa: el perjuicio proviene de alterarlo, por lo que el juez ordena mantenerlo. Ambas surgen del inciso 2 del Art.~230 CPCC.
}
\cqrow{¿Por qué se considera a la prohibición de innovar la medida cautelar matriz?}{Porque permite cautelar no sólo bienes sino también personas y derechos, según los hechos del caso concreto.
}
}
\medskip
\cornellblocksum{
\cqrow{¿Qué es la cautela autónoma y por qué surge?}{Es la petición al juez de suspender los efectos de un acto administrativo hasta que se resuelva el recurso administrativo interpuesto. Surge porque, según el Art.~12 de la Ley 19.549, el acto goza de presunción de legitimidad y los recursos no suspenden su ejecución, lo que puede causar un perjuicio irreparable. Fue una creación pretoriana (antes se usaba el amparo, que la Corte consideró no idóneo).
}
\cqrow{¿Qué cambió con la Ley 26.854?}{Para las medidas contra el Estado, el trámite deja de ser inaudita parte (el juez convoca a la entidad, que presenta un informe) y, para suspender un acto estatal (Art.~13), se exige acreditar perjuicios graves de imposible reparación posterior y haber solicitado antes la suspensión en sede administrativa, con negativa o silencio de 5 días.
}
}{El sistema asegurativo genérico es residual y subsidiario: comprende la inhibición general de bienes (Art.~228, indeterminada, sobre bienes registrables, sin preferencia ni ejecución forzada, por 5 años), la medida cautelar genérica (Art.~232, la medida más apta según las circunstancias), la prohibición de innovar (Art.~230, con facetas innovativa y no innovativa, medida matriz) y la cautela autónoma (suspensión de efectos de actos administrativos, hoy condicionada por la Ley 26.854 con trámite bilateralizado y requisitos agravados).
}

\chapter{Ámbito excepcional: tutela anticipada}\label{cap:tutela}

\section{Qué determina que un ámbito sea excepcional}

El ámbito excepcional se da cuando, a través de la petición y concesión de la cautela, se adelanta en todo o en parte aquello que luego será decidido en la sentencia de mérito. La pretensión cautelar se superpone total o parcialmente con la pretensión sustancial (la invocada en la demanda, que es el objeto del proceso). Lo que se peticiona o se concede cautelarmente es total o parcialmente igual a lo que se pretende que se reconozca al dictar sentencia.

\begin{definition}{Tutela anticipada}
Sistema cautelar en virtud del cual la jurisdicción dicta una medida para proteger y asegurar anticipadamente (de modo \textbf{provisorio}), total o parcialmente, aquello que después se decidirá de modo \textbf{definitivo} al dictar sentencia.
\end{definition}

\section{¿El juez no estaría prejuzgando?}

Cabe preguntarse si es válido este sistema o si el juez estaría prejuzgando al otorgarlo antes de tiempo. La clave está en no confundir el carácter \textbf{provisorio} de lo que se decide cautelarmente con el carácter \textbf{definitivo} de lo que se decide en la sentencia.

La superposición es posible porque la tutela es provisoria: para otorgarla el juez sólo exige que se acredite verosimilitud en el derecho y, en virtud del peligro en la demora y de la urgencia, protege y asegura. Luego, en el debate que se desarrolla a lo largo del proceso (defensa del afectado por la medida, pruebas que se produzcan), al dictar sentencia el juez puede resolver en el mismo sentido, pero ya con carácter definitivo, o en sentido contrario, y entonces queda sin efecto lo que se concedió provisionalmente.

\medskip
\renewcommand{\arraystretch}{1.3}
\noindent\begin{tabularx}{\textwidth}{@{}>{\raggedright\arraybackslash}p{2.6cm} Y Y Y@{}}
\toprule
\textbf{Momento} & \textbf{Decisión} & \textbf{Estándar de prueba} & \textbf{Carácter} \\
\midrule
Al otorgar la tutela anticipada &
Protege provisionalmente lo que será luego objeto de sentencia. &
Verosimilitud en el derecho y peligro en la demora. &
\textbf{Provisorio:} puede modificarse. \\
\addlinespace
Al dictar sentencia definitiva &
Decide con pleno conocimiento del caso. &
Certeza, basada en la prueba producida. &
\textbf{Definitivo:} cosa juzgada. \\
\bottomrule
\end{tabularx}
\renewcommand{\arraystretch}{1}

\medskip
Cuando el juez otorga la tutela anticipada de modo provisorio hay \textbf{cautela}; cuando decide en la sentencia definitiva sobre esa misma pretensión hay \textbf{condena}.

\section{Vías de materialización de la tutela anticipada}

La tutela anticipada se presenta a través de dos vías principales:

\begin{enumerate}
    \item \textbf{Sistemas cautelares (medidas cautelares):} mediante la prohibición de innovar (Art.~230 CPCC), la medida cautelar genérica (Art.~232 CPCC) o cualquier medida cautelar. El juez actúa dentro del ámbito excepcional por la superposición de pretensiones.
    \item \textbf{Subsistemas expresamente creados por el legislador:} normas sustanciales o procesales que regulan explícitamente la anticipación de la tutela para situaciones específicas.
\end{enumerate}

\section{Vías legales específicas de tutela anticipada}

\subsection{Alimentos provisorios (Art.~544 CCyC)}

\begin{quote}
\small
\textbf{Artículo 544 --- Código Civil y Comercial de la Nación.} \textit{Desde el principio de la causa o en el transcurso de ella el juez puede decretar la prestación de alimentos provisionales.}
\end{quote}

\begin{example}{Juicio de alimentos}
Una madre, en representación de sus hijos menores, demanda a su cónyuge, padre de los menores, para que se fije una cuota alimentaria a favor de los hijos, porque el progenitor demandado no abona ninguna suma por ese concepto. Dado el tiempo que demoran estos procesos y el perjuicio que supondría que durante años los menores no contaran con cuota alguna, el legislador creó el sistema de tutela anticipada del Art.~544 CCyC, denominado \textbf{alimentos provisorios}.

El juez puede otorgar de modo provisional un alimento a lo largo del proceso. Esa pretensión es igual, de algún modo, a la pretensión sustancial que el juez debe decidir en la sentencia; sin embargo, es plenamente viable y válida porque los alimentos provisionales son provisorios. Al dictar sentencia, luego de conocer las pruebas de las partes, el juez podrá fijar un alimento similar al provisional, uno considerablemente mayor o menor, o rechazar la demanda.
\end{example}

\subsection{Entrega anticipada del inmueble en el desalojo (Art.~680 bis CPCC)}

\begin{quote}
\small
\textbf{Artículo 680 bis --- Código Procesal Civil y Comercial de la Nación.} \textit{En los casos en que la acción de desalojo se dirija contra el intruso, en cualquier estado del juicio, después de trabada la litis, el juez puede imponer la inmediata entrega del inmueble, si el derecho invocado fuere verosímil, previa caución brindada por el actor por los eventuales daños y perjuicios que se pudieran irrogar.}
\end{quote}

\begin{example}{Proceso de desalojo contra el intruso}
La pretensión sustancial principal en el proceso de desalojo es la desocupación y entrega del inmueble. El Art.~680 bis permite que se adelante cautelarmente aquello que también es objeto de la sentencia de mérito: la desocupación y entrega anticipada del inmueble. Es provisorio, porque luego el juez tomará una decisión definitiva al dictar sentencia.

Por eso el Art.~680 bis exige verosimilitud en el derecho y que se conceda una caución (contracautela) para responder por los eventuales daños y perjuicios que pudiera generar al demandado una sentencia posterior en contra.
\end{example}

\section{Jurisprudencia de la Corte Suprema}

\subsection{Camacho Acosta (CSJN 320:1633): faceta innovativa}

\textbf{Hechos del caso.} El señor Camacho Acosta era un trabajador que prestaba servicios en una empresa. A raíz de un accidente de trabajo le amputaron el brazo izquierdo. En ese momento no estaba vigente ninguna póliza de aseguradora de riesgos de trabajo, por lo que promovió contra la empresa un juicio por daños y perjuicios reclamando todos los daños provenientes del accidente.

\textbf{Pretensión cautelar.} Camacho Acosta necesitaba colocar en su brazo izquierdo la prótesis con cierta cercanía en el tiempo a la amputación, porque si esperaba a la sentencia correría el riesgo de que la prótesis ya no pudiera ser colocada (por rechazo anatómico) y, además, no contar con ella le dificultaba reinsertarse en el mercado laboral. Peticionó una medida \textbf{innovativa} dentro del Art.~230 CPCC: que se le otorgara el dinero para acceder a la compra de la prótesis.

\textbf{Superposición de pretensiones.} El otorgamiento del dinero para la prótesis se superponía \emph{parcialmente} con la pretensión sustancial (la indemnización integral), lo que ubica el caso en el ámbito excepcional: tutela anticipada.

\textbf{Resolución.} El juez de primera instancia y la Sala J de la Cámara Civil rechazaron el pedido. Al llegar el expediente a la Corte por recurso extraordinario, la CSJN hizo lugar a la medida innovativa, reconociendo la tutela anticipada, al reunirse en el caso verosimilitud en el derecho, peligro en la demora y contracautela, y para evitar el daño irreparable que podía producirse si recién al dictarse la sentencia pudiera acceder a la prótesis.

\begin{important}
Lo único que agrega la Corte respecto de la tutela anticipada es que en estos casos se requiere \textbf{mayor prudencia} de los jueces al valorar los presupuestos y los hechos del caso concreto. Queda librado al juez, según su interpretación y valoración de los hechos, otorgar o no esta tutela judicial efectiva.
\end{important}

\subsubsection{Inversión del contradictorio}

Lo típico de estos sistemas es que se \textbf{invierte el contradictorio}: la jurisdicción primero protege y luego conoce. Se tutela primero para evitar el daño irreparable, y luego viene la etapa del contradictorio y la decisión definitiva, con certeza, al dictar sentencia. Si primero no se protegiera y se procediera a conocer, la sentencia podría ser estéril y el reconocimiento de derechos llegaría tarde.

\subsection{Clarín (CSJN 333:1885): faceta no innovativa}

\textbf{Contexto.} Hace aproximadamente diez años se dictó la Ley 26.522, conocida como ``Ley de Medios'', cuya finalidad era evitar la monopolización de los distintos medios de comunicación y telecomunicaciones. Las empresas que contaban con una gran cantidad de licencias debían desinvertir y vender parte de ellas. Esto afectaba de modo muy importante al Grupo Clarín, que debía iniciar, dentro del plazo legal, un proceso de desinversión.

\textbf{Proceso principal.} El Grupo Clarín planteó un proceso en el que discutió la constitucionalidad de los artículos 41 y 161 de la ley, que eran los que lo obligaban a desinvertir.

\textbf{Pretensión cautelar.} Solicitó una tutela anticipada dentro de la prohibición de innovar (Art.~230 CPCC), en su faceta \textbf{no innovativa}: que no se le aplicara la ley y no se alterara la situación existente hasta que se resolviera la pretensión sustancial sobre la constitucionalidad de los artículos. Si no se concedía y la sentencia declaraba la inconstitucionalidad, llegaría tarde: el Grupo ya habría tenido que desinvertir.

\textbf{Resolución y enseñanza.} La CSJN hizo lugar a la prohibición de innovar en su faceta no innovativa, y la ley no se aplicó mientras duró el proceso. Sin embargo, en la sentencia definitiva la Corte resolvió que los artículos 41 y 161 \textbf{eran constitucionales}, decisión contraria a lo protegido cautelarmente.
% TODO: contenido incompleto o ambiguo en las notas originales (la cita del docente sobre este fallo contiene una frase confusa: "el juez consideró autor necesario otorgar esta tutela").

El caso permite ver, por un lado, cómo opera la faceta de no innovar de la prohibición y, por otro, con claridad, el \textbf{carácter provisorio} de este tipo de medidas: se otorgó la tutela para evitar un daño muy difícil de reparar y, una vez producido el conocimiento a lo largo del proceso, se decidió que las normas eran constitucionales.

\section{Debate doctrinario sobre la regulación de la tutela anticipada}

\subsection{La posición crítica: regulación independiente en las provincias}

Una parte de la doctrina entiende que la tutela anticipada \textbf{no es un sistema cautelar} y que no resulta adecuado que cobre vida a través de la prohibición de innovar del Art.~230. Sostiene que debe regularse expresamente de manera distinta de las medidas cautelares, porque adelantar total o parcialmente, de modo provisorio y cautelar, lo que será objeto de la sentencia de mérito implicaría un instituto distinto.

Esta posición llevó a varias provincias ---como San Luis, San Juan y Chaco--- a regular la tutela anticipada de modo independiente en sus códigos procesales, con requisitos más exigentes. La crítica expresada es que ya no bastaría acreditar los presupuestos que exige la Corte (verosimilitud en el derecho, peligro en la demora, contracautela y mayor prudencia del juez): se establecen numerosos y rígidos requisitos adicionales que entorpecen el acceso del justiciable a esta tutela.

\subsection{Artículo 403 del Proyecto de reforma del CPCC: tutela anticipada de urgencia}

El proyecto de reforma del Código Procesal Civil y Comercial de la Nación regula en su artículo 403 la llamada ``tutela anticipada de urgencia'', con requisitos notablemente más exigentes:

\medskip
\renewcommand{\arraystretch}{1.3}
\noindent\begin{tabularx}{\textwidth}{@{}>{\raggedright\arraybackslash}p{2.8cm} Y Y@{}}
\toprule
\textbf{Requisito} & \textbf{Régimen actual (CSJN, Camacho Acosta)} & \textbf{Proyecto, Art.~403} \\
\midrule
Presupuesto temporal &
Peligro en la demora. &
Urgencia (estándar más elevado). \\
\addlinespace
Tipo de perjuicio &
Perjuicio inminente o irreparable. &
Frustración del derecho y daño irreparable. \\
\addlinespace
Prueba del derecho &
Verosimilitud (probabilidad). &
Debe ofrecer pruebas para acreditar la probabilidad del derecho. \\
\addlinespace
Contracautela &
Sí. &
Sí. \\
\addlinespace
Materias alcanzadas &
Cualquier tipo de derecho. &
Sólo prestaciones dinerarias. \\
\addlinespace
Trámite &
Inaudita parte: el afectado se defiende después. &
Con audiencia previa: primero se cita al demandado y luego se resuelve. \\
\bottomrule
\end{tabularx}
\renewcommand{\arraystretch}{1}

\medskip
Se critica la limitación a prestaciones dinerarias como incomprensible: se pierde una tutela judicial efectiva para situaciones vinculadas con otros derechos fundamentales y no sólo con cuestiones netamente dinerarias.

También se critica el trámite con audiencia previa por la demora que puede implicar. Si el demandado aún no se presentó en el proceso, debe notificárselo en su domicilio real (por ejemplo, Bahía Blanca, cuando el proceso tramita en la Ciudad de Buenos Aires); a veces es necesario diligenciar varias cédulas a través de la oficina de notificaciones, lo que puede demorar meses, para otorgar una tutela que en general requiere cierta urgencia.

\subsection{Principios vulnerados según la crítica}

Si la decisión es provisoria, si el afectado puede luego defenderse, si existe una contracautela como garantía por los daños y perjuicios, y si el instituto funciona del modo visto como medida cautelar dentro de la doctrina de la Corte, se cuestiona por qué dictar regulaciones que atentan contra principios reconocidos en tratados de derechos humanos con jerarquía constitucional:

\begin{itemize}
    \item \textbf{Tutela judicial efectiva:} consagrada en los tratados internacionales de derechos humanos con jerarquía constitucional.
    \item \textbf{Principio pro homine:} debe estarse siempre a la interpretación más favorable a la persona.
    \item \textbf{Principio de progresividad:} una vez que los justiciables adquieren cierto nivel de protección, el legislador no puede regular de modo regresivo en su contra.
\end{itemize}

Según esta posición, los sistemas que regulan la tutela anticipada de urgencia de este modo podrían ser tachados de \textbf{inconstitucionales y anticonvencionales}, porque el legislador crea un sistema apartándose de los principios de tutela judicial efectiva y de progresividad.

\cornelltitle{ámbito excepcional y tutela anticipada}
\cornellblock{
\cqrow{¿Qué define al ámbito excepcional de los sistemas cautelares?}{La superposición total o parcial de la pretensión cautelar con la pretensión sustancial (el objeto del proceso principal), a diferencia del ámbito normal, donde ambas son distintas.
}
\cqrow{¿Qué es la tutela anticipada?}{Es el sistema cautelar por el cual la jurisdicción protege y asegura anticipadamente y de modo provisorio, en todo o en parte, lo que luego se decidirá definitivamente en la sentencia de mérito.
}
\cqrow{¿Por qué no hay prejuzgamiento si se adelanta lo que decidirá la sentencia?}{Porque lo decidido cautelarmente es provisorio: se otorga con verosimilitud en el derecho y peligro en la demora, y luego, tras el debate y la prueba, la sentencia puede resolver en el mismo sentido con carácter definitivo o en sentido contrario, dejando sin efecto lo concedido.
}
\cqrow{¿Qué diferencia hay entre cautela y condena?}{Si el juez otorga la tutela anticipada de modo provisorio hay cautela; si decide en la sentencia definitiva sobre esa misma pretensión hay condena, que es definitiva y con autoridad de cosa juzgada.
}
\cqrow{¿Cuáles son las vías de materialización de la tutela anticipada?}{(1) Los sistemas cautelares: prohibición de innovar (Art.~230), medida cautelar genérica (Art.~232) u otra medida cautelar. (2) Subsistemas creados por el legislador: alimentos provisorios (Art.~544 CCyC) y entrega anticipada del inmueble en el desalojo contra el intruso (Art.~680 bis CPCC).
}
\cqrow{¿Qué resolvió la Corte en Camacho Acosta (CSJN 320:1633)?}{Hizo lugar a una medida innovativa (Art.~230 CPCC) como tutela anticipada: se anticipó el dinero para la prótesis del trabajador al que se le amputó el brazo, por verosimilitud, peligro en la demora y contracautela. Agregó que estos casos exigen mayor prudencia del juez al valorar presupuestos y hechos.
}
\cqrow{¿Qué resolvió la Corte en Clarín (CSJN 333:1885) y qué demuestra?}{Hizo lugar a la prohibición de no innovar: no se aplicó la Ley 26.522 al Grupo Clarín mientras duró el proceso. En la sentencia definitiva declaró constitucionales los artículos 41 y 161, lo que muestra el carácter provisorio de la tutela anticipada y la faceta no innovativa de la prohibición de innovar.
}
}
\medskip
\cornellblocksum{
\cqrow{¿Qué significa la inversión del contradictorio?}{En la tutela anticipada la jurisdicción primero protege y luego conoce, para evitar un daño irreparable o una sentencia estéril por tardía; el contradictorio y la decisión definitiva vienen después.
}
\cqrow{¿Qué críticas se formulan a la regulación independiente de la tutela anticipada?}{Que las provincias que la regulan y el Art.~403 del proyecto agregan requisitos numerosos y rígidos (urgencia, pruebas, sólo prestaciones dinerarias, audiencia previa) que entorpecen el acceso a la tutela y pueden demorar meses. Se afirma que vulneran la tutela judicial efectiva, el principio pro homine y la progresividad, y podrían ser inconstitucionales y anticonvencionales.
}
}{El ámbito excepcional se configura cuando la pretensión cautelar se superpone total o parcialmente con la sustancial; esto es la tutela anticipada, válida porque es provisoria y se otorga con verosimilitud y peligro en la demora, con contracautela y mayor prudencia judicial. Se materializa por vías cautelares (prohibición de innovar, medida genérica) y por subsistemas legales (alimentos provisorios, desalojo anticipado). La Corte aplicó la faceta innovativa en Camacho Acosta y la no innovativa en Clarín, donde la sentencia final fue contraria a lo cautelarmente protegido. La regulación independiente con requisitos agravados es criticada por contrariar la tutela judicial efectiva, el principio pro homine y la progresividad.
}

\chapter{El embargo}\label{cap:embargo}

\section{El embargo como medida cautelar por antonomasia}

\begin{definition}{Embargo}
La doctrina ha dicho generalmente que el embargo es la \textbf{medida cautelar por antonomasia}: es la más usada y la más habitual, y consiste en la \textbf{afectación e individualización de un bien} ---mueble o inmueble, registrable o no--- al resultado de un proceso.
\end{definition}

El bien así individualizado queda vinculado directamente con el proceso, queda a disposición del juez y opera como \textbf{prenda concreta} para garantizar la efectividad del resultado del juicio.

Ese bien se encuentra afectado ante la eventualidad de que se concrete la realización de bienes, es decir, la necesidad de sacar a la venta bienes para pagar la deuda reclamada en el proceso. Será el bien fundamental a la hora de procederse a la subasta pública como modo de garantizar el pago de la deuda.

\section{Tipos de embargo}

El ordenamiento reconoce tres tipos de embargo, que se distinguen según el momento procesal en que se traban:

\begin{table}[H]
\centering
\small
\begin{tabularx}{\textwidth}{>{\raggedright\arraybackslash}p{2.6cm} >{\raggedright\arraybackslash}p{3.6cm} >{\raggedright\arraybackslash}X}
\toprule
\textbf{Tipo} & \textbf{Momento} & \textbf{Normativa / características} \\
\midrule
\textbf{Embargo preventivo} & Antes del proceso o durante su tramitación & Arts.\ 209 a 212 CPCCN. Eminentemente cautelar. Exige los tres recaudos: verosimilitud del derecho, peligro en la demora y contracautela suficiente. \\
\addlinespace
\textbf{Embargo ejecutivo} & Una vez admitido el título en el juicio ejecutivo & Se traba cuando el juez reconoce el título como admisible y libra la orden de intimación de pago. Individualiza el bien que será ejecutado si no se paga. \\
\addlinespace
\textbf{Embargo ejecutorio} & En el proceso de ejecución de sentencia & Se dicta cuando la sentencia está confirmada y se avanza en la concreción del pago. Es el paso previo a la subasta del bien. Art.\ 551 CPCCN (sentencia de trance y remate). \\
\bottomrule
\end{tabularx}
\caption{Tipos de embargo}
\end{table}

\subsection{El embargo preventivo}

El embargo preventivo es el eminentemente cautelar. Puede trabarse aun antes del proceso o durante la secuela del proceso, siempre que se reúnan los tres recaudos desarrollados en el capítulo~\ref{cap:presupuestos}: \textbf{verosimilitud del derecho}, \textbf{peligro en la demora} y \textbf{contracautela} adecuada o suficiente.

Los artículos de referencia son los siguientes:
\begin{itemize}
    \item \textbf{Art.\ 209 CPCCN:} establece los supuestos de procedencia del embargo preventivo.
    \item \textbf{Art.\ 210 CPCCN:} amplía los supuestos y complementa el artículo anterior.
    \item \textbf{Art.\ 211 CPCCN:} procedimiento para la traba del embargo.
    \item \textbf{Art.\ 212 CPCCN:} hipótesis específicas derivadas de la secuela del proceso.
\end{itemize}

\subsection{El embargo ejecutivo}

En el trámite del proceso de ejecución existe una primera etapa, que es el análisis del título. No siempre el título trae aparejada, por sí mismo, la ejecución: a veces es necesario preparar la vía ejecutiva. En esos casos procede el embargo preventivo, porque todavía no hay un título ejecutivo declarado admisible.

Una vez que el juez advierte que se trata de un título \textit{prima facie} hábil ---ya sea porque se preparó la vía ejecutiva o porque objetivamente reúne todos los recaudos propios del título---, el título se considera admisible y se libra la orden de intimación de pago. A partir de ese momento procede el \textbf{embargo ejecutivo}.

Por lo tanto, en un juicio ejecutivo puede haber un embargo preventivo trabado antes del inicio del proceso o con el inicio del proceso, si este no comienza directamente con la intimación de pago porque es necesario preparar la vía ejecutiva. Una vez preparada la vía ejecutiva y admitido el título, el embargo pasa a ser ejecutivo, ya sea porque el preventivo se transforma en ejecutivo, ya sea porque se traba directamente como embargo ejecutivo.

\subsection{El embargo ejecutorio}

El embargo ejecutorio es el que se traba en el marco del proceso de ejecución de sentencia. Se dicta cuando una sentencia ya ha sido confirmada por el superior y se empieza a avanzar en la concreción del pago derivado de una sentencia condenatoria. Con sentencia firme y ejecutoriada se admite la traba de un embargo que ya no es preventivo ---porque el derecho está consolidado y reconocido--- ni ejecutivo ---porque no se traba en el trámite de ejecución de un título ejecutivo---, sino \textbf{ejecutorio}, propio del ámbito de la ejecución de sentencia.

\begin{important}
\textbf{Art.\ 551 CPCCN (sentencia de trance y remate).} A partir de su dictado, los embargos anteriores se transforman en ejecutorios por efecto de la ley. La transformación opera automáticamente, muchas veces sin disposición concreta del tribunal, por la etapa del proceso en que se encuentra la causa. El embargo ejecutorio es el paso previo a la subasta del bien.
\end{important}

\section{Efectos de la traba del embargo}

El embargo es una medida procesal que se dicta en el proceso, pero que tiene efectos regulados en los códigos sustanciales. Trae consecuencias civiles y penales, que se ponen en juego en caso de incumplimiento de las obligaciones del depositario de la cosa embargada.

\begin{important}
\textbf{Consecuencias extralegales del embargo.} El depositario del bien embargado tiene la obligación de conservar la cosa, no disminuir su valor y ponerla a disposición del tribunal cuando sea requerido. El Código Penal tipifica delitos vinculados al incumplimiento de las obligaciones del depositario.
\end{important}

Los efectos fundamentales son los siguientes:
\begin{itemize}
    \item El embargo individualiza el bien, que queda afectado a la garantía y vinculado al proceso.
    \item El embargo \textbf{no produce la desapropiación} del bien por parte del afectado. El afectado puede usar y gozar de la cosa, en tanto no la haga perder valor.
    \item Tratándose de bienes registrables, el embargo \textbf{no produce la imposibilidad de enajenación} del bien: el bien puede ser vendido.
\end{itemize}

\subsection{El principio \textit{nemo plus iuris} y el embargo}

Por el principio \textit{nemo plus iuris} ---nadie puede transmitir sobre una cosa un derecho mayor al que posee sobre ella---, quien posee un bien embargado lo posee en tanto se encuentra embargado. Si lo transfiere, lo transmite embargado, y quien lo adquiere adquiere sobre la cosa el mismo derecho que tenía su propietario.

\subsection{El plenario «Sertop» --- Doctrina judicial obligatoria}

\begin{important}
\textbf{Plenario «Sertop»} (Cámara Nacional de Apelaciones en lo Civil). Constituye doctrina judicial obligatoria en la Ciudad Autónoma de Buenos Aires.

\textbf{Regla:} el adquirente de un bien embargado asume la totalidad de la deuda ---no solo el monto por el cual se trabó la medida---, incluyendo ampliaciones, intereses y costas del proceso.

En el fuero \textbf{comercial} este criterio no rige: la transmisión se limita al monto de la cautela.
\end{important}

Ante la duda que genera el criterio que se empleará en este tipo de transmisiones, difícilmente alguien esté dispuesto a adquirir una cosa embargada. Por ello, el embargo suele operar como una suerte de traba a la libre circulación de los bienes. Generalmente quienes adquieren bienes embargados especulan con ello, pagando un precio significativamente menor al de mercado, o tienen algún tipo de vinculación con el sujeto embargado.

\section{Supuestos de procedencia del embargo preventivo (arts.\ 209 y 210 CPCCN)}

La enumeración legal es meramente enunciativa y orientativa, no taxativa: las situaciones que habilitan el embargo pueden ser muchas más, y no es necesario que deban encuadrarse en alguno de los presupuestos que menciona el código.

Los supuestos típicos del art.\ 209 CPCCN son:
\begin{enumerate}[label=\textbf{\arabic*.}]
    \item \textbf{Deudor sin domicilio en la República} (inciso 1°). Se presume el peligro en la demora, porque puede ser dificultoso encontrarle bienes para responder.
    \item \textbf{Reconocimiento de la deuda a través de instrumento público.} Se presume reunida la verosimilitud y el peligro en la demora.
    \item \textbf{Instrumento privado cuyas firmas son abonadas por dos testigos} (información sumaria). El aval de dos testigos que declaran haber presenciado la firma o conocer la rúbrica del deudor permite dictar la medida sin mayores miramientos.
    \item \textbf{Contratos bilaterales} instrumentados públicamente o en instrumento privado con dos testigos (información sumaria), que avalan que la prestación a cargo del acreedor fue cumplida y resta cumplir la contraprestación.
\end{enumerate}

\subsection{Hipótesis procesales del artículo 212 CPCCN}

El art.\ 212 describe tres hipótesis en las que la secuela del proceso determina la procedencia del embargo preventivo:
\begin{itemize}
    \item \textbf{Inciso 1° (remisión al art.\ 63):} declaración de rebeldía del demandado. El hecho de que la parte sea declarada rebelde da motivo formal para la traba del embargo preventivo.
    \item \textbf{Inciso 2°:} confesión de la existencia de la deuda en audiencia confesional, o declaración de confeso ficto (por no comparecer a la audiencia habiendo sido citado). En ambos casos procede el embargo sin mayores recaudos.
    \item \textbf{Inciso 3°:} existencia de sentencia definitiva condenatoria, aunque esté apelada. Si el juez de primera instancia condenó al pago, el embargo preventivo está plenamente configurado; la eventual revisión por la Cámara no impide la medida.
\end{itemize}

\begin{note}
Cuando el demandado no contesta la demanda, por efecto del art.\ 356 se presume la veracidad de los hechos descriptos, lo que también habilita la traba. % TODO: contenido incompleto o ambiguo en las notas originales (las notas aluden a «hechos ilícitos descriptos»)
\end{note}

La existencia de una sentencia judicial condenatoria, aunque se encuentre apelada, otorga una mayor verosimilitud del derecho. La Cámara puede revisar y revocar la sentencia; sin embargo, ya existe un juez de primera instancia que consideró que el deudor debe pagar la suma reclamada, y es lógico que en ese contexto se dicte el embargo preventivo como modo de garantizar el cumplimiento de la sentencia ya dictada.

\section{Bienes inembargables --- Artículo 219 CPCCN}

\begin{enumerate}[label=\textbf{Inciso \arabic*°:}, leftmargin=2.4cm]
    \item El lecho cotidiano del deudor, su esposa e hijos; las ropas y bienes muebles de indispensable uso. Según el criterio judicial, si el deudor tiene cinco televisores puede reconocerse la inembargabilidad de uno, y los demás quedan embargables. También son inembargables las herramientas u objetos que el deudor utiliza para el ejercicio de su profesión u oficio.
    \item Los sepulcros, salvo que la deuda sea por el saldo del precio de compra del sepulcro o por bienes adquiridos para su construcción.
    \item Los bienes declarados inembargables por leyes especiales (jubilaciones, pensiones, salarios, etc.).
\end{enumerate}

\subsection{Inembargabilidad del salario}

El trabajador tiene en su salario un porcentaje de inembargabilidad establecido por la Ley de Contrato de Trabajo, que establece que solo puede ser embargado hasta el \textbf{20\,\%} de sus ingresos.

\begin{example}{Límite del 20\,\% y deuda alimentaria}
Si el salario del deudor es de \$30.000, el 20\,\% equivale a \$6.000. Sin embargo, si el deudor tiene tres hijos menores y la madre reclama alimentos, \$6.000 puede resultar insuficiente. En ese supuesto, el juez puede excepcionalmente elevar el límite de embargabilidad al 40\,\% o 50\,\% del ingreso, dejando de lado el límite de la Ley de Contrato de Trabajo, siempre como modo de garantizar la subsistencia del trabajador.
\end{example}

Un criterio similar rige para jubilaciones y pensiones, que en principio son inembargables, salvo en caso de deudas alimentarias, donde el juez puede establecer el porcentaje embargable garantizando la subsistencia del jubilado o pensionado.

\section{Caducidad del embargo --- Artículo 207 CPCCN}

En materia de embargo rige, como en general en todas las medidas cautelares, lo establecido por el art.\ 207 del código procesal sobre la caducidad de la traba.

\begin{formula}{Caducidad (art.\ 207 CPCCN)}
\begin{itemize}
    \item \textbf{Plazo:} 10 días desde la traba de la medida para iniciar el proceso principal.
    \item En procesos sujetos a mediación previa, el inicio de la mediación interrumpe el plazo de caducidad (equivale al inicio del proceso).
    \item \textbf{Embargos registrales} (Registro del Automotor): caducan a los \textbf{cinco años}. Debe pedirse la reinscripción dentro de ese plazo para mantener subsistente la medida.
\end{itemize}
\end{formula}

\cornelltitle{el embargo}
\cornellblock{
\cqrow{¿Qué es el embargo y por qué se lo llama «medida cautelar por antonomasia»?}{Es la medida cautelar más usada y habitual. Consiste en la afectación e individualización de un bien ---mueble o inmueble, registrable o no--- al resultado de un proceso. El bien queda vinculado al proceso, a disposición del juez, y opera como prenda concreta para garantizar el resultado del juicio y, llegado el caso, la subasta pública.
}
\cqrow{¿Cuáles son los tres tipos de embargo?}{
}
\cqrow{¿Qué efectos produce la traba del embargo?}{Individualiza el bien y lo afecta a la garantía. No desapropia al afectado, que puede usar y gozar de la cosa mientras no la haga perder valor. En bienes registrables no impide la enajenación. El depositario debe conservar la cosa, no disminuir su valor y ponerla a disposición del tribunal; su incumplimiento tiene consecuencias civiles y penales.
}
\cqrow{¿Puede venderse un bien embargado? ¿Cuáles son las consecuencias?}{Sí, pero por \textit{nemo plus iuris} se transmite embargado y el adquirente obtiene el mismo derecho que tenía el propietario. Según el plenario «Sertop» (doctrina obligatoria en la Ciudad Autónoma de Buenos Aires), el adquirente asume la totalidad de la deuda, incluyendo ampliaciones, intereses y costas; en el fuero comercial la transmisión se limita al monto de la cautela. Por ello difícilmente alguien adquiere bienes embargados y el embargo opera como traba a la libre circulación.
}
\cqrow{¿Cuáles son los supuestos de procedencia del embargo preventivo?}{La enumeración de los arts.\ 209 y 210 es enunciativa, no taxativa. Art.\ 209: deudor sin domicilio en la República; reconocimiento de deuda por instrumento público; instrumento privado con firmas abonadas por dos testigos; contratos bilaterales con prestación del acreedor cumplida. Art.\ 212: rebeldía, confesión o confeso ficto, y sentencia condenatoria aunque esté apelada.
}
\cqrow{¿Cuáles son los principales bienes inembargables?}{Art.\ 219 CPCCN: el lecho cotidiano y las ropas y bienes muebles de indispensable uso; los sepulcros (con la salvedad indicada); y los bienes declarados inembargables por leyes especiales (jubilaciones, pensiones, salarios). El salario es embargable hasta el 20\,\%, límite que el juez puede elevar excepcionalmente en deudas alimentarias.
}
}
\medskip
\cornellblocksum{
\cqrow{¿Cuándo caduca el embargo?}{Art.\ 207 CPCCN: a los 10 días desde la traba si no se inicia el proceso principal (la mediación previa interrumpe el plazo). Los embargos registrales sobre automotores caducan a los cinco años, salvo reinscripción dentro de ese plazo. \\




%==============================================================
}
}{El embargo individualiza un bien y lo afecta al proceso sin desapropiar al deudor. Se clasifica según el estadio procesal en preventivo, ejecutivo y ejecutorio. El adquirente de un bien embargado lo recibe con la afectación (\textit{nemo plus iuris}; plenario «Sertop»). Rigen límites de embargabilidad (art.\ 219, límite del 20\,\% del salario) y la caducidad del art.\ 207.
}

\chapter[Secuestro e intervención judicial]{Secuestro e intervención judicial}\label{cap:secuestro}

\section{El secuestro}
%==============================================================

\subsection{Concepto y ámbito de aplicación}

El secuestro está orientado, como medida cautelar, a desapoderar al deudor de la disposición material y jurídica sobre una cosa de su propiedad.

\begin{important}
El secuestro solo opera respecto de bienes \textbf{muebles}. No existe secuestro de bienes inmuebles. Puede recaer sobre bienes muebles registrables (autos, motos) y no registrables. Los animales (ganado, caballos de carrera) también pueden ser objeto de secuestro.
\end{important}

\subsection{Tipos de secuestro}

\subsubsection{Secuestro como medida cautelar propiamente dicha (art.\ 221 CPCCN)}

Es excepcional y no tan habitual como el embargo. De acuerdo con las características del bien, las del deudor o la situación en que se encuentra el bien, en algunas hipótesis procede el secuestro como único modo de garantizar la intangibilidad de la cosa.

Procede cuando la cosa puede ser deteriorada por el uso continuo del deudor, o cuando el deudor ha dado muestras de falta de cuidado y ha permitido que el bien pierda su valor.

\begin{example}{Procedencia del secuestro como cautelar}
\textbf{Ejemplo 1.} La cosa es una obra de arte almacenada en un depósito con humedad, que no recibe los cuidados indispensables para su preservación. Es preferible el secuestro al embargo para evitar el deterioro del bien.

\textbf{Ejemplo 2.} Se trata de aparatología para realizar estudios técnicos o científicos que requiere un mantenimiento que el deudor no realiza, generando un deterioro progresivo.
\end{example}

\subsubsection{Secuestro como sanción}

La doctrina reconoce que el secuestro a veces es una sanción derivada de la conducta del deudor o de quien incumple una obligación procesal. El código lo prevé en dos hipótesis:

\textbf{Art.\ 128 CPCCN --- Secuestro por retención de expediente.} Cuando un letrado o un auxiliar del tribunal (escribano, perito) retira un expediente en préstamo y no lo devuelve en el plazo fijado por el juez:
\begin{enumerate}
    \item Se intima a devolverlo.
    \item Si no lo devuelve en el plazo de la intimación, se dispone el secuestro del expediente por oficial de justicia, con allanamiento de morada si fuera necesario.
\end{enumerate}
Es una hipótesis de secuestro como sanción por incumplimiento de la obligación de restituir.

\textbf{Art.\ 323 CPCCN, incisos 1° a 5° --- Diligencias preliminares.} En las diligencias preliminares puede pedirse a la parte demandada que exhiba un objeto, un inmueble, un testamento o los títulos de una cosa vendida. Si se incumple la intimación de exhibición, una de las sanciones previstas es el secuestro de la cosa, título o documento de que se trate.

\subsubsection{Secuestro como paso previo a la subasta}

Una vez que el bien es embargado en el marco de un juicio de ejecución y es necesario proceder a su venta en pública subasta, porque ya se agotaron los pasos procesales previos, el secuestro actúa como presupuesto de la subasta pública de los bienes embargados.

Si el bien es un inmueble, no corresponde el secuestro, porque este no procede respecto de todo lo que sea inmueble. Los bienes muebles, sean registrables o no, tienen como paso previo a la subasta el secuestro.

\cornelltitle{el secuestro}
\cornellblocksum{
\cqrow{¿Qué es el secuestro y sobre qué bienes recae?}{Es una medida cautelar orientada a desapoderar al deudor de la disposición material y jurídica sobre una cosa de su propiedad. Solo recae sobre bienes muebles, registrables o no, incluidos los animales. No existe secuestro de inmuebles.
}
\cqrow{¿Cuándo procede el secuestro como medida cautelar (art.\ 221 CPCCN)?}{Excepcionalmente, cuando por las características del bien, del deudor o de la situación del bien es el único modo de garantizar la intangibilidad de la cosa: si puede deteriorarse por el uso continuo del deudor o si el deudor ha mostrado falta de cuidado y permitió que pierda valor (por ejemplo, una obra de arte en un depósito con humedad).
}
\cqrow{¿En qué hipótesis el secuestro opera como sanción?}{En dos: art.\ 128 CPCCN, cuando un letrado o auxiliar del tribunal no devuelve un expediente tras la intimación (secuestro por oficial de justicia, con allanamiento si fuera necesario); y art.\ 323 CPCCN, cuando se incumple la intimación a exhibir en diligencias preliminares.
}
\cqrow{¿Qué función cumple el secuestro respecto de la subasta?}{Es presupuesto de la subasta pública de los bienes muebles embargados en un juicio de ejecución. Los inmuebles no se secuestran. \\




%==============================================================
}
}{El secuestro desapodera al deudor de la disposición material y jurídica de bienes muebles. Es excepcional como cautelar (art.\ 221), opera como sanción (arts.\ 128 y 323) y constituye paso previo a la subasta de los muebles embargados. No procede sobre inmuebles.
}

\section{La intervención judicial}
%==============================================================

\subsection{Concepto y fundamento (arts.\ 222 a 227 CPCCN)}

\begin{definition}{Intervención judicial}
Consiste en la designación, por parte del juez, de un \textbf{auxiliar externo al tribunal} para que se involucre en la administración o gestión de un ente colectivo, un patrimonio o una empresa. El interventor \textbf{no tiene facultades de gobierno ni de disposición de bienes}. Su función es garantizar una saludable administración o brindar información al juez.
\end{definition}

\subsection{Tipos de intervención judicial}

\subsubsection{Interventor recaudador (art.\ 223 CPCCN)}

El interventor recaudador es un auxiliar del tribunal designado por el juez, cuya misión es recaudar fondos del giro habitual de la administración de una empresa, un comercio o un patrimonio.

Suele darse cuando la parte deudora no tiene un patrimonio significativo pero tiene un comercio o consultorio. El interventor concurre cotidianamente al establecimiento y, del dinero que ingresa, retiene una suma para afectarla al pago de la deuda reclamada en el proceso.

\begin{example}{Empresa de transporte de ómnibus}
Una empresa de ómnibus de larga distancia que cobra sus ingresos por ventanilla (venta de pasajes) tiene una deuda de \$20.000. Embargar un ómnibus y sacarlo a subasta resulta un despropósito, porque el gasto de la subasta supera la deuda.

\textbf{Solución:} se designa un interventor a ventanilla que retiene del ingreso cotidiano la suma correspondiente, logrando así el cumplimiento de la deuda de forma proporcional y eficiente.

En el caso de empresas de colectivos con tarjeta SUBE, se puede directamente comunicar por oficio a la administración de la gestión de las tarjetas para que retengan lo cobrado.
\end{example}

\subsubsection{Interventor informante (art.\ 224 CPCCN)}

Está orientado a que el interventor que concurre al comercio de la empresa de la parte afectada dé noticia al juez acerca de la manera en que se lleva a cabo la administración, el patrimonio que posee la empresa, el estado de deudas y todo otro dato que el juez necesite. Aporta así información valiosa al proceso sobre la gestión de un comercio, empresa o patrimonio, que puede ser necesaria para decisiones que deban tomarse en el proceso.

\begin{example}{Sucesión con empresa}
En un juicio sucesorio, el patrimonio del causante incluye una empresa o comercio del que era titular. Para distribuir la herencia entre los herederos es necesario contar con información precisa sobre ingresos, egresos, bienes y deudas de ese patrimonio. En esos casos la intervención del informante es necesaria para «poner blanco sobre negro» acerca del alcance de los bienes a distribuir y los ingresos que genera el patrimonio.
\end{example}

\subsubsection{Veedor (art.\ 115 Ley 19.550 --- Ley de Sociedades Comerciales)}

En doctrina se discute qué es el veedor a diferencia del informante. Prácticamente es lo mismo: la Ley de Sociedades denomina veedor al informante, y la diferencia tiene que ver más con una cuestión terminológica que con el alcance de la gestión.

En ambos casos ---interventor informante del art.\ 224 CPCCN o veedor del art.\ 115 de la Ley 19.550--- el alcance de su gestión y la información que debe brindar está establecido con mayor precisión por lo que dispone el juez en la resolución que ordena la intervención que por lo preestablecido en las normas del código.

\subsubsection{Interventor administrador o co-administrador}

El interventor gestiona la administración ---en casos puntuales que el juez comercial establece--- o la cogestiona con administradores propios de la empresa. Se da un paso más: ya no es un informante, sino que realiza cuestiones vinculadas con la administración cotidiana.

Suele suceder cuando la empresa está en estado prefalencial y es necesario evitar que, por maniobras de administración, se propicie la insolvencia a través del retiro de dinero por parte de los socios u otras maniobras perjudiciales.

\begin{important}
\textbf{Límites del interventor.} La intervención no permite la disposición de bienes ni el gobierno de la empresa. Solo en el caso del administrador o co-administrador está habilitada la administración. Los actos de disposición o gobierno deben ser siempre adoptados por el juez, con intervención del síndico, cumpliendo los pasos procesales del proceso falencial. El administrador nunca tiene facultades para actos de disposición sin disposición expresa del juez.
\end{important}

\subsubsection{Quiénes pueden ser designados interventores}

Se aconseja que sean personas idóneas o técnicas en el negocio de que se trate. Generalmente se los busca entre administradores de empresa, contadores o incluso abogados para ciertas gestiones, como la de recaudación en un comercio.

Los abogados, contadores y licenciados en administración de empresas que pretenden ofrecer sus servicios para este tipo de intervenciones se anotan en los listados que elabora anualmente la Cámara Civil y la Cámara Comercial, presentando su currículum y poniéndose a disposición de los tribunales.

\cornelltitle{la intervención judicial}
\cornellblocksum{
\cqrow{¿Qué es la intervención judicial?}{Es la designación por el juez de un auxiliar externo al tribunal para que se involucre en la administración o gestión de un ente colectivo, un patrimonio o una empresa, a fin de garantizar una saludable administración o brindar información al juez.
}
\cqrow{¿Qué facultades tiene el interventor judicial?}{No tiene facultades de gobierno ni de disposición de bienes. Solo el administrador o co-administrador está habilitado para la administración. Los actos de disposición o gobierno los adopta siempre el juez, con intervención del síndico en el proceso falencial.
}
\cqrow{¿Qué hace el interventor recaudador (art.\ 223)?}{Recauda fondos del giro habitual de una empresa, comercio o patrimonio, reteniendo del dinero que ingresa una suma para afectarla al pago de la deuda. Es útil cuando el deudor no tiene patrimonio significativo pero sí un comercio o consultorio (por ejemplo, el interventor a ventanilla en una empresa de ómnibus).
}
\cqrow{¿Qué diferencia hay entre el interventor informante y el veedor?}{Prácticamente ninguna: el informante del art.\ 224 CPCCN y el veedor del art.\ 115 de la Ley 19.550 difieren más por la terminología que por el alcance de su gestión, que está determinado sobre todo por la resolución judicial que ordena la intervención.
}
\cqrow{¿Cuándo actúa un interventor administrador o co-administrador?}{En casos puntuales que establece el juez comercial, cuando la empresa está en estado prefalencial y es necesario evitar maniobras de administración que propicien la insolvencia. Realiza tareas de administración cotidiana, sola o junto con los administradores de la empresa.
}
\cqrow{¿Quiénes pueden ser interventores?}{Personas idóneas o técnicas en el negocio (administradores de empresa, contadores, abogados), inscriptas en los listados que elaboran anualmente la Cámara Civil y la Cámara Comercial. \\




%==============================================================
}
}{La intervención judicial coloca a un auxiliar del tribunal en la gestión de una empresa o patrimonio, sin facultades de gobierno ni de disposición. Sus variantes son el recaudador (art.\ 223), el informante o veedor (art.\ 224 CPCCN; art.\ 115 Ley 19.550) y el administrador o co-administrador; los actos de disposición corresponden siempre al juez.
}

\chapter[Anotación de litis y prohibición de contratar]{Anotación de litis y prohibición de contratar}\label{cap:anotacion}

\section{La anotación de litis}
%==============================================================

\subsection{Concepto y fundamento}

Este tema se conecta con el principio del tercero adquirente de buena fe a título oneroso, estudiado en Derecho Civil I. Entre los efectos de la nulidad de actos jurídicos se encuentra volver las cosas hacia atrás: ante la nulidad, todo vuelve al \textit{statu quo ante}. Si el acto jurídico era la compraventa o la donación de un inmueble y es declarado nulo, el inmueble vuelve al patrimonio de quien lo enajenó nulamente.

\begin{important}
\textbf{Art.\ 392 CCyCN --- Tercero adquirente de buena fe a título oneroso.} Cuando existe un tercero adquirente de buena fe a título oneroso, la ley privilegia su situación: este tercero, que no fue parte en el acto declarado nulo y no tenía nada que ver con él, no tiene por qué cargar con las consecuencias de las operaciones viciadas realizadas anteriormente respecto del bien. Por lo tanto, los efectos de la nulidad no se extienden a su respecto.
\end{important}

\begin{definition}{Anotación de litis}
Es una medida cautelar que, adoptada siempre sobre bienes \textbf{registrables}, publicita \textit{erga omnes} (con alcance a todos) la existencia de un proceso en trámite en el que se discute la validez o el alcance de una registración existente respecto de ese bien.

\textbf{Consecuencia:} al ser pública la existencia del juicio, nadie puede invocar luego la condición de adquirente de buena fe, porque nadie puede alegar buena fe cuando sabe que hay un litigio en trámite respecto del bien.
\end{definition}

Quien adquiere el bien en esos términos sabe que adquiere un bien en riesgo de un litigio que puede resolverse en favor de uno u otro y que, según cómo salga el juicio, tal vez todo vuelva atrás y el bien deba devolverse a un anterior propietario que lo enajenó en virtud de un acto viciado.

\subsection{Requisitos de procedencia}

\begin{itemize}
    \item Solo procede contra bienes \textbf{registrables}.
    \item Debe haber una \textbf{demanda iniciada}. A diferencia del embargo, que puede trabarse antes del proceso, la anotación de litis exige en principio que exista juicio ya iniciado (se admite que el inicio de la mediación equivale al inicio del proceso).
    \item El juicio debe tratarse de una pretensión que, de prosperar, implicaría la modificación de un asiento registral respecto del bien (ya sea de titularidad u otra cuestión vinculada a lo que se anota registralmente).
\end{itemize}

\subsection{Supuestos típicos de procedencia}

\begin{table}[H]
\centering
\small
\begin{tabularx}{\textwidth}{>{\raggedright\arraybackslash}p{3.2cm} >{\raggedright\arraybackslash}p{3.0cm} >{\raggedright\arraybackslash}X}
\toprule
\textbf{Tipo de juicio} & \textbf{Normativa} & \textbf{Particularidad} \\
\midrule
\textbf{Usucapión / prescripción adquisitiva} & Art.\ 1905 CCyCN & El juez ordena la anotación de oficio (sin pedido de parte) al disponer el traslado de la demanda. \\
\addlinespace
\textbf{Expropiación} & Arts.\ 24 y 14 Ley 21.499 & Una de las primeras medidas al iniciarse el proceso es la anotación de litis, para impedir que el propietario venda el bien y luego un tercero invoque buena fe frente a la expropiación. \\
\addlinespace
\textbf{Escrituración} & CPCCN & Busca modificar el asiento registral para que el adquirente por boleto pase a ser titular registral. \\
\addlinespace
\textbf{Nulidad de escritura} & Art.\ 392 CCyCN & Supuesto clásico de vicios del consentimiento o actos irregulares que intentan cancelar un asiento registral. \\
\addlinespace
\textbf{Simulación} & CCyCN & Caso típico de procedencia de la anotación de litis. \\
\addlinespace
\textbf{Interdicto de adquirir} & Art.\ 609 CPCCN & Hipótesis propia de la anotación de litis. \\
\bottomrule
\end{tabularx}
\caption{Supuestos típicos de anotación de litis}
\end{table}

\subsection{Extinción de la anotación de litis}

La anotación de litis \textbf{no se rige por el art.\ 207 CPCCN}. Se extingue por:
\begin{enumerate}
    \item \textbf{Cumplimiento de la sentencia:} si la demanda prospera, el cumplimiento de la sentencia determina la modificación registral buscada y concluye la anotación.
    \item \textbf{Desestimación del reclamo:} cuando el proceso concluye por sentencia desestimatoria, se agota la anotación de litis. Debe comunicarse al registro.
    \item \textbf{Modos anormales de terminación:} caducidad de instancia, desistimiento o transacción. Cualquiera de ellos trae aparejada la extinción de la anotación.
\end{enumerate}

\cornelltitle{la anotación de litis}
\cornellblocksum{
\cqrow{¿Qué es la anotación de litis y qué protege?}{Es una medida cautelar sobre bienes registrables que publicita erga omnes la existencia de un proceso en el que se discute la validez o el alcance de una registración. Al ser pública, impide que alguien invoque luego buena fe como adquirente, ya que nadie puede alegar buena fe sabiendo que hay un litigio en trámite sobre el bien.
}
\cqrow{¿Qué relación tiene con el art.\ 392 CCyCN?}{El art.\ 392 privilegia al tercero adquirente de buena fe a título oneroso, a quien no se extienden los efectos de la nulidad. La anotación de litis neutraliza esa buena fe al hacer pública la existencia del juicio.
}
\cqrow{¿Cuáles son los requisitos de procedencia?}{Que se trate de un bien registrable; que haya demanda iniciada (el inicio de la mediación equivale al inicio del proceso); y que la pretensión, de prosperar, implique la modificación de un asiento registral.
}
\cqrow{¿En qué supuestos procede la anotación de litis de oficio?}{En la usucapión (art.\ 1905 CCyCN), donde el juez la ordena sin pedido de parte al disponer el traslado de la demanda. En la expropiación (arts.\ 24 y 14 Ley 21.499) es una de las primeras medidas al iniciarse el proceso.
}
\cqrow{¿Qué otros juicios son supuestos típicos?}{Escrituración, nulidad de escritura (art.\ 392 CCyCN), simulación e interdicto de adquirir (art.\ 609 CPCCN).
}
\cqrow{¿Cómo se extingue la anotación de litis?}{No se rige por el art.\ 207 CPCCN. Se extingue por cumplimiento de la sentencia, por desestimación del reclamo (con comunicación al registro) o por modos anormales de terminación del proceso (caducidad de instancia, desistimiento o transacción). \\




%==============================================================
}
}{La anotación de litis opera sobre bienes registrables como mecanismo de publicidad que neutraliza la buena fe del tercero adquirente. Exige juicio iniciado cuya pretensión implique modificar un asiento registral, es ordenada de oficio en usucapión, es una de las primeras medidas en expropiación y se extingue por cumplimiento de la sentencia, desestimación o modos anormales de terminación.
}

\section{La prohibición de contratar}
%==============================================================

\begin{note}
La prohibición de innovar (art.\ 230 CPCCN) se analiza en el capítulo~\ref{cap:generico}, donde es objeto de una explicación más detenida. Aquí se desarrolla directamente la prohibición de contratar (art.\ 231 CPCCN).
\end{note}

\subsection{Concepto y procedencia (art.\ 231 CPCCN)}

Es una cautela que en determinados casos resulta procedente, orientada a evitar que, con relación a un bien que es objeto del proceso o que se encuentra embargado en el proceso, se realicen actos jurídicos que puedan disminuir o condicionar la cautela o el trámite del cumplimiento de la sentencia.

\begin{example}{Ejecución hipotecaria}
En el marco de un juicio de ejecución hipotecaria, además del embargo del bien, suele adoptarse la prohibición de contratar para evitar que el deudor hipotecario introduzca personas al inmueble mediante un contrato de locación o comodato, dificultando el avance de las medidas de cumplimiento de la sentencia en la etapa de subasta. La prohibición de contratar puede notificarse también al locatario existente, para hacerle saber que no podrá prorrogarse ni ampliarse el contrato próximo a vencer.
\end{example}

\subsection{Notificación y publicidad}

La prohibición de contratar debe:
\begin{itemize}
    \item Si se trata de un bien registrable, asentarse en el registro para su publicidad \textit{erga omnes}.
    \item Notificarse al titular de la cosa o propietario.
    \item Poder notificarse también a cualquier tercero que se sospeche que pueda ser la persona que contrate con el propietario.
\end{itemize}

\subsection{Caducidad especial (art.\ 231 CPCCN)}

\begin{important}
La prohibición de contratar tiene una caducidad especial, más breve que la del art.\ 207: \textbf{cinco (5) días}. Este plazo corre desde el \textbf{dictado} de la medida (no desde su traba). Solo se aplica si la prohibición se adopta \textbf{antes del inicio del proceso}. Si se adopta durante el curso del proceso ---incluso en la etapa conclusional--- no rige este plazo especial.
\end{important}

\cornelltitle{la prohibición de contratar}
\cornellblocksum{
\cqrow{¿Para qué sirve la prohibición de contratar?}{Para evitar que, respecto de un bien objeto del proceso o embargado en él, se realicen actos jurídicos que puedan disminuir o condicionar la cautela o el cumplimiento de la sentencia. Por ejemplo, en una ejecución hipotecaria evita que el deudor introduzca personas al inmueble mediante locación o comodato.
}
\cqrow{¿Cómo se notifica y se publicita?}{Se asienta en el registro si el bien es registrable (publicidad erga omnes), se notifica al titular o propietario y puede notificarse a cualquier tercero que se sospeche pueda contratar con él, como el locatario existente, a quien se le hace saber que no podrá prorrogarse ni ampliarse el contrato.
}
\cqrow{¿Qué caducidad especial tiene y cuándo se aplica?}{Cinco días, contados desde el dictado de la medida (no desde su traba), y solo si se adopta antes del inicio del proceso. Si se adopta durante el proceso, incluso en la etapa conclusional, no rige ese plazo especial. \\




%==============================================================
}
}{La prohibición de contratar protege la integridad del bien durante el proceso evitando actos jurídicos que condicionen la cautela o la ejecución de la sentencia. Debe publicitarse y notificarse, y tiene una caducidad especial de cinco días desde su dictado cuando es previa al inicio del proceso.
}

\chapter[Protección de personas y familia]{Protección de personas y medidas cautelares en materia de familia}\label{cap:proteccion}

%==============================================================

\section{La protección de personas y el artículo 234 CPCCN}

El análisis de esta medida en el ámbito procesal obedece a una tradición antigua, que incluía en el antiguo art.\ 234 del código procesal una serie de potestades cautelares del juez referidas a la protección de personas. Uno de sus incisos daba amplias facultades al juez para tomar medidas de cambio de guarda y de disposición respecto de niños, en hipótesis en que se observara o se sospechara que estaban sufriendo algún tipo de maltrato, violencia o falta de cuidado.

\begin{important}
\textbf{Ley 26.061 --- Sistema Nacional de Protección Integral de Niñas, Niños y Adolescentes.} Esta ley \textbf{derogó} el inciso del art.\ 234 CPCCN que otorgaba facultades cautelares a los jueces respecto de niños. Como consecuencia, sacó de la órbita del Poder Judicial toda la gestión de cuidado y protección de situaciones de maltrato, abuso o descuido de niñas, niños y adolescentes, y las colocó en el ámbito de la \textbf{administración pública}, a través de las Defensorías de Niños, Niñas y Adolescentes.

Los llamados a intervenir ahora son la escuela, el hospital, los centros de salud, la policía y cualquier organismo que detecte una situación de riesgo.
\end{important}

Lo que sí procede y debe ser judicializado, de acuerdo con los parámetros de la ley, es el \textbf{control de legalidad} de las medidas que la administración ---las defensorías--- adopta respecto del cuidado o el abrigo de situaciones de riesgo de los niños. En tanto el órgano de la administración adopte medidas excepcionales de separación del niño de su ámbito familiar primario, estas deben ser supervisadas y controladas en cuanto a su legalidad, para que no se prolonguen en el tiempo más allá de lo necesario.

\textbf{Estado actual del art.\ 234 CPCCN.} El grueso de la actividad protectora que poseía con anterioridad ha quedado transferida a la administración pública a través del Sistema Nacional de Protección Integral de Niños, Niñas y Adolescentes (Ley 26.061). El art.\ 234 solo conserva una función tuitiva de los jueces en materia de \textbf{mayores de edad con restricciones a la capacidad por problemáticas de salud mental}.

\section{Potestades cautelares en materia de violencia}

\subsection{Ley 24.417 --- Ley de Violencia Familiar}

Establece una serie de medidas cautelares que los jueces pueden adoptar para erradicar cualquier tipo de situación de agresión padecida por personas mayores de edad (hombres, mujeres y niños involucrados con personas mayores de edad). % TODO: contenido incompleto o ambiguo en las notas originales (redacción del ámbito de aplicación)
Incluye medidas de distanciamiento, prohibición de acercamiento, fijación provisional de alimentos, regímenes de comunicación y cuidado personal de los niños.

\subsection{Ley 26.485 --- Ley de Protección de Violencia contra la Mujer}

Refuerza y amplía las medidas de la Ley 24.417 en materia de perspectiva de género. Otorga al sistema judicial potestades cautelares de amplísima gama, que permiten tomar medidas de distanciamiento, prohibición de acercamiento, prestación alimentaria, fijación de regímenes de comunicación entre padres e hijos, establecimiento del cuidado personal de los niños y medidas patrimoniales vinculadas a la distribución de bienes.

\section{Medidas cautelares en el proceso de familia}

En materia de derecho de familia, la exigencia de los presupuestos de verosimilitud del derecho, peligro en la demora y contracautela está superada por el vínculo familiar que existe entre las personas involucradas.

\begin{table}[H]
\centering
\small
\begin{tabularx}{\textwidth}{>{\raggedright\arraybackslash}p{3.2cm} >{\raggedright\arraybackslash}X}
\toprule
\textbf{Presupuesto} & \textbf{Aplicación en procesos de familia} \\
\midrule
\textbf{Verosimilitud del derecho} & Surge del propio vínculo familiar (conyugal, de pareja, filial). No es necesario evaluar más que la existencia del vínculo como presupuesto que satisface este recaudo. \\
\addlinespace
\textbf{Peligro en la demora} & Está dado \textit{in re ipsa} (en la propia naturaleza de las cosas). La urgencia de cautelar el cuidado de los hijos, los alimentos provisorios, la vivienda y el régimen de comunicación es evidente por los derechos fundamentales comprometidos. \\
\addlinespace
\textbf{Contracautela} & Prácticamente en la generalidad de los casos no existe contracautela exigible. La caución juratoria está implícita y ni siquiera se pide, ya que se torna innecesaria. \\
\bottomrule
\end{tabularx}
\caption{Presupuestos cautelares en los procesos de familia}
\end{table}

\subsection{Las cautelares de familia en el Código Civil y Comercial de la Nación}

Las cautelares de familia se encuentran en su gran mayoría insertas en el código de fondo, el Código Civil y Comercial de la Nación. Esto no es algo nuevo ni nació con la reforma del año 2015, sino que ya venía con el código anterior.

El legislador siempre ha tenido en cuenta que ciertos aspectos del conflicto familiar compete al legislador nacional garantizarlos, y no ha querido dejarlos librados a lo que pueda establecer el legislador local provincial. Por eso plasmó en el código de fondo normas eminentemente operativas vinculadas con medidas cautelares propias del derecho de familia, que abarcan:
\begin{itemize}
    \item La atribución del hogar familiar.
    \item La garantía de subsistencia alimentaria provisoria mientras se define en un proceso de conocimiento las prestaciones alimentarias para los hijos menores y, eventualmente, entre cónyuges.
    \item La definición de la manera en que los padres ejercerán el cuidado personal de los hijos.
    \item El régimen de comunicación de los hijos con sus padres.
    \item Las medidas de carácter patrimonial vinculadas con la distribución de bienes como consecuencia de la extinción del régimen matrimonial.
\end{itemize}

\begin{note}
\textbf{Lectura recomendada.} Se recomienda como lectura obligatoria el artículo publicado en el sitio web \url{www.jorgearojas.com.ar}: «Las medidas cautelares en el Código Civil y Comercial de la Nación», del Dr.\ Jorge A.\ Rojas. Analiza las medidas cautelares nominadas e innominadas del CCyCN, las heredadas del código anterior, las absolutamente novedosas y aspectos de tutelas anticipatorias contempladas en el Código Civil y Comercial.
\end{note}

\cornelltitle{protección de personas y familia}
\cornellblocksum{
\cqrow{¿Qué cambió respecto de la protección de niños con la Ley 26.061?}{La ley derogó el inciso del art.\ 234 CPCCN que daba facultades cautelares a los jueces sobre niños. La gestión de protección ante maltrato, abuso o descuido pasó a la administración pública (Defensorías de Niños, Niñas y Adolescentes), con intervención de escuelas, hospitales, centros de salud, policía y otros organismos. Al Poder Judicial le queda el control de legalidad de las medidas excepcionales de separación del niño de su ámbito familiar, para que no se prolonguen más de lo necesario.
}
\cqrow{¿Qué función conserva hoy el art.\ 234 CPCCN?}{Solo una función tuitiva de los jueces respecto de mayores de edad con restricciones a la capacidad por problemáticas de salud mental.
}
\cqrow{¿Qué medidas permiten las leyes 24.417 y 26.485?}{Otorgan al Poder Judicial potestades cautelares amplias en materia de violencia familiar y de género: distanciamiento, prohibición de acercamiento, alimentos provisorios, regímenes de comunicación y cuidado personal de los niños y, en la Ley 26.485, también medidas patrimoniales vinculadas a la distribución de bienes. La Ley 26.485 refuerza y amplía la Ley 24.417 con perspectiva de género.
}
\cqrow{¿Qué particularidades tienen los presupuestos de las cautelares en familia?}{La verosimilitud surge del propio vínculo familiar; el peligro en la demora está dado \textit{in re ipsa} por los derechos fundamentales comprometidos; y en general no hay contracautela exigible, pues la caución juratoria está implícita.
}
\cqrow{¿Dónde se regulan las cautelares de familia y qué abarcan?}{En su gran mayoría en el Código Civil y Comercial de la Nación (ya ocurría con el código anterior), con normas operativas sobre atribución del hogar familiar, alimentos provisorios, cuidado personal de los hijos, régimen de comunicación y medidas patrimoniales tras la extinción del régimen matrimonial. \\




%--------------------------------------------------------------
% BIBLIOGRAFÍA (solo la referencia que figura en el apunte)
%--------------------------------------------------------------
}
}{Con la Ley 26.061 la adopción de medidas de protección de niños se trasladó a la administración pública, reservando a los jueces el control de legalidad. Las leyes 24.417 y 26.485 otorgan al Poder Judicial amplias potestades cautelares en violencia familiar y de género. En los procesos de familia los presupuestos cautelares se satisfacen por el propio vínculo, por el peligro en la demora \textit{in re ipsa} y por la ausencia de contracautela exigible.
}

\chapter*{Bibliografía}
\addcontentsline{toc}{chapter}{Bibliografía}
\markboth{Bibliografía}{Bibliografía}
\begin{itemize}
    \item Rojas y Falcón. \textit{Cómo se hace un alegato}. Ed. Abeledo Perrot.
    \item Rojas, Jorge A. «Las medidas cautelares en el Código Civil y Comercial de la Nación». Disponible en: \url{www.jorgearojas.com.ar}.
\end{itemize}

\end{document}
